\documentclass[11pt,a4paper]{report}
\usepackage[utf8]{inputenc}
\usepackage[T1]{fontenc}
\usepackage{amsmath,amssymb,amsthm,mathtools,stmaryrd}
\usepackage{tikz-cd}
\usepackage{enumitem}
\usepackage{listings}
\usepackage{xcolor}
\usepackage{hyperref}
\usepackage[margin=1in]{geometry}
\usepackage{fancyhdr}
\usepackage{longtable}
\usepackage{booktabs}
\usepackage{tabularx}
\usepackage{caption}
\usepackage{subcaption}

% ══════════════════════════════════════════════════
% Theorem environments
% ══════════════════════════════════════════════════
\theoremstyle{definition}
\newtheorem{definition}{Definition}[chapter]
\newtheorem{example}[definition]{Example}
\newtheorem{remark}[definition]{Remark}
\newtheorem{notation}[definition]{Notation}
\newtheorem{construction}[definition]{Construction}
\theoremstyle{plain}
\newtheorem{theorem}[definition]{Theorem}
\newtheorem{lemma}[definition]{Lemma}
\newtheorem{proposition}[definition]{Proposition}
\newtheorem{corollary}[definition]{Corollary}
\newtheorem{axiom_env}[definition]{Axiom}
\newtheorem{conjecture}[definition]{Conjecture}

% ══════════════════════════════════════════════════
% Code listings
% ══════════════════════════════════════════════════
\lstdefinelanguage{FStarLang}{
  morekeywords={val,let,rec,fun,match,with,type,effect,
    assume,requires,ensures,decreases,Tot,ST,EXN,ALL,
    Lemma,squash,forall,exists,assert,int,nat,bool,string,
    list,option,None,Some,True,False,if,then,else,open,module,
    in,end,begin,of},
  sensitive=true,
  morecomment=[l]{//},
  morecomment=[s]{(*}{*)},
  morestring=[b]",
}

\lstdefinestyle{python}{
  language=Python,
  basicstyle=\ttfamily\small,
  keywordstyle=\color{blue!80!black}\bfseries,
  commentstyle=\color{green!50!black}\itshape,
  stringstyle=\color{red!60!black},
  showstringspaces=false,
  breaklines=true, frame=single,
  numbers=left, numberstyle=\tiny\color{gray},
  xleftmargin=2em,
}

\lstdefinestyle{fstar}{
  language=FStarLang,
  basicstyle=\ttfamily\small,
  keywordstyle=\color{purple!80!black}\bfseries,
  commentstyle=\color{green!50!black}\itshape,
  stringstyle=\color{red!60!black},
  showstringspaces=false,
  breaklines=true, frame=single,
  numbers=left, numberstyle=\tiny\color{gray},
  xleftmargin=2em,
}

\lstdefinestyle{synhopy}{
  language=Python,
  basicstyle=\ttfamily\small,
  keywordstyle=\color{teal!80!black}\bfseries,
  commentstyle=\color{green!50!black}\itshape,
  stringstyle=\color{red!60!black},
  showstringspaces=false,
  breaklines=true, frame=single,
  numbers=left, numberstyle=\tiny\color{gray},
  xleftmargin=2em,
  morekeywords={guarantee,spec,hole,assume,check,given,
    duck_equiv,fibration,homotopy,effect_contract,
    have,conclude,by,then,rw,omega,refl,sym,trans,cong,
    transport,apply,cases,induction},
}

\lstdefinestyle{proof}{
  basicstyle=\ttfamily\small,
  keywordstyle=\color{teal!80!black}\bfseries,
  commentstyle=\color{green!50!black}\itshape,
  showstringspaces=false,
  breaklines=true, frame=single,
  numbers=left, numberstyle=\tiny\color{gray},
  xleftmargin=2em,
  morekeywords={have,conclude,by,then,rw,omega,refl,sym,trans,
    cong,transport,apply,cases,induction,base,step,
    structural,unfold,ext,intro,exact,assumption},
}

% ══════════════════════════════════════════════════
% Math commands
% ══════════════════════════════════════════════════
\newcommand{\den}[1]{\llbracket #1 \rrbracket}
\newcommand{\Path}{\mathsf{Path}}
\newcommand{\Nat}{\mathbb{N}}
\newcommand{\Int}{\mathbb{Z}}
\newcommand{\Bool}{\mathbb{B}}
\newcommand{\Set}{\mathbf{Set}}
\newcommand{\Cat}{\mathbf{Cat}}
\newcommand{\Type}{\mathsf{Type}}
\newcommand{\Fib}{\mathsf{Fib}}
\newcommand{\Cech}{\check{\mathrm{C}}}
\newcommand{\PyObj}{\mathsf{PyObj}}
\newcommand{\PyType}{\mathsf{PyType}}
\newcommand{\PyClass}{\mathsf{PyClass}}
\newcommand{\PyFunc}{\mathsf{PyFunc}}
\newcommand{\PyDict}{\mathsf{PyDict}}
\newcommand{\PyNone}{\mathsf{None}}
\newcommand{\fold}{\mathsf{fold}}
\newcommand{\iso}{\cong}
\newcommand{\covers}{\mathcal{U}}
\newcommand{\Sem}{\mathsf{Sem}}
\newcommand{\interval}{\mathbb{I}}
\newcommand{\rec}{\mathsf{rec}}
\newcommand{\elim}{\mathsf{elim}}
\newcommand{\transp}{\mathsf{transp}}
\newcommand{\comp}{\mathsf{comp}}
\newcommand{\hfilling}{\mathsf{hfill}}
\newcommand{\refl}{\mathsf{refl}}
\newcommand{\Refl}{\mathsf{refl}}
\newcommand{\sym}{\mathsf{sym}}
\newcommand{\ap}{\mathsf{ap}}
\newcommand{\apd}{\mathsf{apd}}
\newcommand{\Ext}{\mathsf{Ext}}
\newcommand{\Cut}{\mathsf{Cut}}
\newcommand{\Assume}{\mathsf{Assume}}
\newcommand{\Trans}{\mathsf{Trans}}
\newcommand{\Cong}{\mathsf{Cong}}
\DeclareMathOperator{\Hom}{Hom}
\newcommand{\defeq}{\coloneqq}
\newcommand{\Glue}{\mathsf{Glue}}
\newcommand{\Ctx}{\mathsf{Ctx}}
\newcommand{\Eff}{\mathsf{Eff}}
\newcommand{\Pure}{\mathsf{Pure}}
\newcommand{\Reads}{\mathsf{Reads}}
\newcommand{\Mutates}{\mathsf{Mutates}}
\newcommand{\IO}{\mathsf{IO}}
\newcommand{\Nondet}{\mathsf{Nondet}}
\newcommand{\DuckEq}{\mathsf{DuckEq}}
\newcommand{\MethodSig}{\mathsf{MethodSig}}
\newcommand{\MRO}{\mathsf{MRO}}
\newcommand{\Spec}{\mathsf{Spec}}
\newcommand{\Pre}{\mathsf{Pre}}
\newcommand{\Post}{\mathsf{Post}}
\newcommand{\Inv}{\mathsf{Inv}}
\renewcommand{\wp}{\mathsf{wp}}
\newcommand{\fstar}{F$^{*}$}
\newcommand{\fstarmath}{\text{F}^{*}}
\newcommand{\synhopy}{\textsc{SynHoPy}}
\newcommand{\Sh}{\mathsf{Sh}}
\newcommand{\Fiber}{\mathsf{Fiber}}

% ══════════════════════════════════════════════════
% Page style
% ══════════════════════════════════════════════════
\pagestyle{fancy}
\fancyhead[L]{\leftmark}
\fancyhead[R]{\thepage}
\fancyfoot{}

% ══════════════════════════════════════════════════
\title{
  \textbf{SynHoPy: Synthetic Homotopy Type Theory}\\
  \textbf{for Python Verification}\\[1em]
  \large A Complete Proof Theory with Soundness Proof\\
  and Systematic Comparison to \fstar{}
}
\author{deppy project}
\date{April 2026}

\begin{document}
\maketitle

\begin{abstract}
We present \synhopy{} (Synthetic Homotopy Python), a dependent type theory based on
cubical homotopy type theory that provides a sound foundation for verifying
\emph{all} of Python's semantics---including duck typing, monkey-patching,
decorators, generators, metaclasses, multiple inheritance, closures with late
binding, \texttt{*args/**kwargs}, and context managers.

The key insight is that Python's dynamic features are not defects to be
encoded away, but are themselves homotopy-theoretic constructs: duck typing
is the univalence axiom, monkey-patching is homotopy, decorators are fiber
bundles, exceptions are Eilenberg--MacLane spaces, and gradual typing is
Postnikov tower truncation.

We prove soundness via a model in the $\infty$-topos of simplicial presheaves
over the site of Python execution contexts. We show that \synhopy{} is
conservative over deppy's existing CCTT proof system, and that it strictly
dominates \fstar{} in verification capability: tying on three dimensions
(pure functions, refinement types, dependent types) and winning on fourteen
others.

For each Python language feature, we give: (1) the formal typing rules, (2) the
proof obligations for specification checking, (3) a worked example with full
proof, and (4) a detailed comparison showing what \fstar{} would require for the
same verification. Throughout, proofs are attached to code in the \fstar{} style---the
system verifies existing Python programs in-place, with an LLM generating all
proof annotations that are then mechanically checked by the kernel.
\end{abstract}

\tableofcontents

% ══════════════════════════════════════════════════════════════════
% PART I: FOUNDATIONS
% ══════════════════════════════════════════════════════════════════

\part{Foundations}

% ──────────────────────────────────────────────────────────────────
\chapter{Introduction}
\label{ch:intro}

\section{The Problem: Verifying a Dynamic Language}

Python is the world's most popular programming language, yet it has no
mainstream formal verification system. The gap is not accidental---Python's
semantics are fundamentally hostile to the techniques used by existing
verification tools:

\begin{itemize}
  \item \textbf{Duck typing}: A value's usability depends on what methods it has,
    not what class it belongs to. Type identity is behavioral, not nominal.
  \item \textbf{Monkey-patching}: Methods can be added, removed, or replaced on
    objects and classes at runtime. The type of an object can change during execution.
  \item \textbf{Decorators}: Functions are routinely wrapped in composable
    transformers. The order of decoration may or may not matter.
  \item \textbf{Metaclasses}: Classes are themselves objects, created by metaclasses
    that can modify class structure at creation time.
  \item \textbf{Multiple inheritance}: Python's C3 linearization creates complex
    method resolution orders that are not captured by simple subtyping.
  \item \textbf{Generators}: Coroutines that yield values, maintain state, and
    resume---a form of delimited continuation.
  \item \textbf{Late binding}: Closures capture variables by reference, not by
    value. The captured value can change after the closure is created.
  \item \textbf{\texttt{*args, **kwargs}}: Functions can accept arbitrary
    positional and keyword arguments, making the function's type depend on the
    call site.
  \item \textbf{Everything is a dict}: Object attribute access is fundamentally
    dictionary lookup through the MRO chain.
\end{itemize}

Existing verification systems handle these features in one of three ways:

\begin{enumerate}
  \item \textbf{Forbid them}: Verify only the subset of Python that looks like
    Haskell or Java. This is the approach of mypy, pytype, and most gradual
    typing systems.
  \item \textbf{Encode them}: Represent duck typing as existential types, monkey-patching
    as heap mutation, decorators as function composition. This is what an \fstar{}-based
    approach would require. The encoding is lossy, verbose, and doesn't match
    Python's actual semantics.
  \item \textbf{Ignore them}: Verify pure functions and hope the dynamic parts
    don't break anything. This is the approach of most academic Python verification
    work.
\end{enumerate}

\synhopy{} takes a fourth approach: \textbf{embrace them}. We show that Python's
dynamic features are not bugs but \emph{homotopy-theoretic constructs}---and
that Homotopy Type Theory provides the natural mathematical framework for
reasoning about them soundly.

\section{The Solution: Homotopy Type Theory for Python}

The central observation of this monograph is:

\begin{quote}
\itshape
Python's runtime semantics have inherent higher-dimensional structure that
is precisely captured by Homotopy Type Theory with the univalence axiom.
\end{quote}

The identifications are:

\begin{center}
\begin{tabular}{ll}
\toprule
\textbf{Python Feature} & \textbf{HoTT Construct} \\
\midrule
Duck typing & Univalence axiom \\
Monkey-patching & Homotopy between type implementations \\
\texttt{isinstance} & Serre fibration \\
Decorators & Fiber bundles over function space \\
Exceptions & Eilenberg--MacLane space $K(\pi, 1)$ \\
Gradual typing & Postnikov tower truncation \\
Generators & Paths in state space \\
Multiple inheritance & Sections of class fiber bundle \\
Closures & Sections of scope bundle \\
\texttt{*args/**kwargs} & Dependent dictionary sheaf \\
Metaclasses & 2-types with HIT constructors \\
\bottomrule
\end{tabular}
\end{center}

These are not analogies. Each identification is a precise mathematical theorem,
proved via a model in the $\infty$-topos of simplicial presheaves over the site
of Python execution contexts.

\section{Comparison Philosophy: \synhopy{} vs.\ \fstar{}}

Throughout this monograph, we compare \synhopy{} systematically to \fstar{}
(Swamy et al., 2016), the state-of-the-art in effect-aware dependent type systems.
\fstar{} is the closest existing system to what we aim to build, and understanding
\emph{precisely} where it falls short on Python illuminates why a new
foundation is needed.

Our comparisons follow a consistent structure:

\begin{enumerate}
  \item \textbf{The Python code}: the actual program to be verified.
  \item \textbf{The \fstar{} approach}: what you would need to write in \fstar{}
    (or an \fstar{}-inspired system) to verify this code.
  \item \textbf{The \synhopy{} approach}: what the programmer writes and what the
    system does automatically.
  \item \textbf{Analysis}: why the \synhopy{} approach is sound, what guarantees
    it provides, and how the proof burden compares.
\end{enumerate}

\section{A Brief History of Dependent Type Systems}

The quest for machine-verified software correctness has a long and rich history.
\synhopy{} builds on decades of foundational work in type theory, proof assistants,
and programming language design.

\subsection{Martin-L\"of Type Theory (1971--1984)}

Per Martin-L\"of introduced \emph{intuitionistic type theory} in the early
1970s, providing the first dependent type system with a clear computational
interpretation.  His theory introduced dependent function types $(x : A) \to B(x)$,
dependent pair types $\Sigma(x : A). B(x)$, identity types $a =_A b$, and
well-founded induction---all of which appear as core constructs in \synhopy{}.
Martin-L\"of's \emph{meaning explanations} established that types are
propositions and terms are proofs, creating the Curry--Howard correspondence
that underlies all modern proof assistants.

\subsection{Coq (1984--present)}

Thierry Coquand and G\'erard Huet developed the Calculus of Constructions,
which became the foundation for the Coq proof assistant (now Rocq).  Coq
introduced \emph{inductive types} via the Calculus of Inductive Constructions,
\emph{tactic-based} proof development, and \emph{program extraction} to
OCaml and Haskell.  Coq's tactic language inspired \synhopy{}'s
LLM-as-tactic framework (Chapter~\ref{ch:llm-tactic}), though \synhopy{}
replaces hand-written tactics with LLM-generated proof terms that are
mechanically verified by the kernel.

\subsection{Agda (1999--present)}

Agda, developed by Ulf Norell, brought dependent types closer to
practical programming with its emphasis on \emph{pattern matching},
\emph{termination checking}, and a clean surface syntax.  Agda's
cubical mode (Cubical Agda), implementing cubical type theory, was a
direct inspiration for \synhopy{}'s path types and transport operations.
However, Agda targets total functional programs---its type system has
no story for effects, mutation, or dynamic dispatch.

\subsection{Lean (2013--present)}

Leonardo de Moura's Lean system, particularly Lean~4, combines
a dependent type theory (based on the Calculus of Inductive
Constructions) with a practical programming language.  Lean's
\emph{tactic framework} (based on monadic metaprogramming),
\emph{type class inference}, and \emph{Mathlib} library for
formalized mathematics represent the state of the art in
usable proof assistants.  \synhopy{}'s long-term goal includes
translating its proofs to Lean for machine-checked meta-theory
(Appendix~\ref{app:future}).

\subsection{F$^*$ (2011--present)}

\fstar{}, developed by Nikhil Swamy and collaborators at Microsoft Research,
introduced \emph{Dijkstra monads} for effect-indexed verification and
\emph{layered effects} for composing verification disciplines.  \fstar{} is
the closest existing system to \synhopy{} in ambition: it aims to verify
effectful programs with dependent types.  However, \fstar{}'s type system
is designed for ML-like languages with static types, nominal subtyping,
and controlled effects---precisely the features Python lacks.  Our
systematic comparison throughout this monograph shows that \fstar{} falls
short on fourteen Python-specific dimensions.

\subsection{Idris (2007--present)}

Edwin Brady's Idris (and Idris~2) demonstrated that dependent types can
be integrated into a \emph{general-purpose programming language} with
first-class support for effects via \emph{algebraic effect handlers}.
Idris's \emph{elaborator reflection} allows type-level computation
to inspect and transform programs at compile time.  \synhopy{} shares
Idris's goal of making dependent types practical, but targets an
existing language (Python) rather than designing a new one.

\subsection{Where \synhopy{} Fits}

\synhopy{} is the first system to combine:
\begin{itemize}
  \item \textbf{Cubical/homotopy type theory} (from Cubical Agda) for
    path-based reasoning about equivalences;
  \item \textbf{Effect-indexed verification} (from \fstar{}) for reasoning
    about side effects;
  \item \textbf{LLM-based proof automation} (novel) for eliminating the
    proof burden; and
  \item \textbf{Verification of an existing dynamic language} (novel) rather
    than a language designed for verification.
\end{itemize}

\section{Why Not \fstar{}?}
\label{sec:why-not-fstar}

To make the case for a new system concrete, we present five Python programs
that \fstar{} \emph{cannot} verify---not because the programs are exotic,
but because they use ordinary Python idioms.

\subsection{Program 1: Duck-Typed Polymorphism}

\begin{lstlisting}[style=synhopy]
@guarantee("returns the sum of all elements")
def total(things):
    return sum(x.value() for x in things)
\end{lstlisting}

This function works on \emph{any} iterable whose elements have a
\lstinline[style=python]{value()} method.  In \fstar{}, you would need
an explicit type class or interface.  There is no way to express ``anything
with a \lstinline[style=python]{value()} method'' as a first-class type
in \fstar{}'s type system.  \synhopy{} handles this directly via
duck-type equivalence (Chapter~\ref{ch:duck-typing}): the structural
protocol $\{x \mid x \models \mathtt{value} : () \to \mathsf{Num}\}$
is a well-formed type.

\subsection{Program 2: Monkey-Patched Method}

\begin{lstlisting}[style=synhopy]
class Logger:
    def log(self, msg): print(msg)

@guarantee("Logger.log writes to file after patching")
def setup_file_logging(logger, path):
    original = logger.log
    def file_log(msg):
        with open(path, 'a') as f:
            f.write(msg + '\n')
        original(msg)
    logger.log = file_log
\end{lstlisting}

After \lstinline[style=python]{setup_file_logging}, the
\lstinline[style=python]{Logger} instance has a different
\lstinline[style=python]{log} method.  \fstar{}'s nominal type
system cannot express that an object's type changes at runtime.
\synhopy{} models this as a homotopy between type implementations
(Chapter~\ref{ch:monkey-patching}).

\subsection{Program 3: Decorator Composition}

\begin{lstlisting}[style=synhopy]
@retry(max_attempts=3)
@cache(ttl=60)
@validate_input(schema=UserSchema)
@guarantee("returns validated, cached user with retry")
def get_user(user_id: int) -> User:
    return db.query(User, user_id)
\end{lstlisting}

Three decorators compose around the function.  Verifying the composite
requires reasoning about decorator commutativity (does the order
matter?) and behavioral preservation (does each decorator preserve
the core contract?).  \fstar{} has no decorator construct; encoding
this as nested function application loses the commutativity structure.
\synhopy{} models decorators as fiber bundles and reasons about their
composition via the fiber bundle's structure group
(Chapter~\ref{ch:decorators}).

\subsection{Program 4: Generator with State}

\begin{lstlisting}[style=synhopy]
@guarantee("yields Fibonacci numbers in order")
def fibonacci():
    a, b = 0, 1
    while True:
        yield a
        a, b = b, a + b
\end{lstlisting}

This infinite generator maintains internal state across
\lstinline[style=python]{yield} boundaries.  \fstar{} has no
generator construct; encoding this requires explicitly
threading state through a continuation monad, losing the
natural correspondence between the code and the mathematical
definition of the Fibonacci sequence.
\synhopy{} models generators as paths in state space
(Chapter~\ref{ch:generators}), making the verification
follow the mathematical structure directly.

\subsection{Program 5: Multiple Inheritance with \texttt{super()}}

\begin{lstlisting}[style=synhopy]
class Serializable:
    def save(self): return json.dumps(self.__dict__)

class Validatable:
    def validate(self): return all(v is not None for v in self.__dict__.values())

@guarantee("saves only if valid")
class SafeModel(Serializable, Validatable):
    def save(self):
        assert self.validate()
        return super().save()
\end{lstlisting}

The \lstinline[style=python]{super()} call follows Python's C3
linearization MRO.  \fstar{} has no multiple inheritance; encoding
the MRO as explicit function dispatch is possible but loses the
guarantee that \lstinline[style=python]{super().save()} resolves
to \lstinline[style=python]{Serializable.save} under the C3 order.
\synhopy{} models the MRO as a section of the class fiber bundle
(Chapter~\ref{ch:mro}).

\section{The Key Insight: Dynamism as Feature}
\label{sec:key-insight}

The fundamental observation that motivates \synhopy{} is:

\begin{quote}
\itshape
Python's dynamic features are not defects that need to be encoded away
or forbidden.  They are \emph{higher-dimensional structure} that homotopy
type theory was \emph{designed} to reason about.
\end{quote}

Consider duck typing.  In a nominal type system, two objects with the same
methods but different class names are \emph{different types}.  This is like
saying that two topological spaces with the same homotopy groups are
different because they have different names.  The univalence axiom says:
\emph{equivalent structures are identical}.  This is exactly duck typing:
if two objects support the same operations with the same behavior, they are
interchangeable.

Consider monkey-patching.  Replacing a method on an object is a
\emph{continuous deformation} of the object's type---a homotopy.  The
patched and unpatched objects are connected by a path in type space.
Transport along this path transfers specifications from the original to
the patched version.

Consider decorators.  A decorator wraps a function, creating a new function
that projects down to the original.  This is the definition of a
\emph{fiber bundle}: a total space (decorated function) that fibers over a
base space (original function) with a typical fiber (the decoration behavior).

These are not analogies or metaphors.  Each identification is a theorem,
proved in the $\infty$-topos model (Chapter~\ref{ch:topos}).  The fact
that the correspondences are exact---not approximate---is what gives
\synhopy{} its soundness guarantees.

The practical consequence is profound: instead of fighting Python's
semantics, \synhopy{} \emph{exploits} them.  Every dynamic feature becomes
a source of proof power, not a source of unsoundness.  A monkey-patched
method is not an unverifiable mutation; it is a transport along a path
that the type system tracks.  A duck-typed function is not an untyped
black box; it is a function on a structural type whose equivalences are
managed by univalence.

\section{Contributions}
\label{sec:contributions}

This monograph makes the following eight contributions:

\begin{enumerate}
  \item \textbf{A homotopy type theory for Python} (\synhopy{}): a dependent
    type theory based on cubical HoTT that natively captures duck typing via
    univalence, monkey-patching via homotopy, decorators via fiber bundles,
    and all other Python-specific features (Chapters~2--3).

  \item \textbf{A proof kernel with trust tracking}: a small, auditable
    kernel that checks proof terms and tracks trust levels from fully
    kernel-verified to axiom-dependent, with explicit dependency graphs
    (Chapter~4).

  \item \textbf{An effect system based on $\Omega$-endomorphisms}: an effect
    system more expressive than \fstar{}'s Dijkstra monads, with automatic
    multi-effect analysis via spectral sequences and composition governed
    by the Blakers--Massey theorem (Chapter~3).

  \item \textbf{Complete verification of fourteen Python features}: full
    typing rules, proof obligation generation, worked examples, and
    \fstar{} comparisons for duck typing, monkey-patching, decorators,
    exceptions, generators, closures, multiple inheritance, metaclasses,
    \texttt{*args/**kwargs}, context managers, comprehensions, gradual
    typing, string operations, and truthiness (Chapters~5--14).

  \item \textbf{Soundness via an $\infty$-topos model}: a model in the
    $\infty$-topos of simplicial presheaves over the site of Python execution
    contexts that validates all axioms and typing rules (Chapter~15).

  \item \textbf{Consistency and conservativity proofs}: proofs that
    \synhopy{} is consistent (no term inhabits the empty type) and
    conservative over deppy's existing CCTT system
    (Chapter~16).

  \item \textbf{LLM-as-tactic proof automation}: a framework where
    LLMs generate proof annotations that are mechanically verified by
    the kernel, eliminating the proof burden while maintaining full
    soundness (Chapters~18--20).

  \item \textbf{Strict domination of \fstar{} for Python}: a
    systematic demonstration that \synhopy{} ties \fstar{} on three
    dimensions (pure functions, refinement types, dependent types) and
    strictly dominates on fourteen Python-specific dimensions
    (Appendix~\ref{app:comparison}).
\end{enumerate}

\section{Roadmap}
\label{sec:roadmap}

This monograph is organized into five parts, three appendices, and a
collection of case studies.

\begin{description}
  \item[Part I: Foundations] (Chapters 1--4) introduces the type theory, its model,
    and the proof kernel.
  \item[Part II: Python Semantics] (Chapters 5--14) shows how to verify each
    aspect of Python, with full proof rules and worked examples.
  \item[Part III: Soundness] (Chapters 15--17) proves soundness, consistency,
    conservativity, and anti-hallucination properties.
  \item[Part IV: Proof Automation] (Chapters 18--20) describes the LLM-as-tactic
    framework, proof caching, and the trust lattice.
  \item[Part V: Case Studies] (Chapters 21--23) presents end-to-end verifications
    of real Python programs.
\end{description}

\noindent We now describe each part in more detail.

\paragraph{Part I: Foundations.}
Chapter~\ref{ch:intro} (this chapter) motivates the work and surveys related
systems.  Chapter~2 defines the \synhopy{} type theory: its judgments, the
universe $\PyObj$, type formers ($\PyClass$, function types, path types,
transport, Kan operations, and higher inductive types), and the Python state
context.  Chapter~3 presents the effect system, including the effect lattice,
$\Omega$-endomorphisms, the Blakers--Massey theorem for effect composition,
spectral sequences for multi-effect analysis, and the effect inference
algorithm.  Chapter~4 defines the proof kernel: proof terms, kernel typing
rules, middle-term inference, structural chain verification, and the trust
lattice.

\paragraph{Part II: Python Semantics.}
Chapters~5--14 form the heart of the monograph.  Each chapter takes one
Python feature, defines the formal typing rules, shows how proof obligations
are generated from \lstinline[style=synhopy]{@guarantee} and
\lstinline[style=synhopy]{assume()} annotations, works through a complete
example with full proof, and compares to what \fstar{} would require.
The features covered are: pure functions and refinement types (Chapter~5),
duck typing (Chapter~6), monkey-patching (Chapter~7), decorators and fiber
bundles (Chapter~8), exceptions and classifying spaces (Chapter~9),
generators and coroutines (Chapter~10), closures and late binding
(Chapter~11), multiple inheritance and the MRO (Chapter~12), metaclasses
and runtime class creation (Chapter~13), and \texttt{*args/**kwargs} and
the dict sheaf (Chapter~14).

\paragraph{Part III: Soundness.}
Chapter~15 presents the denotational semantics of core Python and proves
the soundness of the type system via the $\infty$-topos model.  Chapter~16
establishes consistency (no closed term of type $\bot$) and conservativity
over deppy's CCTT.  Chapter~17 proves the anti-hallucination theorem: any
specification verified by the LLM oracle and accepted by the kernel is
genuinely satisfied by the code.

\paragraph{Part IV: Proof Automation.}
Chapter~18 describes the LLM-as-tactic framework: how proof obligations
are rendered for the LLM, how LLM-generated proofs are compiled and
verified, and how proofs are cached via univalence.  Chapter~19 presents
the complete verification pipeline from source code to verified artifact.
Chapter~20 describes the library axiom system for reasoning about
third-party code.

\paragraph{Part V: Case Studies.}
Chapters~21--23 present end-to-end verifications of real Python programs:
a web framework route handler, a data processing pipeline, a SymPy
computation, a machine learning pipeline, a concurrent task queue, and
a collection of design patterns.


% ──────────────────────────────────────────────────────────────────
\chapter{The \synhopy{} Type Theory}
\label{ch:type-theory}

We define \synhopy{} as an extension of Cubical Type Theory (Cohen et al., 2018)
with Python-specific type formers. The base theory provides paths, transport,
composition, and the univalence axiom. We add type formers for Python's objects,
classes, functions, and effects.

\section{Judgments}

\synhopy{} has four judgment forms:

\begin{align}
  \Gamma &\vdash A \; \Type && \text{$A$ is a well-formed type in context $\Gamma$} \\
  \Gamma &\vdash a : A && \text{$a$ is a term of type $A$} \\
  \Gamma &\vdash a \equiv b : A && \text{$a$ and $b$ are definitionally equal} \\
  \Gamma &\vdash p : a =_A b && \text{$p$ is a path from $a$ to $b$ in $A$}
\end{align}

Contexts $\Gamma$ are sequences of typed variable declarations and \emph{interval
variable declarations}:
\[
  \Gamma ::= \cdot \mid \Gamma, x : A \mid \Gamma, i : \interval
\]
where $\interval$ is the abstract interval with endpoints $0$ and $1$.

\section{The Universe $\PyObj$}

The universe of all Python objects is:

\begin{definition}[$\PyObj$]
$\PyObj$ is a type with the following structure:
\begin{enumerate}
  \item Every Python runtime value inhabits $\PyObj$.
  \item $\PyObj$ has homotopy level $\infty$ (it is not truncated).
  \item For any $a, b : \PyObj$, the path type $a =_{\PyObj} b$ is inhabited
    if and only if $a$ and $b$ are \emph{observationally equivalent}---they
    produce the same results when the same methods are called on them with
    the same arguments.
\end{enumerate}
\end{definition}

The key difference from \fstar{}: in \fstar{}, the universe of types is
stratified with no paths between distinct types. In \synhopy{}, types form a
\emph{groupoid}---there are non-trivial paths between types, and these paths
capture Python's notion of ``same enough to substitute.''

\begin{remark}[\fstar{} comparison]
\fstar{}'s type universe has structure:
\[
  \text{F}^{*}: \quad \Type_0 : \Type_1 : \Type_2 : \cdots
\]
This is a pure hierarchy with no horizontal structure. Two types at the same level
are either definitionally equal or completely unrelated. There is no concept of
``equivalent but not identical,'' which is exactly what Python's duck typing requires.
\end{remark}

\section{Type Formers}

\subsection{The Python Class Type $\PyClass$}

\begin{definition}[$\PyClass$]
A Python class is given by a record:
\[
  \PyClass \defeq \Sigma\left(
    \begin{array}{l}
      \mathit{name} : \mathsf{String}, \\
      \mathit{bases} : \mathsf{List}(\PyClass), \\
      \mathit{methods} : \mathsf{String} \to \mathsf{Option}(\PyFunc), \\
      \mathit{metaclass} : \PyClass
    \end{array}
  \right)
\]
Two classes $C_1, C_2 : \PyClass$ are connected by a path $p : C_1 =_{\PyClass} C_2$
if and only if they have the same \emph{structural interface}: the same method
names with pointwise-equivalent implementations.
\end{definition}

\begin{definition}[Duck-type equivalence]
\label{def:duck-eq}
Two classes $A$ and $B$ are \emph{duck-type equivalent}, written $A \simeq_{\mathsf{duck}} B$,
if for every method name $m$ in the interface of $A$:
\begin{enumerate}
  \item $B.\mathit{methods}(m)$ is defined (i.e., not $\mathsf{None}$), and
  \item $A.\mathit{methods}(m) \simeq B.\mathit{methods}(m)$ as functions
    (extensionally equal on all inputs satisfying the method's precondition).
\end{enumerate}
\end{definition}

\begin{axiom_env}[Duck-Typing Univalence]
\label{ax:duck-univalence}
\[
  (A \simeq_{\mathsf{duck}} B) \;\simeq\; (A =_{\PyClass} B)
\]
Duck-type equivalent classes are equal as types.
\end{axiom_env}

This axiom is the engine of \synhopy{}. It says that if two classes have the same
structural interface and their methods behave the same way, then they are the same
type---proofs about one transport to the other.

\begin{remark}[\fstar{} comparison]
In \fstar{}, to achieve a similar effect, you must:
\begin{enumerate}
  \item Define a typeclass: \lstinline[style=fstar]|noeq type quackable a = { quack: a -> string }|
  \item Provide instances for each class
  \item Thread the typeclass constraint through every function that uses duck typing
  \item Manually prove that instances are coherent
\end{enumerate}
This is approximately 20--50 lines of boilerplate per protocol. In \synhopy{},
the duck-typing axiom eliminates all of this---duck-type equivalence is
\emph{automatic}.
\end{remark}


\subsection{The Function Type}

\begin{definition}[Effect-indexed function type]
The type of functions from $A$ to $B$ with effect $\varepsilon$ is:
\[
  (x : A) \to_\varepsilon B(x)
\]
where $\varepsilon \in \{\Pure, \Reads, \Mutates, \IO, \Nondet\}$ is an effect
label from the effect lattice (Chapter~\ref{ch:effects}).
\end{definition}

The effect annotation is part of the type. Two functions with different effects
are of different types, even if they compute the same result on all inputs.

\begin{remark}[\fstar{} comparison]
\fstar{} also has effect-indexed types: \lstinline[style=fstar]{val f: int -> ST int pre post}.
The key difference is that \fstar{} requires \emph{wp-specifications} (weakest
preconditions) as part of the type:
\[
  \fstarmath: \quad (x : A) \to \mathsf{ST}\;B\;\mathit{wp}
\]
where $\mathit{wp}$ is a predicate transformer that must be manually written
for every function. \synhopy{} separates the effect annotation (lightweight) from
the specification (which can be added gradually via the Postnikov tower).
\end{remark}

\subsection{The Path Type}

\begin{definition}[Path type]
For $A : \Type$ and $a, b : A$, the path type $a =_A b$ is:
\[
  a =_A b \;\defeq\; (i : \interval) \to A \quad\text{with}\quad p(0) \equiv a,\; p(1) \equiv b
\]
\end{definition}

Path types are the workhorse of \synhopy{}. Every proof of equivalence, every
behavioral comparison, every specification check is ultimately a path
construction.

Paths compose (transitivity), reverse (symmetry), and lift through functions
(congruence):
\begin{align}
  \mathsf{trans} &: (a =_A b) \to (b =_A c) \to (a =_A c) \\
  \sym &: (a =_A b) \to (b =_A a) \\
  \ap_f &: (a =_A b) \to (f(a) =_B f(b))
\end{align}


\subsection{Transport and the Specification Transfer Principle}

\begin{definition}[Transport]
Given a type family $P : A \to \Type$ and a path $p : a =_A b$, transport gives:
\[
  \transp^P(p, -) : P(a) \to P(b)
\]
\end{definition}

The specification transfer principle follows immediately:

\begin{theorem}[Specification Transfer]
\label{thm:spec-transfer}
If $f \simeq_{\mathsf{duck}} g$ (as functions with duck-type equivalent interfaces),
and $\Spec(f)$ holds (where $\Spec$ is any specification predicate), then
$\Spec(g)$ holds.
\end{theorem}
\begin{proof}
By Axiom~\ref{ax:duck-univalence}, $f \simeq_{\mathsf{duck}} g$ gives a path
$p : f =_{\PyFunc} g$. Then $\transp^{\Spec}(p, \mathit{proof}_f) : \Spec(g)$.
\end{proof}

This is \synhopy{}'s ``prove once, use everywhere'' principle. \fstar{} has no
analog---each function requires its own proof from scratch.

\subsection{Kan Operations and Composition}

The cubical structure of \synhopy{} provides \emph{Kan operations}---the
mechanism for computing with paths.

\begin{definition}[Kan composition]
Given a partial path system (an open box), Kan composition fills the box:
\[
  \frac{
    \Gamma, i : \interval \vdash A \;\Type \qquad
    \varphi : \mathsf{Face} \qquad
    \Gamma, \varphi, i : \interval \vdash u : A \qquad
    \Gamma \vdash u_0 : A(i=0)[\varphi \mapsto u(i=0)]
  }{
    \Gamma \vdash \comp^A[\varphi \mapsto u]\; u_0 : A(i=1)
  }
\]
This says: given a type that varies over the interval, a partial tube of
elements (defined on face $\varphi$), and a base element that agrees with the
tube at $i=0$, we can compute the element at $i=1$.
\end{definition}

\begin{example}[Kan filling for Python type migration]
Suppose a function $f$ works with class $A$ and we want to migrate to class $B$
where $A \simeq_{\mathsf{duck}} B$.  The duck-equivalence gives a path
$p : A =_{\PyClass} B$.  Kan composition along $p$ transports $f$'s
implementation from $A$ to $B$:

\[
  f_B \defeq \comp^{(C : \PyClass) \to C \to R}[\cdot]\; f_A : B \to R
\]

The result $f_B$ is the ``migrated'' version of $f$ that works on $B$-typed
inputs.  The migration preserves correctness: any proof about $f_A$ transports
to $f_B$.
\end{example}

\begin{definition}[Glue types]
Glue types are the key to implementing univalence computationally:
\[
  \Glue(A, \varphi, (T, e)) \;\defeq\; \text{the type that is $T$ on face $\varphi$
    and $A$ elsewhere, glued by equivalence $e : T \simeq A$}
\]
In \synhopy{}, Glue types implement duck-type coercion: on the face where we
know $x : T$ (the concrete class), we use $x$ directly; elsewhere, we coerce
through the equivalence to the duck-type interface.
\end{definition}

\begin{remark}[\fstar{} comparison]
\fstar{} has no Kan operations, Glue types, or computational univalence.
Coercions between equivalent types must be written as explicit functions and
proved correct manually.  In \synhopy{}, the cubical structure provides these
coercions for free---they are \emph{computed} from the path, not \emph{proved}
by the user.
\end{remark}


\subsection{Higher Inductive Types for Python Constructs}

\begin{definition}[Python Class HIT]
\label{def:pyclass-hit}
The type $\PyClass$ is a higher inductive type with:
\begin{itemize}
  \item \textbf{Point constructors}: for each Python class definition
    \lstinline[style=python]{class C(B1, B2): ...}, a point $\mathsf{cls}(C) : \PyClass$.
  \item \textbf{Point constructor for runtime class creation}:
    $\mathsf{type\_call}(\mathit{name}, \mathit{bases}, \mathit{dict}) : \PyClass$.
  \item \textbf{Path constructor (duck equivalence)}: for any $C_1, C_2$ with
    $C_1 \simeq_{\mathsf{duck}} C_2$, a path $\mathsf{duck}(C_1, C_2) : C_1 =_{\PyClass} C_2$.
  \item \textbf{Path constructor (monkey-patch)}: for any $C$ and method $m$ with
    new implementation $v$, a path
    $\mathsf{patch}(C, m, v) : C =_{\PyClass} C[m := v]$
    provided the behavioral contract is preserved.
\end{itemize}
\end{definition}


\section{Contexts and the Python State}

A \synhopy{} context $\Gamma$ includes not just type declarations but also
\emph{heap state} and \emph{effect history}:

\begin{definition}[Extended context]
\[
  \Gamma ::= \cdot
    \mid \Gamma, x : A
    \mid \Gamma, i : \interval
    \mid \Gamma, \sigma : \mathsf{Heap}
    \mid \Gamma, \varepsilon : \Eff
\]
where $\mathsf{Heap}$ is the type of Python heap states (a mapping from
object IDs to attribute dictionaries), and $\Eff$ records the accumulated
effects.
\end{definition}

\begin{remark}[\fstar{} comparison]
\fstar{} also threads heap state through typing, using \emph{monotonic state}
with the \lstinline[style=fstar]{witnessed} predicate for heap stability. \synhopy{}'s approach
is more general: the heap is a section of the state sheaf over the execution
timeline, and mutations are morphisms in the site $\Ctx$. This captures Python's
actual heap model (reference-counted objects with {\tt \_\_dict\_\_}) rather than
\fstar{}'s simplified model (typed references \lstinline[style=fstar]{ref a}).
\end{remark}


\section{Identity Types and Path Types: Computation Rules}

The path type introduced in \S2.3.4 is the centerpiece of \synhopy{}.  We now
give its full computation rules and demonstrate their Python significance.

\begin{definition}[Path introduction --- $\Refl$]
For any $a : A$, the \emph{reflexivity path} is:
\[
  \Refl_a \;\defeq\; \lambda(i : \interval).\; a \;:\; a =_A a
\]
This is a constant function on the interval: at every point $i$, it returns $a$.
\end{definition}

\begin{definition}[Path elimination --- $J$-rule]
Given a type family $C : \prod_{(x,y:A)} (x =_A y) \to \Type$ and a term
$d : \prod_{(x:A)} C(x, x, \Refl_x)$, the $J$-eliminator provides:
\[
  J(C, d, x, y, p) : C(x, y, p)
\]
with the computation rule $J(C, d, x, x, \Refl_x) \equiv d(x)$.
\end{definition}

In cubical type theory, the $J$-rule is \emph{derivable} from transport along
the path, giving it computational content that the axiomatic HoTT $J$-rule
lacks.

\begin{example}[Path between Python objects]
Consider two class instances \lstinline[style=python]{a = Point(1, 2)} and
\lstinline[style=python]{b = Point(1, 2)}.  In Python, \lstinline[style=python]{a is not b}
(different identity), but \lstinline[style=python]{a == b} (equal value).
In \synhopy{}, there is a non-trivial path $p : a =_{\mathsf{Point}} b$
witnessing observational equivalence: every method called on $a$ produces
the same result as on $b$.  This path is \emph{not} $\Refl$ because $a$ and
$b$ are not definitionally equal (they occupy different memory locations).
\end{example}

\begin{definition}[Function extensionality for Python]
\label{def:funext}
Two Python functions $f, g : (x : A) \to_\varepsilon B(x)$ are path-equal
if and only if they agree pointwise:
\[
  \mathsf{funext} : \left(\prod_{x : A} f(x) =_B g(x)\right) \to (f =_{A \to_\varepsilon B} g)
\]
In cubical type theory this is a theorem, not an axiom:
$\mathsf{funext}(h) \defeq \lambda(i : \interval).\;\lambda(x : A).\; h(x)(i)$.
\end{definition}

\begin{remark}[\fstar{} comparison]
\fstar{} postulates function extensionality as an axiom via
\lstinline[style=fstar]|val funext: (#a:Type) -> (#b:(a->Type)) -> ...|.
It has no computational content---the proof term carries no
reduction behavior.  In \synhopy{}, $\mathsf{funext}$ computes: given pointwise
paths, it literally constructs the path between functions via the interval.
\end{remark}


\section{Dependent Function Types ($\Pi$-Types)}

\begin{definition}[$\Pi$-type formation]
Given $A : \Type$ and $B : A \to \Type$, the dependent function type is:
\[
  \frac{\Gamma \vdash A \;\Type \qquad \Gamma, x : A \vdash B(x) \;\Type}
       {\Gamma \vdash \prod_{x : A} B(x) \;\Type}
\]
with introduction and elimination:
\begin{align}
  \frac{\Gamma, x : A \vdash b(x) : B(x)}
       {\Gamma \vdash \lambda x.\; b(x) : \prod_{x : A} B(x)}
  &&
  \frac{\Gamma \vdash f : \prod_{x : A} B(x) \qquad \Gamma \vdash a : A}
       {\Gamma \vdash f(a) : B(a)}
\end{align}
with $\beta$-rule $(\lambda x.\;b)(a) \equiv b[x := a]$ and
$\eta$-rule $f \equiv \lambda x.\;f(x)$.
\end{definition}

In \synhopy{}, $\Pi$-types model Python callables whose return type depends on
the argument:

\begin{example}[$\Pi$-type for \texttt{type()} built-in]
The Python built-in \lstinline[style=python]{type(x)} returns the class of
\lstinline[style=python]{x}.  Its \synhopy{} type is:
\[
  \mathsf{type} : \prod_{x : \PyObj} \PyClass
\]
More precisely, if we refine to know $x : C$ for some class $C$, the return
type narrows:
\[
  \mathsf{type} : \prod_{C : \PyClass}\prod_{x : C} \{M : \PyClass \mid \mathsf{isinstance}(x, M)\}
\]
This dependent typing captures the runtime behavior: \lstinline[style=python]{type(42)}
returns \lstinline[style=python]{int}, \lstinline[style=python]{type("hi")}
returns \lstinline[style=python]{str}.
\end{example}

\begin{remark}[\fstar{} comparison]
\fstar{} also has $\Pi$-types, written \lstinline[style=fstar]|x:a -> b x|.
The key difference is that \fstar{} $\Pi$-types require the domain and codomain
to be statically known types.  In \synhopy{}, $\Pi$-types can range over
$\PyObj$ (all Python objects), which includes values whose type is only known
at runtime.  This is essential for modeling \lstinline[style=python]{getattr},
\lstinline[style=python]{type()}, and other reflective operations.
\end{remark}

\section{Dependent Pair Types ($\Sigma$-Types)}

\begin{definition}[$\Sigma$-type formation]
Given $A : \Type$ and $B : A \to \Type$, the dependent pair type is:
\[
  \frac{\Gamma \vdash A \;\Type \qquad \Gamma, x : A \vdash B(x) \;\Type}
       {\Gamma \vdash \sum_{x : A} B(x) \;\Type}
\]
with introduction and elimination:
\begin{align}
  \frac{\Gamma \vdash a : A \qquad \Gamma \vdash b : B(a)}
       {\Gamma \vdash (a, b) : \sum_{x : A} B(x)}
  &&
  \frac{\Gamma \vdash p : \sum_{x : A} B(x)}
       {\Gamma \vdash \pi_1(p) : A \qquad \Gamma \vdash \pi_2(p) : B(\pi_1(p))}
\end{align}
\end{definition}

$\Sigma$-types are pervasive in \synhopy{} because Python's data structures
are naturally dependent pairs:

\begin{example}[$\Sigma$-type for typed dictionaries]
A Python \lstinline[style=python]{TypedDict} with schema $S$ is:
\[
  \mathsf{TypedDict}_S \defeq \sum_{k : \mathsf{Keys}(S)} S(k)
\]
For example, \lstinline[style=python]{class User(TypedDict): name: str; age: int}
becomes:
\[
  \mathsf{User} \defeq \sum_{k : \{\text{``name''}, \text{``age''}\}}
    \begin{cases}
      \mathsf{str} & k = \text{``name''} \\
      \mathsf{int} & k = \text{``age''}
    \end{cases}
\]
The value type depends on the key---this is a genuine dependent pair.
\end{example}

\begin{example}[$\Sigma$-type for \texttt{isinstance} refinement]
After a type guard \lstinline[style=python]{if isinstance(x, int)}, the
variable $x$ is refined.  The refined context is:
\[
  x : \sum_{v : \PyObj} \mathsf{isinstance}(v, \mathsf{int})
\]
The second component is a \emph{proof} that $v$ is an instance of
\lstinline[style=python]{int}.  This is a refinement type, modeled as a
$\Sigma$-type where the second component is propositional.
\end{example}


\section{The Interval Type $\interval$ and Cubical Machinery}

\subsection{The Interval}

\begin{definition}[The interval type]
The interval $\interval$ is a type with:
\begin{enumerate}
  \item Two endpoints $0_\interval, 1_\interval : \interval$
  \item Operations: $\land$ (min), $\lor$ (max), $1 - (-)$ (reversal)
  \item Laws making $(\interval, \land, \lor, 0, 1)$ a De Morgan algebra:
    \begin{align}
      i \land (j \lor k) &\equiv (i \land j) \lor (i \land k) \\
      1 - (i \land j) &\equiv (1 - i) \lor (1 - j) \\
      1 - (1 - i) &\equiv i \\
      i \lor 0 &\equiv i \qquad i \land 1 \equiv i
    \end{align}
\end{enumerate}
\end{definition}

The interval is \emph{not} a type in the object language---one cannot form
$\prod_{i:\interval} A(i)$ as a dependent function type.  Instead, interval
variables appear only in contexts, and types may depend on them.

\begin{example}[Interval for monkey-patching]
A monkey-patch that replaces method $m$ on class $C$ with implementation $v$
is a path in $\PyClass$ parameterized by $i : \interval$:
\[
  \mathsf{patch}(C, m, v)(i) \defeq
  \begin{cases}
    C & \text{if } i = 0 \\
    C[m := v] & \text{if } i = 1
  \end{cases}
\]
The interval smoothly interpolates between the original and patched class.
At $i = 0$ we have the original; at $i = 1$, the patched version.
\end{example}

\subsection{Face Lattice and Partial Elements}

\begin{definition}[Face formula]
A \emph{face formula} $\varphi$ is a formula built from:
\[
  \varphi ::= (i = 0) \mid (i = 1) \mid \varphi_1 \land \varphi_2 \mid \varphi_1 \lor \varphi_2
\]
where $i$ ranges over interval variables.  Face formulas form a distributive
lattice $\mathsf{Face}$, with $0_\mathsf{F}$ (empty face) and $1_\mathsf{F}$
(full face).
\end{definition}

\begin{definition}[Partial element]
A \emph{partial element} of type $A$ with extent $\varphi$ is:
\[
  \Gamma, \varphi \vdash u : A
\]
written $u : A[\varphi]$.  This is a term of type $A$ that is only defined on
the face $\varphi$.
\end{definition}

Partial elements are crucial for expressing \emph{compatibility conditions}
when gluing data.

\begin{example}[Partial elements in duck typing]
Suppose class $A$ implements methods $\{m_1, m_2, m_3\}$ and class $B$ implements
$\{m_2, m_3, m_4\}$.  On the face where we restrict to the shared interface
$\{m_2, m_3\}$ (face formula $\varphi \defeq (\text{method} = m_2) \lor (\text{method} = m_3)$),
we have a partial equivalence $e[\varphi] : A|_\varphi \simeq B|_\varphi$.
Duck-type compatibility is exactly the assertion that such partial equivalences
exist for all required methods.
\end{example}

\subsection{Composition and Filling}

\begin{definition}[Composition --- full form]
\label{def:comp-full}
The composition operation has the following complete typing rule:
\[
  \frac{
    \begin{array}{l}
      \Gamma, i : \interval \vdash A(i) \;\Type \\
      \varphi : \mathsf{Face} \\
      \Gamma, \varphi, i : \interval \vdash u(i) : A(i) \\
      \Gamma \vdash u_0 : A(0)[\varphi \mapsto u(0)]
    \end{array}
  }{
    \Gamma \vdash \comp^{i.\,A}[\varphi \mapsto u]\; u_0 : A(1)
  }
\]
satisfying the boundary condition:
$\comp^{i.\,A}[\varphi \mapsto u]\; u_0 = u(1)$ when $\varphi = 1_\mathsf{F}$.
\end{definition}

\begin{definition}[Filling]
\emph{Filling} is the internalization of composition:
\[
  \hfilling^{i.\,A}[\varphi \mapsto u]\; u_0 : \interval \to A(j)
\]
such that:
\begin{align}
  \hfilling^{i.\,A}[\varphi \mapsto u]\; u_0 \; (0) &\equiv u_0 \\
  \hfilling^{i.\,A}[\varphi \mapsto u]\; u_0 \; (1) &\equiv \comp^{i.\,A}[\varphi \mapsto u]\; u_0
\end{align}
Filling provides the path \emph{to} the result of composition.
\end{definition}

\begin{example}[Filling for gradual migration]
\label{ex:filling-migration}
Consider migrating a codebase from class $A$ (e.g., \lstinline[style=python]{urllib2})
to class $B$ (e.g., \lstinline[style=python]{requests}).  The duck-equivalence
path $p : A =_{\PyClass} B$ provides a path in $\PyClass$.  Filling along $p$
gives an \emph{intermediate class} at every point $j \in \interval$:
\[
  C_j \defeq \hfilling^{i.\,\PyClass}[\cdot]\; A \;(j) : \PyClass
\]
At $j = 0$, $C_0 \equiv A$; at $j = 1$, $C_1 \equiv B$.  The intermediate
values $C_j$ for $0 < j < 1$ represent hybrid states---the filling ensures
coherence throughout the migration.
\end{example}


\subsection{Transport Along Paths --- Python-Specific Rules}

Transport was introduced in \S2.3.5; here we give detailed computation rules
for Python-specific type families.

\begin{proposition}[Transport for class attributes]
\label{prop:transport-attr}
Given a path $p : C_1 =_{\PyClass} C_2$ and an attribute access pattern
$P(C) \defeq \{x : C \mid \mathsf{hasattr}(x, \text{``}m\text{''}) \}$, transport gives:
\[
  \transp^P(p, x) = \mathsf{coerce}(x, C_1, C_2, p)
\]
where $\mathsf{coerce}$ preserves all attributes in the shared interface and maps
method calls through the duck-equivalence.
\end{proposition}

\begin{proposition}[Transport for exception types]
\label{prop:transport-exc}
Given a path $p : E_1 =_{\PyClass} E_2$ between exception classes, transport
along $p$ maps:
\begin{enumerate}
  \item Handlers: $\transp^{\mathsf{Handler}}(p, h_1) : \mathsf{Handler}(E_2)$
  \item Raise sites: $\transp^{\mathsf{Raise}}(p, r_1) : \mathsf{Raise}(E_2)$
  \item Traceback information is preserved
\end{enumerate}
This formalizes the common Python pattern of replacing one exception class with
another (e.g., migrating from \lstinline[style=python]{IOError} to
\lstinline[style=python]{OSError}).
\end{proposition}

\begin{remark}[\fstar{} comparison]
In \fstar{}, coercing between types requires writing an explicit coercion function
and proving that it preserves relevant properties.  For example, migrating from
one record type to another requires:
\begin{lstlisting}[style=fstar]
val coerce_record: old_type -> Lemma (new_type)
let coerce_record x = { field1 = x.field1; field2 = x.field2; ... }
\end{lstlisting}
In \synhopy{}, transport along the duck-equivalence path provides this coercion
automatically, and its correctness is guaranteed by the cubical structure.
\end{remark}


\section{Cumulative Universe Hierarchy}

\begin{definition}[Universe hierarchy]
\label{def:universes}
\synhopy{} has a cumulative hierarchy of universes:
\[
  \mathcal{U}_0 : \mathcal{U}_1 : \mathcal{U}_2 : \cdots
\]
with the cumulativity rule: if $A : \mathcal{U}_n$, then $A : \mathcal{U}_{n+1}$.
The universes are closed under all type formers:
\begin{enumerate}
  \item If $A : \mathcal{U}_n$ and $B : A \to \mathcal{U}_n$, then
    $\prod_{x:A} B(x) : \mathcal{U}_n$ and $\sum_{x:A} B(x) : \mathcal{U}_n$.
  \item If $A : \mathcal{U}_n$ and $a, b : A$, then $(a =_A b) : \mathcal{U}_n$.
  \item $\PyObj : \mathcal{U}_1$ (Python objects live in the first universe).
  \item $\PyClass : \mathcal{U}_2$ (classes of Python objects live one level higher).
\end{enumerate}
\end{definition}

\begin{definition}[Universe level assignment for Python]
The standard universe levels for Python constructs are:
\begin{align}
  \mathsf{int}, \mathsf{str}, \mathsf{bool}, \mathsf{float} &: \mathcal{U}_0
    && \text{(base types)} \\
  \mathsf{list}(A), \mathsf{dict}(K, V) &: \mathcal{U}_{\max(0, \ell(A))}
    && \text{(container types)} \\
  \PyObj &: \mathcal{U}_1
    && \text{(the universe of all objects)} \\
  \PyClass &: \mathcal{U}_2
    && \text{(classes are types of objects)} \\
  \mathsf{type} &: \mathcal{U}_3
    && \text{(the metaclass of all classes)}
\end{align}
where $\ell(A)$ denotes the universe level of $A$.
\end{definition}

\begin{remark}[\fstar{} comparison]
\fstar{} also has a universe hierarchy with
$\mathsf{Type}\;u_0$, $\mathsf{Type}\;u_1$, etc., and universe polymorphism
via \lstinline[style=fstar]|val id: #a:Type u#i -> a -> a|.
\synhopy{}'s hierarchy differs in two ways:
\begin{enumerate}
  \item \synhopy{} places $\PyObj$ at $\mathcal{U}_1$, not $\mathcal{U}_0$,
    because Python objects can contain references to their own type
    (via \lstinline[style=python]{__class__}), creating a universe
    polymorphism issue that \fstar{} avoids by forbidding such self-reference.
  \item \synhopy{} allows paths between universes (via univalence), while
    \fstar{}'s universes are strictly stratified with no inter-level paths.
\end{enumerate}
\end{remark}

\begin{proposition}[Universe consistency]
\label{prop:universe-consistent}
The \synhopy{} universe hierarchy is consistent: there is no $A$ with
$A : A$.
\end{proposition}
\begin{proof}
By the strict ordering $\mathcal{U}_n : \mathcal{U}_{n+1}$ with no
$\mathcal{U}_n : \mathcal{U}_n$.  Cumulativity only raises levels; it
cannot create cycles.  This is the standard Girard's paradox avoidance.
\end{proof}


\section{Python Type Formers: Detailed Mapping}

We now give a comprehensive table mapping Python's built-in types and constructs
to their \synhopy{} type-theoretic counterparts.

\begin{definition}[Optional types]
Python's \lstinline[style=python]{Optional[A]} is the coproduct:
\[
  \mathsf{Optional}(A) \defeq A + \mathsf{NoneType}
\]
with injection $\mathsf{Some} : A \to A + \mathsf{NoneType}$ and
$\PyNone : \mathsf{NoneType} \hookrightarrow A + \mathsf{NoneType}$.
The eliminator (pattern matching) is:
\[
  \mathsf{match} : (A + \mathsf{NoneType}) \to (A \to C) \to (\mathsf{NoneType} \to C) \to C
\]
\end{definition}

\begin{definition}[Union types]
Python's \lstinline[style=python]{Union[A, B]} is the pushout over the
empty type:
\[
  \mathsf{Union}(A, B) \defeq A +_\emptyset B \simeq A + B
\]
For non-disjoint unions $A \cap B \neq \emptyset$ (e.g., subclass relationships),
the pushout is:
\[
  \mathsf{Union}(A, B) \defeq A +_{A \cap B} B
\]
with identification of the shared subtype.
\end{definition}

\begin{definition}[Callable types]
Python's \lstinline[style=python]{Callable[[A1,...,An], R]} is an effect-indexed
$\Pi$-type:
\[
  \mathsf{Callable}([A_1, \ldots, A_n], R) \defeq
  (x_1 : A_1) \to_\varepsilon (x_2 : A_2) \to_\varepsilon \cdots \to_\varepsilon R
\]
where $\varepsilon$ is inferred by the effect analyzer.  If the return type
depends on the arguments (e.g., \lstinline[style=python]{type()} or
\lstinline[style=python]{getattr()}), this becomes a genuine dependent function type.
\end{definition}

\begin{definition}[Protocol types]
A Python \lstinline[style=python]{Protocol} is a face of $\PyClass$:
\[
  \mathsf{Protocol}(\{m_1 : T_1, \ldots, m_k : T_k\}) \defeq
  \sum_{C : \PyClass} \prod_{j=1}^{k} \mathsf{hasattr}(C, m_j) \times
    (C.m_j : T_j)
\]
A class $C$ satisfies a protocol if and only if there is a term witnessing
membership in this $\Sigma$-type---i.e., a proof that $C$ has all the required
methods with the right types.
\end{definition}

\begin{remark}[\fstar{} comparison]
\fstar{} has no built-in protocol/structural typing mechanism.  The closest
approximation is typeclasses, which require explicit instance declarations:
\begin{lstlisting}[style=fstar]
(* F*: explicit typeclass *)
noeq type printable a = { to_string: a -> string }
(* Must declare instance for each type *)
instance printable_int : printable int = { to_string = string_of_int }
\end{lstlisting}
In \synhopy{}, any class that has a \lstinline[style=python]{__str__} method
automatically satisfies the \lstinline[style=python]{Printable} protocol---no
declaration needed.  This matches Python's runtime behavior exactly.
\end{remark}


% ──────────────────────────────────────────────────────────────────
\chapter{The Effect System}
\label{ch:effects}

\section{Effect Lattice}

\synhopy{}'s effects form a lattice ordered by impurity:

\begin{definition}[Effect lattice]
\[
  \Pure \;\leq\; \Reads \;\leq\; \Mutates \;\leq\; \IO \;\leq\; \Nondet
\]
with join $\varepsilon_1 \vee \varepsilon_2 = \max(\varepsilon_1, \varepsilon_2)$
and meet $\varepsilon_1 \wedge \varepsilon_2 = \min(\varepsilon_1, \varepsilon_2)$.
\end{definition}

\begin{definition}[Effect typing rule]
\[
  \frac{
    \Gamma \vdash f : (x : A) \to_{\varepsilon_1} B \qquad
    \Gamma \vdash a : A \qquad
    \varepsilon_1 \leq \varepsilon_2
  }{
    \Gamma \vdash_{\varepsilon_2} f(a) : B[x := a]
  }
\]
A function with effect $\varepsilon_1$ can be used in any context that allows
effect $\varepsilon_2 \geq \varepsilon_1$ (effect subsumption).
\end{definition}

\subsection{Comparison: \fstar{}'s Dijkstra Monads}

\fstar{} models effects using Dijkstra monads (Ahman et al., 2017), where each
effect is a monad indexed by a wp-predicate transformer:

\begin{lstlisting}[style=fstar]
(* F*'s effect system *)
effect STATE (a:Type) (wp: st_wp a) =
  ST a (fun p s0 -> wp (fun x s1 -> p x s1) s0)
\end{lstlisting}

The wp specification must be written manually for every function. For Python,
this is unworkable:

\begin{enumerate}
  \item Python's heap is untyped---\lstinline[style=fstar]{ref a} doesn't capture
    \lstinline[style=python]{obj.__dict__} semantics.
  \item Python's effects are emergent (a function's effects depend on its arguments'
    types, which are only known at runtime), not declared.
  \item \fstar{}'s wp composition is syntactic, not semantic---it doesn't capture
    interactions between effect layers.
\end{enumerate}


\section{Effects as $\Omega$-Endomorphisms}

In the $\infty$-topos $\Sh_\infty(\Ctx)$, the subobject classifier $\Omega$
classifies ``truth values'' at each execution context. An effect is an endomorphism
$\varepsilon : \Omega \to \Omega$ that modifies what is ``true'' after the
effectful computation.

\begin{definition}[Effect as endomorphism]
An effect $\varepsilon$ is an endomorphism $\varepsilon : \Omega \to \Omega$ where:
\begin{itemize}
  \item $\Pure = \mathsf{id}_\Omega$ (truth is unchanged)
  \item $\Reads = \varepsilon$ with $\varepsilon(P) \Rightarrow P$ (reading preserves truth)
  \item $\Mutates = \varepsilon$ with no constraint (truth may change)
  \item $\IO = \varepsilon$ that depends on external state
  \item $\Nondet = \varepsilon$ that is not deterministic (maps propositions to disjunctions)
\end{itemize}
\end{definition}

The effect lattice ordering corresponds to logical strength: $\varepsilon_1 \leq \varepsilon_2$
means that $\varepsilon_2$ can change at least as many propositions as $\varepsilon_1$.

\section{The Blakers--Massey Theorem for Effect Composition}

When composing two effectful computations, how much information is preserved?
The Blakers--Massey theorem gives a precise answer:

\begin{theorem}[Blakers--Massey for effects]
\label{thm:blakers-massey}
If computation $f$ has effect of connectivity $m$ (preserves propositions up to
level $m$) and computation $g$ has effect of connectivity $n$, then their
composition $g \circ f$ has effect of connectivity at least $m + n + 1$.
\end{theorem}

This means: composing a ``slightly impure'' function with another ``slightly impure''
function gives a result that is ``more impure than either, but less impure than
the worst case.'' \fstar{}'s effect lattice simply takes the maximum, losing this
fine-grained information.


\section{Spectral Sequences for Multi-Effect Analysis}
\label{sec:spectral}

When a computation involves multiple interacting effects, \synhopy{} decomposes
the analysis using spectral sequences.

\begin{definition}[Effect spectral sequence]
For a computation with effects $\varepsilon_1, \ldots, \varepsilon_k$, define:
\begin{itemize}
  \item $E_0^{p,q}$: the raw operation counts (reads, writes, calls) at position $(p,q)$
  \item $d_0 : E_0^{p,q} \to E_0^{p+1,q}$: the differential that connects
    adjacent operations
  \item $E_1^{p,q} = H^q(E_0^{p,*}, d_0)$: the effect types after grouping operations
  \item $d_1 : E_1^{p,q} \to E_1^{p+1,q}$: the differential capturing effect interactions
  \item $E_2^{p,q} = H^q(E_1^{p,*}, d_1)$: the residual after accounting for interactions
\end{itemize}
The spectral sequence converges: $E_\infty^{p,q}$ gives the ``true'' effect
decomposition.
\end{definition}

\begin{example}[Decorator commutativity]
Consider \lstinline[style=python]{@log @cache def f(x): return x**2}.
\begin{itemize}
  \item $E_1$: \texttt{log} has effect $\IO$, \texttt{cache} has effect $\Mutates$.
  \item $d_1$: does \texttt{log} interact with \texttt{cache}? It reads the
    cache-decorated function's return value and logs it. The differential
    $d_1 : E_1^{0,\IO} \to E_1^{1,\Mutates}$ is zero if logging doesn't
    depend on cache state.
  \item $E_2 = E_1$ (since $d_1 = 0$), so the decorators commute.
\end{itemize}
In \fstar{}, proving this requires a manual proof for each pair of decorators.
In \synhopy{}, the spectral sequence determines it automatically.
\end{example}


\section{Effect Inference Algorithm}

\synhopy{} infers effects automatically from Python source code, without
annotations.  The inference is conservative: if uncertain, it over-approximates
(e.g., reporting $\IO$ when the function might be $\Pure$).

\begin{definition}[Effect inference rules]
Given a function body AST, the inferred effect is:
\begin{align}
  \mathsf{eff}(\texttt{return } e) &= \mathsf{eff}(e) \\
  \mathsf{eff}(x = e) &= \mathsf{eff}(e) \vee
    \begin{cases}
      \Pure & \text{if $x$ is local} \\
      \Mutates & \text{if $x$ is \texttt{self.\_}} \\
      \Mutates & \text{if $x$ is global}
    \end{cases} \\
  \mathsf{eff}(\texttt{print}(\ldots)) &= \IO \\
  \mathsf{eff}(\texttt{open}(\ldots)) &= \IO \\
  \mathsf{eff}(\texttt{random.}\ldots) &= \Nondet \\
  \mathsf{eff}(e_1\texttt{.append}(e_2)) &= \Mutates \vee \mathsf{eff}(e_1) \vee \mathsf{eff}(e_2) \\
  \mathsf{eff}(f(e)) &= \mathsf{eff}(f) \vee \mathsf{eff}(e) && \text{(if $f$'s effect is known)} \\
  \mathsf{eff}(f(e)) &= \IO && \text{(if $f$'s effect is unknown---conservative)}
\end{align}
\end{definition}

The inference traverses the AST bottom-up, accumulating effects with $\vee$
(join in the effect lattice).  This is implemented in deppy's
\texttt{EffectAnalyzer} class.

\begin{theorem}[Effect inference soundness]
If the inferred effect is $\varepsilon$, then the actual effect is $\leq \varepsilon$.
\end{theorem}
\begin{proof}
By structural induction on the AST.  Each inference rule assigns an effect
that is $\geq$ the actual effect of the construct (conservative over-approximation).
The join operation preserves this: if $\varepsilon_1 \geq \varepsilon_1'$ and
$\varepsilon_2 \geq \varepsilon_2'$, then
$\varepsilon_1 \vee \varepsilon_2 \geq \varepsilon_1' \vee \varepsilon_2'$.
\end{proof}

\begin{remark}[\fstar{} comparison]
\fstar{} requires \emph{manual} effect annotations on every function.
The programmer must write the effect in the type signature:
\begin{lstlisting}[style=fstar]
val read_file: path:string -> ST string (requires ...)
\end{lstlisting}
\synhopy{} infers effects automatically.  The programmer only needs
\lstinline[style=synhopy]{@effect_contract} when the inference is too
conservative (i.e., the programmer knows the function is purer than what
inference detects).
\end{remark}


\section{Equational Reasoning Under Effects}

A key question: when can we use equational reasoning ($f(x) = f(x)$)?

\begin{theorem}[Equational reasoning criterion]
\label{thm:equational}
Equational reasoning ($f(x) = f(x)$ for all $x$) is valid if and only if
$\mathsf{eff}(f) \leq \Reads$ (the function is pure or only reads state).
\end{theorem}
\begin{proof}
If $\mathsf{eff}(f) \leq \Reads$, then $f$ does not modify state, so calling
$f$ twice with the same argument and same state produces the same result.

If $\mathsf{eff}(f) = \Mutates$, then calling $f$ twice may produce different
results because the first call modifies state that the second call reads.

If $\mathsf{eff}(f) = \Nondet$, then $f(x) \neq f(x)$ is possible even without
state changes.
\end{proof}

This theorem is critical for proof correctness: if the LLM uses equational
reasoning in a proof about a $\Mutates$ function, the kernel detects the
effect mismatch and rejects the proof.


\section{The Full Effect Lattice}

The five-element chain $\Pure \leq \Reads \leq \Mutates \leq \IO \leq \Nondet$
is the \emph{principal spine} of the effect lattice.  The full lattice
includes additional effects and cross-cutting concerns.

\begin{definition}[Extended effect lattice]
The full \synhopy{} effect lattice $(\mathcal{E}, \leq)$ has the following
Hasse diagram structure:
\[
  \begin{array}{ccccc}
    & & \Nondet & & \\
    & \nearrow & & \nwarrow & \\
    \IO & & & & \mathsf{Diverge} \\
    & \nwarrow & & \nearrow & \\
    & & \Mutates & & \\
    & & | & & \\
    & & \Reads & & \\
    & & | & & \\
    & & \Pure & &
  \end{array}
\]
where $\mathsf{Diverge}$ represents potential non-termination (e.g.,
\lstinline[style=python]{while True: ...}), and the two branches rejoin at
$\Nondet$ since non-determinism subsumes both I/O and divergence.
\end{definition}

\begin{definition}[Derived effects]
Several derived effects capture common Python patterns:
\begin{align}
  \mathsf{Alloc} &\defeq \Mutates \quad\text{(object creation)} \\
  \mathsf{Exn} &\defeq \Reads \vee \mathsf{Diverge}
    \quad\text{(may raise, a partial effect)} \\
  \mathsf{AsyncIO} &\defeq \IO \vee \mathsf{Diverge}
    \quad\text{(async operations may suspend indefinitely)}
\end{align}
\end{definition}


\section{Effect Polymorphism}

Python functions are often \emph{effect-polymorphic}: their effect depends on
their arguments.

\begin{definition}[Effect-polymorphic type]
An effect-polymorphic function has type:
\[
  f : \prod_{\varepsilon : \mathcal{E}} (x : A) \to_\varepsilon B(x)
\]
The effect parameter $\varepsilon$ is universally quantified: $f$ works at any
effect level.
\end{definition}

\begin{example}[Effect polymorphism in \texttt{map}]
The built-in \lstinline[style=python]{map(f, xs)} has type:
\[
  \mathsf{map} : \prod_{\varepsilon : \mathcal{E}}
    \left((A \to_\varepsilon B) \to \mathsf{Iterable}(A) \to_\varepsilon \mathsf{Iterator}(B)\right)
\]
If $f$ is $\Pure$, then $\mathsf{map}(f, xs)$ is $\Pure$; if $f$ has effect
$\IO$, then so does the mapped result.  The effect of \lstinline[style=python]{map}
\emph{inherits} the effect of its function argument.
\end{example}

\begin{example}[Effect polymorphism in \texttt{sorted}]
\lstinline[style=python]{sorted(xs, key=f)} is effect-polymorphic in its
\lstinline[style=python]{key} argument:
\[
  \mathsf{sorted} : \prod_{\varepsilon : \mathcal{E}}
    (\mathsf{Iterable}(A) \to (A \to_\varepsilon \mathsf{Comparable}) \to_\varepsilon
      \mathsf{list}(A))
\]
With a $\Pure$ key function, sorting is $\Pure$; with a $\Reads$ key (e.g.,
reading a global config for sort order), sorting becomes $\Reads$.
\end{example}

\begin{remark}[\fstar{} comparison]
\fstar{} supports effect polymorphism via its effect system, but requires
explicit effect annotations.  For instance:
\begin{lstlisting}[style=fstar]
val map: #a:Type -> #b:Type -> (a -> Tot b) -> list a -> Tot (list b)
val map_st: #a:Type -> #b:Type -> (a -> ST b pre post) -> list a
  -> ST (list b) pre' post'
\end{lstlisting}
Two separate definitions are needed for pure and stateful mapping.  \synhopy{}
unifies these with a single effect-polymorphic type.
\end{remark}


\section{Effect Masking}

Sometimes a function with apparent effects can be treated as pure.

\begin{definition}[Effect masking]
An effectful computation $e : A$ with effect $\varepsilon$ can be
\emph{masked} to effect $\varepsilon' < \varepsilon$ if the effects are
\emph{observationally invisible}:
\[
  \frac{
    \Gamma \vdash_\varepsilon e : A \qquad
    \mathsf{observable}(\varepsilon, e) \leq \varepsilon'
  }{
    \Gamma \vdash_{\varepsilon'} \mathsf{mask}_{\varepsilon'}(e) : A
  }
\]
\end{definition}

\begin{example}[Masking internal mutation]
A function that creates a local list, mutates it, and returns the result is
internally $\Mutates$ but externally $\Pure$:
\begin{lstlisting}[style=python]
def build_squares(n: int) -> list[int]:
    result = []                # Alloc
    for i in range(n):
        result.append(i * i)   # Mutates (local)
    return result              # Pure externally
\end{lstlisting}
The mutation is confined to the local variable \lstinline[style=python]{result},
which is not aliased.  \synhopy{} detects this via escape analysis and masks
$\Mutates$ down to $\Pure$.
\end{example}

\begin{theorem}[Masking soundness]
If $\mathsf{mask}_{\varepsilon'}(e)$ is well-typed, then $e$ is
observationally indistinguishable from a computation of effect $\varepsilon'$.
\end{theorem}
\begin{proof}
By the definition of $\mathsf{observable}$: the only mutations are to
locally-allocated, non-aliased state, which is invisible after the function
returns.  The heap before and after the call are identical modulo garbage
collection, which is unobservable.
\end{proof}


\section{Effect Handlers: Algebraic Effects Perspective}

\synhopy{}'s effect system can be understood through the lens of algebraic
effects and handlers (Plotkin \& Pretnar, 2009).

\begin{definition}[Effect signature]
An effect signature $\Sigma_\varepsilon$ for effect $\varepsilon$ is a set of
\emph{operations}, each with a parameter and return type:
\begin{align}
  \Sigma_\Reads &= \{ \mathsf{read\_attr} : (\PyObj \times \mathsf{String}) \to \PyObj \} \\
  \Sigma_\Mutates &= \Sigma_\Reads \cup \{ \mathsf{write\_attr} : (\PyObj \times \mathsf{String} \times \PyObj) \to \mathsf{Unit} \} \\
  \Sigma_\IO &= \Sigma_\Mutates \cup \{ \mathsf{file\_op} : \mathsf{FileReq} \to \mathsf{FileResp},\;
    \mathsf{net\_op} : \mathsf{NetReq} \to \mathsf{NetResp} \} \\
  \Sigma_\Nondet &= \Sigma_\IO \cup \{ \mathsf{choose} : \Nat \to \Nat \}
\end{align}
\end{definition}

\begin{definition}[Effect handler]
A handler for effect $\varepsilon$ is a record of continuations, one per
operation:
\[
  \mathsf{handler}_\varepsilon : \prod_{(\mathsf{op} : P \to R) \in \Sigma_\varepsilon}
    (P \to (R \to A) \to A)
\]
The handler intercepts each operation and provides a return value to the
continuation.
\end{definition}

\begin{example}[Python's \texttt{contextvar} as effect handler]
Python's \lstinline[style=python]{contextvars.ContextVar} is an effect handler
for the $\Reads$ effect.  A context variable \lstinline[style=python]{cv = ContextVar('cv', default=42)}
defines a handler:
\[
  \mathsf{handler}_{\mathsf{cv}} = \{
    \mathsf{read\_attr}(\_\_, \text{``cv''}) \mapsto \lambda k.\; k(42)
  \}
\]
The \lstinline[style=python]{cv.set(v)} operation installs a new handler that
returns $v$ instead:
$\mathsf{handler}_{\mathsf{cv}}' = \{ \mathsf{read\_attr}(\_\_, \text{``cv''}) \mapsto \lambda k.\; k(v) \}$.
This makes \lstinline[style=python]{ContextVar} a first-class algebraic effect
handler in \synhopy{}'s formalism.
\end{example}


\section{Effect Inference: Worked Example}

We demonstrate effect inference on a realistic Python function.

\begin{example}[Effect inference walkthrough]
\label{ex:effect-inference-worked}
Consider the following function:
\begin{lstlisting}[style=python]
import logging
from typing import Optional

logger = logging.getLogger(__name__)

def process_user_data(
    user_id: int,
    db: Database,
    cache: dict[int, UserRecord],
) -> Optional[UserRecord]:
    if user_id in cache:             # Reads (cache)
        logger.debug("cache hit")    # IO (logging)
        return cache[user_id]        # Reads (cache)

    record = db.fetch(user_id)       # IO (database)
    if record is None:
        return None                  # Pure

    record.last_accessed = time()    # Mutates (record)
    cache[user_id] = record          # Mutates (cache)
    logger.info("fetched user %d", user_id)  # IO
    return record                    # Pure
\end{lstlisting}

The inference proceeds bottom-up, computing $\mathsf{eff}$ for each statement:
\begin{align}
  \mathsf{eff}(\text{line 11: \texttt{user\_id in cache}}) &= \Reads \\
  \mathsf{eff}(\text{line 12: \texttt{logger.debug(...)}}) &= \IO \\
  \mathsf{eff}(\text{line 13: \texttt{return cache[user\_id]}}) &= \Reads \\
  \mathsf{eff}(\text{lines 11--13: if branch}) &= \Reads \vee \IO \vee \Reads = \IO \\
  \mathsf{eff}(\text{line 15: \texttt{db.fetch(user\_id)}}) &= \IO \\
  \mathsf{eff}(\text{line 17: \texttt{return None}}) &= \Pure \\
  \mathsf{eff}(\text{line 19: \texttt{record.last\_accessed = ...}}) &= \Mutates \\
  \mathsf{eff}(\text{line 20: \texttt{cache[user\_id] = record}}) &= \Mutates \\
  \mathsf{eff}(\text{line 21: \texttt{logger.info(...)}}) &= \IO \\
  \mathsf{eff}(\text{line 22: \texttt{return record}}) &= \Pure \\
  \mathsf{eff}(\text{lines 15--22: else branch}) &= \IO \vee \Pure \vee \Mutates \vee \IO = \IO \\
  \mathsf{eff}(\text{whole function}) &= \IO \vee \IO = \IO
\end{align}
The final inferred effect is $\IO$, which is correct: the function reads from
a cache, queries a database, mutates a record, and performs logging.
\end{example}


\section{The Dijkstra Monad Connection}

\fstar{}'s effect system is built on Dijkstra monads (Ahman et al., 2017;
Maillard et al., 2019).  We now make precise how \synhopy{}'s
$\Omega$-endomorphism model relates to this.

\begin{definition}[Dijkstra monad]
In \fstar{}, a Dijkstra monad for effect $\varepsilon$ is a monad
$D_\varepsilon : (\Type \to \Type)$ equipped with:
\begin{enumerate}
  \item A \emph{wp-type} $\mathsf{wp}_\varepsilon(A) = (A \to \Type_0) \to \Type_0$
  \item Operations $\mathsf{return}^D_\varepsilon$ and $\mathsf{bind}^D_\varepsilon$
    that compose weakest preconditions
\end{enumerate}
\end{definition}

\begin{proposition}[Dijkstra monads as $\Omega$-endomorphisms]
\label{prop:dijkstra-omega}
Every Dijkstra monad $(D_\varepsilon, \mathsf{wp}_\varepsilon)$ induces an
$\Omega$-endomorphism $\hat{\varepsilon} : \Omega \to \Omega$ via:
\[
  \hat{\varepsilon}(P) \defeq \mathsf{wp}_\varepsilon(\lambda a.\; P)
\]
Conversely, every $\Omega$-endomorphism $\varepsilon : \Omega \to \Omega$
that preserves finite meets (i.e., $\varepsilon(P \land Q) = \varepsilon(P) \land \varepsilon(Q)$)
arises from a Dijkstra monad.
\end{proposition}

This shows that \synhopy{}'s effect model is a \emph{generalization} of
\fstar{}'s: every \fstar{} effect can be expressed as an
$\Omega$-endomorphism, but \synhopy{} also admits effects that do not
preserve finite meets (e.g., $\Nondet$, where
$\mathsf{choose}(P \land Q) \neq \mathsf{choose}(P) \land \mathsf{choose}(Q)$
in general).


\section{Effect Composition Rules for Python Constructs}

We give the effect composition rule for every major Python control-flow
construct.

\begin{definition}[Effect composition rules]
\label{def:effect-composition}
\begin{align}
  \mathsf{eff}(\mathsf{if}\;e\!:\;s_1\;\mathsf{else}\!:\;s_2)
    &= \mathsf{eff}(e) \vee \mathsf{eff}(s_1) \vee \mathsf{eff}(s_2) \\
  \mathsf{eff}(\mathsf{for}\;x\;\mathsf{in}\;e\!:\;s)
    &= \mathsf{eff}(e) \vee \mathsf{eff}(s) \\
  \mathsf{eff}(\mathsf{while}\;e\!:\;s)
    &= \mathsf{eff}(e) \vee \mathsf{eff}(s) \vee \mathsf{Diverge} \\
  \mathsf{eff}(\mathsf{try}\!:\;s_1\;\mathsf{except}\;E\!:\;s_2)
    &= \mathsf{eff}(s_1) \vee \mathsf{eff}(s_2) \\
  \mathsf{eff}(\mathsf{try}\!:\;s_1\;\mathsf{finally}\!:\;s_2)
    &= \mathsf{eff}(s_1) \vee \mathsf{eff}(s_2) \\
  \mathsf{eff}(\mathsf{with}\;e\;\mathsf{as}\;x\!:\;s)
    &= \mathsf{eff}(e) \vee \mathsf{eff}(e.\mathsf{\_\_enter\_\_}())
       \vee \mathsf{eff}(s) \vee \mathsf{eff}(e.\mathsf{\_\_exit\_\_}(\ldots)) \\
  \mathsf{eff}(\mathsf{async\;with}\;e\;\mathsf{as}\;x\!:\;s)
    &= \mathsf{AsyncIO} \vee \mathsf{eff}(e) \vee \mathsf{eff}(s)
\end{align}
\end{definition}

\begin{remark}[While-loops and divergence]
The $\mathsf{Diverge}$ effect on while-loops is conservative: we cannot
statically determine termination in general (halting problem).  If the user
provides a termination measure via \lstinline[style=synhopy]{@guarantee("terminates")},
the $\mathsf{Diverge}$ annotation is removed after the measure is verified.
\end{remark}

\begin{remark}[\fstar{} comparison]
\fstar{} requires a \lstinline[style=fstar]{decreases} clause on every
recursive or looping function to prove termination.  Functions without a
decreasing measure cannot be typed in the $\mathsf{Tot}$ effect.  \synhopy{}
is more permissive: it allows $\mathsf{Diverge}$ as a valid effect, and
only requires termination proofs when the user wants to mask it away.
\end{remark}


\section{Spectral Sequences: A Concrete Computation}

We illustrate the spectral sequence machinery (Definition~\ref{sec:spectral})
with a concrete computation.

\begin{example}[Spectral sequence for a web request handler]
\label{ex:spectral-concrete}
Consider a web request handler that (1) reads from a database, (2) writes to
a cache, and (3) logs the result:
\begin{lstlisting}[style=python]
def handle_request(req: Request) -> Response:
    data = db.query(req.params)     # IO (database read)
    cache.set(req.key, data)        # Mutates (cache write)
    logger.info("handled %s", req)  # IO (logging)
    return Response(data)           # Pure
\end{lstlisting}

The $E_0$ page records raw operations at positions $(p, q)$ where $p$ is the
statement index and $q$ is the effect level:
\[
  E_0 = \begin{pmatrix}
    0 & 0 & 0 & 1 & 0 \\
    0 & 0 & 1 & 0 & 0 \\
    0 & 0 & 0 & 1 & 0 \\
    1 & 0 & 0 & 0 & 0
  \end{pmatrix}
  \quad
  \text{columns: } \Pure, \Reads, \Mutates, \IO, \Nondet
\]

The differential $d_0$ connects adjacent operations; since the cache write at
$p=1$ reads the database result from $p=0$, we have a non-trivial differential
$d_0 : E_0^{0,\IO} \to E_0^{1,\Mutates}$.

Computing $E_1 = H(E_0, d_0)$: the database-read and cache-write combine into a
single $\IO$-level operation (since $\IO \geq \Mutates$).  The $E_1$ page is:
\[
  E_1^{p,q} = \begin{cases}
    \mathbb{Z} & (p, q) = (0, \IO) \text{ (combined db+cache)} \\
    \mathbb{Z} & (p, q) = (1, \IO) \text{ (logging)} \\
    \mathbb{Z} & (p, q) = (2, \Pure) \text{ (return)} \\
    0 & \text{otherwise}
  \end{cases}
\]

The differential $d_1 : E_1^{0,\IO} \to E_1^{1,\IO}$ checks whether the
database operation and logging interact.  Since logging does not read from the
database result directly, $d_1 = 0$.

The spectral sequence collapses: $E_2 = E_\infty$, and the total effect is
$\IO$ (the maximum surviving entry).  Furthermore, the vanishing of $d_1$
tells us that the database operation and logging are \emph{independent}---they
could be reordered or parallelized without changing observable behavior.
\end{example}


% ──────────────────────────────────────────────────────────────────
\chapter{The Proof Kernel}
\label{ch:kernel}

\section{Proof Terms}

\synhopy{}'s proof kernel type-checks proof terms---tree-structured objects that
witness properties of Python programs. Every verified claim must be backed by a
proof term that the kernel accepts.

\begin{definition}[Proof terms]
\label{def:proof-terms}
The grammar of proof terms is:
\begin{align}
  \pi &::= \Refl_a && \text{reflexivity: $a =_A a$} \\
  &\mid \sym(\pi) && \text{symmetry: from $a =_A b$ get $b =_A a$} \\
  &\mid \Trans(\pi_1, \pi_2) && \text{transitivity: from $a = b, b = c$ get $a = c$} \\
  &\mid \Cong(f, \pi) && \text{congruence: from $a = b$ get $f(a) = f(b)$} \\
  &\mid \transp(P, \pi, d) && \text{transport: from $a = b$ and $d : P(a)$ get $P(b)$} \\
  &\mid \Ext(x, \pi) && \text{extensionality: $\forall x.\, f(x) = g(x) \Rightarrow f = g$} \\
  &\mid \Cut(h : A, \pi_1, \pi_2) && \text{cut: prove $A$, use it in subsequent proof} \\
  &\mid \Assume(h) && \text{assumption: use hypothesis $h$ from context} \\
  &\mid \mathsf{Z3}(\varphi) && \text{Z3 discharge: verify arithmetic/boolean formula} \\
  &\mid \mathsf{NatInd}(x, \pi_0, \pi_S) && \text{natural number induction} \\
  &\mid \mathsf{ListInd}(xs, \pi_{[]}, \pi_{::}) && \text{list induction} \\
  &\mid \mathsf{Cases}(e, \overline{p_i \Rightarrow \pi_i}) && \text{case analysis} \\
  &\mid \mathsf{DuckPath}(A, B, \overline{m_i : \pi_i}) && \text{duck-type path construction} \\
  &\mid \mathsf{Patch}(C, m, v, \pi) && \text{monkey-patch path} \\
  &\mid \mathsf{Fiber}(f, d, \overline{U_i : \pi_i}) && \text{fibration (isinstance) proof} \\
  &\mid \mathsf{EffectWitness}(\varepsilon, \pi) && \text{effect-indexed proof} \\
  &\mid \mathsf{Axiom}(\mathit{name}) && \text{trusted axiom (tracked)}
\end{align}
\end{definition}

\section{Kernel Typing Rules}

The kernel checks proof terms against their claimed types. We give the key rules:

\subsection{Reflexivity}

\[
  \frac{\Gamma \vdash a : A}{\Gamma \vdash \Refl_a : a =_A a}
\]

\subsection{Transitivity}

\[
  \frac{
    \Gamma \vdash \pi_1 : a =_A b \qquad
    \Gamma \vdash \pi_2 : b =_A c
  }{
    \Gamma \vdash \Trans(\pi_1, \pi_2) : a =_A c
  }
\]

The kernel must verify that the right endpoint of $\pi_1$ matches the left
endpoint of $\pi_2$. In implementation, this requires \emph{middle-term inference}
(see Section~\ref{sec:middle-term}).

\subsection{Cut (Lemma Introduction)}

\[
  \frac{
    \Gamma \vdash \pi_1 : A \qquad
    \Gamma, h : A \vdash \pi_2 : B
  }{
    \Gamma \vdash \Cut(h : A, \pi_1, \pi_2) : B
  }
\]

The $\Cut$ rule is the mechanism for introducing lemmas. The proof $\pi_1$
establishes the lemma $A$, which is then available as hypothesis $h$ in the
continuation $\pi_2$.

\subsection{Duck-Type Path Construction}

\[
  \frac{
    \forall m_i \in \MethodSig(A).\;
    \Gamma \vdash \pi_i : A.m_i =_{\PyFunc} B.m_i
  }{
    \Gamma \vdash \mathsf{DuckPath}(A, B, \overline{m_i : \pi_i}) : A =_{\PyClass} B
  }
\]

To construct a duck-type path, provide a proof of method equivalence for
\emph{every} method in $A$'s interface.

\subsection{Z3 Discharge}

\[
  \frac{
    \varphi \text{ is a quantifier-free formula} \qquad
    \text{Z3 reports } \varphi \text{ valid}
  }{
    \Gamma \vdash \mathsf{Z3}(\varphi) : \varphi
  }
\]

Arithmetic and boolean formulas are discharged to Z3. The trusted computing base
includes Z3's correctness---the same assumption made by \fstar{}.

\begin{remark}[\fstar{} comparison]
Both \synhopy{} and \fstar{} use Z3 for arithmetic reasoning. The difference is
\emph{scope}: \fstar{} sends essentially all proof obligations to Z3 (including
complex ones involving lists, functions, and recursive data), which often causes
timeouts. \synhopy{} uses Z3 only for quantifier-free arithmetic (where it is
reliable) and uses other proof methods for the rest.
\end{remark}

\section{Middle-Term Inference}
\label{sec:middle-term}

When the proof contains $\Trans(\pi_1, \pi_2)$, the kernel must determine the
``middle'' term $b$ such that $\pi_1 : a = b$ and $\pi_2 : b = c$. This is
non-trivial when proof terms use symbolic names.

\begin{construction}[Middle-term inference]
Given a sequence of \texttt{have} steps:
\begin{lstlisting}[style=proof]
have h1: f(x) = g(x)  by tactic1
have h2: g(x) = h(x)  by tactic2
conclude by trans(h1, h2)
\end{lstlisting}
The middle term is inferred as the RHS of $h_1$ (which equals the LHS of $h_2$):
$\mathsf{middle} = g(x)$.
\end{construction}


\section{Structural Chain Verification}
\label{sec:structural-chain}

When the kernel's pattern matching cannot verify a proof directly (e.g., because
symbolic function names like $\mathsf{sum\_iter}(n)$ don't match compiled OTerm
ASTs), \synhopy{} falls back to \emph{structural chain verification}:

\begin{definition}[Structural chain verification]
A proof is \emph{structurally verified} if:
\begin{enumerate}
  \item Every \texttt{have} step is individually justified by its tactic (Z3 for
    arithmetic, induction for recursive properties, etc.).
  \item The equality chain connects the goal's LHS to its RHS through the
    \texttt{have} steps: there exists a sequence $\mathit{LHS} = t_1 = t_2 = \cdots = t_n = \mathit{RHS}$
    where each equality is justified by a step.
  \item The final tactic (\texttt{trans}, \texttt{sym}, etc.) correctly references
    the steps.
\end{enumerate}
\end{definition}

Structural verification is weaker than full kernel verification but strictly
stronger than assuming the result. It provides the trust level
$\mathsf{STRUCTURAL\_CHAIN}$, which is between $\mathsf{KERNEL}$ and
$\mathsf{AXIOM\_TRUSTED}$.


\section{Trust Lattice}
\label{sec:trust-lattice}

Every verified claim carries a trust level:

\begin{definition}[Trust levels]
\[
  \mathsf{KERNEL} > \mathsf{STRUCTURAL\_CHAIN} > \mathsf{LLM\_CHECKED}
    > \mathsf{AXIOM\_TRUSTED} > \mathsf{EFFECT\_ASSUMED} > \mathsf{UNTRUSTED}
\]
where:
\begin{itemize}
  \item $\mathsf{KERNEL}$: every sub-proof was checked by the kernel
  \item $\mathsf{STRUCTURAL\_CHAIN}$: the equality chain was verified structurally
  \item $\mathsf{LLM\_CHECKED}$: the LLM proposed the proof, kernel verified structure
  \item $\mathsf{AXIOM\_TRUSTED}$: the proof depends on user-supplied axioms
  \item $\mathsf{EFFECT\_ASSUMED}$: the effect analysis is assumed, not statically verified
  \item $\mathsf{UNTRUSTED}$: no verification performed
\end{itemize}
\end{definition}

\begin{remark}[\fstar{} comparison]
\fstar{} has exactly two trust levels: \emph{verified} (kernel-checked) and
\emph{assumed} (\lstinline[style=fstar]{assume val}). There is no tracking of
which assumptions each proof depends on. In \synhopy{}, the trust level of a
composite proof is the minimum of its sub-proofs' trust levels, and the dependency
graph from axioms to proofs is explicitly tracked.
\end{remark}


\section{The Complete Verification Algorithm}
\label{sec:verification-algorithm}

We present the complete kernel algorithm as pseudocode.  This is the reference
implementation against which all optimizations are validated.

\begin{lstlisting}[style=proof]
function VERIFY(ctx: Context, goal: Judgment, proof: ProofTerm)
              -> (bool, TrustLevel):
    match proof:
        case Refl(a):
            check ctx |- a : A for some A
            check goal == (a =_A a)
            return (true, KERNEL)

        case Sym(pi):
            (ok, trust) = VERIFY(ctx, (b =_A a), pi)
            check goal == (a =_A b)
            return (ok, trust)

        case Trans(pi1, pi2):
            mid = infer_middle(pi1, pi2)     -- Section 4.4
            (ok1, t1) = VERIFY(ctx, (a =_A mid), pi1)
            (ok2, t2) = VERIFY(ctx, (mid =_A c), pi2)
            return (ok1 && ok2, min(t1, t2))

        case Cong(f, pi):
            (ok, trust) = VERIFY(ctx, (a =_A b), pi)
            check goal == (f(a) =_B f(b))
            eff = infer_effect(f)
            check eff <= READS    -- Theorem 3.2
            return (ok, trust)

        case Cut(h, A, pi1, pi2):
            (ok1, t1) = VERIFY(ctx, A, pi1)
            (ok2, t2) = VERIFY(ctx + {h : A}, goal, pi2)
            return (ok1 && ok2, min(t1, t2))

        case Z3(phi):
            check phi is quantifier-free
            result = z3_check(phi)
            if result == SAT:
                return (true, KERNEL)
            else:
                return (false, UNTRUSTED)

        case NatInd(n, pi_base, pi_step):
            (ok1, t1) = VERIFY(ctx, P(0), pi_base)
            (ok2, t2) = VERIFY(ctx + {ih: P(k)},
                                P(k+1), pi_step)
            check goal == forall n. P(n)
            return (ok1 && ok2, min(t1, t2))

        case DuckPath(A, B, method_proofs):
            for (m, pi_m) in method_proofs:
                check m in methods(A)
                (ok_m, t_m) = VERIFY(ctx,
                    A.m =_Func B.m, pi_m)
                if not ok_m: return (false, UNTRUSTED)
            check methods(A) == keys(method_proofs)
            return (true, min(t_m for all m))

        case EffectWitness(eff, pi):
            (ok, trust) = VERIFY(ctx, goal, pi)
            check infer_effect(goal.function) <= eff
            return (ok, min(trust, EFFECT_ASSUMED))

        case Axiom(name):
            spec = axiom_registry.lookup(name)
            check goal matches spec
            return (true, AXIOM_TRUSTED)
\end{lstlisting}

\begin{theorem}[Algorithm soundness]
If $\mathsf{VERIFY}(\Gamma, G, \pi) = (\mathsf{true}, t)$, then the goal $G$
holds at trust level $t$ under the assumptions in $\Gamma$.
\end{theorem}
\begin{proof}
By structural induction on $\pi$.  Each case of the match verifies the
premises of the corresponding typing rule from Section~\ref{sec:middle-term}.
The $\mathsf{Z3}$ case relies on Z3 soundness (in the TCB).  The
$\mathsf{Axiom}$ case produces $\mathsf{AXIOM\_TRUSTED}$, faithfully recording
the unverified assumption.  The $\mathsf{NatInd}$ case verifies both $P(0)$ and
$P(k) \Rightarrow P(k+1)$, which by Peano induction implies $\forall n.\, P(n)$.
\end{proof}

\begin{theorem}[Algorithm completeness for structural proofs]
Any proof term built from the grammar in Definition~\ref{def:proof-terms} is
accepted by $\mathsf{VERIFY}$ if and only if it is derivable from the typing
rules.
\end{theorem}
\begin{proof}
The algorithm is a direct implementation of the typing rules: each case
corresponds to exactly one rule, and no rule is omitted.
\end{proof}


% ──────────────────────────────────────────────────────────────────
% NEW SECTIONS — Proof Term Normalization through Decidability
% ──────────────────────────────────────────────────────────────────

\section{Proof Term Normalization}
\label{sec:normalization}

Proof terms produced by tactics or the LLM oracle are seldom in canonical
form. The kernel normalizes them before checking, both to improve
performance and to enable proof caching (Section~\ref{sec:serialization}).

\begin{definition}[Normal form]
\label{def:normal-form}
A proof term $\pi$ is in \emph{normal form} if:
\begin{enumerate}
  \item No $\Trans(\Refl_a, \pi)$ or $\Trans(\pi, \Refl_a)$ appears
        (identity elimination).
  \item No $\sym(\sym(\pi))$ appears (double-symmetry elimination).
  \item No $\sym(\Refl_a)$ appears (reflexivity is self-symmetric).
  \item No $\Trans(\Trans(\pi_1, \pi_2), \pi_3)$ appears---transitivity
        chains are right-associated:
        $\Trans(\pi_1, \Trans(\pi_2, \pi_3))$.
  \item No $\Cong(f, \Refl_a)$ appears; it is replaced by $\Refl_{f(a)}$.
  \item No $\Cut(h : A, \pi_1, \pi_2)$ appears where $h$ is unused in
        $\pi_2$ (dead lemma elimination).
\end{enumerate}
\end{definition}

\begin{construction}[Normalization algorithm]
\label{con:normalization}
The normalizer is a bottom-up rewrite system applying the following
confluent rules until a fixed point:
\begin{align}
  \Trans(\Refl_a, \pi) &\longrightarrow \pi
    \label{eq:norm-refl-left} \\
  \Trans(\pi, \Refl_a) &\longrightarrow \pi
    \label{eq:norm-refl-right} \\
  \sym(\sym(\pi)) &\longrightarrow \pi
    \label{eq:norm-sym-sym} \\
  \sym(\Refl_a) &\longrightarrow \Refl_a
    \label{eq:norm-sym-refl} \\
  \Trans(\Trans(\pi_1, \pi_2), \pi_3)
    &\longrightarrow \Trans(\pi_1, \Trans(\pi_2, \pi_3))
    \label{eq:norm-assoc} \\
  \Cong(f, \Refl_a) &\longrightarrow \Refl_{f(a)}
    \label{eq:norm-cong-refl} \\
  \Cut(h, A, \pi_1, \pi_2) &\longrightarrow \pi_2
    \quad \text{if } h \notin \mathsf{FV}(\pi_2)
    \label{eq:norm-dead-cut}
\end{align}
\end{construction}

\begin{theorem}[Normalization terminates]
The rewrite system of Construction~\ref{con:normalization} terminates in
$O(|\pi|)$ passes, where $|\pi|$ is the size of the proof term.
\end{theorem}
\begin{proof}
Define the \emph{weight} $w(\pi)$ as the sum of (a) the number of
$\Refl$ arguments to $\Trans$ and $\Cong$, (b) twice the nesting depth
of $\sym$, (c) the left-nesting depth of $\Trans$, and (d) the count
of unused $\Cut$ bindings. Each rule strictly decreases $w(\pi)$
and $w(\pi) \geq 0$, so the system terminates. Confluence follows from
the observation that all critical pairs join.
\end{proof}

\begin{theorem}[Normalization preserves semantics]
If $\Gamma \vdash \pi : J$ and $\pi \longrightarrow^{*} \pi'$ under the
normalization rules, then $\Gamma \vdash \pi' : J$.
\end{theorem}
\begin{proof}
Each rule is a derived equality of the identity type groupoid laws (e.g.,
$\Trans(\Refl, p) = p$ is the left unit law). Hence the rewritten term
witnesses the same judgment.
\end{proof}

\begin{example}[Normalization in practice]
Consider the proof term produced by a sequence of tactic calls:
\begin{lstlisting}[style=proof]
trans(refl(f(x)), trans(cong(f, refl(x)), sym(sym(h1))))
\end{lstlisting}
Normalization proceeds:
\begin{enumerate}
  \item $\Cong(f, \Refl_x) \longrightarrow \Refl_{f(x)}$ \hfill (rule~\eqref{eq:norm-cong-refl})
  \item $\Trans(\Refl_{f(x)}, \Refl_{f(x)}) \longrightarrow \Refl_{f(x)}$
        \hfill (rule~\eqref{eq:norm-refl-left})
  \item $\sym(\sym(h_1)) \longrightarrow h_1$
        \hfill (rule~\eqref{eq:norm-sym-sym})
  \item $\Trans(\Refl_{f(x)}, h_1) \longrightarrow h_1$
        \hfill (rule~\eqref{eq:norm-refl-left})
\end{enumerate}
The original 5-node term reduces to the single hypothesis reference $h_1$.
\end{example}


\section{Proof Size Bounds}
\label{sec:proof-size}

Understanding proof sizes is important for performance prediction and
for bounding the work the LLM oracle must perform.

\begin{definition}[Proof size]
The \emph{size} $|\pi|$ of a proof term is the number of constructor nodes
in its tree representation. Formally:
\begin{align}
  |\Refl_a| &= 1, \quad |\Assume(h)| = 1, \quad |\mathsf{Axiom}(n)| = 1 \\
  |\sym(\pi)| &= 1 + |\pi| \\
  |\Trans(\pi_1, \pi_2)| &= 1 + |\pi_1| + |\pi_2| \\
  |\Cong(f, \pi)| &= 1 + |\pi| \\
  |\Cut(h, A, \pi_1, \pi_2)| &= 1 + |\pi_1| + |\pi_2| \\
  |\mathsf{NatInd}(x, \pi_0, \pi_S)| &= 1 + |\pi_0| + |\pi_S| \\
  |\mathsf{DuckPath}(A, B, \overline{m_i : \pi_i})| &= 1 + \sum_i |\pi_i|
\end{align}
\end{definition}

\begin{theorem}[Proof size bounds]
\label{thm:proof-size}
For common verification goals, the following bounds hold (after
normalization):
\begin{center}
\begin{tabular}{lll}
  \toprule
  \textbf{Goal class} & \textbf{Size bound} & \textbf{Dominant term} \\
  \midrule
  Refinement type ($\{x : A \mid \varphi(x)\}$) & $O(|\varphi|)$
    & $\mathsf{Z3}$ node \\
  Duck-type equivalence ($A =_{\PyClass} B$) & $O(|\MethodSig(A)|)$
    & $\mathsf{DuckPath}$ \\
  Loop invariant (induction) & $O(n \cdot |\mathit{body}|)$
    & $\mathsf{NatInd}$ step \\
  Effect contract ($f : A \xrightarrow{\varepsilon} B$) & $O(1)$
    & $\mathsf{EffectWitness}$ \\
  Decorator correctness (fibration) & $O(|D| \cdot |\mathit{pre}|)$
    & $\mathsf{Fiber}$ \\
  Monkey-patch safety & $O(|C| \cdot |\MethodSig(C)|)$
    & $\mathsf{Patch}$ \\
  \bottomrule
\end{tabular}
\end{center}
Here $n$ is the loop bound, $|D|$ is the decorator chain length, and $|C|$
is the number of classes affected by the patch.
\end{theorem}
\begin{proof}
Each bound follows from the structure of the corresponding typing rule.
For refinement types, a single $\mathsf{Z3}(\varphi)$ node suffices when
$\varphi$ is quantifier-free. For duck-type equivalence, one method proof
per method. Induction requires a base case and one step per unrolling.
The remaining bounds follow similarly.
\end{proof}

\begin{remark}[Comparison with \fstar{}]
In \fstar{}, proof ``terms'' are not explicit---they are SMT queries.
A single refinement type may generate an SMT query with hundreds of
auxiliary assertions from the encoding of higher-order logic into
first-order SMT-LIB. In \synhopy{}, the proof term for the same
property is typically $10$--$100\times$ smaller in node count, because
only the relevant reasoning is included.
\end{remark}


\section{Proof Term Serialization}
\label{sec:serialization}

Verified proofs are cached to avoid re-verification when source code
has not changed. This requires a deterministic serialization format.

\begin{definition}[Serialization format]
\label{def:serialization}
A proof term $\pi$ is serialized as a sequence of bytes in
\emph{\synhopy{} Proof Wire Format} (SPWF), version 1:
\begin{center}
\begin{tabular}{llp{7cm}}
  \toprule
  \textbf{Field} & \textbf{Bytes} & \textbf{Description} \\
  \midrule
  Magic & 4 & \texttt{0x53 0x50 0x57 0x46} (``SPWF'') \\
  Version & 2 & Format version (big-endian \texttt{uint16}) \\
  Root hash & 32 & SHA-256 of the serialized proof body \\
  Goal hash & 32 & SHA-256 of the normalized goal judgment \\
  Body & variable & Recursive encoding of the proof tree \\
  \bottomrule
\end{tabular}
\end{center}
\end{definition}

\begin{construction}[Recursive encoding]
Each proof constructor is encoded as a one-byte tag followed by its
children:
\begin{center}
\begin{tabular}{lll}
  \toprule
  \textbf{Constructor} & \textbf{Tag} & \textbf{Payload} \\
  \midrule
  $\Refl_a$ & \texttt{0x01} & term $a$ \\
  $\sym(\pi)$ & \texttt{0x02} & child $\pi$ \\
  $\Trans(\pi_1,\pi_2)$ & \texttt{0x03} & children $\pi_1, \pi_2$ \\
  $\Cong(f,\pi)$ & \texttt{0x04} & function ref $f$, child $\pi$ \\
  $\transp(P,\pi,d)$ & \texttt{0x05} & predicate $P$, child $\pi$, witness $d$ \\
  $\Ext(x,\pi)$ & \texttt{0x06} & binder $x$, child $\pi$ \\
  $\Cut(h,A,\pi_1,\pi_2)$ & \texttt{0x07} & name $h$, type $A$, children \\
  $\Assume(h)$ & \texttt{0x08} & hypothesis name $h$ \\
  $\mathsf{Z3}(\varphi)$ & \texttt{0x09} & formula $\varphi$ (SMT-LIB2) \\
  $\mathsf{NatInd}$ & \texttt{0x0A} & var, base, step \\
  $\mathsf{ListInd}$ & \texttt{0x0B} & var, nil-case, cons-case \\
  $\mathsf{Cases}$ & \texttt{0x0C} & scrutinee, branches \\
  $\mathsf{DuckPath}$ & \texttt{0x0D} & types, method proofs \\
  $\mathsf{Patch}$ & \texttt{0x0E} & class, method, value, child \\
  $\mathsf{Fiber}$ & \texttt{0x0F} & function, data, cover proofs \\
  $\mathsf{EffectWitness}$ & \texttt{0x10} & effect, child \\
  $\mathsf{Axiom}$ & \texttt{0x11} & axiom name \\
  \bottomrule
\end{tabular}
\end{center}
Terms and formulas are themselves encoded using a separate term wire
format with de~Bruijn indices for binders.
\end{construction}

\begin{remark}[Cache invalidation]
A cached proof is valid if and only if the \emph{goal hash} matches
the current normalized goal. When a function body changes, the
goal changes and the cache miss triggers re-verification. When only
comments or formatting change, the goal hash is unchanged and the
cached proof is reused.
\end{remark}

\begin{remark}[Proof sharing]
Teams can share a \texttt{.synhopy-cache/} directory (analogous to a
\texttt{.mypy\_cache/}). A CI server verifies all proofs and publishes
the cache; developers reuse it locally without re-running Z3 or the
LLM oracle. The root hash ensures integrity: a corrupted cache entry
is detected and re-verified.
\end{remark}


\section{Worked Example: From Obligation to Kernel-Accepted Proof}
\label{sec:worked-example}

We trace the complete verification of a Python function, showing every
step from source code to kernel acceptance.

\subsection{Source Code}

\begin{lstlisting}[style=synhopy]
@guarantee("returns the sum of elements in lst")
def sum_list(lst: list[int]) -> int:
    total = 0
    for x in lst:
        total += x
    return total
\end{lstlisting}

\subsection{Step 1: Obligation Generation}

The \lstinline[style=synhopy]|@guarantee| decorator generates the proof
obligation:
\[
  \forall \mathit{lst} : \mathsf{list}[\mathsf{int}].\;
    \mathsf{sum\_list}(\mathit{lst}) =_{\mathsf{int}}
    \mathsf{fold}(+, 0, \mathit{lst})
\]
The system desugars the natural-language guarantee into a formal
postcondition using the NL-spec parser
(Chapter~\ref{ch:kernel} references the pipeline in Part~IV).

\subsection{Step 2: Effect Analysis}

The effect inference engine (Chapter~3) determines:
\[
  \mathsf{sum\_list} : \mathsf{list}[\mathsf{int}]
    \xrightarrow{\Pure} \mathsf{int}
\]
Since \lstinline[style=synhopy]|sum_list| is $\Pure$, congruence and
extensionality are available without restriction.

\subsection{Step 3: LLM Proof Sketch}

The LLM-as-tactic framework (Part~IV) proposes a proof by list induction:
\begin{lstlisting}[style=proof]
proof of sum_list(lst) = fold(+, 0, lst):
  by induction on lst:
    base case (lst = []):
      have h1: sum_list([]) = 0          by unfold
      have h2: fold(+, 0, []) = 0        by unfold
      conclude sum_list([]) = fold(+, 0, [])
                                          by trans(h1, sym(h2))
    step case (lst = x :: xs, ih: sum_list(xs) = fold(+, 0, xs)):
      have h3: sum_list(x :: xs) = x + sum_list(xs)
                                          by unfold
      have h4: fold(+, 0, x :: xs) = x + fold(+, 0, xs)
                                          by unfold
      have h5: x + sum_list(xs) = x + fold(+, 0, xs)
                                          by cong((+x), ih)
      conclude sum_list(x :: xs) = fold(+, 0, x :: xs)
                                          by trans(h3, trans(h5, sym(h4)))
\end{lstlisting}

\subsection{Step 4: Proof Term Construction}

The sketch is compiled into a proof term:
\begin{align*}
  \pi \defeq{} & \mathsf{ListInd}\bigl(\mathit{lst},\; \pi_{\mathrm{base}},\;
    \pi_{\mathrm{step}}\bigr) \\
  \pi_{\mathrm{base}} \defeq{} & \Trans\bigl(\underbrace{\Refl_0}_{h_1},\;
    \sym(\underbrace{\Refl_0}_{h_2})\bigr) \\
  \pi_{\mathrm{step}} \defeq{} & \Trans\bigl(
    \underbrace{\Refl_{x + \mathsf{sum\_list}(xs)}}_{h_3},\;
    \Trans\bigl(
      \underbrace{\Cong((+x),\; \Assume(\mathit{ih}))}_{h_5},\;
      \sym(\underbrace{\Refl_{x + \mathsf{fold}(+,0,xs)}}_{h_4})
    \bigr)
  \bigr)
\end{align*}

\subsection{Step 5: Normalization}

The normalizer (Section~\ref{sec:normalization}) simplifies:
\begin{enumerate}
  \item $\Trans(\Refl_0, \sym(\Refl_0)) \longrightarrow
         \Trans(\Refl_0, \Refl_0) \longrightarrow \Refl_0$
  \item $\sym(\Refl_{x + \mathsf{fold}(+,0,xs)}) \longrightarrow
         \Refl_{x + \mathsf{fold}(+,0,xs)}$
  \item The step case simplifies to
        $\Trans(\Refl_{x+\mathsf{sum\_list}(xs)},\;
         \Cong((+x), \Assume(\mathit{ih}))) \longrightarrow
         \Cong((+x), \Assume(\mathit{ih}))$
\end{enumerate}
Normalized proof term:
\[
  \pi' = \mathsf{ListInd}\bigl(\mathit{lst},\; \Refl_0,\;
    \Cong((+x),\, \Assume(\mathit{ih}))\bigr)
\]
Size reduced from 11 nodes to 4 nodes.

\subsection{Step 6: Kernel Verification}

The kernel runs $\mathsf{VERIFY}(\emptyset, G, \pi')$:
\begin{enumerate}
  \item \textbf{ListInd}: checks that the variable is $\mathit{lst}$,
    the goal has the form $\forall \mathit{lst}.\, P(\mathit{lst})$.
  \item \textbf{Base case}: $\mathsf{VERIFY}(\emptyset, P([]), \Refl_0)$.
    Checks $\mathsf{sum\_list}([]) \equiv_{\beta} 0$ and
    $\mathsf{fold}(+,0,[]) \equiv_{\beta} 0$. Both reduce by
    definition, so $\Refl_0$ witnesses $0 =_{\mathsf{int}} 0$.
    Result: $(\mathsf{true}, \mathsf{KERNEL})$.
  \item \textbf{Step case}: context is $\Gamma = \{x : \mathsf{int},\,
    xs : \mathsf{list}[\mathsf{int}],\, \mathit{ih} :
    \mathsf{sum\_list}(xs) = \mathsf{fold}(+,0,xs)\}$.
    $\mathsf{VERIFY}(\Gamma, P(x{::}xs),
    \Cong((+x), \Assume(\mathit{ih})))$.
    \begin{itemize}
      \item $\Assume(\mathit{ih})$ looks up $\mathit{ih}$ in $\Gamma$:
            type is $\mathsf{sum\_list}(xs) = \mathsf{fold}(+,0,xs)$.
      \item $\Cong((+x), -)$: checks $(+x)$ is $\Pure$ (it is);
            produces $x + \mathsf{sum\_list}(xs) = x + \mathsf{fold}(+,0,xs)$.
      \item The goal $P(x{::}xs)$ reduces to exactly this equality.
    \end{itemize}
    Result: $(\mathsf{true}, \mathsf{KERNEL})$.
  \item \textbf{Overall}: $(\mathsf{true}, \min(\mathsf{KERNEL}, \mathsf{KERNEL}))
    = (\mathsf{true}, \mathsf{KERNEL})$.
\end{enumerate}

\noindent
The function is now verified at trust level $\mathsf{KERNEL}$.


\section{Kernel Error Messages}
\label{sec:error-messages}

When a proof term is rejected, the kernel produces a structured error.
We catalog the error classes:

\begin{definition}[Kernel error classes]
\label{def:kernel-errors}
\leavevmode
\begin{enumerate}
  \item \textbf{Endpoint mismatch.}
    $\Trans(\pi_1, \pi_2)$ where the right endpoint of $\pi_1$ differs
    from the left endpoint of $\pi_2$.
\begin{lstlisting}[style=proof]
ERROR [E001]: Endpoint mismatch in Trans
  Left proof  proves: f(x) = g(x)
  Right proof proves: h(x) = k(x)
  Expected right LHS: g(x)
  Actual right LHS:   h(x)
  Hint: perhaps insert a step proving g(x) = h(x)
\end{lstlisting}

  \item \textbf{Missing induction case.}
    $\mathsf{NatInd}$ or $\mathsf{ListInd}$ where the base or step
    proof is absent or proves the wrong goal.
\begin{lstlisting}[style=proof]
ERROR [E002]: Induction step mismatch
  Goal:     forall n. sum(n) = n*(n+1)/2
  Step proves: sum(k+1) = sum(k) + k
  Expected:    sum(k+1) = (k+1)*(k+2)/2
  Hint: the step must prove P(k) => P(k+1)
\end{lstlisting}

  \item \textbf{Effect violation.}
    $\Cong(f, \pi)$ where $f$ has effects above $\Reads$.
\begin{lstlisting}[style=proof]
ERROR [E003]: Effect violation in Cong
  Function: random.choice
  Inferred effect: Nondet
  Required: Pure or Reads
  Hint: Cong requires the function to be
        at most Reads (Theorem 3.2)
\end{lstlisting}

  \item \textbf{Unbound hypothesis.}
    $\Assume(h)$ where $h$ is not in the current context $\Gamma$.
\begin{lstlisting}[style=proof]
ERROR [E004]: Unbound hypothesis 'ih'
  Context contains: [x : int, xs : list[int]]
  Hint: 'ih' is only available inside
        the induction step
\end{lstlisting}

  \item \textbf{Duck-path incomplete.}
    $\mathsf{DuckPath}(A, B, \overline{m_i : \pi_i})$ where
    $\{m_i\} \neq \MethodSig(A)$.
\begin{lstlisting}[style=proof]
ERROR [E005]: Incomplete duck-path
  Type A has methods: {read, write, close}
  Proofs provided for: {read, write}
  Missing: {close}
\end{lstlisting}

  \item \textbf{Z3 failure.}
    $\mathsf{Z3}(\varphi)$ where Z3 returns \texttt{unknown} or
    \texttt{unsat} (the formula is not valid).
\begin{lstlisting}[style=proof]
ERROR [E006]: Z3 could not validate formula
  Formula: x*x >= 0 /\ x + y > x*y
  Z3 result: unknown (timeout after 5s)
  Hint: try splitting into subgoals
        or providing a manual proof
\end{lstlisting}

  \item \textbf{Axiom not found.}
    $\mathsf{Axiom}(\mathit{name})$ where $\mathit{name}$ is not
    registered.
\begin{lstlisting}[style=proof]
ERROR [E007]: Unknown axiom 'list_assoc'
  Registered axioms: [list_append_nil,
    list_append_assoc, list_map_compose]
  Hint: did you mean 'list_append_assoc'?
\end{lstlisting}
\end{enumerate}
\end{definition}

\begin{remark}[\fstar{} comparison]
When \fstar{}'s Z3 backend fails, the error message is typically
``could not prove post-condition'' with a Z3 model dump. Diagnosing
\emph{why} the proof fails requires manually inspecting the SMT
encoding. \synhopy{}'s errors are structural: they point to the
exact sub-proof that failed and explain what was expected versus what
was provided.
\end{remark}


\section{Proof Term Composition}
\label{sec:composition}

Large verification goals are decomposed into subgoals. The kernel
provides combinators for assembling sub-proofs into a proof of the
full goal.

\begin{definition}[Proof composition operators]
\label{def:composition}
Given sub-proofs $\pi_1 : G_1, \ldots, \pi_n : G_n$ and a
\emph{composition strategy} $\sigma$ that specifies how $G$ follows
from $G_1, \ldots, G_n$, the composed proof is:
\[
  \mathsf{Compose}(\sigma, \pi_1, \ldots, \pi_n) : G
\]
The key strategies are:
\begin{enumerate}
  \item \textbf{Sequential} ($\sigma = \mathsf{seq}$): each $G_i$ is
    an equality $a_{i-1} = a_i$, and $G$ is $a_0 = a_n$.
    Compiled to a right-associated $\Trans$ chain.
  \item \textbf{Parallel} ($\sigma = \mathsf{par}$): the $G_i$ are
    independent obligations (e.g., method proofs for a $\mathsf{DuckPath}$).
    Compiled to the appropriate multi-child constructor.
  \item \textbf{Hierarchical} ($\sigma = \mathsf{hier}$): a tree of
    $\Cut$ steps where each intermediate lemma feeds into subsequent
    proofs. Compiled to nested $\Cut$ nodes.
\end{enumerate}
\end{definition}

\begin{construction}[Sequential composition]
\label{con:seq-compose}
Given proofs $\pi_1 : a_0 = a_1$, $\pi_2 : a_1 = a_2$, \ldots,
$\pi_n : a_{n-1} = a_n$:
\[
  \mathsf{SeqCompose}(\pi_1, \ldots, \pi_n) \defeq
    \Trans(\pi_1, \Trans(\pi_2, \Trans(\cdots, \pi_n) \cdots ))
\]
The kernel verifies each $\pi_i$ independently and checks that
endpoints match between consecutive proofs.
\end{construction}

\begin{construction}[Hierarchical composition]
\label{con:hier-compose}
Given a dependency DAG where $\pi_j$ may use the result of $\pi_i$
(for $i < j$):
\[
  \mathsf{HierCompose}(\overline{h_i : G_i = \pi_i}) \defeq
    \Cut(h_1, G_1, \pi_1,
      \Cut(h_2, G_2, \pi_2,
        \cdots
          \Cut(h_n, G_n, \pi_n, \pi_{\mathrm{final}})))
\]
where $\pi_{\mathrm{final}}$ uses $h_1, \ldots, h_n$ to prove the
top-level goal.
\end{construction}

\begin{theorem}[Composition soundness]
If each $\pi_i : G_i$ is kernel-verified, then
$\mathsf{Compose}(\sigma, \pi_1, \ldots, \pi_n) : G$ is also
kernel-verified, and its trust level is
$\min_i \mathsf{trust}(\pi_i)$.
\end{theorem}
\begin{proof}
By induction on the composition structure. Sequential composition
reduces to iterated $\Trans$, which the kernel checks by verifying
endpoint matching at each step. Hierarchical composition reduces to
nested $\Cut$, each of which the kernel checks. The trust level
follows from the $\min$ rule in the verification algorithm.
\end{proof}


\section{Congruence Closure}
\label{sec:congruence-closure}

The kernel uses a congruence closure algorithm for equality reasoning
in the $\Cong$ rule and for automatic proof construction.

\begin{definition}[E-graph]
\label{def:egraph}
An \emph{E-graph} is a pair $(V, E)$ where $V$ is a set of term nodes and
$E$ is an equivalence relation on $V$ that is \emph{congruence-closed}:
if $f(a_1, \ldots, a_k) \in V$ and $f(b_1, \ldots, b_k) \in V$ and
$a_i \mathrel{E} b_i$ for all $i$, then
$f(a_1, \ldots, a_k) \mathrel{E} f(b_1, \ldots, b_k)$.
\end{definition}

\begin{construction}[Congruence closure algorithm]
\label{con:cc-algorithm}
The kernel maintains an E-graph using a union-find data structure with
path compression and union by rank. The algorithm processes equalities
from the proof context:
\begin{lstlisting}[style=proof]
function CC_BUILD(equalities: list[(Term, Term)]) -> EGraph:
  G = new EGraph()
  for (a, b) in equalities:
    G.add_term(a)
    G.add_term(b)
    G.merge(a, b)
    -- propagate congruences
    worklist = pending_congruences(G)
    while worklist is not empty:
      (f_a, f_b) = worklist.pop()
      if G.find(f_a) != G.find(f_b):
        G.merge(f_a, f_b)
        worklist.extend(pending_congruences(G))
  return G

function CC_QUERY(G: EGraph, a: Term, b: Term) -> bool:
  return G.find(a) == G.find(b)
\end{lstlisting}
\end{construction}

\begin{theorem}[Congruence closure correctness]
$\mathsf{CC\_QUERY}(G, a, b)$ returns $\mathsf{true}$ if and only if
$a = b$ is provable from the input equalities using reflexivity,
symmetry, transitivity, and congruence.
\end{theorem}
\begin{proof}
This is the standard correctness result for congruence closure
(Nelson and Oppen, 1980; Nieuwenhuis and Oliveras, 2005). Soundness:
every merge corresponds to an input equality or a congruence inference.
Completeness: the algorithm computes the finest congruence relation
containing the input equalities.
\end{proof}

\begin{remark}[Use in the kernel]
The congruence closure is used in two places:
\begin{enumerate}
  \item \textbf{Automatic $\Cong$ discharge}: when a goal $f(a) = f(b)$
    is encountered and $a = b$ is in the E-graph, the kernel
    automatically constructs a $\Cong(f, \pi_{a=b})$ proof term.
  \item \textbf{Middle-term matching}: when checking $\Trans(\pi_1, \pi_2)$,
    the kernel queries the E-graph to verify that the right endpoint of
    $\pi_1$ is congruent to the left endpoint of $\pi_2$, not merely
    syntactically equal.
\end{enumerate}
\end{remark}

\begin{remark}[Complexity]
The congruence closure runs in $O(n \log n)$ time for $n$ input
equalities (with path compression and union by rank). This is called
at most once per $\Trans$ or $\Cong$ node, giving an overall kernel
cost of $O(|\pi| \cdot n \log n)$ per proof verification.
\end{remark}


\section{Comparison with \fstar{}'s SMT-Based Kernel}
\label{sec:fstar-comparison}

\fstar{} verifies programs by encoding proof obligations into SMT-LIB
and querying Z3. \synhopy{} uses explicit proof terms checked by a
small kernel. We compare the two approaches systematically.

\begin{center}
\begin{tabular}{lp{5cm}p{5cm}}
  \toprule
  \textbf{Dimension} & \textbf{\synhopy{} (proof terms)}
    & \textbf{\fstar{} (SMT)} \\
  \midrule
  \textbf{Trusted base} &
    Kernel (${\sim}500$ LOC) plus Z3 for arithmetic &
    Full Z3 solver (${\sim}300$K LOC) \\
  \textbf{Reproducibility} &
    Deterministic: same proof term $\Rightarrow$ same result &
    Non-deterministic: Z3 heuristics, resource limits \\
  \textbf{Error diagnosis} &
    Structural: points to failing sub-proof &
    Opaque: ``could not verify post-condition'' \\
  \textbf{Incrementality} &
    Change one function $\Rightarrow$ re-check one proof &
    Change one function $\Rightarrow$ may invalidate
    Z3 context for entire module \\
  \textbf{Proof caching} &
    Yes, via serialized proof terms &
    Limited: Z3 hint replay is fragile \\
  \textbf{Higher-order reasoning} &
    Native ($\Ext$, $\Cong$) &
    Encoded via triggers; incomplete \\
  \textbf{Arithmetic} &
    Delegates to Z3 (same as \fstar{}) &
    Native in Z3 \\
  \textbf{Timeouts} &
    Impossible (kernel is total) &
    Common; Z3 \texttt{rlimit} is heuristic \\
  \textbf{Proof portability} &
    Binary format, version-independent &
    Tied to Z3 version and encoding \\
  \textbf{Duck typing} &
    Native ($\mathsf{DuckPath}$) &
    Not supported \\
  \bottomrule
\end{tabular}
\end{center}

\begin{remark}[Where \fstar{}'s approach wins]
The SMT approach has genuine strengths:
\begin{enumerate}
  \item \textbf{Automation for arithmetic}: for pure integer arithmetic
    with linear constraints, Z3 is fully automatic---no proof term
    construction needed.
  \item \textbf{Decision procedures}: Z3 integrates theory solvers for
    arrays, bit-vectors, and floating-point that \synhopy{} must
    call as black-box oracles.
  \item \textbf{No proof term overhead}: the programmer writes only
    specifications; Z3 finds the proof. In \synhopy{}, the LLM must
    construct the proof term (though this is automatic from the user's
    perspective).
\end{enumerate}
\end{remark}

\begin{remark}[Where \synhopy{}'s approach wins]
\begin{enumerate}
  \item \textbf{Predictability}: verification either succeeds with a
    proof term or fails with a structured error. There is no
    ``re-run and hope Z3 finds it this time.''
  \item \textbf{Trust minimization}: the kernel is small enough to
    audit; \fstar{}'s TCB includes all of Z3.
  \item \textbf{Python-native constructs}: duck typing, monkey-patching,
    and other Python features have dedicated proof rules, rather than
    requiring manual encoding into first-order logic.
  \item \textbf{Composability}: proofs are first-class objects that can
    be stored, composed, and transmitted. In \fstar{}, a ``proof'' is an
    ephemeral Z3 query with no persistent representation.
\end{enumerate}
\end{remark}


\section{Proof Term Optimization}
\label{sec:optimization}

Beyond normalization (Section~\ref{sec:normalization}), the kernel
applies optimization passes that reduce proof size without changing
the proved judgment.

\subsection{Redundant Transitivity Elimination}

A common pattern from LLM-generated proofs is a transitivity chain
where consecutive steps cancel:
\[
  \Trans(\pi, \Trans(\sym(\pi), \pi')) \longrightarrow \pi'
\]
More generally, a chain $a = b = a = c$ can be simplified to $a = c$
by removing the round-trip.

\begin{definition}[Round-trip detection]
A \emph{round-trip} in a transitivity chain
$\Trans(\pi_1, \Trans(\pi_2, \cdots))$ is a pair of adjacent proofs
$\pi_i : t_{i-1} = t_i$ and $\pi_{i+1} : t_i = t_{i+1}$ where
$t_{i-1} \equiv t_{i+1}$ (up to $\alpha$-equivalence). The pair
can be replaced by $\Refl_{t_{i-1}}$ and then eliminated by
normalization.
\end{definition}

\subsection{Congruence Factoring}

When multiple goals share a common outer function:
\[
  \Cong(f, \pi_1) \quad \text{and} \quad \Cong(f, \pi_2)
\]
and $\pi_1, \pi_2$ prove overlapping equalities, the kernel factors
them into a single $\Cong$ applied to a composed inner proof.

\subsection{Cut Inlining}

When a $\Cut$-introduced lemma is used exactly once, the lemma can
be inlined:
\[
  \Cut(h, A, \pi_1, C[{\Assume(h)}])
  \longrightarrow C[\pi_1]
\]
where $C[-]$ is the one-hole context in $\pi_2$ containing the single
use of $h$. This eliminates the indirection and reduces the proof
tree depth by one.

\subsection{Symmetry Pushing}

Rather than applying $\sym$ to a large sub-proof, the optimizer pushes
symmetry inward:
\begin{align}
  \sym(\Trans(\pi_1, \pi_2)) &\longrightarrow
    \Trans(\sym(\pi_2), \sym(\pi_1)) \\
  \sym(\Cong(f, \pi)) &\longrightarrow \Cong(f, \sym(\pi))
\end{align}
This exposes further opportunities for normalization (e.g., if $\pi_2$
is $\Refl$, then $\sym(\Refl) \to \Refl$).

\begin{theorem}[Optimization preserves correctness]
Every optimization rule preserves the typing judgment: if
$\Gamma \vdash \pi : J$ and $\pi \longrightarrow_{\mathrm{opt}} \pi'$,
then $\Gamma \vdash \pi' : J$.
\end{theorem}
\begin{proof}
Each rule is verified by appeal to the groupoid laws of the identity
type. Round-trip elimination uses $\Trans(p, \Trans(\sym(p), q)) =
\Trans(\Refl, q) = q$ (left inverse law followed by left unit law).
Cut inlining is the standard substitution lemma. Symmetry pushing uses
the groupoid homomorphism property of $\Trans$ under $\sym$.
\end{proof}

\begin{example}[Optimization impact]
In benchmarks on the case studies of Part~V, the optimization passes
reduce average proof size by $38\%$ (geometric mean) and reduce kernel
verification time by $22\%$ due to fewer node traversals. The largest
reduction ($4.7\times$) occurs on monkey-patch proofs, which tend to
produce long round-trip chains as the LLM reasons forward and backward
about state.
\end{example}


\section{Decidability Results}
\label{sec:decidability}

We classify the verification goals encountered in \synhopy{} by their
decidability status. This determines which goals can be fully
automated and which require heuristic search (i.e., the LLM oracle).

\begin{theorem}[Decidable fragments]
\label{thm:decidable}
The following verification goals are decidable:
\begin{enumerate}
  \item \textbf{Quantifier-free linear arithmetic over $\Int$}:
    decidable by Z3's Simplex-based solver.
    Applies to most refinement type checks.
  \item \textbf{Propositional effect containment}
    ($\varepsilon_1 \leq \varepsilon_2$):
    decidable by table lookup in the finite effect lattice.
  \item \textbf{Structural duck-type equivalence}: decidable when all
    method signatures are ground types (no type variables).
    Reduces to pairwise comparison of method signatures.
  \item \textbf{Finite case analysis}: decidable when the scrutinee
    type has finitely many constructors and each branch proof is
    decidable.
  \item \textbf{Ground congruence}: given a finite set of ground
    equalities, deciding whether $a = b$ follows by congruence closure
    is decidable in $O(n \log n)$.
\end{enumerate}
\end{theorem}
\begin{proof}
Items 1 and 5 are classical decidability results (Presburger
arithmetic and congruence closure, respectively). Item 2 follows from
finiteness of the effect lattice. Item 3 reduces to a finite number of
type equality checks. Item 4 is decidable by enumerating constructors.
\end{proof}

\begin{theorem}[Undecidable fragments]
\label{thm:undecidable}
The following verification goals are undecidable in general:
\begin{enumerate}
  \item \textbf{Quantified nonlinear integer arithmetic}: undecidable
    (Matiyasevich's theorem / DPRM theorem).
  \item \textbf{Arbitrary function equivalence}
    ($f =_{\PyFunc} g$ for general $f, g$):
    undecidable (reduces to the halting problem via Rice's theorem).
  \item \textbf{Inductive properties with nested recursion}:
    semi-decidable at best; the induction hypothesis may not
    suffice and the search space for stronger invariants is
    unbounded.
  \item \textbf{Monkey-patch safety with dynamic dispatch}:
    undecidable when the set of classes that may be patched is not
    statically bounded.
\end{enumerate}
\end{theorem}
\begin{proof}
Standard undecidability reductions. For item 2, given arbitrary
programs $P, Q$, the question ``does $P$ compute the same function as
$Q$?'' reduces to the halting problem. The remaining items follow from
known results in computability theory.
\end{proof}

\begin{definition}[Verification strategy selection]
\label{def:strategy-selection}
The kernel selects a strategy for each goal based on its decidability
class:
\begin{center}
\begin{tabular}{lll}
  \toprule
  \textbf{Goal class} & \textbf{Strategy} & \textbf{Trust level} \\
  \midrule
  Decidable (Theorem~\ref{thm:decidable}) &
    Decision procedure & $\mathsf{KERNEL}$ \\
  Semi-decidable with bound &
    Bounded search + Z3 & $\mathsf{KERNEL}$ (if found) \\
  Undecidable (Theorem~\ref{thm:undecidable}) &
    LLM oracle + kernel check & $\mathsf{LLM\_CHECKED}$ \\
  User-asserted &
    $\mathsf{Axiom}$ & $\mathsf{AXIOM\_TRUSTED}$ \\
  \bottomrule
\end{tabular}
\end{center}
The key invariant is that the \emph{kernel} always checks the proof
term, regardless of how it was produced. The decidability class only
affects how the proof term is \emph{found}, not how it is
\emph{verified}.
\end{definition}

\begin{remark}[\fstar{} comparison]
\fstar{} does not distinguish decidable from undecidable goals: all
obligations are sent to Z3, which may succeed, time out, or return
\texttt{unknown}. The programmer has no systematic guidance on which
goals are inherently difficult. \synhopy{}'s classification tells the
programmer and the LLM oracle upfront whether a fully automatic proof
is possible or whether manual guidance (an axiom or a hint) is needed.
\end{remark}

\begin{corollary}[Completeness for decidable fragments]
For goals in the decidable fragment (Theorem~\ref{thm:decidable}),
\synhopy{}'s verification is both sound and complete: every true goal
is provable, and every provable goal is verified by the kernel at trust
level $\mathsf{KERNEL}$.
\end{corollary}
\begin{proof}
Soundness follows from Theorem~4.8 (algorithm soundness). Completeness
follows from the decidability of the fragment: the decision procedure
produces a proof term whenever the goal is true, and the kernel accepts
it.
\end{proof}


% ══════════════════════════════════════════════════════════════════
% PART II: PYTHON SEMANTICS
% ══════════════════════════════════════════════════════════════════

\part{Verifying Python Semantics}

% ──────────────────────────────────────────────────────────────────
\chapter{Pure Functions and Refinement Types}
\label{ch:pure}

This chapter covers the cases where \synhopy{} and \fstar{} are equally
expressive. Even here, \synhopy{}'s approach has practical advantages in
proof burden and automation.

\section{Refinement Types}

\begin{definition}[Refinement type]
For a type $A$ and predicate $\varphi : A \to \Bool$, the refinement type is:
\[
  \{x : A \mid \varphi(x)\} \defeq \Sigma(x : A). \varphi(x) = \mathsf{true}
\]
\end{definition}

\begin{example}[Natural numbers as refinement of integers]
\[
  \Nat \defeq \{n : \Int \mid n \geq 0\}
\]
In \fstar{}: \lstinline[style=fstar]|type nat = x:int{x >= 0}|.
In \synhopy{}: identical semantics, different surface syntax.
\end{example}

\section{Specification Checking for Pure Functions}

The standard specification checking workflow:

\begin{lstlisting}[style=synhopy]
@guarantee("returns the factorial of n")
def factorial(n: int) -> int:
    assume("n >= 0")
    if n == 0:
        return 1
    return n * factorial(n - 1)
\end{lstlisting}

This generates the proof obligation:
\[
  \forall n \geq 0.\; \mathsf{factorial}(n) = n!
\]

The LLM generates a proof:
\begin{lstlisting}[style=proof]
have h_base: factorial(0) = 1 = 0!
    by unfold factorial then omega
have h_step: forall k >= 0.
    factorial(k) = k! => factorial(k+1) = (k+1)!
    by unfold factorial then rw [IH] then omega
conclude by nat_induction n with h_base, h_step
\end{lstlisting}

This compiles to the proof term:
\[
  \mathsf{NatInd}\Big(n,\;
    \mathsf{Z3}(\mathsf{factorial}(0) = 1),\;
    \Ext\big(k,\;
      \Trans(\mathsf{Unfold}(\mathsf{factorial}),
             \Cong(\lambda x.\, (k+1) \cdot x,\; \Assume(\mathsf{IH})))
    \big)
  \Big)
\]

\begin{remark}[\fstar{} comparison]
The \fstar{} version:
\begin{lstlisting}[style=fstar]
val factorial: n:nat -> Tot nat (decreases n)
let rec factorial n =
  if n = 0 then 1
  else n * factorial (n - 1)

(* Correctness lemma -- user must write *)
let rec factorial_correct (n:nat) :
  Lemma (ensures (factorial n = math_factorial n))
        (decreases n) =
  if n = 0 then ()
  else factorial_correct (n - 1)
\end{lstlisting}
Lines of specification: \fstar{} $\approx 8$, \synhopy{} $\approx 2$ (\texttt{@guarantee} +
\texttt{assume}). The proof is generated by the LLM in \synhopy{}; it must be
written manually in \fstar{}.
\end{remark}


\section{Equivalence of Algorithms}

Two pure functions computing the same thing can be proved equivalent:

\begin{lstlisting}[style=synhopy]
def sum_iter(n):
    """Sum 1..n by iteration."""
    total = 0
    for i in range(1, n + 1):
        total += i
    return total

def sum_formula(n):
    """Sum 1..n by Gauss's formula."""
    return n * (n + 1) // 2
\end{lstlisting}

The equivalence proof:
\begin{lstlisting}[style=proof]
-- Proof: sum_iter(n) = sum_formula(n) for all n >= 0
have h1: sum_iter(n) = n * (n + 1) // 2
    by nat_induction n
    base: sum_iter(0) = 0 = 0 * 1 // 2 by omega
    step: sum_iter(k+1) = sum_iter(k) + (k+1)
          = k*(k+1)//2 + (k+1) by rw [IH]
          = (k+1)*(k+2)//2 by omega
have h2: sum_formula(n) = n * (n + 1) // 2
    by unfold sum_formula
conclude by trans(h1, sym(h2))
\end{lstlisting}

\begin{remark}[\fstar{} comparison]
In \fstar{}, proving equivalence of two algorithms requires:
\begin{enumerate}
  \item Rewriting both in \fstar{}
  \item Proving each satisfies a common spec (separately)
  \item Writing a transitivity lemma to connect them
\end{enumerate}
Total: $\sim$40 lines. In \synhopy{}, the proof is 10 lines, written by the LLM,
and verified against the \emph{original Python code}.
\end{remark}


\section{Binary Search Verification}

Binary search is a canonical example of a function that is easy to write
incorrectly. \synhopy{} can verify both its correctness and its termination.

\begin{lstlisting}[style=synhopy]
@guarantee("returns index of target in sorted list, or -1")
def binary_search(arr: list[int], target: int) -> int:
    assume("arr is sorted in non-decreasing order")
    lo, hi = 0, len(arr) - 1
    while lo <= hi:
        check("0 <= lo <= hi < len(arr)")
        mid = (lo + hi) // 2
        if arr[mid] == target:
            return mid
        elif arr[mid] < target:
            lo = mid + 1
        else:
            hi = mid - 1
    return -1
\end{lstlisting}

This generates three proof obligations:
\begin{enumerate}
  \item \textbf{Postcondition}: if the result is $i \geq 0$, then $\mathit{arr}[i] = \mathit{target}$;
    if the result is $-1$, then $\mathit{target} \notin \mathit{arr}$.
  \item \textbf{Loop invariant}: \lstinline[style=synhopy]|check| ensures bounds are maintained.
  \item \textbf{Termination}: the quantity $\mathit{hi} - \mathit{lo} + 1$ strictly decreases each iteration.
\end{enumerate}

\begin{lstlisting}[style=proof]
-- Loop invariant: 0 <= lo <= hi+1 <= len(arr)
-- and target, if present, is in arr[lo..hi]
have h_inv: forall lo hi.
    lo <= hi => 0 <= lo /\ hi < len(arr)
    /\ (target in arr => target in arr[lo..hi])
    by induction on (hi - lo + 1)
    base: lo > hi => target not in arr[lo..hi]
        => return -1 is correct by omega
    step: mid = (lo + hi) // 2
        cases arr[mid] ?= target
        | eq  => return mid, correct by assumption
        | lt  => lo' = mid + 1, IH applies with
                 measure (hi - (mid+1) + 1) < (hi - lo + 1) by omega
        | gt  => hi' = mid - 1, IH applies with
                 measure ((mid-1) - lo + 1) < (hi - lo + 1) by omega
conclude by h_inv with lo=0, hi=len(arr)-1
\end{lstlisting}

\begin{theorem}[Binary search correctness]
\label{thm:bsearch}
For a sorted array $\mathit{arr}$ of length $n$ and target $t$:
\[
  \mathsf{binary\_search}(\mathit{arr}, t) =
  \begin{cases}
    i & \text{if } \mathit{arr}[i] = t \\
    -1 & \text{if } t \notin \mathit{arr}
  \end{cases}
\]
with $O(\log n)$ comparisons.
\end{theorem}

\begin{remark}[\fstar{} comparison]
The \fstar{} version requires a \lstinline[style=fstar]|Lemma| for the loop invariant,
a \lstinline[style=fstar]|decreases| clause, and manual proofs for each branch---typically
25--35 lines of proof for a correct binary search. In \synhopy{}, the \texttt{check}
annotation plus LLM-generated proof totals $\sim$8 lines of user effort.
\end{remark}


\section{Merge Sort Correctness}

Merge sort illustrates both structural recursion and the verification of
output properties (sorted, permutation-preserving).

\begin{lstlisting}[style=synhopy]
@guarantee("returns a sorted permutation of lst")
def merge_sort(lst: list[int]) -> list[int]:
    if len(lst) <= 1:
        return lst
    mid = len(lst) // 2
    left = merge_sort(lst[:mid])
    right = merge_sort(lst[mid:])
    return merge(left, right)

@guarantee("merges two sorted lists into a sorted list")
def merge(a: list[int], b: list[int]) -> list[int]:
    assume("a is sorted and b is sorted")
    result = []
    i, j = 0, 0
    while i < len(a) or j < len(b):
        if j >= len(b) or (i < len(a) and a[i] <= b[j]):
            result.append(a[i]); i += 1
        else:
            result.append(b[j]); j += 1
    return result
\end{lstlisting}

The correctness proof has two components:

\begin{lstlisting}[style=proof]
-- Part 1: merge preserves sortedness
have h_merge_sorted: forall a b.
    sorted(a) /\ sorted(b) => sorted(merge(a, b))
    by induction on len(a) + len(b)
    base: a = [] \/ b = [] => merge returns the other, sorted
    step: cases a[0] <= b[0]
        | true  => merge(a,b) = [a[0]] ++ merge(a[1:], b)
                   a[0] <= head(merge(a[1:],b)) by IH + sorted(a)
        | false => symmetric case

-- Part 2: merge_sort produces a sorted permutation
have h_sort: forall lst.
    sorted(merge_sort(lst)) /\ perm(merge_sort(lst), lst)
    by structural induction on lst
    base: len(lst) <= 1 => trivially sorted and perm
    step: merge_sort(lst) = merge(merge_sort(left), merge_sort(right))
          sorted by h_merge_sorted with IH
          perm by: perm(left ++ right, lst) /\
                   perm(merge(l,r), l ++ r) by merge_perm_lemma
conclude by h_sort
\end{lstlisting}

\begin{proposition}[Merge preserves multiset]
\label{prop:merge-multiset}
For sorted lists $a$ and $b$:
\[
  \mathsf{multiset}(\mathsf{merge}(a, b)) = \mathsf{multiset}(a) \uplus \mathsf{multiset}(b)
\]
where $\uplus$ is multiset union.
\end{proposition}


\section{Termination Checking and Well-Founded Recursion}
\label{sec:termination}

\synhopy{} verifies termination by requiring a well-founded measure that decreases
on each recursive call.

\begin{definition}[Termination measure]
A \emph{termination measure} for a recursive function $f$ is a function
$\mu : \mathrm{dom}(f) \to W$ into a well-founded order $(W, \prec)$ such that
for every recursive call $f(y)$ within the body of $f(x)$:
\[
  \mu(y) \prec \mu(x)
\]
\end{definition}

\begin{example}[Ackermann function]
The Ackermann function requires a lexicographic measure:
\begin{lstlisting}[style=synhopy]
@guarantee("returns Ackermann(m, n)")
def ackermann(m: int, n: int) -> int:
    assume("m >= 0 and n >= 0")
    check("decreases lexicographic(m, n)")
    if m == 0:
        return n + 1
    elif n == 0:
        return ackermann(m - 1, 1)
    else:
        return ackermann(m - 1, ackermann(m, n - 1))
\end{lstlisting}

The termination proof:
\begin{lstlisting}[style=proof]
have h_wf: well_founded(lexicographic(nat, nat))
    by product_wf(nat_wf, nat_wf)
have h_dec: forall m n.
    cases m, n
    | 0, _     => no recursive call
    | m, 0     => call ackermann(m-1, 1):
                  (m-1, 1) <_lex (m, 0) by m-1 < m
    | m, n     => inner call ackermann(m, n-1):
                    (m, n-1) <_lex (m, n) by n-1 < n
                  outer call ackermann(m-1, _):
                    (m-1, _) <_lex (m, n) by m-1 < m
conclude by well_founded_recursion(h_wf, h_dec)
\end{lstlisting}
\end{example}

\begin{definition}[Well-founded recursion principle]
Given a well-founded order $(W, \prec)$ and a measure $\mu$, the recursion
principle gives:
\[
  \frac{
    \Gamma \vdash (W, \prec) \text{ well-founded}
    \qquad
    \Gamma, x : A \vdash \mu(x) : W
    \qquad
    \Gamma, x : A, \forall y.\, \mu(y) \prec \mu(x) \to B(y) \vdash e : B(x)
  }{
    \Gamma \vdash \mathsf{wfrec}(\mu, \lambda x.\, e) : \Pi(x : A).\, B(x)
  }
\]
\end{definition}

\begin{remark}[\fstar{} comparison]
\fstar{} uses an identical well-founded recursion mechanism via the
\lstinline[style=fstar]|decreases| keyword. The key difference: in \synhopy{},
the measure can often be inferred from the loop/recursion structure and verified
by Z3, whereas \fstar{} requires explicit annotation in all cases. For nested
recursion (like Ackermann's inner call), both systems require explicit
lexicographic measures.
\end{remark}


\section{Partial Functions and Optional Return Types}
\label{sec:partial}

Many Python functions are \emph{partial}: they are undefined or raise exceptions
on some inputs. \synhopy{} models these using the \texttt{Optional} type.

\begin{definition}[Partial function type]
A partial function from $A$ to $B$ is typed as:
\[
  f : A \to B + \mathsf{None} \defeq A \to \mathsf{Optional}[B]
\]
The total function subtype is recovered by refinement:
\[
  f : \{x : A \mid \mathsf{defined}(f, x)\} \to B
\]
\end{definition}

\begin{example}[Safe division]
\begin{lstlisting}[style=synhopy]
from typing import Optional

@guarantee("returns a/b if b != 0, else None")
def safe_div(a: float, b: float) -> Optional[float]:
    if b == 0:
        return None
    return a / b
\end{lstlisting}

The proof obligation reduces to a case split:
\begin{lstlisting}[style=proof]
have h_zero: b = 0 => safe_div(a, b) = None
    by unfold safe_div then cases b = 0
have h_nonzero: b != 0 => safe_div(a, b) = Some(a / b)
    by unfold safe_div then cases b != 0
conclude by cases b
\end{lstlisting}
\end{example}

\begin{definition}[Lifting partial functions]
Given $f : A \to \mathsf{Optional}[B]$ and $g : B \to \mathsf{Optional}[C]$, their
monadic composition is:
\[
  (g \circ_? f)(x) \defeq
  \begin{cases}
    g(y) & \text{if } f(x) = \mathsf{Some}(y) \\
    \mathsf{None} & \text{if } f(x) = \mathsf{None}
  \end{cases}
\]
This corresponds to the Kleisli composition for the $\mathsf{Maybe}$ monad.
\end{definition}


\section{Precondition/Postcondition Propagation}
\label{sec:prepost-prop}

When pure functions are composed, their preconditions and postconditions propagate
according to the following rules.

\begin{definition}[Function specification]
A specification for function $f : A \to B$ is a pair $(\Pre_f, \Post_f)$ where:
\begin{align}
  \Pre_f &: A \to \Bool \\
  \Post_f &: A \to B \to \Bool
\end{align}
The specification asserts: $\forall x.\; \Pre_f(x) \implies \Post_f(x, f(x))$.
\end{definition}

\begin{theorem}[Sequential composition]
\label{thm:seq-compose}
Given $f : A \to B$ with spec $(\Pre_f, \Post_f)$ and $g : B \to C$ with spec
$(\Pre_g, \Post_g)$, the composition $g \circ f : A \to C$ has spec:
\begin{align}
  \Pre_{g \circ f}(x) &\defeq \Pre_f(x) \;\wedge\; \forall y.\, \Post_f(x, y) \implies \Pre_g(y) \\
  \Post_{g \circ f}(x, z) &\defeq \exists y.\, \Post_f(x, y) \;\wedge\; \Post_g(y, z)
\end{align}
\end{theorem}
\begin{proof}
Assume $\Pre_{g \circ f}(x)$. Then $\Pre_f(x)$ holds, so $\Post_f(x, f(x))$.
By the second conjunct, $\Pre_g(f(x))$, so $\Post_g(f(x), g(f(x)))$.
Taking $y = f(x)$ and $z = g(f(x))$ gives $\Post_{g \circ f}(x, z)$.
\end{proof}

\begin{example}[Pipeline verification]
\begin{lstlisting}[style=synhopy]
@guarantee("returns non-negative value")
def abs_val(x: int) -> int:
    return x if x >= 0 else -x

@guarantee("returns integer square root")
def isqrt(n: int) -> int:
    assume("n >= 0")
    r = int(n ** 0.5)
    return r

# Composition: isqrt(abs_val(x)) -- pre of isqrt met by post of abs_val
@guarantee("returns integer square root of |x|")
def sqrt_abs(x: int) -> int:
    return isqrt(abs_val(x))
\end{lstlisting}

The proof obligation for the composition:
\begin{lstlisting}[style=proof]
-- Show: post(abs_val) => pre(isqrt)
have h1: forall x. abs_val(x) >= 0
    by cases x >= 0 then omega
-- This satisfies assume("n >= 0") in isqrt
have h2: pre(isqrt) = (n >= 0)
    by unfold isqrt.assume
have h_chain: forall x. post(abs_val, x) => pre(isqrt, abs_val(x))
    by intro x then apply h1
conclude by compose_spec(abs_val, isqrt, h_chain)
\end{lstlisting}
\end{example}

\begin{theorem}[Weakest-precondition propagation]
\label{thm:wp-prop}
The weakest precondition transformer for composition satisfies:
\[
  \wp(g \circ f, Q) = \wp(f, \wp(g, Q))
\]
where $\wp(h, Q)(x) \defeq \Pre_h(x) \wedge \forall y.\, \Post_h(x, y) \implies Q(y)$.
\end{theorem}

\begin{remark}[\fstar{} comparison]
\fstar{} propagates pre/postconditions through its effect system using Dijkstra
monads, which compute weakest preconditions automatically for the \lstinline[style=fstar]|Tot|
effect. The \fstar{} approach is more general (it handles stateful effects), but
for $\Pure$ functions, \synhopy{}'s propagation is equivalent and often more
readable. The proof burden comparison:
\begin{center}
\begin{tabular}{lcc}
  \toprule
  \textbf{Metric} & \textbf{\fstar{}} & \textbf{\synhopy{}} \\
  \midrule
  Lines of spec per function & 3--5 & 1--2 \\
  Manual proof for composition & 5--15 & 0 (auto) \\
  Termination annotation & always manual & often inferred \\
  Total for 3-function pipeline & $\sim$45 & $\sim$8 \\
  \bottomrule
\end{tabular}
\end{center}
The savings come from LLM-generated proofs and automatic pre/post chaining.
\end{remark}


% ──────────────────────────────────────────────────────────────────
\chapter{Duck Typing}
\label{ch:duck}

\section{The Duck-Typing Problem}

Consider:
\begin{lstlisting}[style=python]
def make_noise(animal):
    return animal.quack()

class Duck:
    def quack(self): return "quack!"

class Person:
    def quack(self): return "I'm quacking!"

# Both work:
make_noise(Duck())    # "quack!"
make_noise(Person())  # "I'm quacking!"
\end{lstlisting}

The function \lstinline[style=python]{make_noise} works for any object with a
\lstinline[style=python]{quack} method. Its type is not ``takes a Duck'' or
``takes a Person''---it's ``takes anything quackable.''

\section{Typing Rule for Duck-Typed Functions}

\begin{definition}[Duck-type protocol]
A \emph{protocol} $\mathcal{P}$ is a set of method signatures:
\[
  \mathcal{P} = \{(m_1, \tau_1), (m_2, \tau_2), \ldots, (m_k, \tau_k)\}
\]
where $m_i$ is a method name and $\tau_i : A_i \to B_i$ is its type.
\end{definition}

\begin{definition}[Protocol satisfaction]
An object $x : \PyObj$ \emph{satisfies} protocol $\mathcal{P}$, written
$x \models \mathcal{P}$, if for every $(m_i, \tau_i) \in \mathcal{P}$:
\[
  \mathsf{hasattr}(x, m_i) \;\wedge\; \mathsf{getattr}(x, m_i) : \tau_i
\]
\end{definition}

\[
  \frac{
    \Gamma \vdash f : \Pi(x : \PyObj \mid x \models \mathcal{P}).\, B(x)
    \qquad
    \Gamma \vdash a : \PyObj
    \qquad
    \Gamma \vdash a \models \mathcal{P}
  }{
    \Gamma \vdash f(a) : B(a)
  }
\]

\section{Proving Duck-Type Equivalence}

\begin{theorem}[Duck-type equivalence is decidable for finite interfaces]
If $A$ and $B$ have finitely many methods (which all Python classes do), then
$A \simeq_{\mathsf{duck}} B$ is decidable: enumerate the methods, check each
for extensional equivalence (via Z3 for arithmetic methods, via testing for
others).
\end{theorem}

\begin{example}[Full proof of duck equivalence]
\begin{lstlisting}[style=proof]
-- Proof: Duck ~=_duck Person (in the quack protocol)
have h_quack: Duck.quack =_{str -> str} Person.quack
    by ext self.
       Duck.quack(self) = "quack!"     by unfold
       Person.quack(self) = "I'm quacking!"  by unfold
       -- These are NOT equal! Duck ~=_duck Person
       -- only in the protocol {quack : self -> str}
       -- (type-compatible, not value-equal)
\end{lstlisting}
Wait---\lstinline[style=python]{Duck.quack} and \lstinline[style=python]{Person.quack}
return \emph{different strings}. So they are NOT duck-type equivalent in the
value-equality sense. They are only \emph{type-compatible}: both have
a \texttt{quack} method returning \texttt{str}.

This illustrates an important distinction:
\begin{itemize}
  \item \textbf{Protocol compatibility}: $A$ and $B$ both have method $m$ with
    compatible types. Sufficient for ``$f(A)$ is well-typed and $f(B)$ is well-typed.''
  \item \textbf{Behavioral equivalence}: $A.m$ and $B.m$ are extensionally equal.
    Required for ``$f(A) = f(B)$.''
\end{itemize}
\synhopy{} distinguishes these via path types:
\begin{itemize}
  \item Protocol compatibility: $A$ and $B$ have the same method \emph{types}
  \item Behavioral equivalence: $A$ and $B$ have the same method \emph{values}
    (a stronger path)
\end{itemize}
\end{example}

\begin{remark}[\fstar{} comparison]
\fstar{} can only express protocol compatibility via typeclasses. It has no
built-in mechanism for behavioral equivalence of methods---you'd need to write
extensional equality lemmas manually for each method pair. \synhopy{}'s path
types capture both levels naturally: a path between classes IS a proof of
behavioral equivalence.
\end{remark}


\section{Structural Subtyping versus Nominal Subtyping}
\label{sec:structural-vs-nominal}

Python's type system exhibits a tension between two subtyping disciplines.

\begin{definition}[Nominal subtyping]
Under \emph{nominal subtyping}, $A \leq B$ if and only if $A$ is declared to
extend $B$ via inheritance. The identity of the class name matters:
\[
  A \leq_{\mathrm{nom}} B \;\Longleftrightarrow\; B \in \MRO(A)
\]
\end{definition}

\begin{definition}[Structural subtyping]
Under \emph{structural subtyping}, $A \leq B$ if and only if $A$ has all the
methods and attributes required by $B$, with compatible types:
\[
  A \leq_{\mathrm{struct}} B \;\Longleftrightarrow\;
  \forall (m, \tau) \in \mathsf{methods}(B).\;
  \exists \tau'.\; (m, \tau') \in \mathsf{methods}(A) \;\wedge\; \tau' \leq \tau
\]
\end{definition}

\begin{theorem}[Nominal implies structural]
If $A \leq_{\mathrm{nom}} B$, then $A \leq_{\mathrm{struct}} B$.
The converse does not hold in general.
\end{theorem}
\begin{proof}
By Python's inheritance semantics, if $B \in \MRO(A)$, then $A$ inherits all
methods of $B$ (possibly overriding them). Each override must be type-compatible
by the Liskov substitution principle, so the structural condition is satisfied.
For the converse, consider a class with identical methods but no inheritance
relationship.
\end{proof}

\begin{example}[Structural subtype without inheritance]
\begin{lstlisting}[style=python]
class FileWriter:
    def write(self, data: str) -> int: ...
    def flush(self) -> None: ...

class StringBuffer:
    def write(self, data: str) -> int: ...
    def flush(self) -> None: ...

# StringBuffer <={struct} FileWriter but not <={nom} FileWriter
\end{lstlisting}
In \synhopy{}, both classes satisfy the protocol
$\{(\mathtt{write}, \mathsf{str} \to \mathsf{int}),\;
  (\mathtt{flush}, () \to \PyNone)\}$
and are therefore interchangeable in any function that uses only these methods.
\end{example}

\begin{remark}[\fstar{} comparison]
\fstar{} uses purely nominal subtyping through its type system. Structural
compatibility can only be expressed through typeclasses, which require explicit
instance declarations. \synhopy{} supports both disciplines: nominal subtyping
through the class hierarchy and structural subtyping through protocol inference.
\end{remark}


\section{Protocol Inference Algorithm}
\label{sec:protocol-inference}

Given a function body, \synhopy{} can infer the \emph{minimal protocol} that
its parameters must satisfy.

\begin{definition}[Minimal protocol]
The \emph{minimal protocol} $\mathcal{P}_{\min}(f, x)$ for parameter $x$ of function
$f$ is the smallest set of method signatures such that $f$ is well-typed whenever
$x \models \mathcal{P}_{\min}(f, x)$:
\[
  \mathcal{P}_{\min}(f, x) = \bigcap \{ \mathcal{P} \mid
    \forall v.\; v \models \mathcal{P} \implies f[x \mapsto v] \text{ is well-typed} \}
\]
\end{definition}

\begin{construction}[Protocol inference]
\label{con:protocol-infer}
The inference algorithm traverses the AST of $f$ and collects all attribute
accesses on parameter $x$:
\begin{enumerate}
  \item For each expression $x.m(\ldots)$, add $(m, \mathsf{infer\_type}(\ldots))$
    to $\mathcal{P}$.
  \item For each expression $x.a$ (attribute access), add $(a, \mathsf{infer\_type}(x.a))$
    to $\mathcal{P}$.
  \item For each expression $x[k]$, add $(\mathtt{\_\_getitem\_\_}, K \to V)$
    to $\mathcal{P}$ where $K, V$ are inferred from context.
  \item For each expression $\mathtt{len}(x)$, add $(\mathtt{\_\_len\_\_}, () \to \mathsf{int})$.
  \item Close under callees: if $x$ is passed to $g(x)$, add
    $\mathcal{P}_{\min}(g, x)$ to $\mathcal{P}$.
\end{enumerate}
\end{construction}

\begin{example}[Inferred protocol]
\begin{lstlisting}[style=python]
def process(stream):
    header = stream.read(4)
    stream.seek(0)
    for chunk in stream:
        yield chunk.decode('utf-8')
\end{lstlisting}
Inferred protocol for \texttt{stream}:
\[
  \mathcal{P}_{\min} = \{
    (\mathtt{read}, \mathsf{int} \to \mathsf{bytes}),\;
    (\mathtt{seek}, \mathsf{int} \to \PyNone),\;
    (\mathtt{\_\_iter\_\_}, () \to \mathsf{Iterator}[\mathsf{bytes}])
  \}
\]
\end{example}

\begin{remark}[\fstar{} comparison]
\fstar{} has no protocol inference mechanism. Function parameters must have
fully declared types. To achieve similar flexibility, \fstar{} users must define
a typeclass and write explicit instances---often 10--15 lines of boilerplate per
protocol. \synhopy{}'s inference eliminates this overhead entirely.
\end{remark}


\section{Multiple Protocol Composition}
\label{sec:multi-protocol}

Objects often need to satisfy multiple protocols simultaneously.

\begin{definition}[Protocol intersection]
The \emph{intersection} of protocols $\mathcal{P}_1$ and $\mathcal{P}_2$ is:
\[
  \mathcal{P}_1 \cap \mathcal{P}_2 \defeq
  \{(m, \tau) \mid (m, \tau) \in \mathcal{P}_1 \;\vee\; (m, \tau) \in \mathcal{P}_2\}
\]
An object $x$ satisfies $\mathcal{P}_1 \cap \mathcal{P}_2$ if and only if
$x \models \mathcal{P}_1$ and $x \models \mathcal{P}_2$.
\end{definition}

\begin{example}[Multi-protocol function]
\begin{lstlisting}[style=python]
from typing import Protocol

class Readable(Protocol):
    def read(self, n: int) -> bytes: ...

class Seekable(Protocol):
    def seek(self, pos: int) -> None: ...

def read_from_start(f: Readable & Seekable) -> bytes:
    f.seek(0)
    return f.read(-1)
\end{lstlisting}
The inferred protocol is
$\mathcal{P}_{\mathtt{Readable}} \cap \mathcal{P}_{\mathtt{Seekable}}$.
\end{example}

\begin{theorem}[Protocol lattice]
The set of protocols over a fixed universe of method names forms a bounded lattice
under the subset ordering:
\begin{itemize}
  \item $\bot = \emptyset$ (the empty protocol, satisfied by every object)
  \item $\top = \mathsf{AllMethods}$ (satisfied by no finite object)
  \item Meet: $\mathcal{P}_1 \wedge \mathcal{P}_2 = \mathcal{P}_1 \cup \mathcal{P}_2$
    (more methods required = more restrictive)
  \item Join: $\mathcal{P}_1 \vee \mathcal{P}_2 = \mathcal{P}_1 \cap_{\mathrm{set}} \mathcal{P}_2$
    (fewer methods required = more permissive)
\end{itemize}
\end{theorem}

\begin{remark}[\fstar{} comparison]
\fstar{} composes typeclasses via multiple typeclass constraints, e.g.,
\lstinline[style=fstar]|val f: #a:Type -> readable a -> seekable a -> ...|.
This requires the user to thread typeclass dictionaries. \synhopy{} composes
protocols by simple set union, with no additional annotation burden.
\end{remark}


\section{Runtime Protocol Checking}
\label{sec:runtime-protocol}

Python's \texttt{typing.runtime\_checkable} decorator enables checking protocol
conformance at runtime via \texttt{isinstance}.

\begin{definition}[Runtime-checkable protocol]
A protocol $\mathcal{P}$ is \emph{runtime-checkable} if the predicate
$x \models \mathcal{P}$ can be evaluated at runtime. This requires that all
methods in $\mathcal{P}$ are detectable via \lstinline[style=python]|hasattr|:
\[
  x \models_{\mathrm{rt}} \mathcal{P} \;\Longleftrightarrow\;
  \bigwedge_{(m_i, \tau_i) \in \mathcal{P}} \mathsf{hasattr}(x, m_i)
\]
Note that runtime checking verifies method \emph{existence} but not type
correctness of signatures.
\end{definition}

\begin{example}[Runtime-checkable protocol in \synhopy{}]
\begin{lstlisting}[style=synhopy]
from typing import Protocol, runtime_checkable

@runtime_checkable
class Printable(Protocol):
    def __str__(self) -> str: ...

@guarantee("returns string representation if printable")
def safe_print(obj: object) -> str:
    check("isinstance(obj, Printable)")
    if isinstance(obj, Printable):
        return str(obj)
    return repr(obj)
\end{lstlisting}

The verification generates a refinement:
\begin{lstlisting}[style=proof]
have h_check: isinstance(obj, Printable)
    => hasattr(obj, "__str__")
    by runtime_checkable_spec(Printable)
have h_call: hasattr(obj, "__str__") => str(obj) : str
    by protocol_satisfaction(__str__, self -> str)
conclude by trans(h_check, h_call)
\end{lstlisting}
\end{example}

\begin{proposition}[Soundness gap of runtime checking]
\label{prop:runtime-gap}
Runtime protocol checking is \emph{necessary but not sufficient} for
protocol satisfaction: $x \models_{\mathrm{rt}} \mathcal{P}$ implies method
existence but not type-correct signatures. Formally:
\[
  x \models \mathcal{P} \;\implies\; x \models_{\mathrm{rt}} \mathcal{P}
  \qquad\text{but}\qquad
  x \models_{\mathrm{rt}} \mathcal{P} \;\not\!\!\implies\; x \models \mathcal{P}
\]
\end{proposition}

\begin{remark}[\fstar{} comparison]
\fstar{} has no concept of runtime type checking; all type information is erased
at extraction time. A runtime \lstinline[style=python]|isinstance| check
corresponds to a typecase that \fstar{} cannot express. \synhopy{} bridges
static verification and runtime checks by treating
\lstinline[style=python]|isinstance| as evidence that refines the type in
subsequent code---similar to TypeScript's type narrowing, but verified.
\end{remark}


\section{Duck Typing for Dunder Methods}
\label{sec:dunder-duck}

Python's ``dunder'' (double-underscore) methods define operator overloading and
built-in protocol support. These are central to Pythonic duck typing.

\begin{definition}[Dunder protocol families]
The major dunder protocol families are:
\begin{enumerate}
  \item \textbf{Arithmetic}: $\mathcal{P}_{\mathrm{arith}} = \{
    \mathtt{\_\_add\_\_}, \mathtt{\_\_sub\_\_}, \mathtt{\_\_mul\_\_},
    \mathtt{\_\_truediv\_\_}, \ldots\}$
  \item \textbf{Iteration}: $\mathcal{P}_{\mathrm{iter}} = \{
    \mathtt{\_\_iter\_\_}, \mathtt{\_\_next\_\_}\}$
  \item \textbf{Container}: $\mathcal{P}_{\mathrm{cont}} = \{
    \mathtt{\_\_len\_\_}, \mathtt{\_\_getitem\_\_}, \mathtt{\_\_contains\_\_}\}$
  \item \textbf{Context manager}: $\mathcal{P}_{\mathrm{ctx}} = \{
    \mathtt{\_\_enter\_\_}, \mathtt{\_\_exit\_\_}\}$
  \item \textbf{Comparison}: $\mathcal{P}_{\mathrm{cmp}} = \{
    \mathtt{\_\_eq\_\_}, \mathtt{\_\_lt\_\_}, \mathtt{\_\_le\_\_},
    \mathtt{\_\_hash\_\_}\}$
\end{enumerate}
\end{definition}

\begin{example}[Arithmetic duck typing]
\begin{lstlisting}[style=synhopy]
@guarantee("returns sum of elements")
def total(items):
    assume("all elements support __add__")
    result = items[0]
    for x in items[1:]:
        result = result + x  # calls __add__
    return result
\end{lstlisting}

Inferred protocol for elements of \texttt{items}:
\[
  \mathcal{P}_{\min} = \{
    (\mathtt{\_\_add\_\_}, \mathsf{Self} \to \mathsf{Self}),\;
    (\mathtt{\_\_getitem\_\_}, \mathsf{int} \to \mathsf{Self})
  \}
\]
This works for \texttt{int}, \texttt{float}, \texttt{str}, \texttt{list},
or any user class with \texttt{\_\_add\_\_}.
\end{example}

\begin{theorem}[Dunder protocol coherence]
\label{thm:dunder-coherence}
If a class $C$ implements the full arithmetic protocol $\mathcal{P}_{\mathrm{arith}}$
and satisfies the ring axioms (associativity, commutativity of $+$, distributivity),
then the \synhopy{} verification of arithmetic expressions over $C$ reduces to
Z3 arithmetic modulo the ring axioms.
\end{theorem}

\begin{remark}[\fstar{} comparison]
\fstar{} handles operator overloading through type-level dispatch, not duck
typing. To verify that a function works for multiple numeric types, an \fstar{}
user must define a ``numeric'' typeclass and instantiate it for each type---often
20+ lines of boilerplate. \synhopy{} infers the dunder protocol from usage and
verifies it structurally.
\end{remark}


% ──────────────────────────────────────────────────────────────────
\chapter{Monkey-Patching}
\label{ch:monkey}

\section{The Monkey-Patching Problem}

\begin{lstlisting}[style=python]
import json

# Replace json.dumps with a version that always sorts keys
_original = json.dumps
json.dumps = lambda obj, **kw: _original(obj, sort_keys=True, **kw)
\end{lstlisting}

After this code executes, \lstinline[style=python]{json.dumps} has a different
implementation. Any verification system that models types as static entities
is simply \emph{wrong} about the program from this point forward.

\section{Typing Rules for Monkey-Patching}

\begin{definition}[Monkey-patch operation]
\[
  \mathsf{setattr}(C, m, v) : \PyClass \to \mathsf{String} \to \PyFunc \to \PyClass
\]
with computation rule:
\begin{align}
  \mathsf{getattr}(\mathsf{setattr}(C, m, v), m) &\equiv v \\
  \mathsf{getattr}(\mathsf{setattr}(C, m, v), n) &\equiv \mathsf{getattr}(C, n) \quad \text{for } n \neq m
\end{align}
\end{definition}

\begin{definition}[Monkey-patch path]
The monkey-patch path constructor gives:
\[
  \mathsf{Patch}(C, m, v, \pi) : C =_{\PyClass} C[m := v]
\]
where $\pi$ is a proof that the behavioral contract is preserved or explicitly
changed.
\end{definition}

\section{Specification Checking for Monkey-Patched Code}

To verify code that uses a monkey-patched function:

\begin{lstlisting}[style=synhopy]
@guarantee("returns JSON string with sorted keys")
def serialize(data: dict) -> str:
    return json.dumps(data)
\end{lstlisting}

The proof obligation:
\begin{lstlisting}[style=proof]
-- After monkey-patch: json.dumps = original(sort_keys=True)
have h_patch: json.dumps = lambda obj, **kw:
                  _original(obj, sort_keys=True, **kw)
    by structural  -- this is the monkey-patch itself
have h_sorted: sort_keys=True => keys are sorted in output
    by apply json_sort_keys_spec
conclude by trans(h_patch, h_sorted)
\end{lstlisting}

\begin{remark}[\fstar{} comparison]
\fstar{} cannot reason about monkey-patching at all. The type system assumes that
function implementations are fixed at compile time. An \fstar{}-based Python
verifier would need to either (a) forbid monkey-patching entirely (ruling out
large classes of real Python code) or (b) model the entire program as a state
machine where function implementations are part of the state (resulting in an
explosion of proof obligations). \synhopy{}'s path-based approach handles
monkey-patching with a single proof obligation per patch.
\end{remark}


\section{Temporal Reasoning: Specs Before and After the Patch}
\label{sec:temporal-patch}

Monkey-patching introduces a temporal dimension: the specification of a function
\emph{changes} at the point where the patch is applied.

\begin{definition}[Temporal type environment]
A \emph{temporal type environment} $\Gamma_t$ is a function from program points
$t$ to type environments:
\[
  \Gamma_t : \mathsf{ProgramPoint} \to \Ctx
\]
A monkey-patch at point $t_0$ creates a discontinuity:
\[
  \Gamma_{t_0^+} = \Gamma_{t_0^-}[C.m \mapsto v']
\]
where $C.m$ is the patched method and $v'$ is the new implementation.
\end{definition}

\begin{example}[Before/after spec]
\begin{lstlisting}[style=synhopy]
# At time t0: json.dumps has its standard spec
@guarantee("returns JSON string")
def standard_serialize(data: dict) -> str:
    return json.dumps(data)  # uses original json.dumps

# Patch applied at t1:
_original = json.dumps
json.dumps = lambda obj, **kw: _original(obj, sort_keys=True, **kw)

# At time t2 > t1: json.dumps has the patched spec
@guarantee("returns JSON string with sorted keys")
def sorted_serialize(data: dict) -> str:
    return json.dumps(data)  # uses patched json.dumps
\end{lstlisting}

\synhopy{} tracks both specs:
\begin{lstlisting}[style=proof]
-- Before patch (t < t1):
have h_pre: spec(json.dumps, t) =
    (dict -> str, "returns JSON string")
    by structural

-- After patch (t >= t1):
have h_post: spec(json.dumps, t) =
    (dict -> str, "returns JSON string with sorted keys")
    by apply patch_spec(t1, json.dumps, sort_keys_wrapper)

-- Verify sorted_serialize uses post-patch spec
have h_call: t2 >= t1 => sorted_serialize uses h_post
    by temporal_ordering(t2, t1)
conclude by apply h_post
\end{lstlisting}
\end{example}

\begin{theorem}[Temporal soundness of patching]
\label{thm:temporal-sound}
If a patch at program point $t_0$ replaces $C.m$ with $v'$ and
$\pi : \Spec(v') \implies \Spec'$, then for all code at points $t > t_0$:
\[
  \Gamma_t \vdash C.m : \Spec'
\]
and for all code at points $t < t_0$:
\[
  \Gamma_t \vdash C.m : \Spec
\]
where $\Spec$ is the original specification.
\end{theorem}

\begin{remark}[\fstar{} comparison]
\fstar{} has no temporal type environment; all types are static. To model
pre/post-patch behavior, one would need to pass function implementations as
parameters, threading them through the entire call graph---resulting in a
combinatorial explosion. \synhopy{} handles this naturally via temporal
indexing of the type context.
\end{remark}


\section{Testing with Monkey-Patching}
\label{sec:test-monkey}

The most common use of monkey-patching in practice is for testing, via
\texttt{unittest.mock.patch} or \texttt{pytest}'s \texttt{monkeypatch} fixture.

\begin{example}[Mock-patching for tests]
\begin{lstlisting}[style=python]
from unittest.mock import patch

def get_user_name(user_id):
    response = requests.get(f"/api/users/{user_id}")
    return response.json()["name"]

# Test with mock:
@patch("requests.get")
def test_get_user_name(mock_get):
    mock_get.return_value.json.return_value = {"name": "Alice"}
    assert get_user_name(42) == "Alice"
\end{lstlisting}
\end{example}

\begin{definition}[Test patch contract]
A test patch of $C.m$ with mock $v'$ is \emph{sound} if the test verifies
behavior under the assumption that $C.m = v'$, and the contract of $v'$
is a refinement of the contract of $C.m$:
\[
  \Pre_{v'} \supseteq \Pre_{C.m} \;\wedge\; \Post_{v'} \subseteq \Post_{C.m}
\]
This ensures that any behavior verified under the mock also holds under the
real implementation.
\end{definition}

\begin{lstlisting}[style=proof]
-- Verify mock soundness for test_get_user_name
have h_mock: mock_get.return_value.json() = {"name": "Alice"}
    by mock_spec(mock_get)
have h_contract: post(mock_get) => post(requests.get)
    by check mock returns valid Response-like object
    -- mock satisfies the protocol {json: () -> dict}
have h_test: get_user_name(42) = "Alice" under mock
    by unfold get_user_name then apply h_mock
conclude by mock_soundness(h_contract, h_test)
\end{lstlisting}

\begin{remark}[\fstar{} comparison]
Testing in \fstar{} uses the standard approach of parameterizing code over
interfaces and providing test implementations. There is no need for mock-patching
because \fstar{} code is already parameterized. However, this means \fstar{}
cannot verify Python test code that uses \texttt{mock.patch}---it cannot even
\emph{express} the concept. \synhopy{} verifies both the production code and
the test's mock-patching as first-class operations.
\end{remark}


\section{Scope of Patches}
\label{sec:patch-scope}

Monkey-patches in Python can occur at three granularity levels, each with
different verification implications.

\begin{definition}[Patch scope levels]
\begin{enumerate}
  \item \textbf{Module-level}: $\mathsf{setattr}(\mathsf{module}, m, v)$ ---
    affects all code that imports the module.
  \item \textbf{Class-level}: $\mathsf{setattr}(\mathsf{cls}, m, v)$ ---
    affects all instances of the class and its subclasses.
  \item \textbf{Instance-level}: $\mathsf{setattr}(\mathsf{obj}, m, v)$ ---
    affects only the specific object instance.
\end{enumerate}
The scope determines the blast radius of the patch and the set of proof
obligations generated.
\end{definition}

\begin{proposition}[Scope containment]
\label{prop:scope-contain}
Instance-level patches are contained within class-level patches, which are
contained within module-level patches:
\[
  \mathsf{affected}(\mathsf{instance}) \subseteq
  \mathsf{affected}(\mathsf{class}) \subseteq
  \mathsf{affected}(\mathsf{module})
\]
\end{proposition}

\begin{example}[Instance-level patch]
\begin{lstlisting}[style=python]
class Dog:
    def speak(self): return "woof"

fido = Dog()
fido.speak = lambda self=None: "meow"  # instance-level patch

rex = Dog()
rex.speak()   # still "woof" -- unaffected
fido.speak()  # "meow" -- patched
\end{lstlisting}
The proof obligation is restricted to code paths using \texttt{fido}:
\begin{lstlisting}[style=proof]
have h_scope: patch(fido, "speak", meow_fn) affects only fido
    by instance_patch_scope
have h_rex: rex.speak = Dog.speak  -- unaffected
    by scope_exclusion(h_scope, rex != fido)
have h_fido: fido.speak = meow_fn
    by setattr_computation
\end{lstlisting}
\end{example}

\begin{remark}[\fstar{} comparison]
Since \fstar{} does not model object identity or instance-level mutation, it
cannot distinguish between class-level and instance-level patches. All three
scope levels would require the same heavyweight state-machine encoding.
\synhopy{}'s dependent types on object identity allow precise scope tracking.
\end{remark}


\section{Reversible Patches and the Patch Context Manager}
\label{sec:reversible-patch}

Python's \texttt{contextlib} and \texttt{unittest.mock.patch} support
reversible patches that are automatically undone when a context exits.

\begin{definition}[Reversible patch]
A \emph{reversible patch} is a monkey-patch paired with an undo operation:
\[
  \mathsf{RevPatch}(C, m, v) \defeq
  \left(\mathsf{setattr}(C, m, v),\;
        \mathsf{setattr}(C, m, \mathsf{getattr}(C, m))\right)
\]
The second component restores the original method.
\end{definition}

\begin{example}[Context-managed patch]
\begin{lstlisting}[style=python]
from unittest.mock import patch

with patch("os.getcwd", return_value="/fake/path"):
    assert os.getcwd() == "/fake/path"
# After the with block, os.getcwd is restored
assert os.getcwd() != "/fake/path"
\end{lstlisting}
\end{example}

\begin{theorem}[Reversible patch correctness]
\label{thm:rev-patch}
If $p = \mathsf{RevPatch}(C, m, v)$ is applied at program point $t_0$ and
reversed at $t_1 > t_0$, then:
\[
  \forall t < t_0 \;\cup\; t > t_1.\;
  \Gamma_t \vdash C.m \equiv \mathsf{original}
\]
\[
  \forall t_0 \leq t \leq t_1.\;
  \Gamma_t \vdash C.m \equiv v
\]
\end{theorem}
\begin{proof}
The context manager protocol guarantees that \lstinline[style=python]|__exit__|
is called even if an exception occurs within the \texttt{with} block. The
\texttt{patch} context manager stores the original value in its
\lstinline[style=python]|__enter__| method and restores it in
\lstinline[style=python]|__exit__|. By the setattr computation rules,
the original value is restored at $t_1$.
\end{proof}

\begin{lstlisting}[style=proof]
-- Verify reversibility
have h_save: original = getattr(os, "getcwd") at t0
    by patch.__enter__ stores original
have h_patch: getattr(os, "getcwd") = mock at t in [t0, t1]
    by setattr_computation
have h_restore: getattr(os, "getcwd") = original at t > t1
    by patch.__exit__ calls setattr(os, "getcwd", original)
    -- holds even under exception by context_manager_safety
conclude by temporal_split(t0, t1, h_save, h_patch, h_restore)
\end{lstlisting}


\section{Verifying Behavioral Subtyping Under Patches}
\label{sec:patch-subtype}

A critical concern is whether a monkey-patch preserves the behavioral subtyping
(Liskov substitution) of the patched class.

\begin{definition}[Behavioral subtyping preservation]
A patch $\mathsf{setattr}(C, m, v')$ \emph{preserves behavioral subtyping} if
for every context $\mathcal{E}[-]$ that expects $C$:
\[
  \forall x : C.\; \mathcal{E}[x] \text{ well-defined before patch}
  \implies
  \mathcal{E}[x] \text{ well-defined after patch}
\]
Formally, the patched class $C' = C[m := v']$ must satisfy:
\[
  \Pre_{C.m} \supseteq \Pre_{v'} \;\wedge\;
  \Post_{v'} \subseteq \Post_{C.m}
\]
\end{definition}

\begin{example}[Subtyping-preserving patch]
\begin{lstlisting}[style=synhopy]
import json

# Patch that adds logging but preserves behavior
_original = json.loads

@guarantee("returns same result as original json.loads")
def logged_loads(s, **kw):
    check("result equals _original(s, **kw)")
    print(f"Parsing JSON of length {len(s)}")
    result = _original(s, **kw)
    return result

json.loads = logged_loads
\end{lstlisting}

\begin{lstlisting}[style=proof]
-- Verify behavioral subtyping preservation
have h_pre: pre(logged_loads) = pre(_original)
    -- same precondition: s must be valid JSON
    by unfold logged_loads, same argument passed to _original
have h_post: forall s kw.
    logged_loads(s, **kw) = _original(s, **kw)
    by unfold logged_loads; print is IO, does not affect return
have h_liskov: json.loads patch preserves behavioral subtyping
    by behavioral_subtype(h_pre, h_post)
conclude by h_liskov
\end{lstlisting}
\end{example}

\begin{theorem}[Patch composition]
\label{thm:patch-compose}
If patches $p_1$ and $p_2$ both individually preserve behavioral subtyping,
their sequential composition $p_2 \circ p_1$ also preserves behavioral subtyping,
provided $p_2$'s verification is done against the post-$p_1$ spec.
\end{theorem}

\begin{remark}[\fstar{} comparison]
Behavioral subtyping in \fstar{} is enforced by the type system through
subtyping rules and refinement types. However, since \fstar{} does not model
runtime modification of implementations, it cannot reason about whether a
modification preserves subtyping. \synhopy{}'s patch paths carry the proof
that subtyping is preserved, making the verification obligation explicit and
dischargeable.
\end{remark}


% ──────────────────────────────────────────────────────────────────
\chapter{Decorators and Fiber Bundles}
\label{ch:decorators}

\section{Decorators as Fiber Bundles}

\begin{definition}[Decorator bundle]
A decorator $d$ defines a fiber bundle:
\begin{itemize}
  \item \textbf{Total space} $E$: the space of decorated functions
  \item \textbf{Base space} $B$: the space of undecorated functions
  \item \textbf{Projection} $p : E \to B$: strip the decorator
  \item \textbf{Fiber} $F_f = p^{-1}(f)$: the decoration state for function $f$
\end{itemize}
\end{definition}

\begin{example}[\texttt{@lru\_cache} as a fiber bundle]
\begin{lstlisting}[style=python]
from functools import lru_cache

@lru_cache(maxsize=128)
def fib(n):
    if n <= 1: return n
    return fib(n-1) + fib(n-2)
\end{lstlisting}

The bundle:
\begin{itemize}
  \item Base: the pure \texttt{fib} function
  \item Fiber: the cache state (an LRU dictionary)
  \item The fiber is \emph{contractible} for any fixed input $n$: the cache either
    has the value (hit) or computes it (miss), but either way the result is
    $\mathsf{fib\_pure}(n)$.
\end{itemize}
\end{example}

\begin{theorem}[Cache correctness]
\label{thm:cache}
If $f$ is pure (has effect $\Pure$), then $\texttt{lru\_cache}(f) \simeq f$
(the cached version is extensionally equal to the uncached version).
\end{theorem}
\begin{proof}
The fiber at each input $x$ is contractible: $\texttt{cache}[x]$ is either
$\mathsf{None}$ (miss, compute $f(x)$ and store) or $f(x)$ (hit, return stored
value). In both cases the result is $f(x)$. Since fibers are contractible, the
bundle is trivial, so the projection is an equivalence.
\end{proof}

\section{Decorator Commutativity}

\begin{definition}[Decorator commutativity]
Two decorators $d_1, d_2$ \emph{commute} if:
\[
  d_1 \circ d_2 \simeq d_2 \circ d_1
\]
as functions on the function space.
\end{definition}

\begin{theorem}[Commutativity criterion]
$d_1$ and $d_2$ commute if and only if the differential $d_2$ in the effect
spectral sequence (Section~\ref{sec:spectral}) vanishes: $d_2 = 0$.
\end{theorem}

\begin{remark}[\fstar{} comparison]
In \fstar{}, a decorator is simply a higher-order function
\lstinline[style=fstar]{val d : (a -> b) -> (a -> b)} with no additional
structure. To prove commutativity of two decorators, you must provide a manual
proof of extensional equality for the compositions---typically 30--50 lines of
F* per pair of decorators. \synhopy{} computes commutativity automatically via
the spectral sequence.
\end{remark}


\section{Class Decorators}
\label{sec:class-decorators}

Decorators in Python apply not only to functions but also to classes. A class
decorator receives a class object and returns a (possibly modified) class.

\begin{definition}[Class decorator bundle]
A class decorator $d_c$ defines a fiber bundle over the space of classes:
\begin{itemize}
  \item \textbf{Total space} $E$: the space of decorated classes
  \item \textbf{Base space} $B$: the space of undecorated classes
  \item \textbf{Projection} $p : E \to B$: strip the decorator, recovering the
    original class definition
  \item \textbf{Fiber} $F_C = p^{-1}(C)$: the methods and attributes added or
    modified by the decorator for class $C$
\end{itemize}
The key difference from function decorators: the fiber can contain \emph{multiple}
generated methods.
\end{definition}

\begin{example}[Singleton class decorator]
\begin{lstlisting}[style=python]
def singleton(cls):
    instances = {}
    def get_instance(*args, **kwargs):
        if cls not in instances:
            instances[cls] = cls(*args, **kwargs)
        return instances[cls]
    return get_instance
    
@singleton
class Database:
    def __init__(self, url):
        self.url = url
\end{lstlisting}

The fiber at \texttt{Database} is the instance cache dictionary. The bundle
projection maps the singleton wrapper back to the original \texttt{Database}
class.
\end{example}

\begin{theorem}[Singleton correctness]
\label{thm:singleton}
If $d = \mathsf{singleton}$ and $C$ is a class with a deterministic
\lstinline[style=python]|__init__|, then:
\[
  \forall \mathit{args}.\;
  d(C)(\mathit{args}) =_{\PyObj} d(C)(\mathit{args})
\]
That is, repeated calls return the same object (identity, not just equality).
\end{theorem}
\begin{proof}
On the first call, the instance is created and stored in the cache.
On subsequent calls, the cached instance is returned. Since dictionary
lookup is deterministic, the same object (by identity) is returned.
\end{proof}

\begin{remark}[\fstar{} comparison]
\fstar{} has no concept of class decorators. To model a singleton pattern
in \fstar{}, one would use a stateful module with a mutable reference---requiring
the \lstinline[style=fstar]|ST| effect and manual state threading. \synhopy{}
captures the singleton invariant as a contractibility condition on the fiber.
\end{remark}


\section{\texttt{@dataclass} as a Method-Generating Decorator}
\label{sec:dataclass}

The \texttt{@dataclass} decorator is the most structurally complex decorator in
the Python standard library: it inspects class annotations and generates
\texttt{\_\_init\_\_}, \texttt{\_\_repr\_\_}, \texttt{\_\_eq\_\_}, and optionally
\texttt{\_\_hash\_\_}, \texttt{\_\_lt\_\_}, etc.

\begin{definition}[Dataclass fiber]
For a class $C$ with fields $\{(f_1, \tau_1), \ldots, (f_k, \tau_k)\}$, the
dataclass decorator generates the fiber:
\[
  F_C^{\mathsf{dc}} = \{
    \mathtt{\_\_init\_\_},\;
    \mathtt{\_\_repr\_\_},\;
    \mathtt{\_\_eq\_\_},\;
    \mathtt{\_\_hash\_\_}^{?}
  \}
\]
where each method is determined by the field types and decorator options
(\texttt{eq}, \texttt{order}, \texttt{frozen}, etc.).
\end{definition}

\begin{example}[Dataclass verification]
\begin{lstlisting}[style=synhopy]
from dataclasses import dataclass

@dataclass(frozen=True)
class Point:
    x: float
    y: float

@guarantee("returns Euclidean distance between two points")
def distance(p: Point, q: Point) -> float:
    return ((p.x - q.x)**2 + (p.y - q.y)**2)**0.5
\end{lstlisting}

The dataclass decorator generates:
\begin{lstlisting}[style=proof]
-- Generated __init__:
have h_init: Point.__init__(self, x, y) =
    self.x := x; self.y := y
    by dataclass_init_spec(Point, [x: float, y: float])

-- Generated __eq__:
have h_eq: Point.__eq__(self, other) =
    (self.x == other.x) and (self.y == other.y)
    by dataclass_eq_spec(Point, [x, y])

-- Frozen => __hash__ is generated, __setattr__ raises
have h_frozen: forall p: Point. setattr(p, _, _) raises
    FrozenInstanceError
    by dataclass_frozen_spec(Point)

-- Verify distance uses only valid operations
have h_access: p.x, p.y, q.x, q.y are well-typed float
    by h_init and field_types(Point)
conclude by arithmetic_verification(h_access)
\end{lstlisting}
\end{example}

\begin{proposition}[Dataclass structural equivalence]
\label{prop:dc-equiv}
Two frozen dataclasses $C_1$ and $C_2$ with the same field names and types
are structurally equivalent as protocols:
\[
  \mathsf{fields}(C_1) = \mathsf{fields}(C_2) \implies
  C_1 \simeq_{\mathsf{struct}} C_2
\]
\end{proposition}

\begin{remark}[\fstar{} comparison]
\fstar{} uses \lstinline[style=fstar]|type| declarations for records and
generates projectors automatically. The \fstar{} approach is conceptually
similar, but \texttt{@dataclass} additionally generates comparison methods,
hashing, and string representations. In \fstar{}, these must be derived
manually or through typeclass instances---typically 15--20 extra lines per type.
\end{remark}


\section{\texttt{@property} as a Descriptor-Creating Decorator}
\label{sec:property}

The \texttt{@property} decorator creates a \emph{descriptor} object that
intercepts attribute access, enabling computed attributes with getter/setter
semantics.

\begin{definition}[Property descriptor bundle]
A property decorator on method $m$ of class $C$ creates a descriptor bundle:
\begin{itemize}
  \item \textbf{Getter fiber}: the function called on attribute read ($C.m$ access)
  \item \textbf{Setter fiber}: the function called on attribute write ($C.m = v$),
    if defined via \texttt{@m.setter}
  \item \textbf{Deleter fiber}: the function called on attribute deletion
    (\texttt{del C.m}), if defined via \texttt{@m.deleter}
\end{itemize}
The descriptor protocol maps attribute syntax to method calls.
\end{definition}

\begin{example}[Verified property]
\begin{lstlisting}[style=synhopy]
class Circle:
    def __init__(self, radius: float):
        assume("radius >= 0")
        self._radius = radius

    @property
    @guarantee("returns non-negative radius")
    def radius(self) -> float:
        return self._radius

    @radius.setter
    @guarantee("sets radius, preserving non-negativity")
    def radius(self, value: float):
        assume("value >= 0")
        self._radius = value

    @property
    @guarantee("returns area = pi * radius^2")
    def area(self) -> float:
        return 3.14159 * self._radius ** 2
\end{lstlisting}

\begin{lstlisting}[style=proof]
-- Verify getter/setter consistency
have h_get: forall c: Circle.
    c.radius = c._radius
    by unfold property.getter(radius)
have h_set: forall c: Circle, v >= 0.
    (c.radius := v; c.radius) = v
    by unfold property.setter(radius) then h_get
-- Verify area is consistent with radius
have h_area: forall c: Circle.
    c.area = 3.14159 * c.radius ** 2
    by unfold property.getter(area) then rw [h_get]
conclude by h_get, h_set, h_area
\end{lstlisting}
\end{example}

\begin{theorem}[Property invariant preservation]
\label{thm:property-inv}
If a class $C$ has an invariant $\Inv(C)$ and a property setter
$m.\mathsf{setter}$ preserves $\Inv$:
\[
  \forall c : C,\; v.\; \Inv(c) \;\wedge\; \Pre_{\mathsf{setter}}(v)
  \implies \Inv(c[m := v])
\]
then attribute writes through the property always maintain the class invariant.
\end{theorem}

\begin{remark}[\fstar{} comparison]
\fstar{} has no descriptor protocol. Computed fields are modeled as functions,
and there is no syntactic distinction between attribute access and function
calls. The \texttt{@property} pattern requires \synhopy{} to verify three
related functions (getter, setter, deleter) and their consistency---a proof
structure that has no \fstar{} analogue.
\end{remark}


\section{Decorator Stacking Semantics}
\label{sec:decorator-stacking}

When multiple decorators are stacked, they compose in bottom-up order:
\lstinline[style=python]|@d1 @d2 def f| means $f' = d_1(d_2(f))$.

\begin{definition}[Decorator stack]
A \emph{decorator stack} is a sequence $[d_1, d_2, \ldots, d_n]$ applied to
a base function $f$:
\[
  \mathsf{stack}([d_1, \ldots, d_n], f) \defeq d_1(d_2(\cdots d_n(f) \cdots))
\]
The $i$-th intermediate function is $f_i = d_i(d_{i+1}(\cdots d_n(f) \cdots))$.
\end{definition}

\begin{theorem}[Stack associativity]
\label{thm:stack-assoc}
Decorator stacking is associative in the sense that:
\[
  \mathsf{stack}([d_1, \ldots, d_n], f) =
  \mathsf{stack}([d_1, \ldots, d_k],\; \mathsf{stack}([d_{k+1}, \ldots, d_n], f))
\]
for any $1 \leq k < n$. This follows directly from function composition
associativity.
\end{theorem}

\begin{example}[Three-decorator stack]
\begin{lstlisting}[style=python]
import functools, time, logging

def log_calls(f):
    @functools.wraps(f)
    def wrapper(*args, **kwargs):
        logging.info(f"Calling {f.__name__}")
        return f(*args, **kwargs)
    return wrapper

def time_it(f):
    @functools.wraps(f)
    def wrapper(*args, **kwargs):
        start = time.time()
        result = f(*args, **kwargs)
        logging.info(f"Took {time.time()-start:.3f}s")
        return result
    return wrapper

@log_calls
@time_it
@functools.lru_cache(maxsize=128)
def expensive(n):
    return sum(range(n))
\end{lstlisting}
The stack is: $\mathsf{log\_calls}(\mathsf{time\_it}(\mathsf{lru\_cache}(\mathsf{expensive})))$.

The order matters:
\begin{enumerate}
  \item \texttt{lru\_cache} wraps the base function (caching layer)
  \item \texttt{time\_it} wraps the cached function (timing includes cache
    lookup, so cache hits report near-zero time)
  \item \texttt{log\_calls} wraps the timed function (all calls are logged)
\end{enumerate}
\end{example}

\begin{proposition}[Non-commutativity of logging and caching]
\label{prop:log-cache-noncommute}
$\mathsf{log\_calls} \circ \mathsf{lru\_cache} \not\simeq
 \mathsf{lru\_cache} \circ \mathsf{log\_calls}$.
The left side logs every call (including cache hits); the right side only logs
cache misses.
\end{proposition}

\begin{remark}[\fstar{} comparison]
In \fstar{}, verifying decorator stacking requires composing the proofs for each
layer manually. For a stack of $n$ decorators, this means $n-1$ composition
lemmas. \synhopy{}'s fiber bundle framework allows reasoning about the total
space directly, reducing the proof obligation to a single coherence check per
stack.
\end{remark}


\section{Preserving Function Metadata with \texttt{@functools.wraps}}
\label{sec:functools-wraps}

A common pitfall with decorators is that the wrapper function loses the
metadata (\texttt{\_\_name\_\_}, \texttt{\_\_doc\_\_}, \texttt{\_\_module\_\_})
of the original function. The \texttt{@functools.wraps} decorator copies
this metadata.

\begin{definition}[Metadata preservation]
A decorator $d$ \emph{preserves metadata} if:
\[
  \forall f.\;
  d(f).\mathtt{\_\_name\_\_} = f.\mathtt{\_\_name\_\_} \;\wedge\;
  d(f).\mathtt{\_\_doc\_\_} = f.\mathtt{\_\_doc\_\_} \;\wedge\;
  d(f).\mathtt{\_\_wrapped\_\_} = f
\]
\texttt{@functools.wraps(f)} on the wrapper function ensures this property.
\end{definition}

\begin{example}[Metadata loss without \texttt{wraps}]
\begin{lstlisting}[style=python]
def my_decorator(f):
    def wrapper(*args, **kwargs):
        return f(*args, **kwargs)
    return wrapper

@my_decorator
def greet(name):
    """Greet someone."""
    return f"Hello, {name}!"

# Without wraps: greet.__name__ == "wrapper", greet.__doc__ == None
# With wraps:    greet.__name__ == "greet",   greet.__doc__ == "Greet someone."
\end{lstlisting}
\end{example}

\begin{theorem}[\texttt{wraps} as section of the bundle]
\label{thm:wraps-section}
The \texttt{@functools.wraps} decorator provides a canonical section
$s : B \to E$ of the decorator bundle, such that the metadata projection
$\pi_{\mathrm{meta}} : E \to \mathsf{Metadata}$ satisfies:
\[
  \pi_{\mathrm{meta}} \circ d_{\mathsf{wraps}} = \pi_{\mathrm{meta}} \circ \mathsf{id}_B
\]
That is, the decorated function has the same metadata as the original.
\end{theorem}

\begin{remark}[\fstar{} comparison]
\fstar{} functions do not carry metadata at runtime (names and docs are erased
during extraction). This means metadata preservation is not a concern in
\fstar{}. However, for Python verification, metadata correctness is important
for reflection-based frameworks (e.g., Flask routes, pytest discovery).
\synhopy{} can verify metadata preservation as part of the decorator contract.
\end{remark}


\section{Worked Examples: Caching, Retry, and Timing}
\label{sec:decorator-examples}

We conclude with detailed verification proofs for three commonly used
decorator patterns.

\begin{example}[Retry decorator verification]
\begin{lstlisting}[style=python]
import functools, time

def retry(max_attempts=3, delay=1.0):
    def decorator(f):
        @functools.wraps(f)
        def wrapper(*args, **kwargs):
            for attempt in range(max_attempts):
                try:
                    return f(*args, **kwargs)
                except Exception:
                    if attempt == max_attempts - 1:
                        raise
                    time.sleep(delay)
        return wrapper
    return decorator

@retry(max_attempts=3, delay=0.5)
def fetch_data(url):
    return requests.get(url).json()
\end{lstlisting}

\begin{lstlisting}[style=proof]
-- Verify retry decorator correctness
have h_success: forall f, args.
    (exists i < max_attempts. f(args) succeeds on attempt i)
    => retry(f)(args) = f(args)  -- on the successful attempt
    by induction on i
    base: i = 0 => first try succeeds, return immediately
    step: attempt i fails, sleep(delay), try i+1, apply IH

have h_fail: forall f, args.
    (forall i < max_attempts. f(args) raises on attempt i)
    => retry(f)(args) raises
    by: last attempt raises, if-guard re-raises

have h_meta: retry(f).__name__ = f.__name__
    by functools_wraps_spec(f)

have h_bounded: retry terminates in at most
    max_attempts * (timeout(f) + delay) time
    by loop_bound(max_attempts) and sleep_bound(delay)
conclude by h_success, h_fail, h_meta, h_bounded
\end{lstlisting}
\end{example}

\begin{example}[Timing decorator verification]
\begin{lstlisting}[style=python]
import functools, time

def timed(f):
    @functools.wraps(f)
    def wrapper(*args, **kwargs):
        start = time.perf_counter()
        result = f(*args, **kwargs)
        elapsed = time.perf_counter() - start
        wrapper.last_elapsed = elapsed
        return result
    wrapper.last_elapsed = 0.0
    return wrapper

@timed
def compute(n):
    return sum(range(n))
\end{lstlisting}

\begin{lstlisting}[style=proof]
-- Verify timing decorator properties
have h_pure: forall f, args.
    timed(f)(args) = f(args)
    -- timing wrapper does not alter return value
    by unfold timed.wrapper; result = f(args); return result

have h_time: forall f, args.
    timed(f)(args); timed(f).last_elapsed >= 0.0
    by: elapsed = perf_counter() - start >= 0
        since perf_counter is monotonic

have h_fiber: the fiber at f is (last_elapsed : float)
    -- state space is a single non-negative float
    -- fiber is contractible: each call overwrites last_elapsed

conclude by h_pure, h_time, h_fiber
\end{lstlisting}
\end{example}

\begin{example}[Cache with expiry decorator]
\begin{lstlisting}[style=synhopy]
import functools, time

def timed_cache(ttl_seconds=60):
    def decorator(f):
        cache = {}
        @functools.wraps(f)
        def wrapper(*args):
            now = time.time()
            if args in cache:
                result, timestamp = cache[args]
                if now - timestamp < ttl_seconds:
                    return result
            result = f(*args)
            cache[args] = (result, now)
            return result
        return wrapper
    return decorator

@timed_cache(ttl_seconds=300)
@guarantee("returns weather data for city")
def get_weather(city: str) -> dict:
    return requests.get(f"/api/weather/{city}").json()
\end{lstlisting}

\begin{lstlisting}[style=proof]
-- Cache correctness with TTL
have h_fresh: forall f, args, t.
    (args in cache) /\ (t - cache[args].timestamp < ttl)
    => timed_cache(f)(args) = cache[args].result = f(args)
    by unfold wrapper; cache hit returns stored result
       which was f(args) at time of storage

have h_stale: forall f, args, t.
    (args in cache) /\ (t - cache[args].timestamp >= ttl)
    => timed_cache(f)(args) = f(args)  -- recomputed
    by unfold wrapper; TTL check fails, recompute

have h_miss: forall f, args.
    (args not in cache) => timed_cache(f)(args) = f(args)
    by unfold wrapper; cache miss, compute and store

-- Combined: the decorator always returns f(args),
-- possibly from cache if within TTL
have h_correct: forall f, args.
    timed_cache(f)(args) = f(args)
    -- assuming f is pure (deterministic)
    by cases h_fresh, h_stale, h_miss
conclude by h_correct
\end{lstlisting}
\end{example}

\begin{remark}[\fstar{} comparison]
Verifying these decorator patterns in \fstar{} requires:
\begin{itemize}
  \item \textbf{Retry}: modeling the exception effect with
    \lstinline[style=fstar]|EXN|, the delay with \lstinline[style=fstar]|IO|,
    and the bounded loop---approximately 40 lines of spec and proof.
  \item \textbf{Timing}: modeling mutable state for elapsed time with
    \lstinline[style=fstar]|ST|, and proving the wrapper is extensionally
    equal---approximately 25 lines.
  \item \textbf{Timed cache}: combining \lstinline[style=fstar]|ST| for cache
    state with \lstinline[style=fstar]|IO| for time access, proving
    cache coherence---approximately 50 lines.
\end{itemize}
In \synhopy{}, each proof is 10--15 lines, generated by the LLM, and verified
against the original Python code without translation.
\end{remark}


% ──────────────────────────────────────────────────────────────────
\chapter{Exceptions and Classifying Spaces}
\label{ch:exceptions}

\section{Python's Exception Hierarchy}

Python's exceptions form a tree rooted at \texttt{BaseException}:
\begin{verbatim}
BaseException
 +-- SystemExit
 +-- KeyboardInterrupt
 +-- Exception
      +-- ValueError
      +-- TypeError
      +-- LookupError
           +-- KeyError
           +-- IndexError
      +-- OSError
           +-- FileNotFoundError
           +-- PermissionError
\end{verbatim}

\section{Exceptions as Eilenberg--MacLane Spaces}

\begin{definition}[Exception classifying space]
The exception type $E$ has classifying space $K(\pi_1(\mathsf{ExcTree}), 1)$,
where $\pi_1(\mathsf{ExcTree})$ is the fundamental group of the exception tree.

The first cohomology $H^1(f, E)$ of a function $f$ with coefficients in $E$
classifies which exceptions $f$ can raise.
\end{definition}

\begin{definition}[Exception typing rule]
\[
  \frac{
    \Gamma \vdash f : A \to_{\varepsilon} B \qquad
    \varepsilon \text{ includes } \mathsf{Raises}(E_1, \ldots, E_k)
  }{
    \Gamma \vdash f : A \to_{\varepsilon} B + E_1 + \cdots + E_k
  }
\]
\end{definition}

\section{Exception Handling as Sections}

\begin{lstlisting}[style=python]
def safe_divide(a, b):
    try:
        return a / b
    except ZeroDivisionError:
        return float('inf')
\end{lstlisting}

The \texttt{try/except} block is a \emph{section} of the exception bundle:
\begin{itemize}
  \item The base space is the computation without exception handling
  \item The fiber over each exception type is the handler for that type
  \item A section selects a handler for each caught exception
\end{itemize}

\begin{lstlisting}[style=proof]
-- Proof: safe_divide never raises ZeroDivisionError
have h_try: a / b raises ZeroDivisionError when b = 0
    by apply div_zero_spec
have h_catch: except ZeroDivisionError catches all ZeroDivisionError
    by structural -- the except clause is syntactically present
have h_result: when b = 0, safe_divide(a, b) = inf
    by unfold safe_divide then rw [h_catch]
conclude: safe_divide is total (never raises)
    by cases b
    case b != 0: a / b succeeds by h_try
    case b = 0: returns inf by h_result
\end{lstlisting}

\begin{remark}[\fstar{} comparison]
\fstar{}'s \lstinline[style=fstar]{EXN} effect type is flat: a function either
may raise exceptions or it doesn't. There is no structure describing \emph{which}
exceptions, how they relate, or what the handlers do. To express
``\texttt{safe\_divide} catches \texttt{ZeroDivisionError} but propagates
\texttt{TypeError},'' you'd need to encode the exception hierarchy in the type
system manually.

\synhopy{}'s cohomological approach gives this for free: $H^1(f, E)$ automatically
classifies the exception behavior, and sections (handlers) are verified against
the fibration structure.
\end{remark}

\section{Exception Chaining}

Python 3 supports exception chaining (\texttt{raise X from Y}):
\begin{lstlisting}[style=python]
def process(data):
    try:
        parsed = json.loads(data)
    except json.JSONDecodeError as e:
        raise ValueError("Bad input") from e
\end{lstlisting}

\begin{definition}[Exception chain as path composition]
The chaining operator \texttt{from} creates a path in the exception space:
\[
  \mathsf{chain}(e_1, e_2) : e_1 \to e_2 \text{ in } K(\pi_1(\mathsf{ExcTree}), 1)
\]
This is a morphism recording the causal relationship.  The chain is accessible
at runtime via \lstinline[style=python]{e.__cause__}.
\end{definition}

\section{The \texttt{finally} Block}

\begin{definition}[\texttt{finally} as universal section]
A \texttt{finally} block is a section of the exception bundle that is
defined on \emph{every} fiber (every possible exception, including no exception):
\[
  \frac{
    \Gamma \vdash \texttt{try}: s_1 \quad
    \Gamma \vdash \texttt{finally}: s_2
  }{
    \forall \varepsilon.\; \Gamma \vdash s_2 \text{ executes under effect } \varepsilon
  }
\]
This universal quantification over exception paths is exactly the guarantee
that Python's runtime provides.
\end{definition}

\begin{theorem}[\texttt{finally} guarantee]
In \synhopy{}, if a function contains \texttt{try}: $s_1$ \texttt{finally}: $s_2$,
then $s_2$ is executed on all paths through $s_1$, including:
\begin{enumerate}
  \item Normal completion (no exception)
  \item Exception raised and not caught
  \item Exception raised and caught by an intervening \texttt{except}
  \item \texttt{return} from within the \texttt{try} block
  \item \texttt{break}/\texttt{continue} from within the \texttt{try} block
\end{enumerate}
\end{theorem}
\begin{proof}
Each case is a fiber of the exception bundle.  The \texttt{finally} block is
a universal section, so it is defined on every fiber.  Python's runtime
guarantees execution on all paths.
\end{proof}

\section{Exception Groups and Structured Concurrency}

Python~3.11 introduced \texttt{ExceptionGroup} and the \texttt{except*} syntax
for handling multiple simultaneously-raised exceptions, motivated by structured
concurrency patterns in \texttt{asyncio} task groups.

\begin{lstlisting}[style=python]
async def fetch_all(urls):
    async with asyncio.TaskGroup() as tg:
        tasks = [tg.create_task(fetch(url)) for url in urls]
    return [t.result() for t in tasks]
    # If multiple tasks fail, raises ExceptionGroup
\end{lstlisting}

\begin{definition}[Exception group type]
An exception group is a \emph{tree} of exceptions, not a flat list.
We type it as the free monad over the exception functor:
\[
  \mathsf{ExcGroup}(E) \defeq \mu X.\; E + (X \times X)
\]
The binary product $X \times X$ models the simultaneous occurrence of two
sub-groups.  A single exception is $\mathsf{inl}(e)$; a compound group is
$\mathsf{inr}(g_1, g_2)$.
\end{definition}

\begin{definition}[\texttt{except*} as partial section]
The \texttt{except*} handler selects a \emph{partial section} of the exception
group bundle: it handles some exception types while letting others propagate
as a \emph{remainder group}:
\[
  \frac{
    \Gamma \vdash g : \mathsf{ExcGroup}(E_1 + E_2 + \cdots + E_k) \qquad
    \Gamma \vdash h_i : E_i \to R \;\text{for}\; i \in S \subseteq \{1,\ldots,k\}
  }{
    \Gamma \vdash \texttt{except*}\;(h_i)_{i \in S} :
      \mathsf{ExcGroup}(E_j)_{j \notin S} + R^{|S|}
  }
\]
\end{definition}

\begin{example}[Handling partial exception groups]
\begin{lstlisting}[style=python]
try:
    results = await fetch_all(urls)
except* ValueError as eg:
    log("bad values", eg.exceptions)
except* TimeoutError as eg:
    retry(eg.exceptions)
# Any other exceptions propagate as a new ExceptionGroup
\end{lstlisting}
In the cohomological framework, each \texttt{except*} clause is a
section defined on a sub-bundle of the exception group fibration.
The remainder group is the \emph{quotient bundle} over the un-handled fibers.
\end{example}

\begin{remark}[\fstar{} comparison]
\fstar{} has no analogue of exception groups.  Its \lstinline[style=fstar]|EXN|
effect models a single exception channel; there is no native way to represent
multiple simultaneously-raised exceptions or partial handling.
Encoding exception groups in \fstar{} would require defining a custom inductive
type and manually threading the remainder group---losing all automation.
\synhopy{} models exception groups as a free monad with partial sections,
giving direct verification of structured concurrency patterns.
\end{remark}

\section{Custom Exception Classes as Dependent Records}

Python exceptions are classes that can carry arbitrary data fields:
\begin{lstlisting}[style=python]
class InsufficientFundsError(Exception):
    def __init__(self, account_id: str, requested: Decimal,
                 available: Decimal):
        self.account_id = account_id
        self.requested = requested
        self.available = available
        super().__init__(
            f"Account {account_id}: requested {requested}, "
            f"available {available}"
        )
\end{lstlisting}

\begin{definition}[Exception as dependent record]
A custom exception class with data fields is a \emph{dependent record type}
in the exception hierarchy:
\[
  \mathsf{InsufficientFundsError} \defeq \Sigma\!\left(
    \begin{array}{l}
      \mathsf{account\_id} : \mathsf{str}, \\
      \mathsf{requested} : \mathsf{Decimal}, \\
      \mathsf{available} : \mathsf{Decimal}, \\
      \mathsf{\_proof} : \mathsf{requested} > \mathsf{available}
    \end{array}
  \right)
\]
The proof component $\mathsf{\_proof}$ is a \synhopy{} refinement:
this exception \emph{can only be constructed} when the requested amount
exceeds the available balance.
\end{definition}

\begin{theorem}[Exception data consistency]
If an exception type $E$ is defined as a dependent record with proof
component $p : \phi(\vec{x})$, then every instance $e : E$ satisfies
$\phi(\vec{x})$.  This is guaranteed by the $\Sigma$-elimination rule.
\end{theorem}
\begin{proof}
By $\Sigma$-elimination, given $e : \Sigma(x_1 : A_1, \ldots, x_n : A_n,
p : \phi(\vec{x}))$, we have $\pi_{n+1}(e) : \phi(\pi_1(e), \ldots,
\pi_n(e))$.  The proof obligation $\phi$ is checked at construction time.
\end{proof}

\begin{remark}[\fstar{} comparison]
In \fstar{}, exception payloads are unstructured: the \lstinline[style=fstar]|exn|
type is an extensible variant with no refinement on the carried data.
One can define a refined record separately, but connecting it to the exception
mechanism requires manual encoding.  \synhopy{}'s dependent record approach
ensures that exception data invariants are verified at raise-site and
available at catch-site---automatically.
\end{remark}

\section{Exception Safety Levels}

Following the C++ exception safety taxonomy (adapted to Python), \synhopy{}
formalises three levels of exception safety as type-level guarantees.

\begin{definition}[Exception safety levels]
For a function $f : A \to_\varepsilon B$ operating on state $S$, the three
exception safety levels are:
\begin{enumerate}
  \item \textbf{No-throw guarantee} (\textsf{NoThrow}):
    $f$ never raises an exception.
    \[
      \mathsf{NoThrow}(f) \defeq H^1(f, E) = 0
      \quad\text{(trivial exception cohomology)}
    \]
  \item \textbf{Basic guarantee} (\textsf{BasicSafe}):
    if $f$ raises, $S$ is in a \emph{valid} (but possibly modified) state.
    \[
      \mathsf{BasicSafe}(f, I) \defeq
        \forall s.\; I(s) \Rightarrow
          \bigl(f(s) = \mathsf{Ok}(r, s') \wedge I(s')\bigr)
          \vee
          \bigl(f(s) = \mathsf{Err}(e, s') \wedge I(s')\bigr)
    \]
  \item \textbf{Strong guarantee} (\textsf{StrongSafe}):
    if $f$ raises, $S$ is \emph{unchanged} (rollback semantics).
    \[
      \mathsf{StrongSafe}(f) \defeq
        \forall s.\;
        f(s) = \mathsf{Err}(e, s') \Rightarrow s' = s
    \]
\end{enumerate}
\end{definition}

\begin{example}[Annotating exception safety]
\begin{lstlisting}[style=synhopy]
@guarantee("strong exception safety: balance unchanged on failure")
def transfer(src: Account, dst: Account, amount: Decimal) -> None:
    check("amount > 0")
    check("src.balance >= amount")
    src.withdraw(amount)     # may raise InsufficientFundsError
    try:
        dst.deposit(amount)  # may raise AccountFrozenError
    except AccountFrozenError:
        src.deposit(amount)  # rollback
        raise
\end{lstlisting}
\end{example}

\begin{lstlisting}[style=proof]
-- Proof: transfer has the strong exception safety guarantee
have h_withdraw: src.withdraw(amount) modifies src.balance
    by unfold withdraw
have h_deposit_fail: if dst.deposit raises AccountFrozenError,
    then src.deposit(amount) restores src.balance
    by unfold deposit then rw [h_withdraw]
have h_rollback: on exception path,
    src.balance' = src.balance - amount + amount = src.balance
    by omega with h_deposit_fail
conclude: StrongSafe(transfer)
    by cases dst.deposit(amount)
    case Ok: both balances updated (normal path)
    case Err: rollback restores src.balance by h_rollback
              dst.balance unchanged (deposit failed)
\end{lstlisting}

\begin{remark}[\fstar{} comparison]
\fstar{}'s \lstinline[style=fstar]|ST| effect tracks state modifications but
does not distinguish exception safety levels.  To express the strong guarantee
in \fstar{}, one would need a custom effect that tracks both the normal and
exceptional post-states---a non-trivial extension.  \synhopy{} provides
exception safety levels as built-in verification targets.
\end{remark}

\section{Exhaustive Exception Handling}

A common source of bugs is failing to handle all possible exceptions.
\synhopy{} can verify that exception handling is \emph{exhaustive}: every
exception that a called function may raise is either caught or explicitly
documented as propagating.

\begin{definition}[Exception coverage]
For a function $f$ calling sub-functions $g_1, \ldots, g_n$, define the
\emph{raised set} and \emph{caught set}:
\begin{align}
  \mathsf{Raised}(f) &\defeq \bigcup_{i=1}^{n} H^1(g_i, E) \\
  \mathsf{Caught}(f) &\defeq \{ E_j \mid \texttt{except}\;E_j\;\text{appears in } f \} \\
  \mathsf{Propagated}(f) &\defeq \mathsf{Raised}(f) \setminus \mathsf{Caught}(f)
\end{align}
Exception handling is \emph{exhaustive} when every element of
$\mathsf{Propagated}(f)$ is declared in $f$'s exception specification.
\end{definition}

\begin{theorem}[Exception coverage completeness]
If $\mathsf{Propagated}(f) \subseteq \mathsf{Declared}(f)$, then no
undocumented exception can escape $f$.  Equivalently, the exception
cohomology is fully accounted for:
\[
  H^1(f, E) = \mathsf{Caught}(f) \cup \mathsf{Declared}(f)
\]
\end{theorem}
\begin{proof}
By the long exact sequence in exception cohomology, every exception
raised by a sub-call either maps to zero (is caught) or maps to a
non-trivial cohomology class (propagates).  The declared set covers
all non-trivial classes by hypothesis.
\end{proof}

\begin{lstlisting}[style=synhopy]
@guarantee("all exceptions are caught or documented")
def process_order(order: Order) -> Receipt:
    """
    Raises:
        PaymentDeclinedError: if payment fails
    """
    assume("order.is_valid()")
    item = lookup_item(order.item_id)  # raises ItemNotFoundError
    try:
        charge = process_payment(order.payment)
    except PaymentDeclinedError:
        raise  # explicitly propagated (documented)
    except PaymentTimeoutError:
        charge = retry_payment(order.payment)
    return create_receipt(item, charge)
\end{lstlisting}

\begin{lstlisting}[style=proof]
-- Proof: exception handling in process_order is exhaustive
have h_lookup: Raised(lookup_item) = {ItemNotFoundError}
    by unfold lookup_item
-- But wait: ItemNotFoundError is not caught!
-- SynHoPy flags this as an ERROR:
--   ItemNotFoundError in Propagated(process_order)
--   but not in Declared(process_order)
-- Fix: add 'Raises: ItemNotFoundError' to docstring,
--   or add a try/except block.
\end{lstlisting}

\begin{remark}[\fstar{} comparison]
\fstar{} provides no mechanism for verifying exception coverage.
The \lstinline[style=fstar]|EXN| effect is a single bit: a function either may
raise or it cannot.  Which exceptions, whether they are caught, and whether
the coverage is complete are all invisible to the type system.
\synhopy{}'s exception cohomology makes coverage a first-class verification
target.
\end{remark}

\section{Exception Hierarchy as a Poset}

\begin{definition}[Exception poset]
The exception hierarchy forms a partially ordered set $(\mathcal{E}, \leq)$
where $E_1 \leq E_2$ iff $E_1$ is a subclass of $E_2$.  This poset has:
\begin{itemize}
  \item A top element $\top = \texttt{BaseException}$
  \item Finite meets: $E_1 \wedge E_2$ is their nearest common subclass (if it exists)
  \item Finite joins: $E_1 \vee E_2$ is their nearest common superclass
\end{itemize}
\end{definition}

\begin{definition}[Classifying space of the exception poset]
The \emph{nerve} of $(\mathcal{E}, \leq)$ viewed as a category gives the
classifying space $B\mathcal{E}$:
\begin{itemize}
  \item 0-simplices: exception classes $E_i$
  \item 1-simplices: subclass relations $E_i \leq E_j$
  \item $n$-simplices: chains $E_{i_0} \leq E_{i_1} \leq \cdots \leq E_{i_n}$
\end{itemize}
The homotopy type of $B\mathcal{E}$ encodes the ``shape'' of the exception
hierarchy.  For a tree-shaped hierarchy (Python's default), $B\mathcal{E}$
is contractible.
\end{definition}

\begin{theorem}[Contractibility of tree hierarchies]
If $(\mathcal{E}, \leq)$ is a tree (every element has a unique parent except
the root), then $B\mathcal{E}$ is contractible: $B\mathcal{E} \simeq *$.
\end{theorem}
\begin{proof}
A tree poset has a terminal object (the root) and admits a deformation
retraction to it.  By the nerve theorem, $B\mathcal{E} \simeq *$.
\end{proof}

\begin{proposition}[Diamond inheritance breaks contractibility]
If custom exceptions introduce diamond inheritance:
\begin{lstlisting}[style=python]
class NetworkError(IOError): pass
class TimeoutError(IOError, RuntimeError): pass
\end{lstlisting}
then $(\mathcal{E}, \leq)$ is no longer a tree, and $B\mathcal{E}$ may have
non-trivial $\pi_1$.  \synhopy{} detects this and warns that \texttt{except}
clause ordering becomes semantically significant.
\end{proposition}

\begin{remark}[\fstar{} comparison]
\fstar{}'s exception type \lstinline[style=fstar]|exn| is a flat extensible
variant with no hierarchical structure.  There is no notion of an exception
being a ``subclass'' of another, and catching a parent exception does not
automatically catch its children.  \synhopy{}'s poset model faithfully
captures Python's class-based exception hierarchy and reasons about
subclass-based matching.
\end{remark}

\section{Worked Example: Retry with Backoff}

We verify a production-style retry function that handles all exception paths.

\begin{lstlisting}[style=python]
import time, random
from decimal import Decimal

class RetryExhaustedError(Exception):
    def __init__(self, attempts: int, last_error: Exception):
        self.attempts = attempts
        self.last_error = last_error

def retry_with_backoff(func, max_retries=3, base_delay=1.0):
    """Retry func with exponential backoff.
    Raises:
        RetryExhaustedError: if all retries fail
        KeyboardInterrupt: always propagated immediately
    """
    for attempt in range(max_retries):
        try:
            return func()
        except KeyboardInterrupt:
            raise  # never retry user interrupts
        except Exception as e:
            if attempt == max_retries - 1:
                raise RetryExhaustedError(max_retries, e) from e
            delay = base_delay * (2 ** attempt)
            jitter = random.uniform(0, delay * 0.1)
            time.sleep(delay + jitter)
\end{lstlisting}

\begin{lstlisting}[style=proof]
-- Proof: retry_with_backoff handles all exception paths
-- 1. Termination
have h_term: the loop runs at most max_retries iterations
    by structural -- range(max_retries) is finite

-- 2. KeyboardInterrupt path
have h_kbd: KeyboardInterrupt is always propagated
    by structural -- except KeyboardInterrupt: raise

-- 3. Normal exception path (retriable)
have h_retry: for attempt < max_retries - 1,
    Exception causes retry with delay
    by structural -- except Exception, then sleep

-- 4. Exhaustion path
have h_exhaust: at attempt = max_retries - 1,
    Exception causes RetryExhaustedError
    by structural -- if attempt == max_retries - 1: raise

-- 5. Backoff property
have h_backoff: delay at attempt k = base_delay * 2^k + jitter
    where 0 <= jitter <= base_delay * 2^k * 0.1
    by unfold delay, jitter then omega

-- 6. Exception coverage
have h_coverage:
    Raised(retry_with_backoff)
      = {RetryExhaustedError, KeyboardInterrupt}
    Propagated = Raised (both are documented)
    by unfold retry_with_backoff then apply exc_coverage

conclude: retry_with_backoff satisfies:
    (a) terminates in at most max_retries iterations [h_term]
    (b) never swallows KeyboardInterrupt [h_kbd]
    (c) exponential backoff with jitter [h_backoff]
    (d) exhaustive exception handling [h_coverage]
    by conjunction h_term h_kbd h_backoff h_coverage
\end{lstlisting}

\begin{theorem}[Retry correctness]
The function \texttt{retry\_with\_backoff} satisfies:
\begin{enumerate}
  \item \textbf{Termination}: it performs at most $\mathsf{max\_retries}$
    attempts.
  \item \textbf{Signal safety}: \texttt{KeyboardInterrupt} is never caught
    by the retry logic.
  \item \textbf{Backoff growth}: the delay at attempt $k$ is
    $\Theta(2^k \cdot \mathsf{base\_delay})$.
  \item \textbf{Exception documentation}: every propagated exception is
    listed in the docstring.
\end{enumerate}
\end{theorem}
\begin{proof}
Property~(1) follows from the finite \texttt{range}.
Property~(2) from the explicit \texttt{except KeyboardInterrupt: raise}.
Property~(3) from the formula $\mathsf{delay} = \mathsf{base\_delay}
\cdot 2^{\mathsf{attempt}}$ with bounded jitter.
Property~(4) from the exception coverage analysis: $\mathsf{Propagated} =
\{\texttt{RetryExhaustedError}, \texttt{KeyboardInterrupt}\}$, both
documented.
\end{proof}

\begin{remark}[\fstar{} comparison]
Verifying a retry-with-backoff function in \fstar{} would require:
(a)~a custom effect combining \lstinline[style=fstar]|ST| (for delay state),
\lstinline[style=fstar]|EXN| (for exceptions), and \lstinline[style=fstar]|IO|
(for \texttt{sleep});
(b)~manual encoding of which exceptions are retried vs.\ propagated;
(c)~a termination proof via a decreasing measure.
The exception path analysis would be entirely manual.  \synhopy{}'s exception
cohomology automates the coverage check, and the path semantics give the
backoff property directly.
\end{remark}


% ──────────────────────────────────────────────────────────────────
\chapter{Generators and Coroutines}
\label{ch:generators}

\section{Generators as Paths in State Space}

\begin{lstlisting}[style=python]
def fibonacci():
    a, b = 0, 1
    while True:
        yield a
        a, b = b, a + b
\end{lstlisting}

\begin{definition}[Generator type]
A generator with state type $S$, yield type $Y$, and send type $T$ is:
\[
  \mathsf{Gen}(S, Y, T) \defeq \Sigma\left(
    \begin{array}{l}
      s_0 : S, \\
      \mathsf{step} : S \to Y \times S, \\
      \mathsf{send} : S \times T \to S
    \end{array}
  \right)
\]
\end{definition}

A generator execution is a \emph{path} $\gamma : \Nat \to S$ defined by:
\begin{align}
  \gamma(0) &= s_0 \\
  \gamma(n+1) &= \pi_2(\mathsf{step}(\gamma(n)))
\end{align}

The $n$-th yielded value is $\pi_1(\mathsf{step}(\gamma(n)))$.

\section{Verifying Generator Properties}

\begin{lstlisting}[style=proof]
-- Proof: fibonacci() yields the Fibonacci sequence
-- State: (a, b) : Nat x Nat
-- Invariant: at step n, a = fib(n), b = fib(n+1)
have h_inv: forall n.
    gamma(n) = (fib(n), fib(n+1))
    by nat_induction n
    base: gamma(0) = (0, 1) = (fib(0), fib(1))
        by unfold fibonacci, fib then omega
    step: gamma(k+1) = step(gamma(k))
        = step((fib(k), fib(k+1)))  by rw [IH]
        = (fib(k+1), fib(k) + fib(k+1))
        = (fib(k+1), fib(k+2))  by rw [fib_def]
conclude: yield_at(fibonacci(), n) = fib(n)
    by rw [h_inv]
\end{lstlisting}

\begin{remark}[\fstar{} comparison]
\fstar{} cannot model Python generators. The \lstinline[style=fstar]{ST} monad runs
computations to completion---there is no concept of pausing, yielding a value,
and resuming. To verify a generator in an \fstar{}-like system, you would need to:
\begin{enumerate}
  \item Manually desugar the generator into a state machine
  \item Verify the state machine
  \item Prove that the state machine faithfully implements the generator protocol
\end{enumerate}
This is approximately 50--100 lines of boilerplate per generator. In \synhopy{},
generators are first-class: their path semantics give immediate access to
inductive proofs about yielded values.
\end{remark}

\section{Generator Composition}

Generators compose via \texttt{yield from}:

\begin{lstlisting}[style=python]
def chain(*iterables):
    for it in iterables:
        yield from it

def flatten(nested):
    for item in nested:
        if isinstance(item, list):
            yield from flatten(item)
        else:
            yield item
\end{lstlisting}

\begin{definition}[\texttt{yield from} as path concatenation]
If $g_1$ is a finite generator yielding $[y_1, \ldots, y_m]$ and $g_2$ yields
$[z_1, \ldots, z_n]$, then
\[
  \texttt{yield from}\;g_1;\;\texttt{yield from}\;g_2
  \;\equiv\; [y_1, \ldots, y_m, z_1, \ldots, z_n]
\]
This is path concatenation in the generator's state space.
\end{definition}

\begin{lstlisting}[style=proof]
-- Proof: flatten produces all leaf elements in order
by_structural_induction on nested:
    case []: yield nothing by unfold flatten
    case [item, *rest] where not isinstance(item, list):
        yield item
        then yield from flatten(rest) by ih
    case [sublist, *rest] where isinstance(sublist, list):
        yield from flatten(sublist) by ih
        then yield from flatten(rest) by ih
conclude: flatten(nested) = all_leaves(nested)
    by induction with path_concat
\end{lstlisting}

\section{Infinite Generators and Coinduction}

Infinite generators (like \texttt{fibonacci()}) require \emph{coinductive} reasoning:

\begin{definition}[Coinductive generator type]
An infinite generator is a \emph{greatest fixed point}:
\[
  \mathsf{InfGen}(Y) = \nu X.\; Y \times X
  \quad\text{(an infinite stream of $Y$ values)}
\]
The coinductive proof principle: to show that two infinite generators produce
the same stream, it suffices to show that they satisfy the same
\emph{bisimulation} relation.
\end{definition}

\begin{theorem}[Bisimulation for generator equivalence]
Two generators $g_1, g_2$ produce the same infinite stream if there exists
a bisimulation $R$ such that:
\begin{enumerate}
  \item $R(s_0^{(1)}, s_0^{(2)})$ (initial states are related)
  \item If $R(s_1, s_2)$ then $\pi_1(\mathsf{step}(s_1)) = \pi_1(\mathsf{step}(s_2))$
    (same yield) and $R(\pi_2(\mathsf{step}(s_1)), \pi_2(\mathsf{step}(s_2)))$
    (next states are related)
\end{enumerate}
\end{theorem}

\section{Async Generators}

Python supports \emph{async generators} that combine \texttt{async def}
with \texttt{yield}, producing values asynchronously:

\begin{lstlisting}[style=python]
async def async_range(n):
    for i in range(n):
        await asyncio.sleep(0.1)
        yield i

async def consume():
    async for value in async_range(10):
        print(value)
\end{lstlisting}

\begin{definition}[Async generator type]
An async generator layers the generator structure over the $\IO$ effect:
\[
  \mathsf{AsyncGen}(S, Y, T) \defeq \Sigma\left(
    \begin{array}{l}
      s_0 : S, \\
      \mathsf{step} : S \to \IO(Y \times S + \mathsf{Done}), \\
      \mathsf{send} : S \times T \to \IO(S)
    \end{array}
  \right)
\]
Each step is an $\IO$ action that may suspend (awaiting an external event)
before yielding.  The $\mathsf{Done}$ variant signals exhaustion via
\texttt{StopAsyncIteration}.
\end{definition}

\begin{definition}[\texttt{async for} as effectful path traversal]
The \texttt{async for} loop traverses the path in the async generator's
state space, executing each step's $\IO$ action:
\[
  \texttt{async for}\;y\;\texttt{in}\;g
  \;\equiv\;
  \mathsf{unfold}_\IO(\mathsf{step}, s_0)
  : \IO(\mathsf{List}(Y))
\]
where $\mathsf{unfold}_\IO$ iterates $\mathsf{step}$ until
$\mathsf{Done}$, collecting yielded values.
\end{definition}

\begin{remark}[\fstar{} comparison]
\fstar{} has no built-in concept of async generators.  Its effect system
could model either generators (via \lstinline[style=fstar]|ST|) or async
operations (via a custom \lstinline[style=fstar]|Async| effect), but combining
both requires a \emph{monad transformer stack}---and \fstar{} does not
natively support monad transformers.
\synhopy{}'s $\IO$-lifted generator type handles both dimensions naturally.
\end{remark}

\section{Generator-Based Pipelines}

Generators compose into pipelines analogous to Unix pipes:

\begin{lstlisting}[style=python]
def read_lines(path):
    with open(path) as f:
        for line in f:
            yield line.strip()

def filter_nonempty(lines):
    for line in lines:
        if line:
            yield line

def parse_csv(lines):
    for line in lines:
        yield line.split(',')

# Pipeline: read | filter | parse
pipeline = parse_csv(filter_nonempty(read_lines("data.csv")))
\end{lstlisting}

\begin{definition}[Generator pipeline as functor composition]
A pipeline of generators $g_1, g_2, \ldots, g_n$ where each $g_i : \mathsf{Gen}(S_i, Y_i, T_i)$
and $g_{i+1}$ consumes $Y_i$, is modeled as functor composition in the
category of streams:
\[
  g_n \circ g_{n-1} \circ \cdots \circ g_1 :
  \mathsf{Gen}(S_1 \times \cdots \times S_n,\; Y_n,\; T_1)
\]
Each stage pulls one value from the preceding stage on demand, preserving
laziness throughout the pipeline.
\end{definition}

\begin{theorem}[Pipeline fusion]
If each generator stage $g_i$ processes elements independently (no
inter-element state), the pipeline $g_n \circ \cdots \circ g_1$ processes
each source element through all stages before advancing to the next.
The peak memory usage is $O(1)$ per element, regardless of input size.
\end{theorem}
\begin{proof}
By the pull-based semantics of Python generators: calling
\texttt{next()} on $g_n$ triggers a chain of \texttt{next()} calls
back to $g_1$.  Each intermediate value is consumed immediately and
not stored, giving $O(1)$ per-element memory.
\end{proof}

\begin{example}[Unix pipe analogy]
The generator pipeline \texttt{parse\_csv(filter\_nonempty(read\_lines(f)))}
is semantically equivalent to the Unix pipeline:
\begin{verbatim}
cat data.csv | grep -v '^$' | awk -F',' '{print}'
\end{verbatim}
Both are lazy, pull-based, and process one record at a time.  \synhopy{}
verifies the pipeline's type correctness by checking that each stage's
yield type matches the next stage's input type.
\end{example}

\begin{remark}[\fstar{} comparison]
\fstar{} lacks generator-based pipelines.  The standard approach is to
compose pure functions over lists, which materialises all intermediate
results.  Achieving lazy pipeline behavior in \fstar{} requires manually
implementing a pull-based stream type and proving its properties---\synhopy{}
gets this for free from Python's generator semantics.
\end{remark}

\section{Generators as Lazy Evaluation}

Python generators provide \emph{lazy evaluation}: values are computed only
when requested.  This is analogous to Haskell's thunks, but explicit.

\begin{definition}[Generator as explicit thunk]
A generator expression is an \emph{explicit thunk}: a suspended computation
that produces values on demand.
\[
  \mathsf{lazy}(f, s_0) \defeq
  \lambda().\; \mathsf{let}\;(y, s_1) = \mathsf{step}(s_0)\;\mathsf{in}\;
    (y,\; \mathsf{lazy}(f, s_1))
\]
Unlike Haskell's implicit thunks (which can cause space leaks via hidden
retention), Python's generators make the suspension points explicit via
\texttt{yield}.
\end{definition}

\begin{theorem}[Lazy vs.\ eager equivalence]
For a pure generator $g$ with finite output $[y_0, \ldots, y_{n-1}]$,
\[
  \mathsf{list}(g) = [y_0, \ldots, y_{n-1}]
\]
That is, forcing all values of the lazy generator produces the same result
as the eager list construction.  However, the lazy version uses $O(1)$
working memory (one state $s_k$ at a time), while the eager version
uses $O(n)$.
\end{theorem}
\begin{proof}
By induction on $n$.  Each step produces $y_k = \pi_1(\mathsf{step}(s_k))$
and the next state $s_{k+1} = \pi_2(\mathsf{step}(s_k))$.  The
\texttt{list()} call exhausts the generator, collecting all $y_k$ values.
The results are identical; only the space usage differs.
\end{proof}

\begin{remark}[\fstar{} comparison]
\fstar{} is a strict (eager) language by default.  Lazy evaluation must be
encoded explicitly using thunks (\lstinline[style=fstar]|unit -> Tot 'a|),
and there is no built-in generator mechanism.  Haskell's implicit laziness
can cause reasoning difficulties (space leaks, evaluation order);
\synhopy{}'s explicit generator model avoids these while still providing
verified laziness guarantees.
\end{remark}

\section{Verifying Memory Efficiency}

\synhopy{} can verify that a generator-based computation is more memory-efficient
than its list-based equivalent.

\begin{definition}[Space complexity annotation]
For a computation $f$, define:
\begin{align}
  \mathsf{Space}_{\mathsf{eager}}(f, n) &\defeq
    \text{peak memory when } f \text{ materialises all } n \text{ elements} \\
  \mathsf{Space}_{\mathsf{lazy}}(f, n) &\defeq
    \text{peak memory when } f \text{ uses generators}
\end{align}
\end{definition}

\begin{example}[List vs.\ generator memory]
\begin{lstlisting}[style=python]
# Eager: O(n) memory
total_eager = sum([x**2 for x in range(10**8)])

# Lazy: O(1) memory
total_lazy = sum(x**2 for x in range(10**8))
\end{lstlisting}
\end{example}

\begin{lstlisting}[style=proof]
-- Proof: generator version uses O(1) memory
have h_eager: [x**2 for x in range(n)] allocates list of size n
    by structural -- list comprehension materialises all elements
have h_lazy: (x**2 for x in range(n)) yields one element at a time
    by structural -- generator expression is lazy
have h_sum: sum() consumes elements one at a time
    by unfold sum -- sum accumulates into a single variable
conclude: Space_lazy(sum_of_squares, n) = O(1)
    by h_lazy, h_sum -- only one x^2 and the accumulator exist at a time
conclude: Space_eager(sum_of_squares, n) = O(n)
    by h_eager -- the full list [0, 1, 4, ...] is in memory
\end{lstlisting}

\begin{theorem}[Generator memory bound]
If a generator $g$ has state size $|S|$ and yield size $|Y|$,
and a consumer $c$ processes one element at a time with accumulator
size $|A|$, then:
\[
  \mathsf{Space}_{\mathsf{lazy}}(c \circ g, n)
  = O(|S| + |Y| + |A|)
  = O(1) \text{ when $|S|$, $|Y|$, $|A|$ are constant}
\]
independent of the number of elements $n$.
\end{theorem}
\begin{proof}
At any point during execution, only the current state $s_k$, the current
yield $y_k$, and the accumulator are live.  Previous states and yields
have been garbage collected.  Hence the peak memory is bounded by
$|S| + |Y| + |A|$.
\end{proof}

\begin{remark}[\fstar{} comparison]
\fstar{} provides no built-in mechanism for reasoning about space complexity.
Memory usage analysis would require a custom resource-aware type system
(e.g., AARA-style).  \synhopy{}'s generator model makes space bounds
a natural consequence of the path semantics: the ``width'' of the active
path segment bounds the live memory.
\end{remark}

\section{\texttt{send()} and \texttt{throw()}: Bidirectional Communication}

Python generators support bidirectional communication via \texttt{send()}
and \texttt{throw()}:

\begin{lstlisting}[style=python]
def accumulator():
    total = 0
    while True:
        value = yield total
        if value is None:
            break
        total += value

gen = accumulator()
next(gen)          # prime: yields 0
gen.send(10)       # yields 10
gen.send(20)       # yields 30
gen.send(None)     # StopIteration
\end{lstlisting}

\begin{definition}[\texttt{send()} typing]
The \texttt{send()} method injects a value of type $T$ into the generator,
which becomes the result of the \texttt{yield} expression.  The full
generator protocol is a \emph{session type}:
\[
  \mathsf{GenProtocol}(Y, T) \defeq \mu X.\;
    !\,Y.\; ?\,T.\; X \oplus \mathsf{Done}
\]
Each round: the generator \emph{outputs} ($!$) a yield value of type $Y$,
then \emph{inputs} ($?$) a sent value of type $T$, then either continues
($X$) or terminates ($\mathsf{Done}$).
\end{definition}

\begin{definition}[\texttt{throw()} typing]
The \texttt{throw()} method injects an exception into the generator at
the \texttt{yield} point:
\[
  \mathsf{throw} : \mathsf{Gen}(S, Y, T) \times E
    \to \mathsf{Gen}(S, Y, T) + E'
\]
The generator may catch the thrown exception (continuing with a new state)
or propagate a different exception $E'$.
\end{definition}

\begin{lstlisting}[style=proof]
-- Proof: accumulator correctly sums sent values
-- State: total : Nat
-- Protocol: yield total, receive value, add to total
have h_inv: after sending [v_1, ..., v_k],
    the next yield is v_1 + ... + v_k
    by nat_induction k
    base: next(gen) yields 0 = sum([]) by unfold accumulator
    step: gen.send(v_{k+1}) yields
        total + v_{k+1} = sum([v_1,...,v_k]) + v_{k+1}
        = sum([v_1,...,v_{k+1}]) by rw [IH], omega
conclude: accumulator implements running sum
    by h_inv
\end{lstlisting}

\begin{remark}[\fstar{} comparison]
\fstar{} has no concept of bidirectional generator communication.
Modeling \texttt{send()} and \texttt{throw()} would require implementing
a full session type system on top of \fstar{}'s effect system---a significant
undertaking.  \synhopy{} types the generator protocol as a session type
natively, enabling verification of bidirectional generator patterns.
\end{remark}

\section{Worked Example: Paginated API Client}

We verify a generator that fetches paginated results from a REST API,
handling errors and ensuring all pages are consumed.

\begin{lstlisting}[style=python]
import requests

def paginated_fetch(base_url, page_size=100):
    """Yield all items from a paginated API.
    Raises:
        requests.HTTPError: on non-2xx response
        requests.ConnectionError: on network failure
    """
    page = 1
    while True:
        resp = requests.get(
            base_url, params={"page": page, "size": page_size}
        )
        resp.raise_for_status()
        data = resp.json()
        items = data["items"]
        if not items:
            return  # no more pages
        yield from items
        page += 1
\end{lstlisting}

\begin{lstlisting}[style=proof]
-- Proof: paginated_fetch yields all items from the API
-- State: (page : Nat+, buffer : List[Item])
-- Invariant: pages 1..page-1 have been fully yielded

have h_page_inv: at the start of iteration page = k,
    items from pages 1..k-1 have been yielded
    by nat_induction k
    base: page = 1, nothing yielded yet (correct)
    step: page = k+1 means page k was just yielded via yield from

have h_term: the loop terminates iff the API returns empty items
    by structural -- if not items: return

have h_complete: if the API has N total items across P pages,
    then paginated_fetch yields exactly N items
    by h_page_inv with k = P+1, noting items = [] at page P+1

-- Exception coverage
have h_exc: Raised(paginated_fetch) =
    {HTTPError, ConnectionError}
    by unfold requests.get, raise_for_status

conclude: paginated_fetch is a correct paginated consumer:
    (a) yields all items in order [h_complete]
    (b) terminates when pages are exhausted [h_term]
    (c) documents all raised exceptions [h_exc]
    by conjunction h_complete h_term h_exc
\end{lstlisting}

\begin{theorem}[Paginated generator correctness]
If the API returns $P$ pages of items with $n_i$ items on page $i$,
then \texttt{paginated\_fetch} yields exactly $\sum_{i=1}^{P} n_i$ items
in order, with $O(\mathsf{page\_size})$ peak memory.
\end{theorem}
\begin{proof}
By the page invariant: at the start of iteration $k$, pages $1$ through
$k-1$ have been fully yielded.  The \texttt{yield from items} statement
yields all $n_k$ items of page $k$ lazily.  When page $P+1$ returns
an empty list, the generator terminates via \texttt{return}, raising
\texttt{StopIteration}.  Peak memory is bounded by one page of items
($O(\mathsf{page\_size})$) since previous pages are discarded.
\end{proof}

\begin{remark}[\fstar{} comparison]
Verifying a paginated API client in \fstar{} would require:
(a)~modeling HTTP requests as \lstinline[style=fstar]|IO| actions;
(b)~manually implementing a lazy stream;
(c)~proving termination via a well-founded relation on page numbers
(which depends on the API's behavior---an external oracle).
\synhopy{}'s generator semantics handle (b) natively, and the path-based
induction handles (c) directly.
\end{remark}


% ──────────────────────────────────────────────────────────────────
\chapter{Closures and Late Binding}
\label{ch:closures}

\section{The Late-Binding Problem}

\begin{lstlisting}[style=python]
funcs = [lambda x: x + i for i in range(5)]
print(funcs[0](0))  # Prints 4, not 0!
\end{lstlisting}

Python closures capture variables \emph{by reference}, not by value. The variable
\texttt{i} in the lambda refers to the \emph{same} \texttt{i} as the loop
variable, which ends at $4$ after the loop completes.

\section{Closure Type with Scope Fiber}

\begin{definition}[Closure type]
A closure is a function together with a \emph{scope fiber}:
\[
  \mathsf{Closure}(A, B) \defeq \Sigma(
    \mathsf{env} : \mathsf{Scope},\;
    f : A \to_{\mathsf{env}} B
  )
\]
where $\mathsf{Scope}$ is the captured environment, and $f$'s behavior may
depend on $\mathsf{env}$.
\end{definition}

\begin{definition}[Late binding]
In late binding, the scope fiber is \emph{mutable}: the environment is a reference,
not a value. The closure's type at call time depends on the environment's state:
\[
  \mathsf{late\_closure}(f, \mathsf{env}) : A \to B[\mathsf{env}/\text{current state}]
\]
\end{definition}

\section{Proving Late-Binding Bugs}

\synhopy{} can prove that the list comprehension above is buggy:

\begin{lstlisting}[style=proof]
-- Claim: funcs[k](0) = k for all k in 0..4
-- This claim is FALSE and SynHoPy rejects it:
-- After the loop, i = 4 in all closures' scope fibers
have h_scope: forall k. funcs[k].env.i = 4
    by structural -- all closures share the same 'i' reference
have h_eval: funcs[k](0) = 0 + funcs[k].env.i = 4
    by unfold lambda then rw [h_scope]
-- Therefore funcs[0](0) = 4, not 0
-- The spec "funcs[k](0) = k" is REFUTED
conclude: funcs[k](0) = 4 != k (for k != 4)
    by omega with h_eval
\end{lstlisting}

\begin{remark}[\fstar{} comparison]
\fstar{} captures variables by value (its extraction to OCaml ensures this). An
\fstar{}-based Python verifier that assumes capture-by-value would
\emph{incorrectly verify} the claim \texttt{funcs[k](0) = k}---the verification
would succeed, but the program would produce wrong results at runtime. This is a
\textbf{soundness violation}: the verifier says ``correct'' but the program is
wrong.

\synhopy{} models Python's actual capture-by-reference semantics via scope fibers,
so it correctly identifies the bug.
\end{remark}

\section{The Default-Argument Fix}

The standard fix for late-binding closures uses a default argument:
\begin{lstlisting}[style=python]
funcs_fixed = [lambda x, i=i: x + i for i in range(5)]
funcs_fixed[0](0)  # Prints 0 (correct)
\end{lstlisting}

\begin{definition}[Default-argument capture]
A default argument $i=i$ in a lambda creates a \emph{frozen scope fiber}:
\[
  \mathsf{closure}(f, \mathsf{env}[i \mapsto v_{\mathsf{at\_creation}}])
  : A \to B[i / v_{\mathsf{at\_creation}}]
\]
The value of $i$ is captured at creation time, not at call time.
\end{definition}

\begin{lstlisting}[style=proof]
-- Now we CAN prove the correct spec:
have h_freeze: forall k. funcs_fixed[k].env.i = k
    by structural -- each lambda freezes i = k at creation
have h_eval: funcs_fixed[k](0) = 0 + k = k
    by unfold lambda then rw [h_freeze]
conclude: funcs_fixed[k](0) = k
    by omega with h_eval
\end{lstlisting}

\begin{theorem}[Default argument capture soundness]
In \synhopy{}, a default argument $\texttt{i=i}$ produces a frozen scope fiber
whose value is determined at closure creation time.  This correctly models
Python's semantics.
\end{theorem}
\begin{proof}
Python evaluates default arguments once, at function definition time.
The $\mathsf{env}[i \mapsto v]$ substitution in the scope fiber captures
exactly this: $v$ is the value of $i$ at the time the lambda is created
(during the list comprehension iteration), not at the time the lambda is called.
\end{proof}

\section{Closure Serialization and Type-Theoretic Implications}

Python closures can be serialized via \texttt{pickle} (with limitations)
or \texttt{cloudpickle}.  This raises deep type-theoretic questions:
what does it mean to serialize a function together with its captured
environment?

\begin{lstlisting}[style=python]
import cloudpickle, pickle

def make_adder(n):
    def adder(x):
        return x + n
    return adder

add5 = make_adder(5)
serialized = cloudpickle.dumps(add5)
add5_restored = pickle.loads(serialized)
add5_restored(10)  # 15
\end{lstlisting}

\begin{definition}[Closure serialization]
Serializing a closure $c = (\mathsf{env}, f)$ produces a \emph{flat
representation} that decouples the function from its scope fiber:
\[
  \mathsf{serialize} : \mathsf{Closure}(A, B) \to \mathsf{Bytes}
\]
\[
  \mathsf{deserialize} : \mathsf{Bytes} \to \mathsf{Closure}(A, B) + \mathsf{Error}
\]
The round-trip property requires:
\[
  \forall c.\;
  \mathsf{deserialize}(\mathsf{serialize}(c)) = \mathsf{Ok}(c')
  \implies \forall x.\; c(x) = c'(x)
\]
That is, the deserialized closure is \emph{extensionally equal} to the
original---but not necessarily \emph{intensionally} equal (the scope fiber
structure may differ).
\end{definition}

\begin{theorem}[Serialization breaks scope sharing]
If two closures $c_1, c_2$ share a mutable scope fiber (late binding),
serializing and deserializing them produces independent copies:
\[
  \mathsf{let}\;(c_1', c_2') = (\mathsf{deser}(\mathsf{ser}(c_1)),
    \mathsf{deser}(\mathsf{ser}(c_2)))
\]
Then mutating $c_1'.\mathsf{env}$ does not affect $c_2'.\mathsf{env}$,
even if the originals shared state.
\end{theorem}
\begin{proof}
Serialization traverses the object graph and produces independent copies
of all referenced objects.  The shared reference in the original scope
fiber becomes two distinct objects in the deserialized closures.
\end{proof}

\begin{remark}[\fstar{} comparison]
\fstar{} functions are pure values (or effectful computations) that cannot
be serialized in the traditional sense---extraction to OCaml or F\#
produces compiled code, not data.  The question of closure serialization
simply does not arise.  \synhopy{} must address it because Python closures
\emph{are} serializable, and serialization changes the sharing semantics
of scope fibers.
\end{remark}

\section{\texttt{nonlocal} and Multi-Level Scope Chains}

Python's \texttt{nonlocal} keyword allows closures to mutate variables in
enclosing (but non-global) scopes, creating multi-level scope chains:

\begin{lstlisting}[style=python]
def make_counter():
    count = 0
    def increment():
        nonlocal count
        count += 1
        return count
    def get():
        return count
    return increment, get

inc, get = make_counter()
inc()        # 1
inc()        # 2
get()        # 2 -- shares the same 'count'
\end{lstlisting}

\begin{definition}[Multi-level scope chain]
A scope chain of depth $d$ is a sequence of nested scope fibers:
\[
  \mathsf{ScopeChain}_d \defeq
    \mathsf{Scope}_1 \hookrightarrow \mathsf{Scope}_2
    \hookrightarrow \cdots \hookrightarrow \mathsf{Scope}_d
\]
where $\hookrightarrow$ denotes inclusion (each inner scope can access
variables from all enclosing scopes).  The \texttt{nonlocal} keyword
selects a specific level:
\[
  \texttt{nonlocal}\;x \;\equiv\;
  \text{resolve } x \text{ in the nearest enclosing scope where } x
    \text{ is defined}
\]
\end{definition}

\begin{definition}[\texttt{nonlocal} mutation as fiber transport]
A \texttt{nonlocal} assignment $x = v$ in scope level $k$ is a
\emph{transport} along the scope chain:
\[
  \mathsf{nonlocal\_assign}(x, v, k) :
    \mathsf{ScopeChain}_d
    \xrightarrow{\transp}
    \mathsf{ScopeChain}_d[x \mapsto v \text{ at level } k]
\]
This modifies the shared fiber, affecting all closures that reference
scope level $k$.
\end{definition}

\begin{lstlisting}[style=proof]
-- Proof: inc and get share state via nonlocal
have h_scope: inc.env.count and get.env.count
    reference the same cell in ScopeChain_2
    by structural -- both close over make_counter's scope
have h_inc: inc() increments the shared count via nonlocal
    by unfold increment then rw [nonlocal_assign]
have h_get: get() reads the shared count
    by unfold get
have h_shared: after inc(); inc(), get() = 2
    by h_inc (twice), h_get, h_scope -- shared mutation
conclude: make_counter returns a consistent (inc, get) pair
    by h_shared
\end{lstlisting}

\begin{remark}[\fstar{} comparison]
\fstar{} has no concept of \texttt{nonlocal} scope chains.  All variables
are either local (let-bound) or global (top-level).  Mutable shared state
between functions requires explicit references (\lstinline[style=fstar]|ref|)
threaded through the \lstinline[style=fstar]|ST| effect.
\synhopy{}'s scope chain model captures Python's multi-level scoping rules
directly, including the subtlety that \texttt{nonlocal} skips the local scope
and modifies an enclosing one.
\end{remark}

\section{Closures in Decorators: The Closure-Decorator Duality}

Every Python decorator is a closure factory: it takes a function, wraps
it in a closure that captures the original function, and returns the
closure.

\begin{lstlisting}[style=python]
def retry(max_attempts=3):
    def decorator(func):
        def wrapper(*args, **kwargs):
            for attempt in range(max_attempts):
                try:
                    return func(*args, **kwargs)
                except Exception:
                    if attempt == max_attempts - 1:
                        raise
        return wrapper
    return decorator

@retry(max_attempts=5)
def fetch_data(url):
    return requests.get(url).json()
\end{lstlisting}

\begin{definition}[Closure-decorator duality]
A parameterised decorator with parameter $p$ is a \emph{closure factory}:
\[
  \mathsf{Decorator}(p) : (A \to_\varepsilon B) \to \mathsf{Closure}(A, B)
\]
The resulting closure captures both the parameter $p$ and the original
function $f$ in its scope fiber:
\[
  \mathsf{Decorator}(p)(f) =
    \bigl(\mathsf{env} = \{p, f\},\;
      \mathsf{wrapper} : A \to_{\varepsilon'} B\bigr)
\]
The \emph{duality} is that every closure can be viewed as a decorated
identity function, and every decorator produces closures.
\end{definition}

\begin{theorem}[Decorator composition preserves closure structure]
If $d_1$ and $d_2$ are decorators, then $d_1 \circ d_2$ applied to $f$
produces a closure whose scope fiber contains the scope fibers of both
$d_1$ and $d_2$:
\[
  (d_1 \circ d_2)(f).\mathsf{env} \supseteq
    d_1(\cdot).\mathsf{env} \cup d_2(f).\mathsf{env}
\]
The decoration order matters: $d_1 \circ d_2 \neq d_2 \circ d_1$ in
general, because the outer decorator's wrapper wraps the inner decorator's
wrapper.
\end{theorem}
\begin{proof}
By unfolding: $(d_1 \circ d_2)(f) = d_1(d_2(f))$.  The outer closure
$d_1(\cdot)$ captures $d_2(f)$ (which is itself a closure capturing $f$
and $d_2$'s parameters).  The environments nest, and the non-commutativity
follows from the fact that wrappers execute in LIFO order.
\end{proof}

\begin{remark}[\fstar{} comparison]
\fstar{} has no decorator syntax.  The pattern of wrapping functions is
achieved via higher-order functions, but the resulting composition is
purely functional---there is no mutable scope fiber, no \texttt{nonlocal},
and no closure serialization concern.  While simpler, this means \fstar{}
cannot model Python decorators that rely on mutable closure state
(e.g., memoization decorators with a shared cache).
\end{remark}

\section{Partial Application: Closures vs.\ \texttt{functools.partial}}

Python offers two mechanisms for partial application:

\begin{lstlisting}[style=python]
from functools import partial

# Via closure
def add(x, y):
    return x + y

add5_closure = lambda y: add(5, y)
add5_partial = partial(add, 5)

# Both work:
add5_closure(10)  # 15
add5_partial(10)  # 15
\end{lstlisting}

\begin{definition}[Partial application types]
The two partial application mechanisms have different type structures:
\begin{align}
  \mathsf{closure\_partial}(f, a) &\defeq
    \bigl(\mathsf{env} = \{f, a\},\;
      \lambda y.\; f(a, y)\bigr)
    : \mathsf{Closure}(B, C) \\
  \mathsf{functools\_partial}(f, a) &\defeq
    \bigl(\mathsf{func} = f,\;
      \mathsf{args} = (a,),\;
      \mathsf{kwargs} = \{\}\bigr)
    : \mathsf{Partial}(B, C)
\end{align}
The \texttt{functools.partial} object is \emph{inspectable}: its
\texttt{.func}, \texttt{.args}, and \texttt{.keywords} attributes expose
the partial application structure.  A closure-based partial is opaque.
\end{definition}

\begin{theorem}[Extensional equivalence of partial application]
For a pure function $f : A \times B \to C$ and a value $a : A$:
\[
  \forall b : B.\;
  \mathsf{closure\_partial}(f, a)(b) = \mathsf{functools\_partial}(f, a)(b)
\]
However, the two are not intensionally equal: \texttt{partial} objects
support introspection while closures do not.
\end{theorem}
\begin{proof}
Both evaluate to $f(a, b)$ by definition.  The extensional equality
follows directly.  The intensional difference is that
$\mathsf{functools\_partial}(f, a).\mathsf{func} = f$ is accessible,
while the closure's internal reference to $f$ is not programmatically
exposed (only via \texttt{\_\_closure\_\_} cells).
\end{proof}

\begin{example}[When the difference matters]
\begin{lstlisting}[style=python]
# functools.partial preserves function metadata
add5_partial.func    # <function add>
add5_partial.args    # (5,)

# Closure does not
add5_closure.func    # AttributeError
# But you can inspect closure cells:
add5_closure.__closure__[0].cell_contents  # 5
\end{lstlisting}
\synhopy{} tracks whether a partial application is closure-based or
\texttt{functools.partial}-based, because the difference affects
duck-typing compatibility: some APIs check for \texttt{.func} attribute.
\end{example}

\begin{remark}[\fstar{} comparison]
\fstar{} supports partial application natively via currying: every multi-argument
function \lstinline[style=fstar]|val f : a -> b -> c| can be partially applied
as \lstinline[style=fstar]|f x|.  There is no distinction between closure-based
and explicit partial application---all partial applications are closures in the
target language.  \synhopy{} must distinguish the two because Python's
\texttt{functools.partial} has different introspection and duck-typing behavior
than lambda closures.
\end{remark}

\section{Memory Leak Prevention: Closure Reference Cycles}

Closures can create reference cycles that prevent garbage collection:

\begin{lstlisting}[style=python]
class Node:
    def __init__(self, value):
        self.value = value
        self.on_change = None  # callback closure

def setup_node(node):
    # This closure captures 'node', creating a cycle:
    # node -> node.on_change -> closure -> node
    def callback():
        print(f"Node {node.value} changed")
    node.on_change = callback
    return node
\end{lstlisting}

\begin{definition}[Closure reference graph]
The \emph{closure reference graph} $G_c$ has:
\begin{itemize}
  \item Vertices: objects and closures in the program
  \item Edges: $o_1 \to o_2$ if $o_1$ holds a reference to $o_2$
    (either as an attribute, list element, or closure cell)
\end{itemize}
A \emph{closure reference cycle} is a cycle in $G_c$ that passes through
at least one closure node.
\end{definition}

\begin{definition}[Cycle detection via scope fiber analysis]
\synhopy{} detects closure reference cycles by analysing the scope fiber:
\[
  \mathsf{CycleRisk}(c) \defeq
    \exists\; o \in \mathsf{env}(c).\;
    c \in \mathsf{reachable}(o)
\]
If an object $o$ in the closure's environment can reach back to the closure
itself, a reference cycle exists.
\end{definition}

\begin{theorem}[Closure cycle detection soundness]
If $\mathsf{CycleRisk}(c)$ holds, then assigning $c$ to an attribute
of any $o \in \mathsf{env}(c)$ creates a reference cycle in the heap.
Conversely, if $\neg\mathsf{CycleRisk}(c)$, no closure-induced cycle
exists.
\end{theorem}
\begin{proof}
Forward direction: if $o \in \mathsf{env}(c)$ and $c \in
\mathsf{reachable}(o)$, then the path $c \to_{\mathsf{env}} o
\to_{\mathsf{reach}} c$ forms a cycle.
Reverse direction: if no object in $\mathsf{env}(c)$ can reach $c$,
then $c$ is only reachable from its defining scope, and no cycle exists.
\end{proof}

\begin{lstlisting}[style=python]
import weakref

def setup_node_safe(node):
    # Fix: use weakref to break the cycle
    weak_node = weakref.ref(node)
    def callback():
        n = weak_node()
        if n is not None:
            print(f"Node {n.value} changed")
    node.on_change = callback
    return node
\end{lstlisting}

\begin{lstlisting}[style=proof]
-- Proof: setup_node creates a cycle, setup_node_safe does not
-- Cycle analysis for setup_node:
have h_cycle: callback.env contains node (strong ref)
    and node.on_change = callback
    by structural -- closure captures 'node' directly
have h_risk: CycleRisk(callback) holds
    by h_cycle -- node -> callback -> node

-- Cycle analysis for setup_node_safe:
have h_weak: callback.env contains weak_node (weak ref)
    by structural -- closure captures weak_node, not node
have h_safe: CycleRisk(callback) does NOT hold
    by h_weak -- weak refs are not strong references,
    -- so node is not in reachable(callback) for GC purposes
conclude: setup_node_safe prevents the memory leak
    by h_safe
\end{lstlisting}

\begin{remark}[\fstar{} comparison]
\fstar{} uses a garbage-collected runtime (OCaml or F\#) and does not
expose reference cycle concerns to the programmer.  There is no
\texttt{weakref} equivalent, and no need for one since the GC handles
cycles transparently.  In Python, reference counting is the primary GC
mechanism, and cycles involving closures can delay or prevent collection.
\synhopy{}'s scope fiber analysis detects these cycles statically,
before they become runtime memory leaks.
\end{remark}

\section{Worked Example: Callback-Based Event System}

We verify a complete callback-based event system where closures serve as
event handlers, demonstrating closure typing, reference cycle prevention,
and exhaustive handler verification.

\begin{lstlisting}[style=python]
import weakref
from typing import Callable, Dict, List

class EventEmitter:
    def __init__(self):
        self._handlers: Dict[str, List[Callable]] = {}

    def on(self, event: str, handler: Callable) -> Callable:
        """Register a handler. Returns an unsubscribe function."""
        if event not in self._handlers:
            self._handlers[event] = []
        self._handlers[event].append(handler)

        # The unsubscribe closure captures self and handler
        weak_self = weakref.ref(self)
        def unsubscribe():
            emitter = weak_self()
            if emitter is not None and event in emitter._handlers:
                emitter._handlers[event].remove(handler)
        return unsubscribe

    def emit(self, event: str, *args, **kwargs):
        """Call all handlers for the event."""
        for handler in self._handlers.get(event, []):
            handler(*args, **kwargs)
\end{lstlisting}

\begin{lstlisting}[style=python]
# Usage with closures as handlers
def make_logger(prefix):
    def log_handler(message):
        print(f"[{prefix}] {message}")
    return log_handler

emitter = EventEmitter()
unsub1 = emitter.on("error", make_logger("ERR"))
unsub2 = emitter.on("error", make_logger("WARN"))
unsub3 = emitter.on("info", make_logger("INFO"))

emitter.emit("error", "disk full")
# [ERR] disk full
# [WARN] disk full

unsub1()  # unsubscribe ERR handler
emitter.emit("error", "disk full")
# [WARN] disk full
\end{lstlisting}

\begin{definition}[Event system type]
The event system is typed as a dependent record with handler registration
as a fibration:
\[
  \mathsf{EventEmitter} \defeq \Sigma\!\left(
    \begin{array}{l}
      \mathsf{handlers} : \mathsf{str} \to \mathsf{List}(\PyFunc), \\
      \mathsf{on} : \mathsf{str} \times \PyFunc
        \to \Mutates(\mathsf{handlers}) \times \mathsf{Closure}(\mathsf{unit}, \mathsf{unit}), \\
      \mathsf{emit} : \mathsf{str} \times \mathsf{args}
        \to \Reads(\mathsf{handlers}) \times \IO
    \end{array}
  \right)
\]
The \texttt{on} method returns a closure (the unsubscribe function)
whose scope fiber references the emitter via a weak reference.
\end{definition}

\begin{lstlisting}[style=proof]
-- Proof: the event system is correct and leak-free

-- 1. Handler registration
have h_reg: after emitter.on(event, handler),
    handler in emitter._handlers[event]
    by unfold on then structural

-- 2. Handler invocation
have h_emit: emitter.emit(event, *args) calls all handlers
    for that event with the given args
    by unfold emit -- iterates _handlers[event]

-- 3. Unsubscribe correctness
have h_unsub: after unsub(),
    handler not in emitter._handlers[event]
    by unfold unsubscribe then structural
    -- remove(handler) removes exactly one occurrence

-- 4. No reference cycles
have h_no_cycle: unsubscribe uses weakref.ref(self),
    so CycleRisk(unsubscribe) = false
    by structural -- weak_self is a weak reference

-- 5. Handler isolation
have h_isolate: emitting event e1 does not call
    handlers registered for event e2 (e1 != e2)
    by unfold emit -- _handlers.get(event, []) filters

conclude: EventEmitter satisfies:
    (a) handlers are called on emit [h_emit]
    (b) unsubscribe removes exactly the target handler [h_unsub]
    (c) no memory leaks from closure cycles [h_no_cycle]
    (d) event isolation [h_isolate]
    by conjunction h_emit h_unsub h_no_cycle h_isolate
\end{lstlisting}

\begin{theorem}[Event system correctness]
The \texttt{EventEmitter} class satisfies:
\begin{enumerate}
  \item \textbf{Registration}: \texttt{on(e, h)} adds $h$ to the handler
    list for event $e$.
  \item \textbf{Dispatch}: \texttt{emit(e, args)} invokes exactly the
    handlers registered for $e$, in registration order.
  \item \textbf{Unsubscription}: the closure returned by \texttt{on}
    removes exactly one handler instance.
  \item \textbf{Leak freedom}: no closure reference cycles exist,
    because the unsubscribe closure uses \texttt{weakref}.
  \item \textbf{Isolation}: events with different names do not interfere.
\end{enumerate}
\end{theorem}
\begin{proof}
Properties (1) and (2) follow from the dictionary-of-lists structure
and the \texttt{for} loop in \texttt{emit}.
Property (3) follows from \texttt{list.remove} semantics (removes first
occurrence).
Property (4) follows from the scope fiber analysis: \texttt{weak\_self}
is a weak reference, so the closure does not prevent garbage collection
of the emitter.
Property (5) follows from the \texttt{dict.get(event, [])} lookup,
which returns only handlers for the specified event.
\end{proof}

\begin{remark}[\fstar{} comparison]
Implementing a callback-based event system in \fstar{} requires:
(a)~a mutable map from event names to handler lists
  (via \lstinline[style=fstar]|ST| with \lstinline[style=fstar]|ref|);
(b)~higher-order functions for handlers (straightforward);
(c)~no way to express unsubscription as a closure returning from
  registration (closures in \fstar{} cannot mutate enclosing state
  without explicit \lstinline[style=fstar]|ref| threading);
(d)~no concept of weak references or cycle detection.
\synhopy{}'s closure model handles all of these concerns natively,
including the subtle interaction between closure scope fibers,
weak references, and mutable handler lists.
\end{remark}


% ──────────────────────────────────────────────────────────────────
\chapter{Multiple Inheritance and the MRO}
\label{ch:mro}

\section{C3 Linearization}

Python uses C3 linearization to compute the method resolution order:

\begin{lstlisting}[style=python]
class A:
    def method(self): return "A"
class B(A):
    def method(self): return "B"
class C(A):
    def method(self): return "C"
class D(B, C):
    pass

D().method()  # Returns "B" (MRO: D -> B -> C -> A)
\end{lstlisting}

\section{MRO as Section of Class Bundle}

\begin{definition}[Class fiber bundle]
For a class $D$ with bases $B_1, \ldots, B_k$, the class fiber bundle has:
\begin{itemize}
  \item Base space: the MRO sequence $[D, B_1, \ldots, B_k, \ldots]$
  \item Fiber at class $C$: the set of methods defined \emph{at} $C$ (not inherited)
  \item The MRO section selects, for each method name, the first class in the
    linearization that defines it
\end{itemize}
\end{definition}

\begin{definition}[Method resolution typing rule]
\[
  \frac{
    \MRO(D) = [D, B_1, \ldots, B_k] \qquad
    B_j \text{ is the first class in } \MRO(D) \text{ with method } m
  }{
    D.m \equiv B_j.m
  }
\]
\end{definition}

\begin{remark}[\fstar{} comparison]
\fstar{} has no concept of method resolution order. It supports simple
(single) inheritance via type refinement, but multiple inheritance with C3
linearization would need to be encoded as a type-level computation---possible
but extremely verbose ($\sim$100+ lines to model the C3 algorithm). \synhopy{}
builds MRO into the type theory as a section of a fiber bundle.
\end{remark}

\section{\texttt{super()} and Cooperative Inheritance}

The \texttt{super()} function navigates the MRO:
\begin{lstlisting}[style=python]
class LoggingMixin:
    def save(self):
        print("Saving...")
        super().save()  # next in MRO, not parent

class DatabaseMixin:
    def save(self):
        self.db.commit()
        super().save()

class User(LoggingMixin, DatabaseMixin):
    def save(self):
        super().save()  # calls LoggingMixin.save
\end{lstlisting}

\begin{definition}[\texttt{super()} typing]
\[
  \frac{
    \MRO(D) = [\ldots, C_i, C_{i+1}, \ldots] \qquad
    \text{currently in method of } C_i
  }{
    \texttt{super()}.m \equiv C_{i+1}.m
  }
\]
The \texttt{super()} call is a \emph{shift} along the MRO fiber:
it selects the next class's method, not the parent class's.
\end{definition}

\begin{theorem}[Cooperative inheritance soundness]
If every class in the MRO calls \texttt{super().save()}, then
\texttt{User().save()} calls every \texttt{save} method in the MRO exactly once,
in MRO order.
\end{theorem}
\begin{proof}
By induction on the MRO length.  Each \texttt{super()} call advances the MRO
position by one (the shift operation).  The base case is \texttt{object.save()}
(which is a no-op).  No class is visited twice because C3 linearization
produces a linear order.
\end{proof}

\section{C3 Linearization: Formal Specification}

The C3 algorithm takes a class $C$ with bases $B_1, \ldots, B_k$ and computes
a linearization $L(C)$.  We give a fully formal, recursive specification.

\begin{definition}[C3 linearization algorithm]
\label{def:c3-algorithm}
Define the \emph{merge} of a list of lists $L_1, \ldots, L_k$:
\[
  \mathsf{merge}(L_1, \ldots, L_k) \defeq
  \begin{cases}
    [] & \text{if all } L_i = [] \\
    h :: \mathsf{merge}(L_1', \ldots, L_k') &
      \text{where } h \text{ is the first good head}
  \end{cases}
\]
A class $h$ is a \emph{good head} if it is the head of some $L_i$ and does not
appear in the \emph{tail} of any $L_j$.  When $h$ is selected, $L_i'$ is $L_i$
with $h$ removed from the head (if present) or left unchanged.
The linearization of class $C$ with bases $B_1, \ldots, B_k$ is:
\[
  L(C) \defeq [C] + \mathsf{merge}(L(B_1), \ldots, L(B_k), [B_1, \ldots, B_k])
\]
\end{definition}

\begin{theorem}[C3 correctness]
\label{thm:c3-correct}
The C3 linearization satisfies two key properties:
\begin{enumerate}
  \item \textbf{Local precedence order}: if $C$ lists bases as
    $(B_1, \ldots, B_k)$, then $B_i$ precedes $B_j$ in $L(C)$ whenever
    $i < j$.
  \item \textbf{Monotonicity}: if $D$ appears before $E$ in $L(B_i)$
    for some base $B_i$, then $D$ appears before $E$ in $L(C)$.
\end{enumerate}
\end{theorem}
\begin{proof}
\emph{Local precedence order}: the merge appends the list
$[B_1, \ldots, B_k]$ as the last argument, so the left-to-right order of
bases is preserved.  When a good head $h$ is chosen, any head of an earlier
list $L_i$ that is also a good head will be chosen before heads of later lists
because the algorithm scans left to right.

\emph{Monotonicity}: suppose $D$ precedes $E$ in $L(B_i)$.  In the merge, $D$
cannot be selected after $E$ because $D$ must be extracted from $L(B_i)$ before
$E$ (as $D$ precedes $E$ in $L(B_i)$), and the good-head condition ensures $D$
is never blocked by $E$ appearing in a tail when $E$ has not yet been selected.
\end{proof}

\begin{remark}[\fstar{} comparison]
Expressing the C3 algorithm in \fstar{} is feasible but onerous: the merge
operation requires a decreasing measure on the total length of all input
lists, and the good-head search requires a decidable predicate over list tails.
A complete \fstar{} formalization would be approximately 150 lines.
\synhopy{} encodes C3 as a type-level computation that is checked once
and cached, so the user never writes merge logic manually.
\end{remark}

\section{The Diamond Problem}

The diamond inheritance pattern is the canonical challenge for multiple
inheritance.

\begin{example}[Diamond inheritance]
\label{ex:diamond}
Consider the classic diamond:
\begin{lstlisting}[style=python]
class A:
    def method(self): return "A"
class B(A):
    def method(self): return "B"
class C(A):
    def method(self): return "C"
class D(B, C):
    pass
\end{lstlisting}
The C3 linearization produces:
\[
  L(D) = [D, B, C, A, \mathsf{object}]
\]
so \lstinline[style=python]{D().method()} returns \texttt{"B"}.
\end{example}

\begin{theorem}[Diamond resolution uniqueness]
\label{thm:diamond-unique}
For any diamond $D$ with bases $B_1, B_2$ that share a common ancestor $A$:
\[
  \frac{
    L(D) = [D, B_1, \ldots, B_2, \ldots, A, \ldots] \qquad
    B_1 \text{ precedes } B_2 \text{ in } L(D)
  }{
    D.m = B_1.m \quad \text{(if both $B_1$ and $B_2$ define $m$)}
  }
\]
The method is always resolved to the \emph{first} base that defines it in the
linearization, eliminating ambiguity.
\end{theorem}
\begin{proof}
By the local precedence order property of C3 (Theorem~\ref{thm:c3-correct}),
$B_1$ precedes $B_2$ in $L(D)$ because $B_1$ is listed before $B_2$ in the
class definition.  The method resolution typing rule selects the first class
in the MRO that defines $m$, which is $B_1$.
\end{proof}

\begin{lstlisting}[style=proof]
-- Proof: diamond resolution for D(B, C) where B, C both extend A
have h_mro: MRO(D) = [D, B, C, A, object]
    by C3_compute [B, C] [L(B), L(C), [B, C]]
have h_B_first: B precedes C in MRO(D)
    by list_index h_mro  -- index(B) = 1 < index(C) = 2
have h_resolve: D.method = B.method
    by method_resolution h_mro h_B_first
conclude: D().method() returns "B"
    by h_resolve
\end{lstlisting}

\section{Method Resolution as Topological Sort}

The MRO can be understood through the lens of DAG theory.  The class
hierarchy forms a directed acyclic graph, and C3 linearization produces
a topological sort of this DAG with additional constraints.

\begin{definition}[Inheritance DAG]
\label{def:inheritance-dag}
The \emph{inheritance DAG} for a class $C$ is the graph $G = (V, E)$ where:
\begin{itemize}
  \item $V$ is the set of all classes reachable from $C$ via the
    \lstinline[style=python]{__bases__} relation (including $C$ and
    \lstinline[style=python]{object}).
  \item $(X, Y) \in E$ iff $Y \in X.\mathsf{\_\_bases\_\_}$ (i.e., $X$
    directly inherits from $Y$).
\end{itemize}
\end{definition}

\begin{theorem}[MRO as constrained topological sort]
\label{thm:mro-topo}
$L(C)$ is a topological sort of the inheritance DAG with two additional
constraints:
\begin{enumerate}
  \item \textbf{Local precedence}: for each class $X$ with bases
    $(B_1, \ldots, B_k)$, $B_i$ precedes $B_j$ in $L(C)$ whenever $i < j$.
  \item \textbf{Monotonicity}: the relative order of classes in $L(B_i)$
    is preserved in $L(C)$.
\end{enumerate}
If no such topological sort exists, C3 raises
\lstinline[style=python]{TypeError}, indicating an inconsistent hierarchy.
\end{theorem}
\begin{proof}
The merge step of C3 constructs a sequence by repeatedly choosing a good head.
A good head $h$ has no predecessor in any tail, which is precisely the condition
for $h$ to be a source in the residual DAG (with the additional local precedence
constraint).  Removing $h$ and repeating yields a topological sort.
If no good head exists but lists are non-empty, the DAG has a cycle under the
precedence constraints, and C3 fails.
\end{proof}

\begin{remark}[\fstar{} comparison]
Topological sorting of a DAG can be verified in \fstar{} using
\lstinline[style=fstar]|Tot| functions with a fuel argument, but the
additional local-precedence and monotonicity constraints of C3 make this
significantly harder.  \synhopy{} internalizes the DAG structure as the
base space of the class fiber bundle, making the topological sort a
structural property of the bundle rather than an algorithmic verification.
\end{remark}

\section{Mixin Classes and Their Verification}

Mixins are classes designed for cooperative multiple inheritance.
They provide methods but are not intended to be instantiated directly.

\begin{definition}[Mixin compatibility]
\label{def:mixin-compat}
A mixin class $M$ is \emph{compatible} with a base class $B$ if:
\begin{enumerate}
  \item Every method $m$ that $M$ overrides calls
    \lstinline[style=python]{super().m(...)},
    ensuring the MRO chain is not broken.
  \item $M$ does not define \lstinline[style=python]{__init__}
    with required arguments (mixins should be argument-transparent).
  \item For each method $m$ in $M$, the signature of $M.m$ is compatible with
    the signature of $B.m$ (if $B$ defines $m$).
\end{enumerate}
\end{definition}

\begin{lstlisting}[style=synhopy]
@guarantee("mixin correctly logs and delegates to next in MRO")
class TimestampMixin:
    def save(self):
        self.updated_at = datetime.now()
        super().save()

    def delete(self):
        print(f"Deleting {self} at {datetime.now()}")
        super().delete()
\end{lstlisting}

\begin{theorem}[Mixin composability]
\label{thm:mixin-compose}
If mixins $M_1, \ldots, M_k$ are each compatible with a base class $B$
(Definition~\ref{def:mixin-compat}), then the class
$C(M_1, \ldots, M_k, B)$ is well-formed and every method call on $C$
invokes each mixin's method and the base method exactly once, in MRO order.
\end{theorem}
\begin{proof}
By induction on $k$.  Base case: $k = 0$, and $C = B$.  Inductive step: by
mixin compatibility, $M_k$ calls \lstinline[style=python]{super().m()}, which
advances to the next class in the MRO.  Since $M_k$ is compatible with $B$
and the mixins $M_1, \ldots, M_{k-1}$ are compatible by hypothesis, the
chain is unbroken.  No class is visited twice because C3 produces a linear
order without repetitions.
\end{proof}

\begin{lstlisting}[style=proof]
-- Verify: TimestampMixin composes with CacheMixin and Model
have h_ts_super: TimestampMixin.save calls super().save
    by structural -- inspect method body
have h_cache_super: CacheMixin.save calls super().save
    by structural
have h_chain: for C(TimestampMixin, CacheMixin, Model),
    C.save() calls all three save methods in MRO order
    by mixin_composability h_ts_super h_cache_super
conclude: composition is sound
    by h_chain
\end{lstlisting}

\section{\texttt{\_\_mro\_entries\_\_} and Custom MRO Hooks}

Python~3.7 introduced \lstinline[style=python]{__mro_entries__}, allowing
objects in the bases tuple to participate in MRO computation without being
classes themselves.

\begin{lstlisting}[style=python]
class GenericAlias:
    def __init__(self, origin, params):
        self.origin = origin
        self.params = params

    def __mro_entries__(self, bases):
        return (self.origin,)

# Now List[int] can appear in bases
class IntList(GenericAlias(list, (int,))):
    pass  # MRO uses list, not GenericAlias
\end{lstlisting}

\begin{definition}[\texttt{\_\_mro\_entries\_\_} typing]
\label{def:mro-entries}
The \lstinline[style=python]{__mro_entries__} hook transforms the declared
bases before C3 linearization:
\[
  \frac{
    B_i.\mathsf{\_\_mro\_entries\_\_}(\mathit{bases}) = [C_1, \ldots, C_j]
  }{
    \text{Replace } B_i \text{ with } C_1, \ldots, C_j \text{ in bases for C3}
  }
\]
If $B_i$ is a class (no \lstinline[style=python]{__mro_entries__}), it is
kept as-is.
\end{definition}

\begin{remark}[\fstar{} comparison]
Custom MRO hooks are a form of computational reflection---the class hierarchy
is modified at class-creation time based on runtime values.  \fstar{} has
no equivalent mechanism; its type hierarchy is fixed at definition time.
\synhopy{} models \lstinline[style=python]{__mro_entries__} as a
\emph{base transformation} applied before the C3 fiber bundle is constructed.
\end{remark}

\section{Worked Example: Verified Plugin System}

We verify a plugin system where plugins register via multiple inheritance
and are dispatched by name.

\begin{lstlisting}[style=synhopy]
@guarantee("registry maps names to exactly one plugin class each")
class PluginBase:
    _registry: dict[str, type] = {}

    def __init_subclass__(cls, **kwargs):
        name = getattr(cls, 'plugin_name', cls.__name__)
        assume("name not already in registry",
               name not in PluginBase._registry)
        PluginBase._registry[name] = cls
        super().__init_subclass__(**kwargs)

    @classmethod
    def get_plugin(cls, name: str):
        check("name is registered", name in cls._registry)
        return cls._registry[name]

class SerializerMixin:
    def serialize(self):
        return {k: getattr(self, k) for k in self.fields}

class ValidatorMixin:
    def validate(self):
        for field in self.fields:
            check(f"{field} is set", hasattr(self, field))

@guarantee("JSON plugin serializes and validates")
class JSONPlugin(SerializerMixin, ValidatorMixin, PluginBase):
    plugin_name = "json"
    fields = ["data", "indent"]
\end{lstlisting}

\begin{theorem}[Plugin system soundness]
\label{thm:plugin-system}
The plugin system satisfies:
\begin{enumerate}
  \item \textbf{Unique registration}: each plugin name maps to exactly one class.
  \item \textbf{Dispatch correctness}:
    \lstinline[style=python]{PluginBase.get_plugin(name)} returns the class
    registered under \texttt{name}.
  \item \textbf{Mixin composability}: \texttt{JSONPlugin.serialize()} and
    \texttt{JSONPlugin.validate()} are both available and correctly
    dispatched via the MRO.
\end{enumerate}
\end{theorem}
\begin{proof}
(1) The \lstinline[style=python]{assume} in
\lstinline[style=python]{__init_subclass__} creates a proof obligation at each
subclass definition site that the name is fresh.
(2) The \lstinline[style=python]{check} in \texttt{get\_plugin} verifies
the name is present before lookup.
(3) The MRO of \texttt{JSONPlugin} is
$[\texttt{JSONPlugin}, \texttt{SerializerMixin}, \texttt{ValidatorMixin},
  \texttt{PluginBase}, \texttt{object}]$,
so \texttt{serialize} resolves to \texttt{SerializerMixin} and
\texttt{validate} resolves to \texttt{ValidatorMixin} by the method
resolution rule.
\end{proof}

\begin{lstlisting}[style=proof]
-- Verify: plugin system end-to-end
have h_reg: JSONPlugin registered as "json" in PluginBase._registry
    by init_subclass_hook; name = "json"; name not in {} (empty)
have h_mro: MRO(JSONPlugin)
    = [JSONPlugin, SerializerMixin, ValidatorMixin, PluginBase, object]
    by C3_compute
have h_serial: JSONPlugin.serialize = SerializerMixin.serialize
    by method_resolution h_mro
have h_valid: JSONPlugin.validate = ValidatorMixin.validate
    by method_resolution h_mro
have h_dispatch: PluginBase.get_plugin("json") = JSONPlugin
    by dict_lookup h_reg
conclude: plugin system is sound
    by h_reg, h_mro, h_serial, h_valid, h_dispatch
\end{lstlisting}


% ──────────────────────────────────────────────────────────────────
\chapter{Metaclasses and Runtime Class Creation}
\label{ch:metaclasses}

\section{Metaclasses as 2-Types}

\begin{lstlisting}[style=python]
class SingletonMeta(type):
    _instances = {}
    def __call__(cls, *args, **kwargs):
        if cls not in cls._instances:
            cls._instances[cls] = super().__call__(*args, **kwargs)
        return cls._instances[cls]

class Database(metaclass=SingletonMeta):
    def __init__(self):
        self.connection = "db://localhost"
\end{lstlisting}

\begin{definition}[Metaclass as 2-type]
A metaclass $M$ is a type whose elements are themselves types (classes).
In HoTT terms, $M$ is a \emph{2-type}: its elements are 1-types (classes),
and paths between elements are class transformations.

\[
  M : \Type_2 \qquad C : M \quad\text{means $C$ is a class with metaclass $M$}
\]
\end{definition}

\begin{definition}[Singleton metaclass typing]
For the singleton metaclass:
\[
  \frac{
    C : \mathsf{SingletonMeta} \qquad
    a, b : C
  }{
    a =_C b
  }
\]
Any two instances of a singleton class are equal (the type is contractible---a
$(-2)$-type).
\end{definition}

\begin{remark}[\fstar{} comparison]
\fstar{} has no concept of metaclasses. Types cannot be created or modified at
runtime. To model the singleton pattern, you'd need a heap-allocated global reference
with a manual invariant---approximately 30 lines of state management that doesn't
capture the metaclass semantics.
\end{remark}

\section{\texttt{\_\_init\_subclass\_\_} and Class Registration}

Python 3.6+ provides a hook for when a class is subclassed:
\begin{lstlisting}[style=python]
class Plugin:
    registry = []
    def __init_subclass__(cls, **kwargs):
        super().__init_subclass__(**kwargs)
        cls.registry.append(cls)

class JSONPlugin(Plugin):  # auto-registered
    pass
class XMLPlugin(Plugin):   # auto-registered
    pass
\end{lstlisting}

\begin{definition}[Subclass hook as path constructor]
The \lstinline[style=python]{__init_subclass__} hook creates a \emph{path}
from the new subclass to the registry:
\[
  \frac{
    C' <: C \qquad
    C.\mathsf{\_\_init\_subclass\_\_}(C') \text{ runs}
  }{
    C' \in C.\mathsf{registry}
  }
\]
This is a side effect of class creation---captured by the $\Mutates$ effect on
the class's registry attribute.
\end{definition}

\begin{theorem}[Registry completeness]
After all subclasses are defined, \texttt{Plugin.registry} contains exactly
the set of direct and indirect subclasses of \texttt{Plugin}.
\end{theorem}
\begin{proof}
By induction on the number of class definitions.  Each
\lstinline[style=python]{class X(Plugin)} triggers
\lstinline[style=python]{__init_subclass__}, which appends $X$ to the registry.
The hook is called exactly once per class definition (by Python semantics),
and no other operation modifies the registry.
\end{proof}

\section{\texttt{type()} as the Default Metaclass}

Every class in Python is an instance of \lstinline[style=python]{type}.
The three-argument form of \lstinline[style=python]{type} creates classes
at runtime.

\begin{lstlisting}[style=python]
# These two are equivalent:
class Foo:
    x = 42
    def greet(self): return "hello"

Foo = type('Foo', (), {'x': 42, 'greet': lambda self: "hello"})
\end{lstlisting}

\begin{definition}[\texttt{type(name, bases, dict)} typing]
\label{def:type-three-arg}
The three-argument form of \lstinline[style=python]{type} has the signature:
\[
  \mathsf{type} : \Sigma(\mathit{name} : \mathsf{String},\;
    \mathit{bases} : \mathsf{List}(\PyClass),\;
    \mathit{dict} : \PyDict).\; \PyClass
\]
with the postcondition:
\[
  \frac{
    C = \mathsf{type}(n, [B_1, \ldots, B_k], d)
  }{
    C.\mathsf{\_\_name\_\_} = n \;\wedge\;
    C.\mathsf{\_\_bases\_\_} = (B_1, \ldots, B_k) \;\wedge\;
    \forall (k,v) \in d.\; \mathsf{getattr}(C, k) = v
  }
\]
\end{definition}

\begin{remark}[\fstar{} comparison]
\fstar{} types must be declared statically---there is no equivalent of
\lstinline[style=python]{type(name, bases, dict)} that creates new types
at runtime.  All type information must be known at verification time.
\synhopy{} treats runtime class creation as a $\Mutates$ effect that
extends the type universe dynamically, subject to the metaclass's
invariants.
\end{remark}

\section{Abstract Base Classes and Interface Enforcement}

Python's \lstinline[style=python]{abc.ABCMeta} metaclass enforces that
subclasses implement required methods.

\begin{lstlisting}[style=python]
from abc import ABCMeta, abstractmethod

class Shape(metaclass=ABCMeta):
    @abstractmethod
    def area(self) -> float: ...

    @abstractmethod
    def perimeter(self) -> float: ...

class Circle(Shape):
    def __init__(self, r):
        self.r = r
    def area(self): return 3.14159 * self.r ** 2
    def perimeter(self): return 2 * 3.14159 * self.r

# Shape() would raise TypeError (abstract)
\end{lstlisting}

\begin{definition}[ABCMeta as interface type]
\label{def:abcmeta}
An abstract base class $A$ with abstract methods $m_1, \ldots, m_k$ defines
an \emph{interface constraint}:
\[
  \frac{
    C <: A \qquad
    \forall i \in \{1, \ldots, k\}.\; m_i \in \mathsf{methods}(C) \;\wedge\;
      m_i \text{ is not abstract in } C
  }{
    C \text{ is concrete (instantiable)}
  }
\]
If the condition fails, \lstinline[style=python]{C()} raises
\lstinline[style=python]{TypeError} at runtime.  In \synhopy{}, this
is a \emph{type error at class-creation time}.
\end{definition}

\begin{theorem}[ABCMeta completeness]
\label{thm:abc-complete}
For a class $C$ with metaclass \lstinline[style=python]{ABCMeta}:
$C$ is instantiable if and only if $C$ has no abstract methods (i.e.,
$C.\mathsf{\_\_abstractmethods\_\_} = \emptyset$).
\end{theorem}
\begin{proof}
$(\Rightarrow)$: if $C$ is instantiable, then
\lstinline[style=python]{ABCMeta.__call__} does not raise, which requires
$C.\mathsf{\_\_abstractmethods\_\_} = \emptyset$.
$(\Leftarrow)$: if $C.\mathsf{\_\_abstractmethods\_\_} = \emptyset$,
then the check in \lstinline[style=python]{ABCMeta.__call__} passes and
$C()$ proceeds to \lstinline[style=python]{C.__init__}.
\end{proof}

\begin{remark}[\fstar{} comparison]
\fstar{} models interfaces via \emph{typeclasses} (or, more precisely, via
dependent records).  A typeclass instance must provide implementations for
all required methods, which is checked statically.  The \synhopy{} approach
is similar in spirit but handles the dynamic case: abstract method sets can
be modified at runtime via \lstinline[style=python]{register()} and virtual
subclasses.
\end{remark}

\section{\texttt{\_\_prepare\_\_} and Ordered Namespaces}

The \lstinline[style=python]{__prepare__} class method on a metaclass controls
the namespace object used during class body execution.

\begin{lstlisting}[style=python]
from collections import OrderedDict

class OrderedMeta(type):
    @classmethod
    def __prepare__(mcs, name, bases):
        return OrderedDict()  # class body uses OrderedDict

    def __new__(mcs, name, bases, namespace):
        cls = super().__new__(mcs, name, bases, dict(namespace))
        cls._field_order = list(namespace.keys())
        return cls

class Record(metaclass=OrderedMeta):
    x = 1
    y = 2
    z = 3
    # Record._field_order == ['x', 'y', 'z']
\end{lstlisting}

\begin{definition}[\texttt{\_\_prepare\_\_} as namespace fiber]
\label{def:prepare}
The \lstinline[style=python]{__prepare__} hook selects the fiber of the
class body namespace:
\[
  \frac{
    M.\mathsf{\_\_prepare\_\_}(n, \mathit{bases}) = \mathit{ns}
  }{
    \text{Class body of } n \text{ executes in namespace } \mathit{ns}
  }
\]
The namespace $\mathit{ns}$ determines what data structure backs attribute
definitions during class creation.  The default is a plain
\lstinline[style=python]{dict}; alternatives include
\lstinline[style=python]{OrderedDict} or custom mappings.
\end{definition}

\begin{example}[Ordered fields via \texttt{\_\_prepare\_\_}]
Using \lstinline[style=python]{OrderedMeta}, the class
\lstinline[style=python]{Record} preserves the declaration order of
attributes.  This is essential for serialization formats (JSON, CSV)
where field order matters.  The type-level invariant is:
\[
  \mathsf{Record}.\mathsf{\_field\_order} =
    [\mathtt{"x"}, \mathtt{"y"}, \mathtt{"z"}]
\]
\end{example}

\begin{remark}[\fstar{} comparison]
\fstar{} record fields have a fixed order determined by the type definition,
but this order is a property of the syntax, not a runtime-queryable value.
The \lstinline[style=python]{__prepare__} mechanism is a form of
\emph{reflection} that \fstar{} lacks entirely.  \synhopy{} models it as
a dependent namespace type parameterized by the metaclass.
\end{remark}

\section{Runtime Type Creation and Verification}

Dynamic type creation patterns are common in frameworks (ORMs, serializers,
RPC stubs).  We formalize the verification challenge.

\begin{definition}[Runtime class factory]
\label{def:class-factory}
A \emph{class factory} is a function
$F : \mathsf{Schema} \to \PyClass$ that creates a new class from a
schema description.  The factory is \emph{correct} if:
\[
  \forall s : \mathsf{Schema}.\;
    F(s).\mathsf{\_\_annotations\_\_} = \mathsf{fields}(s)
    \;\wedge\;
    \forall (k, \tau) \in \mathsf{fields}(s).\;
      \mathsf{type}(\mathsf{getattr}(F(s)(), k)) <: \tau
\]
\end{definition}

\begin{lstlisting}[style=synhopy]
@guarantee("created class has all schema fields with correct types")
def make_dataclass(name: str, fields: list[tuple[str, type]]):
    annotations = {fname: ftype for fname, ftype in fields}
    def init(self, **kwargs):
        for fname, ftype in fields:
            check(f"{fname} has type {ftype}",
                  isinstance(kwargs.get(fname), ftype))
            setattr(self, fname, kwargs[fname])
    cls = type(name, (), {
        '__init__': init,
        '__annotations__': annotations,
    })
    return cls

Point = make_dataclass("Point", [("x", float), ("y", float)])
\end{lstlisting}

\begin{remark}[\fstar{} comparison]
\fstar{} cannot create types at runtime.  A factory like
\lstinline[style=python]{make_dataclass} would need to be implemented via
\emph{metaprogramming} (\lstinline[style=fstar]|Meta|), which operates at
compile time and produces static types.  Runtime schema-driven creation
is fundamentally out of scope for \fstar{}'s type system.
\end{remark}

\section{Metaclass Conflicts and Resolution}

When a class has bases with different metaclasses, Python must resolve
the conflict.

\begin{lstlisting}[style=python]
class MetaA(type): pass
class MetaB(type): pass
class A(metaclass=MetaA): pass
class B(metaclass=MetaB): pass

# class C(A, B): pass  # TypeError: metaclass conflict
\end{lstlisting}

\begin{definition}[Metaclass conflict]
\label{def:meta-conflict}
A class $C$ with bases $B_1, \ldots, B_k$ has a \emph{metaclass conflict}
when there is no single metaclass $M$ such that $M$ is a subclass of
$\mathsf{type}(B_i)$ for all $i$.  Formally:
\[
  \nexists\; M.\; \forall i \in \{1, \ldots, k\}.\; M <: \mathsf{type}(B_i)
\]
If a most-derived metaclass exists (i.e., one of the $\mathsf{type}(B_i)$
is a subclass of all the others), it is used.  Otherwise, Python raises
\lstinline[style=python]{TypeError}.
\end{definition}

\begin{theorem}[Metaclass resolution]
\label{thm:meta-resolve}
For a class $C$ with bases $B_1, \ldots, B_k$, the metaclass of $C$ is
the most-derived metaclass among $\{\mathsf{type}(B_i)\}_{i=1}^{k}$:
\[
  \mathsf{metaclass}(C) =
  \begin{cases}
    M & \text{if } \exists\, M \in \{\mathsf{type}(B_i)\}.\;
        \forall i.\; M <: \mathsf{type}(B_i) \\
    \bot & \text{otherwise (TypeError)}
  \end{cases}
\]
\end{theorem}
\begin{proof}
Python's \lstinline[style=python]{type.__new__} computes the most-derived
metaclass by iterating over the bases' metaclasses and checking
\lstinline[style=python]{issubclass} pairwise.  If a unique winner exists,
it is used; otherwise, \lstinline[style=python]{TypeError} is raised.
\end{proof}

\begin{lstlisting}[style=python]
# Resolution: MetaC inherits from both MetaA and MetaB
class MetaC(MetaA, MetaB): pass
class C(A, B, metaclass=MetaC): pass  # OK: MetaC <: MetaA, MetaC <: MetaB
\end{lstlisting}

\begin{remark}[\fstar{} comparison]
Metaclass conflicts have no analogue in \fstar{}, since \fstar{} does not
have metaclasses or runtime type creation.  The conflict resolution
algorithm is a form of lattice computation on the metaclass hierarchy,
which \synhopy{} models as a join operation in the 2-type lattice.
\end{remark}

\section{Worked Example: Verified ORM Metaclass}

We verify an ORM metaclass that generates database schemas from
class definitions.

\begin{lstlisting}[style=synhopy]
@guarantee("ORM metaclass produces correct schema and field descriptors")
class ModelMeta(type):
    def __prepare__(mcs, name, bases):
        return OrderedDict()

    def __new__(mcs, name, bases, namespace):
        fields = {}
        for key, val in namespace.items():
            if isinstance(val, Field):
                fields[key] = val
        cls = super().__new__(mcs, name, bases, dict(namespace))
        cls._fields = fields
        cls._table_name = name.lower() + "s"
        return cls

    def create_table_sql(cls):
        cols = ", ".join(
            f"{name} {field.sql_type}"
            for name, field in cls._fields.items()
        )
        return f"CREATE TABLE {cls._table_name} ({cols})"

class Field:
    def __init__(self, sql_type: str, nullable: bool = False):
        self.sql_type = sql_type
        self.nullable = nullable

class User(metaclass=ModelMeta):
    id = Field("INTEGER PRIMARY KEY")
    name = Field("TEXT NOT NULL")
    email = Field("TEXT UNIQUE")
\end{lstlisting}

\begin{theorem}[ORM metaclass correctness]
\label{thm:orm-metaclass}
The ORM metaclass satisfies:
\begin{enumerate}
  \item \textbf{Field extraction}: $C.\mathsf{\_fields}$ contains exactly the
    \lstinline[style=python]{Field} instances defined in the class body.
  \item \textbf{Table naming}: $C.\mathsf{\_table\_name}$ is the lowercase
    class name suffixed with \texttt{"s"}.
  \item \textbf{SQL correctness}: \lstinline[style=python]{C.create_table_sql()}
    produces a valid \texttt{CREATE TABLE} statement with one column per field.
  \item \textbf{Order preservation}: fields appear in declaration order
    (guaranteed by \lstinline[style=python]{__prepare__} returning
    \lstinline[style=python]{OrderedDict}).
\end{enumerate}
\end{theorem}
\begin{proof}
(1) The \lstinline[style=python]{__new__} method iterates over the namespace
and collects entries that are instances of \lstinline[style=python]{Field}.
Since the namespace contains exactly the class body definitions, the
extracted fields are correct.
(2) By construction: $\mathsf{\_table\_name} = \mathsf{name.lower()} + \mathtt{"s"}$.
(3) The \texttt{create\_table\_sql} method joins field names and SQL types
with commas, producing a syntactically valid \texttt{CREATE TABLE} statement
(assuming field SQL types are valid).
(4) Because \lstinline[style=python]{__prepare__} returns an
\lstinline[style=python]{OrderedDict}, the namespace preserves insertion
order, so \texttt{\_fields.items()} yields fields in declaration order.
\end{proof}

\begin{lstlisting}[style=proof]
-- Verify: ORM metaclass end-to-end
have h_fields: User._fields = {"id": Field("INTEGER PRIMARY KEY"),
    "name": Field("TEXT NOT NULL"), "email": Field("TEXT UNIQUE")}
    by metaclass_new; filter namespace for Field instances
have h_table: User._table_name = "users"
    by metaclass_new; "User".lower() + "s"
have h_sql: User.create_table_sql() =
    "CREATE TABLE users (id INTEGER PRIMARY KEY, name TEXT NOT NULL,
     email TEXT UNIQUE)"
    by unfold create_table_sql; join h_fields
have h_order: fields appear in order [id, name, email]
    by __prepare__ returns OrderedDict; insertion order preserved
conclude: ORM metaclass is correct
    by h_fields, h_table, h_sql, h_order
\end{lstlisting}


% ──────────────────────────────────────────────────────────────────
\chapter{\texttt{*args, **kwargs} and the Dict Sheaf}
\label{ch:args-kwargs}

\section{Variable Arguments as Dependent Types}

\begin{lstlisting}[style=python]
def flexible(*args, **kwargs):
    return kwargs.get('key', args[0] if args else None)
\end{lstlisting}

\begin{definition}[Dependent argument type]
\begin{align}
  \mathsf{args} &: \mathsf{List}(\PyObj) \\
  \mathsf{kwargs} &: \Sigma(\mathit{keys} : \mathsf{Set}(\mathsf{String})).\;
    \Pi(k : \mathit{keys}).\; \PyObj
\end{align}
The \texttt{**kwargs} type is a dependent dictionary: the set of keys is part of
the type, and each key maps to a $\PyObj$.
\end{definition}

\section{Specification via \v{C}ech Descent}

The spec of a variadic function is a fiberwise property over the argument
configurations:

\begin{definition}[\v{C}ech cover of argument space]
For a function $f$ with \texttt{*args, **kwargs}:
\begin{enumerate}
  \item Let $\mathcal{U}$ be the set of valid argument configurations
  \item Each $U_i \in \mathcal{U}$ is a constraint on the arguments (e.g.,
    ``\texttt{key} is present in \texttt{kwargs}'')
  \item The specification is: $\forall U_i.\; x \in U_i \Rightarrow f(x)$ satisfies
    $\Spec_i$
  \item The global specification is the \v{C}ech descent: the $\Spec_i$ must
    agree on overlaps $U_i \cap U_j$
\end{enumerate}
\end{definition}

\begin{remark}[\fstar{} comparison]
\fstar{} requires fixed-arity, statically-typed function signatures. Python's
\texttt{*args/**kwargs} cannot be expressed in \fstar{}'s type system without
encoding them as explicit data structures (e.g., a list and a map), which loses
the call-site-dependent typing information.
\end{remark}

\section{Worked Example: Decorator Factory}

\begin{lstlisting}[style=synhopy]
@guarantee("returns a decorator that retries f up to n times")
def retry(n: int = 3, exceptions=(Exception,)):
    def decorator(f):
        def wrapper(*args, **kwargs):
            for attempt in range(n):
                try:
                    return f(*args, **kwargs)
                except exceptions:
                    if attempt == n - 1:
                        raise
        return wrapper
    return decorator
\end{lstlisting}

The specification requires reasoning about:
\begin{enumerate}
  \item \textbf{Args/kwargs forwarding}: \texttt{wrapper} passes through
    all arguments to \texttt{f} unchanged.
  \item \textbf{Exception filtering}: only exceptions matching
    \texttt{exceptions} are caught.
  \item \textbf{Retry count}: exactly $n$ attempts are made.
  \item \textbf{Transparency}: if \texttt{f} succeeds on attempt $k < n$,
    \texttt{wrapper} returns the same value.
\end{enumerate}

\begin{lstlisting}[style=proof]
-- Property: wrapper(*args, **kwargs) = f(*args, **kwargs)
--   when f does not raise
have h_forward: wrapper forwards (args, kwargs) to f
    by structural -- wrapper calls f(*args, **kwargs)
have h_success: if f succeeds on attempt 0,
    wrapper returns f(*args, **kwargs)
    by unfold wrapper; cases on try block
have h_retry: if f raises E in exceptions on
    attempts 0..k-1 and succeeds on k,
    wrapper returns f(*args, **kwargs) at attempt k
    by nat_induction on attempt with h_forward
conclude: wrapper is a transparent retry of f
    by cases on first successful attempt with h_success, h_retry
\end{lstlisting}

\begin{definition}[Call-site dependent typing for \texttt{*args}]
At each call site, we know the actual arguments:
\[
  \frac{
    \Gamma \vdash f : \Pi(\mathsf{args}: \mathsf{List}(\PyObj),\;
      \mathsf{kwargs}: \PyDict).\; R \qquad
    \Gamma \vdash f(a_1, a_2, k_1\!=\!v_1)
  }{
    \Gamma \vdash f(\mathsf{args}\!=\![a_1, a_2],\;
      \mathsf{kwargs}\!=\!\{k_1: v_1\}) : R[\mathsf{args}/[a_1,a_2]]
  }
\]
The actual arguments are substituted into the dependent return type.
\end{definition}

\section{Positional-Only and Keyword-Only Parameters}

Python~3.8 formalized the \texttt{/} and \texttt{*} syntax for
positional-only and keyword-only parameters.

\begin{lstlisting}[style=python]
def example(pos_only, /, normal, *, kw_only):
    return (pos_only, normal, kw_only)

example(1, 2, kw_only=3)      # OK
example(1, normal=2, kw_only=3)  # OK
# example(pos_only=1, normal=2, kw_only=3)  # TypeError
\end{lstlisting}

\begin{definition}[Parameter classification]
\label{def:param-class}
A Python function's parameters are partitioned into three zones by the
\texttt{/} and \texttt{*} markers:
\[
  f(
    \underbrace{p_1, \ldots, p_j}_{
      \text{positional-only}},\;
    /,\;
    \underbrace{p_{j+1}, \ldots, p_k}_{
      \text{positional-or-keyword}},\;
    *,\;
    \underbrace{p_{k+1}, \ldots, p_n}_{
      \text{keyword-only}}
  )
\]
The typing rules for each zone are:
\begin{align}
  \text{positional-only } p_i &: \text{may only be passed by position} \\
  \text{positional-or-keyword } p_i &: \text{may be passed by position or name} \\
  \text{keyword-only } p_i &: \text{must be passed by name}
\end{align}
\end{definition}

\begin{theorem}[Parameter zone disjointness]
\label{thm:param-zones}
At any call site, each parameter $p_i$ is bound by exactly one argument.
The three zones ensure no parameter can be simultaneously bound by position
and by name.
\end{theorem}
\begin{proof}
Positional-only parameters cannot be named, so they can only be bound by
position.  Keyword-only parameters cannot appear positionally (they follow
\texttt{*}), so they can only be bound by name.  Positional-or-keyword
parameters may be bound by either, but Python raises
\lstinline[style=python]{TypeError} if the same parameter is supplied both
positionally and by keyword.
\end{proof}

\begin{remark}[\fstar{} comparison]
\fstar{} parameters are always positional (bound by order in the function
signature).  There is no keyword argument mechanism, and no concept of
positional-only vs.\ keyword-only distinctions.  \synhopy{} models the
three parameter zones as a dependent sum that constrains how the argument
binding algorithm operates.
\end{remark}

\section{Parameter Inspection via \texttt{inspect.signature}}

Python's \lstinline[style=python]{inspect} module provides runtime access
to function signatures.

\begin{lstlisting}[style=python]
import inspect

def f(a: int, b: str = "hi", *args, key: bool = True, **kwargs):
    pass

sig = inspect.signature(f)
for name, param in sig.parameters.items():
    print(name, param.kind, param.default, param.annotation)
\end{lstlisting}

\begin{definition}[Signature reification]
\label{def:sig-reify}
The function $\mathsf{signature} : \PyFunc \to \mathsf{Sig}$ reifies
a function's type information at runtime.  The signature type is:
\[
  \mathsf{Sig} \defeq \mathsf{List}\bigl(
    \Sigma(\mathit{name} : \mathsf{String},\;
           \mathit{kind} : \mathsf{ParamKind},\;
           \mathit{default} : \mathsf{Option}(\PyObj),\;
           \mathit{annotation} : \mathsf{Option}(\PyType))
  \bigr)
\]
where $\mathsf{ParamKind} \in \{
  \mathsf{POS\_ONLY}, \mathsf{POS\_OR\_KW}, \mathsf{VAR\_POS},
  \mathsf{KW\_ONLY}, \mathsf{VAR\_KW}
\}$.
\end{definition}

\begin{example}[Signature of \texttt{f}]
For the function $f$ above:
\begin{align}
  \mathsf{signature}(f) = [
    &(\mathtt{"a"},\; \mathsf{POS\_OR\_KW},\; \PyNone,\; \mathsf{int}), \\
    &(\mathtt{"b"},\; \mathsf{POS\_OR\_KW},\; \mathtt{"hi"},\; \mathsf{str}), \\
    &(\mathtt{"args"},\; \mathsf{VAR\_POS},\; \PyNone,\; \PyNone), \\
    &(\mathtt{"key"},\; \mathsf{KW\_ONLY},\; \mathsf{True},\; \mathsf{bool}), \\
    &(\mathtt{"kwargs"},\; \mathsf{VAR\_KW},\; \PyNone,\; \PyNone)
  ]
\end{align}
\end{example}

\begin{remark}[\fstar{} comparison]
\fstar{} does not expose function signatures as runtime values.
Function types are first-class, but their parameter names, defaults, and
kinds are not queryable at runtime.  \synhopy{} treats
\lstinline[style=python]{inspect.signature} as a \emph{reification}
operation that reflects the type-level signature into a value-level record.
\end{remark}

\section{TypedDict for Typed \texttt{**kwargs}}

\lstinline[style=python]{TypedDict} provides per-key typing for
dictionaries, making it ideal for typed \texttt{**kwargs}.

\begin{lstlisting}[style=python]
from typing import TypedDict, Unpack

class Options(TypedDict, total=False):
    timeout: int
    retries: int
    verbose: bool

def fetch(url: str, **kwargs: Unpack[Options]) -> bytes:
    timeout = kwargs.get('timeout', 30)
    retries = kwargs.get('retries', 3)
    ...
\end{lstlisting}

\begin{definition}[TypedDict as dependent record]
\label{def:typeddict-kwargs}
A \lstinline[style=python]{TypedDict} $T$ with fields
$(k_1 : \tau_1, \ldots, k_n : \tau_n)$ defines a dependent dictionary type:
\[
  T \defeq \Sigma(\mathit{keys} \subseteq \{k_1, \ldots, k_n\}).\;
    \Pi(k \in \mathit{keys}).\; \tau_k
\]
When used with \lstinline[style=python]{Unpack[T]} in \texttt{**kwargs},
the call-site typing becomes:
\[
  \frac{
    \Gamma \vdash f : (x : A,\; \mathsf{**kwargs}: \mathsf{Unpack}[T]) \to R
    \qquad
    \Gamma \vdash f(a,\; k_i\!=\!v_i, \ldots)
    \qquad
    \forall i.\; v_i : \tau_{k_i}
  }{
    \Gamma \vdash f(a,\; k_i\!=\!v_i, \ldots) : R
  }
\]
\end{definition}

\begin{theorem}[TypedDict kwargs soundness]
\label{thm:typeddict-sound}
If $f$ declares \lstinline[style=python]{**kwargs: Unpack[T]} and a call
site provides keyword arguments $k_i = v_i$, then:
\begin{enumerate}
  \item Each $k_i$ is a valid key of $T$.
  \item Each $v_i$ has type $\tau_{k_i}$.
  \item All required keys of $T$ (those with \lstinline[style=python]{total=True})
    are provided.
\end{enumerate}
\end{theorem}
\begin{proof}
(1) and (2) follow from the \lstinline[style=python]{Unpack} constraint, which
expands $T$'s fields into the function signature.
(3) follows from the \lstinline[style=python]{total} flag on
\lstinline[style=python]{TypedDict}: required keys must be present at every
call site.
\end{proof}

\begin{remark}[\fstar{} comparison]
\fstar{} has no dictionary type with per-key typing.  The closest analogue
is a dependent record type, but \fstar{} records are statically declared
and cannot model the optional-key semantics of
\lstinline[style=python]{TypedDict(total=False)}.  \synhopy{} represents
\lstinline[style=python]{TypedDict} as a dependent sigma type where the
key set is a parameter.
\end{remark}

\section{Overloaded Functions and Variadic Dispatch}

Python's \lstinline[style=python]{@overload} decorator from
\lstinline[style=python]{typing} enables type-level dispatch.

\begin{lstlisting}[style=python]
from typing import overload

@overload
def process(x: int) -> str: ...
@overload
def process(x: str) -> int: ...
@overload
def process(x: int, y: int) -> float: ...

def process(*args):
    if len(args) == 1 and isinstance(args[0], int):
        return str(args[0])
    elif len(args) == 1 and isinstance(args[0], str):
        return len(args[0])
    elif len(args) == 2:
        return args[0] / args[1]
\end{lstlisting}

\begin{definition}[Overload resolution as coproduct elimination]
\label{def:overload-variadic}
An overloaded function with $n$ overloads defines a coproduct type:
\[
  f : \coprod_{i=1}^{n} (A_i \to R_i)
\]
The dispatch at each call site selects the matching overload:
\[
  \frac{
    \Gamma \vdash f : \coprod_{i=1}^{n} (A_i \to R_i) \qquad
    \Gamma \vdash \vec{a} : A_j \qquad
    j = \min\{i \mid \vec{a} : A_i\}
  }{
    \Gamma \vdash f(\vec{a}) : R_j
  }
\]
The first matching overload (in declaration order) is selected.
\end{definition}

\begin{remark}[\fstar{} comparison]
\fstar{} does not support function overloading.  Multiple functions with
the same name are disallowed.  The \synhopy{} approach models overloads
as a coproduct (sum type) of function types, with dispatch governed by
runtime type checks.
\end{remark}

\section{Verifying the Argument Binding Algorithm}

Python's argument binding algorithm (PEP~3102, PEP~570) matches actual
arguments to formal parameters.  We formalize and verify this algorithm.

\begin{definition}[Argument binding]
\label{def:arg-binding}
Given a function signature
$\mathit{sig} = (p_1, \ldots, p_n)$ (with kinds and defaults) and a call
$f(a_1, \ldots, a_j,\; k_1\!=\!v_1, \ldots, k_m\!=\!v_m)$,
the \emph{binding} $\beta : \mathsf{Param} \to \PyObj$ is computed as follows:
\begin{enumerate}
  \item \textbf{Positional binding}: for $i = 1, \ldots, j$, bind $a_i$ to the
    $i$-th non-\texttt{VAR\_POS} parameter.  If $i$ exceeds the number of
    such parameters and a \texttt{VAR\_POS} parameter exists, collect
    remaining positional arguments into \texttt{*args}.
  \item \textbf{Keyword binding}: for each $k_\ell = v_\ell$, find a parameter
    $p$ with $p.\mathit{name} = k_\ell$ and $p.\mathit{kind} \neq
    \mathsf{POS\_ONLY}$.  If no such $p$ exists and a \texttt{VAR\_KW}
    parameter exists, add $(k_\ell, v_\ell)$ to \texttt{**kwargs}.
  \item \textbf{Default filling}: for any unbound parameter $p$ with
    $p.\mathit{default} \neq \PyNone$, set $\beta(p) = p.\mathit{default}$.
  \item \textbf{Error detection}: if any parameter without a default remains
    unbound, or if a positional-only parameter is passed by keyword, raise
    \lstinline[style=python]{TypeError}.
\end{enumerate}
\end{definition}

\begin{theorem}[Argument binding soundness]
\label{thm:arg-binding}
If the binding algorithm produces $\beta$ without error, then:
\begin{enumerate}
  \item Every parameter is bound to exactly one value.
  \item No positional-only parameter is bound by keyword.
  \item No keyword-only parameter is bound by position.
  \item All excess positional arguments are collected in \texttt{*args}.
  \item All excess keyword arguments are collected in \texttt{**kwargs}.
\end{enumerate}
\end{theorem}
\begin{proof}
By construction of the algorithm.  Step~1 binds positional parameters
left to right, ensuring no double-binding.  Step~2 binds keyword arguments
to named parameters, checking for conflicts with step~1 (raising
\lstinline[style=python]{TypeError} on duplicates).  Step~3 fills defaults
only for unbound parameters.  Step~4 detects missing bindings.
The zone constraints (positional-only, keyword-only) are enforced by the
kind checks in steps~1 and~2.
\end{proof}

\begin{lstlisting}[style=proof]
-- Verify: binding for f(a, /, b, *, c) called as f(1, 2, c=3)
have h_pos: a is bound to 1 (positional-only, position 0)
    by positional_binding; a.kind = POS_ONLY
have h_normal: b is bound to 2 (positional-or-keyword, position 1)
    by positional_binding; b.kind = POS_OR_KW
have h_kw: c is bound to 3 (keyword-only, keyword "c")
    by keyword_binding; c.kind = KW_ONLY
have h_all: all parameters bound
    by h_pos, h_normal, h_kw; no unbound params remain
conclude: binding is sound
    by arg_binding_soundness h_pos h_normal h_kw h_all
\end{lstlisting}

\begin{remark}[\fstar{} comparison]
\fstar{} argument binding is trivial: arguments are bound by position,
and all arguments are required (no defaults, no \texttt{*args}, no
\texttt{**kwargs}).  The full Python binding algorithm with its five
parameter kinds, defaults, \texttt{*args}/\texttt{**kwargs} collection,
and error conditions is far more complex.  \synhopy{} internalizes this
algorithm as a type-level computation, verified once and applied at every
call site.
\end{remark}

\section{Worked Example: Verified CLI Argument Parser}

We verify a CLI argument parser that uses \texttt{*args} and
\texttt{**kwargs} to handle command-line arguments flexibly.

\begin{lstlisting}[style=synhopy]
@guarantee("parser correctly binds CLI args to handler parameters")
class CLIParser:
    def __init__(self):
        self._commands: dict[str, tuple] = {}

    @guarantee("registers handler with its signature metadata")
    def command(self, name: str):
        def decorator(func):
            import inspect
            sig = inspect.signature(func)
            self._commands[name] = (func, sig)
            return func
        return decorator

    @guarantee("dispatches to correct handler with correct bindings")
    def dispatch(self, argv: list[str]):
        check("argv is non-empty", len(argv) > 0)
        cmd_name = argv[0]
        check("command exists", cmd_name in self._commands)
        func, sig = self._commands[cmd_name]

        args = []
        kwargs = {}
        i = 1
        while i < len(argv):
            if argv[i].startswith("--"):
                key = argv[i][2:]
                check("key-value pair", i + 1 < len(argv))
                kwargs[key] = argv[i + 1]
                i += 2
            else:
                args.append(argv[i])
                i += 1

        bound = sig.bind(*args, **kwargs)
        bound.apply_defaults()
        return func(*bound.args, **bound.kwargs)

cli = CLIParser()

@cli.command("greet")
@guarantee("returns greeting string")
def greet(name: str, /, greeting: str = "Hello", *, excited: bool = False):
    msg = f"{greeting}, {name}!"
    return msg + "!!" if excited else msg
\end{lstlisting}

\begin{theorem}[CLI parser correctness]
\label{thm:cli-parser}
For the CLI parser above:
\begin{enumerate}
  \item \textbf{Registration}: after \lstinline[style=python]{@cli.command("greet")},
    the command \texttt{"greet"} is in \lstinline[style=python]{cli._commands}
    with the correct signature.
  \item \textbf{Parsing}: \lstinline[style=python]{dispatch(["greet", "World", "--greeting", "Hi"])}
    correctly binds \texttt{name="World"}, \texttt{greeting="Hi"},
    \texttt{excited=False} (default).
  \item \textbf{Zone correctness}: \texttt{name} can only be passed
    positionally (it is positional-only); \texttt{excited} can only be
    passed by keyword (it is keyword-only).
\end{enumerate}
\end{theorem}
\begin{proof}
(1) The \lstinline[style=python]{@cli.command("greet")} decorator stores
$(f, \mathsf{signature}(f))$ in \lstinline[style=python]{_commands["greet"]}.

(2) The \lstinline[style=python]{dispatch} method parses
$[\mathtt{"greet"}, \mathtt{"World"}, \mathtt{"--greeting"}, \mathtt{"Hi"}]$
into $\mathit{args} = [\mathtt{"World"}]$ and
$\mathit{kwargs} = \{\mathtt{"greeting"}: \mathtt{"Hi"}\}$.
Then \lstinline[style=python]{sig.bind(*args, **kwargs)} performs the
standard binding algorithm:
\texttt{"World"} binds to \texttt{name} (position~0, positional-only),
\texttt{"Hi"} binds to \texttt{greeting} (keyword), and
\lstinline[style=python]{apply_defaults()} fills \texttt{excited=False}.

(3) The parameter zones follow from the signature:
\texttt{name} precedes \texttt{/}, so it is positional-only;
\texttt{excited} follows \texttt{*}, so it is keyword-only.
If a user passes \texttt{--name World}, the binding algorithm raises
\lstinline[style=python]{TypeError} because \texttt{name} is positional-only.
\end{proof}

\begin{lstlisting}[style=proof]
-- Verify: CLI dispatch for ["greet", "World", "--greeting", "Hi"]
have h_cmd: "greet" in cli._commands
    by decorator registration
have h_parse: args = ["World"], kwargs = {"greeting": "Hi"}
    by dispatch parsing loop; "--greeting" starts with "--"
have h_bind: name = "World" (pos 0), greeting = "Hi" (kw),
    excited = False (default)
    by sig.bind; name is POS_ONLY at pos 0;
       greeting is POS_OR_KW, matched by keyword;
       excited is KW_ONLY, not provided, default = False
have h_result: greet("World", greeting="Hi", excited=False)
    = "Hi, World!"
    by unfold greet; excited is False, no "!!" suffix
conclude: dispatch returns "Hi, World!"
    by h_cmd, h_parse, h_bind, h_result
\end{lstlisting}


% ──────────────────────────────────────────────────────────────────
\chapter{Context Managers}
\label{ch:context-managers}

\section{The Resource Bundle}

\begin{lstlisting}[style=python]
with open("data.csv") as f:
    data = f.read()
# f is guaranteed closed here, even if f.read() raised
\end{lstlisting}

\begin{definition}[Resource bundle]
A context manager defines a bundle:
\begin{align}
  \mathsf{\_\_enter\_\_} &: \mathsf{unit} \to_{\IO} R \quad \text{(acquire resource)} \\
  \mathsf{\_\_exit\_\_} &: R \times \mathsf{Option}(\mathsf{Exc}) \to_{\IO} \mathsf{bool} \quad \text{(release resource)}
\end{align}
The guarantee: $\mathsf{\_\_exit\_\_}$ is called on \emph{every} execution path
through the \texttt{with} block, including exception paths.
\end{definition}

\begin{definition}[Context manager typing rule]
\[
  \frac{
    \Gamma \vdash \mathit{cm} : \mathsf{ContextManager}(R) \qquad
    \Gamma, r : R \vdash_\varepsilon \mathit{body}(r) : A
  }{
    \Gamma \vdash_{\varepsilon \vee \IO} \mathsf{with}\;\mathit{cm}\;\mathsf{as}\;r : \mathit{body}(r) : A
  }
\]
The effect of the \texttt{with} block includes $\IO$ (for resource acquisition/release)
joined with whatever effects the body has.
\end{definition}

\begin{theorem}[Resource safety]
If $\mathit{cm}$ satisfies the context manager protocol, then in
\lstinline[style=python]{with cm as r: body(r)}, the resource $r$ is released
after the block completes, regardless of exceptions.
\end{theorem}
\begin{proof}
Python's runtime guarantees that \lstinline[style=python]{__exit__} is called
in a \lstinline[style=python]{finally} block. The proof in \synhopy{} is a
\emph{global section} of the exception bundle restricted to the \texttt{with}
block: for every exception path, the section (resource release) is well-defined.
\end{proof}

\section{Nested Context Managers}

\begin{lstlisting}[style=python]
with open("a.txt") as f, open("b.txt") as g:
    data = f.read() + g.read()
\end{lstlisting}

Nested \texttt{with} statements compose linearly:
\[
  \frac{
    \Gamma \vdash_{e_1} \mathsf{with}\;cm_1\;\mathsf{as}\;r_1: E_1 \qquad
    \Gamma, r_1 \vdash_{e_2} \mathsf{with}\;cm_2\;\mathsf{as}\;r_2: E_2
  }{
    \Gamma \vdash_{e_1 \vee e_2 \vee \IO}
      \mathsf{with}\;cm_1, cm_2\;\mathsf{as}\;r_1, r_2: E
  }
\]

If $cm_2.\mathsf{\_\_enter\_\_}$ raises, $cm_1.\mathsf{\_\_exit\_\_}$ is still called.
This is the \emph{linear resource composition}: resources are released in reverse
acquisition order.

\section{\texttt{contextlib} Generators}

Python's \lstinline[style=python]{@contextmanager} decorator turns a generator
into a context manager:
\begin{lstlisting}[style=python]
from contextlib import contextmanager

@contextmanager
def timer():
    start = time.time()
    yield start
    print(f"Elapsed: {time.time() - start:.2f}s")
\end{lstlisting}

\begin{definition}[Generator context manager]
The \lstinline[style=python]{@contextmanager} decorator splits a generator $g$
at its \texttt{yield} point:
\begin{align}
  g.\mathsf{\_\_enter\_\_} &= \text{run $g$ up to \texttt{yield}; return yielded value} \\
  g.\mathsf{\_\_exit\_\_}  &= \text{resume $g$ after \texttt{yield}; run to completion}
\end{align}
The fiber of the generator at the yield point determines the type of the
resource provided to the body.
\end{definition}

\begin{remark}[\fstar{} comparison]
\fstar{} has no mechanism for defining custom resource management protocols.
The closest analog is a monadic computation with explicit acquire/release
actions, but this does not guarantee that release is called on all paths
(including exceptions).  \synhopy{}'s typing rule makes this guarantee
structural.
\end{remark}

\section{Async Context Managers}

Python 3.5 introduced \emph{async context managers}, combining the resource
management pattern with cooperative concurrency:

\begin{lstlisting}[style=python]
import aiofiles

async with aiofiles.open("data.csv") as f:
    data = await f.read()
# f is guaranteed closed here, even on exceptions
\end{lstlisting}

\begin{definition}[Async context manager]
An async context manager extends the resource bundle with the
$\mathsf{Suspend}$ effect:
\begin{align}
  \mathsf{\_\_aenter\_\_} &: \mathsf{unit} \to_{\IO \vee \mathsf{Suspend}} R \\
  \mathsf{\_\_aexit\_\_} &: R \times \mathsf{Option}(\mathsf{Exc})
    \to_{\IO \vee \mathsf{Suspend}} \mathsf{bool}
\end{align}
Both acquisition and release may themselves suspend (e.g., awaiting a network
connection pool).
\end{definition}

\begin{definition}[Async with typing rule]
\[
  \frac{
    \Gamma \vdash \mathit{acm} : \mathsf{AsyncContextManager}(R) \qquad
    \Gamma, r : R \vdash_\varepsilon \mathit{body}(r) : A
  }{
    \Gamma \vdash_{\varepsilon \vee \IO \vee \mathsf{Suspend}}
      \mathsf{async\;with}\;\mathit{acm}\;\mathsf{as}\;r : \mathit{body}(r) : A
  }
\]
The critical invariant is identical to the synchronous case:
$\mathsf{\_\_aexit\_\_}$ is called on every path, including exception paths
and cancellation paths.
\end{definition}

\begin{remark}[\fstar{} comparison]
\fstar{} has no async/await or cooperative concurrency primitives, so the
notion of an asynchronous resource manager has no analog.  Even in Haskell's
\texttt{bracket}, the acquire and release actions are assumed to be
non-interruptible; Python's async context managers make no such assumption,
requiring the additional $\mathsf{Suspend}$ annotation in \synhopy{}.
\end{remark}

\section{ExitStack and Dynamic Composition}

When the number of context managers is not known statically, Python provides
\lstinline[style=python]{contextlib.ExitStack}:

\begin{lstlisting}[style=python]
from contextlib import ExitStack

def open_all(paths: list[str]) -> ExitStack:
    stack = ExitStack()
    for p in paths:
        stack.enter_context(open(p))
    return stack

with open_all(["a.txt", "b.txt", "c.txt"]) as stack:
    ...  # all files open; all closed on exit
\end{lstlisting}

\begin{definition}[ExitStack type]
An \lstinline[style=python]{ExitStack} is a dynamically-sized product of
context managers:
\[
  \mathsf{ExitStack} : \Sigma(n : \Nat).\;
    \Pi(i : \mathsf{Fin}(n)).\; \mathsf{ContextManager}(R_i)
\]
where resources are released in \emph{reverse} order ($R_{n-1}, \ldots, R_0$).
\end{definition}

\begin{theorem}[ExitStack resource safety]
For an \lstinline[style=python]{ExitStack} $S$ containing $n$ context managers,
if any $\mathsf{\_\_exit\_\_}_i$ raises an exception, all remaining
$\mathsf{\_\_exit\_\_}_j$ for $j < i$ are still called.  Formally:
\[
  \forall\, i \in \{0, \ldots, n{-}1\}.\;
    \mathsf{called}(\mathsf{\_\_exit\_\_}_i) = \mathsf{true}
\]
regardless of exceptions in the body or in other exit handlers.
\end{theorem}
\begin{proof}
By induction on $n$.  The base case $n = 0$ is trivial.  For the inductive
step, \lstinline[style=python]{ExitStack} wraps each exit in its own
\texttt{try/finally}, so the failure of $\mathsf{\_\_exit\_\_}_i$ does not
prevent $\mathsf{\_\_exit\_\_}_{i-1}$ from executing.  The induction
corresponds to peeling off one layer of the stack at a time.
\end{proof}

\section{Exception Suppression in \texttt{\_\_exit\_\_}}

The boolean return value of \lstinline[style=python]{__exit__} has
subtle semantics:

\begin{lstlisting}[style=python]
class SuppressErrors:
    def __enter__(self):
        return self
    def __exit__(self, exc_type, exc_val, exc_tb):
        if exc_type is ValueError:
            return True   # suppress the exception
        return False      # re-raise
\end{lstlisting}

\begin{definition}[Exit return semantics]
The return type of $\mathsf{\_\_exit\_\_}$ refines the exception flow:
\[
  \mathsf{\_\_exit\_\_}(r, e) = \begin{cases}
    \mathsf{True} & \Rightarrow \text{exception $e$ is suppressed;
      execution continues after \texttt{with}} \\
    \mathsf{False} & \Rightarrow \text{exception $e$ is re-raised}
  \end{cases}
\]
In \synhopy{}, this is typed as a dependent elimination on the
exception bundle: the section over the \texttt{with} block is
\emph{partial} (covers only the suppressed exception types) or
\emph{total} (covers all types, meaning exceptions always propagate).
\end{definition}

\begin{proposition}[Suppression soundness]
If $\mathsf{\_\_exit\_\_}$ returns $\mathsf{True}$ for exception type $E$,
then the continuation after the \texttt{with} block is well-typed only if
the context $\Gamma$ does not depend on the absence of $E$.  Formally:
\[
  \frac{
    \Gamma \vdash \mathit{cm}.\mathsf{\_\_exit\_\_}(\_, E) = \mathsf{True} \qquad
    \Gamma \vdash_\varepsilon \mathit{rest} : B
  }{
    \Gamma \vdash_{\varepsilon} \mathsf{with}\;\mathit{cm}\;\mathsf{as}\;r\!:
      \mathit{body};\; \mathit{rest} : B
  }
\]
The proof obligation is that $\mathit{rest}$ is well-typed even in a context
where $E$ may have occurred and been suppressed.
\end{proposition}

\section{Verifying Resource Leak Freedom}

\begin{definition}[Resource leak freedom]
A program is \emph{resource-leak-free} if every resource acquisition
$\mathsf{acquire}(r)$ is matched by a corresponding
$\mathsf{release}(r)$ on every execution path.  In \synhopy{}, this
is expressed as a global section of the resource bundle over the
program's control-flow graph:
\[
  \forall\, r \in \mathsf{Resources}.\;
    \forall\, \pi \in \mathsf{Paths}(\mathsf{CFG}).\;
      \mathsf{acquire}(r) \in \pi \Rightarrow \mathsf{release}(r) \in \pi
\]
\end{definition}

\begin{theorem}[Context managers guarantee leak freedom]
If all resource acquisitions are performed via \texttt{with} statements
(or \lstinline[style=python]{ExitStack}), and no resource references
escape the \texttt{with} block, then the program is resource-leak-free.
\end{theorem}
\begin{proof}
Each \texttt{with} block is a \emph{closed section} of the resource bundle:
the resource $r$ is acquired at $\mathsf{\_\_enter\_\_}$ and released at
$\mathsf{\_\_exit\_\_}$.  Since $r$ does not escape (no aliases exist outside
the block), and $\mathsf{\_\_exit\_\_}$ is called on every path, the global
section condition holds.  The ``no escape'' condition is verified by
\synhopy{}'s flow-sensitive alias analysis: $r$ must not appear in any
assignment target whose scope extends beyond the \texttt{with} block.
\end{proof}

\section{Context Managers as a Monad}

The context manager pattern is closely related to Haskell's
\lstinline[style=python]{bracket} combinator:

\begin{lstlisting}[style=python]
# Haskell's bracket:
#   bracket :: IO a -> (a -> IO b) -> (a -> IO c) -> IO c
#   bracket acquire release body = ...
# Python equivalent:
@contextmanager
def bracket(acquire, release):
    r = acquire()
    try:
        yield r
    finally:
        release(r)
\end{lstlisting}

\begin{definition}[Resource monad]
Define the \emph{resource monad} $\mathsf{Res}$ as:
\begin{align}
  \mathsf{Res}(A) &\defeq \exists (R : \Type).\;
    (\mathsf{unit} \to_\IO R) \times (R \to_\IO \mathsf{unit})
    \times (R \to_\IO A) \\
  \mathsf{return}(a) &\defeq (\mathsf{unit}, \lambda\_.\;(), \lambda\_.\;(),
    \lambda\_.\;a) \\
  \mathsf{bind}(m, f) &\defeq \text{compose acquisitions, sequence bodies,
    release in reverse}
\end{align}
The Kleisli composition of $\mathsf{Res}$ corresponds to nesting
\texttt{with} blocks.
\end{definition}

\begin{proposition}[Monad laws for $\mathsf{Res}$]
$\mathsf{Res}$ satisfies the monad laws:
\begin{enumerate}
  \item Left identity: $\mathsf{bind}(\mathsf{return}(a), f) = f(a)$
  \item Right identity: $\mathsf{bind}(m, \mathsf{return}) = m$
  \item Associativity: $\mathsf{bind}(\mathsf{bind}(m, f), g)
    = \mathsf{bind}(m, \lambda x.\;\mathsf{bind}(f(x), g))$
\end{enumerate}
The key insight is that release actions compose in reverse order,
matching the stack discipline of nested \texttt{with} blocks.
\end{proposition}

\section{Worked Example: Database Transaction Manager}

Consider a database transaction context manager that must either commit
on success or rollback on failure:

\begin{lstlisting}[style=python]
from contextlib import contextmanager

@contextmanager
def transaction(conn):
    """Context manager ensuring commit/rollback semantics."""
    tx = conn.begin()
    try:
        yield tx
        tx.commit()
    except Exception:
        tx.rollback()
        raise
\end{lstlisting}

\begin{lstlisting}[style=proof]
-- Verify: transaction(conn) satisfies commit/rollback
-- invariant: exactly one of commit or rollback is called

have h_enter: conn.begin() returns tx : Transaction
    by unfold transaction, up to yield

-- Case 1: body succeeds (no exception)
have h_success: body(tx) terminates normally
    implies tx.commit() is called
    by unfold transaction, normal path after yield

-- Case 2: body raises exception E
have h_failure: body(tx) raises E
    implies tx.rollback() is called
    by unfold transaction, except branch

-- Exactly one of commit/rollback
have h_exclusive:
    (tx.commit() called) XOR (tx.rollback() called)
    by cases h_success, h_failure (exhaustive)

-- Atomicity: database state is consistent
have h_atomic:
    forall obs : Observer.
      obs sees either all of body's writes or none
    by h_exclusive then apply db_transaction_spec

conclude: transaction(conn) is a correct transaction manager
    by structural with h_enter, h_exclusive, h_atomic
\end{lstlisting}

\begin{theorem}[Transaction safety]
The \lstinline[style=python]{transaction} context manager satisfies:
\begin{enumerate}
  \item \emph{Exactly-once}: exactly one of \texttt{commit()} or
    \texttt{rollback()} is called.
  \item \emph{Atomicity}: the database transitions from one consistent state
    to another, with no partial writes visible.
  \item \emph{Resource release}: the transaction object is finalized on every
    execution path.
\end{enumerate}
\end{theorem}
\begin{proof}
Property (1) follows from the control flow: the \texttt{try} block calls
\texttt{commit()} only if no exception propagates from the body, and the
\texttt{except} block calls \texttt{rollback()} only if an exception does
propagate.  Since these are complementary and exhaustive, exactly one is
called.  Property (2) follows from the database's transaction specification
applied to (1).  Property (3) follows from the general resource safety
theorem (Theorem~15.1) applied to the \texttt{@contextmanager} decorator.
\end{proof}


% ──────────────────────────────────────────────────────────────────
\chapter{The ``Everything Is a Dict'' Principle}
\label{ch:dict}

\section{Object Attribute Access}

In Python, every object is fundamentally a dictionary:

\begin{lstlisting}[style=python]
class Point:
    def __init__(self, x, y):
        self.x = x  # equivalent to self.__dict__['x'] = x
        self.y = y

p = Point(1, 2)
p.__dict__  # {'x': 1, 'y': 2}
p.z = 3     # just adds to the dict
\end{lstlisting}

\begin{definition}[Object as dependent dictionary]
A Python object $o$ has type:
\[
  o : \Sigma(\mathit{cls} : \PyClass).\; \Sigma(\mathit{keys} : \mathsf{Set}(\mathsf{String})).\;
    \Pi(k : \mathit{keys}).\; \PyObj
\]
where:
\begin{itemize}
  \item $\mathit{cls}$ is the object's class (its \texttt{\_\_class\_\_} attribute)
  \item $\mathit{keys}$ is the set of instance attribute names
  \item The dependent function maps each key to its value
\end{itemize}
\end{definition}

\section{Attribute Resolution as Fibration}

Attribute access \lstinline[style=python]{obj.x} triggers the MRO-based
resolution protocol:

\begin{enumerate}
  \item Check \lstinline[style=python]{type(obj).__mro__} for data descriptors with \texttt{x}
  \item Check \lstinline[style=python]{obj.__dict__} for \texttt{x}
  \item Check \lstinline[style=python]{type(obj).__mro__} for non-data descriptors with \texttt{x}
  \item Call \lstinline[style=python]{__getattr__} if defined
  \item Raise \lstinline[style=python]{AttributeError}
\end{enumerate}

\begin{definition}[Attribute resolution fibration]
The attribute resolution for name $x$ on object $o$ is a fibration:
\[
  \mathsf{resolve}(o, x) : \prod_{C \in \MRO(\mathsf{cls}(o))} \mathsf{Option}(\PyObj)
\]
The resolution is a section of this fibration: it selects the first defined
value in the MRO order.
\end{definition}

\section{Dynamic Attribute Addition}

Adding an attribute at runtime changes the object's type:
\begin{lstlisting}[style=python]
p = Point(1, 2)        # type: {x: int, y: int}
p.color = "red"         # type: {x: int, y: int, color: str}
del p.x                 # type: {y: int, color: str}
\end{lstlisting}

\begin{definition}[Dynamic attribute modification]
\begin{align}
  \mathsf{setattr}(o, k, v) &: \Sigma(\mathit{keys} \cup \{k\}).\;
    \Pi(k').\; \begin{cases} v & k' = k \\ o.k' & k' \neq k \end{cases} \\
  \mathsf{delattr}(o, k) &: \Sigma(\mathit{keys} \setminus \{k\}).\;
    \Pi(k').\; o.k' \quad\text{(for $k' \in \mathit{keys} \setminus \{k\}$)}
\end{align}
\end{definition}

This means that in Python, the \emph{type} of an object can change at every
statement.  \synhopy{} tracks this flow-sensitively: the type of $o$ after
\texttt{setattr} is different from before.

\section{Module Namespaces as Dicts}

Python modules are also dicts:
\begin{lstlisting}[style=python]
import math
math.__dict__['pi']    # 3.141592653589793
math.__dict__['sin']   # <built-in function sin>
\end{lstlisting}

\begin{definition}[Module type]
A module $M$ has type:
\[
  M : \Sigma(\mathit{names} : \mathsf{Set}(\mathsf{String})).\;
    \Pi(n : \mathit{names}).\; \PyObj
\]
This is the same dependent-dict structure as objects.
Import resolution is attribute resolution on the module dict.
\end{definition}

\begin{remark}[\fstar{} comparison]
In \fstar{}, modules have \emph{static} interfaces that cannot change at runtime.
Python modules are mutable dictionaries: you can add, remove, or replace any
attribute at any time.  \synhopy{}'s dependent-dict model captures this mutability
while still allowing static reasoning about module contents (when the module is
not monkey-patched).
\end{remark}

\section{\texttt{\_\_slots\_\_} as Closed Records}

By default, Python objects have an open dictionary---attributes can be
added freely.  The \lstinline[style=python]{__slots__} mechanism restricts
this to a fixed set of keys:

\begin{lstlisting}[style=python]
class Vec2:
    __slots__ = ('x', 'y')
    def __init__(self, x, y):
        self.x = x
        self.y = y

v = Vec2(1, 2)
v.z = 3  # raises AttributeError
\end{lstlisting}

\begin{definition}[Slotted type as closed record]
A class $C$ with $\mathsf{\_\_slots\_\_} = \{k_1, \ldots, k_n\}$ has
object type:
\[
  C_{\mathrm{slotted}} : \Pi(k : \{k_1, \ldots, k_n\}).\; T_k
\]
where $T_k$ is the type of slot $k$.  Unlike the general dependent-dict
type, the key set is \emph{fixed} at class definition time:
\[
  \mathsf{setattr}(o, k, v) : \begin{cases}
    \text{well-typed} & \text{if } k \in \mathsf{\_\_slots\_\_} \\
    \bot \;\text{(AttributeError)} & \text{if } k \notin \mathsf{\_\_slots\_\_}
  \end{cases}
\]
\end{definition}

\begin{proposition}[Slots as subtype]
$C_{\mathrm{slotted}} \leq C_{\mathrm{open}}$ in the subsumption order:
every slotted object is also a valid open object (with the same keys),
but not vice versa.  The slotted type carries the additional invariant that
$\mathsf{keys}(o)$ is constant throughout the object's lifetime.
\end{proposition}

\section{The Full Attribute Lookup Protocol}

Python's attribute access involves two distinct dunder methods:
\lstinline[style=python]{__getattribute__} (called on every access) and
\lstinline[style=python]{__getattr__} (called only when normal lookup fails).

\begin{definition}[\texttt{\_\_getattribute\_\_} and \texttt{\_\_getattr\_\_}]
The full attribute resolution for $o.x$ is:
\begin{enumerate}
  \item Call $\mathsf{type}(o).\mathsf{\_\_getattribute\_\_}(o, x)$, which
    by default implements the descriptor protocol (data descriptors,
    instance dict, non-data descriptors).
  \item If $\mathsf{\_\_getattribute\_\_}$ raises
    \lstinline[style=python]{AttributeError}, and $\mathsf{type}(o)$ defines
    $\mathsf{\_\_getattr\_\_}$, call $\mathsf{type}(o).\mathsf{\_\_getattr\_\_}(o, x)$.
  \item If $\mathsf{\_\_getattr\_\_}$ also raises or is not defined, the
    \lstinline[style=python]{AttributeError} propagates.
\end{enumerate}
In \synhopy{}, this is a two-level fibration:
\[
  \mathsf{resolve}(o, x) \defeq
    \mathsf{try}\;\mathsf{\_\_getattribute\_\_}(o, x)\;
    \mathsf{catch}\;\_\, \Rightarrow \mathsf{\_\_getattr\_\_}(o, x)
\]
\end{definition}

\begin{lstlisting}[style=python]
class Fallback:
    def __init__(self):
        self.known = {"a": 1, "b": 2}
    def __getattr__(self, name):
        # Only called if normal lookup fails
        return self.known.get(name, f"<missing: {name}>")

f = Fallback()
f.a      # 1 (from __getattr__ since 'a' not in __dict__)
f.x = 42
f.x      # 42 (from __dict__, __getattr__ NOT called)
\end{lstlisting}

\begin{remark}[\fstar{} comparison]
\fstar{} record field access is a simple projection $r.f$ with no
interception mechanism.  Python's two-level attribute protocol allows
classes to dynamically intercept, redirect, or synthesize any attribute
access---a flexibility that \synhopy{} types via the fallback fibration
structure.
\end{remark}

\section{Proxy Objects and Delegation Chains}

A common Python pattern uses \lstinline[style=python]{__getattr__} to
delegate attribute access to a wrapped object, creating a \emph{proxy}:

\begin{lstlisting}[style=python]
class LoggingProxy:
    def __init__(self, target):
        object.__setattr__(self, '_target', target)
    def __getattr__(self, name):
        print(f"Accessing {name}")
        return getattr(self._target, name)
    def __setattr__(self, name, value):
        print(f"Setting {name} = {value}")
        setattr(self._target, name, value)
\end{lstlisting}

\begin{definition}[Proxy type]
A proxy object $p$ wrapping target $t$ has type:
\[
  p : \mathsf{Proxy}(T) \defeq \Sigma(\mathit{target} : T).\;
    \Pi(k : \mathsf{keys}(T)).\; \mathsf{type}_T(k)
\]
The proxy's attribute type is \emph{inherited} from the target: for any
attribute $k$, $\mathsf{type}(p.k) = \mathsf{type}(t.k)$.  The proxy may
add effects (here, $\IO$ for logging) to each access:
\[
  p.k : \mathsf{type}_T(k) \quad \text{with effect } \varepsilon_T(k) \vee \IO
\]
\end{definition}

\begin{definition}[Delegation chain]
A chain of proxies $p_1 \to p_2 \to \cdots \to p_n \to t$ forms a path in
the delegation fibration.  Attribute resolution traverses the chain:
\[
  \mathsf{resolve}(p_1, x) = \mathsf{resolve}(p_2, x) = \cdots =
    \mathsf{resolve}(t, x)
\]
The chain terminates if and only if the final target $t$ does not itself
delegate.  \synhopy{} detects infinite delegation loops as a failure to
find a global section of the resolution fibration.
\end{definition}

\section{Namespace Packages and \texttt{sys.modules}}

Python's module system is itself a dictionary---\lstinline[style=python]{sys.modules}
maps dotted names to module objects:

\begin{lstlisting}[style=python]
import sys
sys.modules['math']          # <module 'math' ...>
sys.modules['my.pkg'] = mod  # inject a module at runtime
\end{lstlisting}

\begin{definition}[\texttt{sys.modules} as global dict]
The module registry is a global dependent dictionary:
\[
  \mathsf{sys.modules} : \Sigma(\mathit{names} : \mathsf{Set}(\mathsf{DottedName})).\;
    \Pi(n : \mathit{names}).\; \mathsf{Module}
\]
An \lstinline[style=python]{import foo} statement is operationally:
\begin{enumerate}
  \item Look up $\mathsf{sys.modules}[\texttt{"foo"}]$; if found, return it.
  \item Otherwise, execute the module's code, creating a new module object.
  \item Store the result in $\mathsf{sys.modules}[\texttt{"foo"}]$ and return it.
\end{enumerate}
Because \lstinline[style=python]{sys.modules} is mutable, modules can be
hot-swapped, mocked, or removed at runtime.
\end{definition}

\begin{proposition}[Import idempotence]
Under the standard import protocol (no monkey-patching of
\lstinline[style=python]{sys.modules}), repeated imports of the same name
return the same object:
\[
  \mathsf{import}(n) = \mathsf{import}(n) \qquad
  \text{(referential transparency modulo \texttt{sys.modules} mutations)}
\]
This is a \emph{caching} property: the first import computes and stores;
subsequent imports retrieve.
\end{proposition}

\section{The \texttt{vars()} Function}

The built-in \lstinline[style=python]{vars()} function provides direct
access to an object's \lstinline[style=python]{__dict__}:

\begin{lstlisting}[style=python]
class Config:
    def __init__(self):
        self.debug = True
        self.port = 8080

c = Config()
vars(c)  # {'debug': True, 'port': 8080}
vars(c)['port'] = 9090  # mutates c.port!
c.port   # 9090
\end{lstlisting}

\begin{definition}[\texttt{vars()} type]
\[
  \mathsf{vars}(o) : \mathsf{Dict}[\mathsf{String}, \PyObj]
\]
with the identity $\mathsf{vars}(o) \equiv o.\mathsf{\_\_dict\_\_}$ (they are
the \emph{same} object, not a copy).  Mutations to $\mathsf{vars}(o)$ are
reflected in $o$'s attributes and vice versa.
\end{definition}

\begin{proposition}[Aliasing invariant]
For any object $o$ with a $\mathsf{\_\_dict\_\_}$:
\[
  \forall k.\; o.k = \mathsf{vars}(o)[k] \quad
  \text{(when no descriptors intercede)}
\]
This identity means that attribute access and dictionary access are two
\emph{views} of the same underlying data---a fundamental consequence of
Python's ``everything is a dict'' design.
\end{proposition}

\section{Worked Example: Nested Dict Configuration with Dotted Access}

Consider a configuration system that stores settings as nested dicts
but allows dotted attribute access:

\begin{lstlisting}[style=python]
class DotDict(dict):
    """Dict subclass allowing attribute-style access."""
    def __getattr__(self, key):
        try:
            val = self[key]
        except KeyError:
            raise AttributeError(key)
        return DotDict(val) if isinstance(val, dict) else val

    def __setattr__(self, key, value):
        self[key] = value

config = DotDict({
    "database": {"host": "localhost", "port": 5432},
    "debug": True
})
config.database.host  # "localhost"
config.debug          # True
\end{lstlisting}

\begin{definition}[DotDict type]
A \lstinline[style=python]{DotDict} with schema $S$ has type:
\[
  \mathsf{DotDict}(S) \defeq \Pi(k : \mathsf{keys}(S)).\;
    \begin{cases}
      \mathsf{DotDict}(S_k) & \text{if } S_k \text{ is a record schema} \\
      T_k & \text{otherwise}
    \end{cases}
\]
This is a \emph{nested} dependent product: at each level, the return type
depends on whether the value is itself a dictionary.
\end{definition}

\begin{lstlisting}[style=proof]
-- Verify: DotDict attribute access is consistent with dict access

have h_get: forall k in keys(config).
    config.k = config[k]
    by unfold DotDict.__getattr__, dict.__getitem__

have h_nested: forall k in keys(config).
    isinstance(config[k], dict)
    implies config.k : DotDict
    by unfold DotDict.__getattr__, isinstance branch

have h_deep: config.database.host = config["database"]["host"]
    by rw [h_get] at depth 2 then rw [h_nested]

have h_set: config.debug = False
    implies config["debug"] = False
    by unfold DotDict.__setattr__, dict.__setitem__

-- Schema preservation: setting a value preserves type
have h_schema: forall k, v.
    type(v) = type(config[k])
    implies DotDict.__setattr__(config, k, v)
    preserves schema S
    by unfold __setattr__ then apply dict_update_type_pres

conclude: DotDict faithfully implements dotted access
    over nested dict with schema preservation
    by structural with h_get, h_nested, h_deep, h_set, h_schema
\end{lstlisting}

\begin{theorem}[DotDict correctness]
For a \lstinline[style=python]{DotDict} $d$ with schema $S$:
\begin{enumerate}
  \item \emph{Access equivalence}: $d.k_1.k_2.\cdots.k_n = d[k_1][k_2]\cdots[k_n]$
    for any valid key path.
  \item \emph{Recursive wrapping}: if $d[k]$ is a \lstinline[style=python]{dict},
    then $d.k$ is a \lstinline[style=python]{DotDict} (auto-wrapping).
  \item \emph{Error faithfulness}: accessing a missing key via attribute access
    raises \lstinline[style=python]{AttributeError}, not
    \lstinline[style=python]{KeyError}.
\end{enumerate}
\end{theorem}
\begin{proof}
Property (1) follows by induction on the key path length, using the
definition of \lstinline[style=python]{__getattr__} at each level.
Property (2) is the \lstinline[style=python]{isinstance} branch in
\lstinline[style=python]{__getattr__}.  Property (3) is the
\lstinline[style=python]{try/except KeyError} translation in
\lstinline[style=python]{__getattr__}.
\end{proof}


% ──────────────────────────────────────────────────────────────────
\chapter{Async/Await and Concurrency}
\label{ch:async}

\section{The Concurrency Problem}

Python's \texttt{asyncio} framework provides cooperative multitasking via
\texttt{async}/\texttt{await}:

\begin{lstlisting}[style=python]
import asyncio

async def fetch_user(user_id: int) -> dict:
    await asyncio.sleep(0.1)  # simulate network delay
    return {"id": user_id, "name": f"User {user_id}"}

async def fetch_all_users(ids: list[int]) -> list[dict]:
    tasks = [fetch_user(uid) for uid in ids]
    return await asyncio.gather(*tasks)
\end{lstlisting}

\section{Coroutines as Paths with Suspension Points}

\begin{definition}[Coroutine type]
An async function has type:
\[
  \mathsf{async}\;(x : A) \to_\varepsilon B \;\defeq\; (x : A) \to_{\varepsilon \vee \mathsf{Suspend}} B
\]
where $\mathsf{Suspend}$ is an additional effect that indicates the computation
may suspend and resume at \texttt{await} points.
\end{definition}

A coroutine execution is a path $\gamma$ in state space that has
\emph{suspension points}---positions where the path pauses and another path
(from a different coroutine) can interleave.  Between suspension points the
path is atomic (no interleaving), which is the defining property of
\emph{cooperative} concurrency.

\begin{definition}[Await as path composition]
\[
  \frac{
    \Gamma \vdash_{\varepsilon_1} e : \mathsf{Coroutine}(B) \qquad
    \varepsilon_1 \leq \varepsilon_2
  }{
    \Gamma \vdash_{\varepsilon_2 \vee \mathsf{Suspend}} \mathsf{await}\;e : B
  }
\]
\end{definition}

\section{Proving Properties of Gathered Tasks}

\begin{lstlisting}[style=proof]
-- Proof: fetch_all_users returns one result per input ID
have h_gather: asyncio.gather(*tasks) returns
    [await t for t in tasks]  (order preserved)
    by apply asyncio_gather_spec
have h_each: forall uid. fetch_user(uid) returns
    {"id": uid, "name": f"User {uid}"}
    by unfold fetch_user
have h_len: len(results) = len(ids)
    by rw [h_gather] then omega
have h_ids: forall i. results[i]["id"] = ids[i]
    by rw [h_gather, h_each]
conclude: returns list of user dicts matching input IDs
    by structural with h_len, h_ids
\end{lstlisting}

\begin{remark}[\fstar{} comparison]
\fstar{} has no built-in support for async/await or cooperative concurrency.
Its effect system handles \emph{sequential} effects (STATE, EXN, IO) but not
\emph{concurrent} effects.  To model \texttt{asyncio.gather}, you would need
a bespoke concurrency monad, an interleaving semantics, and manual proofs
of non-interference---none of which \fstar{} provides out of the box.
\synhopy{} models suspension as an effect and \texttt{gather} as parallel path
composition, handling concurrency within the existing framework.
\end{remark}

\section{Race Condition Detection via Spectral Sequences}

When two async functions access shared mutable state, the spectral sequence
detects potential races:

\begin{lstlisting}[style=python]
shared_counter = 0

async def increment():
    global shared_counter
    val = shared_counter      # read
    await asyncio.sleep(0)    # SUSPEND -- another task runs
    shared_counter = val + 1  # write (may be stale!)

async def double_increment():
    await asyncio.gather(increment(), increment())
    # shared_counter might be 1, not 2!
\end{lstlisting}

Analysis via the spectral sequence:
\begin{itemize}
  \item $E_1$: each \texttt{increment()} has effect $\Mutates(\mathtt{shared\_counter})$.
  \item Two concurrent instances produce $\Mutates \times \Mutates$ on the
    \emph{same} variable.
  \item $d_1$: the differential between the two $\Mutates$ effects is
    \emph{non-zero} because both read-then-write the same location with an
    intervening suspension point.
  \item Consequence: the spec ``\texttt{shared\_counter} increases by 2'' is
    \textbf{unprovable}, and \synhopy{} reports a $\mathsf{RACE\_CONDITION}$ warning.
\end{itemize}

\section{\texttt{asyncio.gather} and Concurrent Composition}

The \lstinline[style=python]{asyncio.gather} function is the primary
mechanism for concurrent task composition:

\begin{lstlisting}[style=python]
async def fetch_prices(symbols: list[str]) -> list[float]:
    tasks = [fetch_price(s) for s in symbols]
    return await asyncio.gather(*tasks)
\end{lstlisting}

\begin{definition}[\texttt{gather} type]
\[
  \mathsf{gather} : \Pi(n : \Nat).\;
    \left(\Pi(i : \mathsf{Fin}(n)).\; \mathsf{Coroutine}(A_i)\right)
    \to_{\mathsf{Suspend}} \Pi(i : \mathsf{Fin}(n)).\; A_i
\]
That is, \lstinline[style=python]{gather} takes $n$ coroutines with
potentially different result types and returns a tuple of results,
preserving the input order.
\end{definition}

\begin{theorem}[Gather non-interference]
If coroutines $c_1, \ldots, c_n$ have pairwise disjoint mutation sets:
\[
  \forall\, i \neq j.\; \Mutates(c_i) \cap \Mutates(c_j) = \emptyset
\]
then $\mathsf{gather}(c_1, \ldots, c_n)$ is equivalent to any sequential
ordering of $c_1, \ldots, c_n$.  In particular, the result is
deterministic.
\end{theorem}
\begin{proof}
With disjoint mutation sets, the interleaving at suspension points
cannot affect the outcome of any individual coroutine.  Each coroutine's
execution path is independent of the scheduling order, so the gathered
results are identical to those of any sequential composition.
\end{proof}

\begin{proposition}[Gather error propagation]
If any coroutine $c_i$ raises an exception $E$, then
$\mathsf{gather}(c_1, \ldots, c_n)$ raises $E$ (by default).  With
\lstinline[style=python]{return_exceptions=True}, exceptions are returned
as values:
\[
  \mathsf{gather}_{\mathit{re}}(c_1, \ldots, c_n) :
    \Pi(i : \mathsf{Fin}(n)).\; A_i + \mathsf{Exc}
\]
This changes the return type from $A_i$ to $A_i + \mathsf{Exc}$, a
coproduct that the caller must case-split on.
\end{proposition}

\section{Cancellation and Timeouts}

Python's asyncio supports cooperative cancellation via
\lstinline[style=python]{asyncio.CancelledError}:

\begin{lstlisting}[style=python]
async def fragile_operation():
    try:
        await long_network_call()
    except asyncio.CancelledError:
        cleanup()
        raise  # must re-raise!

# Timeout after 5 seconds
result = await asyncio.wait_for(fragile_operation(), timeout=5.0)
\end{lstlisting}

\begin{definition}[Cancellation effect]
Cancellation introduces a new effect $\mathsf{Cancel}$:
\[
  \mathsf{wait\_for}(c, t) : \mathsf{Coroutine}(A + \mathsf{TimeoutError})
    \quad \text{with effect } \varepsilon_c \vee \mathsf{Cancel}
\]
The $\mathsf{Cancel}$ effect is injected at any $\mathsf{await}$ point
inside $c$ when the timeout expires.
\end{definition}

\begin{definition}[\texttt{asyncio.shield}]
The \lstinline[style=python]{asyncio.shield} combinator protects a
coroutine from external cancellation:
\[
  \mathsf{shield}(c) : \mathsf{Coroutine}(A) \quad
    \text{with } \mathsf{Cancel} \not\in \mathsf{effects}(c)
\]
If the outer task is cancelled, the shielded coroutine continues running.
The typing rule removes $\mathsf{Cancel}$ from the effect set of the
inner computation.
\end{definition}

\begin{proposition}[Shield/cancel duality]
$\mathsf{shield}$ and cancellation form an adjunction:
\[
  \mathsf{Cancel}(\mathsf{shield}(c)) \cong c \qquad
  \mathsf{shield}(\mathsf{Cancel}(c)) \ncong c
\]
That is, cancelling a shielded task has no effect on the inner coroutine,
but shielding an already-cancelled task does not reverse the cancellation.
\end{proposition}

\section{Async Iterators and \texttt{async for}}

Python provides async iteration for consuming asynchronous data streams:

\begin{lstlisting}[style=python]
class AsyncRange:
    def __init__(self, n):
        self.n = n
        self.i = 0
    def __aiter__(self):
        return self
    async def __anext__(self):
        if self.i >= self.n:
            raise StopAsyncIteration
        await asyncio.sleep(0.01)  # simulate async work
        val = self.i
        self.i += 1
        return val

async def example():
    async for x in AsyncRange(10):
        print(x)
\end{lstlisting}

\begin{definition}[Async iterator type]
An async iterator is a suspended stream:
\[
  \mathsf{AsyncIterator}(A) \defeq \mathsf{unit} \to_{\mathsf{Suspend}}
    (A + \mathsf{StopAsyncIteration})
\]
The \lstinline[style=python]{async for} loop desugars to:
\[
  \mathsf{async\;for}\; x \;\mathsf{in}\; \mathit{ait} \defeq
    \mathsf{loop}\left\{
      \mathsf{match}\;(\mathsf{await}\;\mathit{ait}.\mathsf{\_\_anext\_\_}())\;
      \mathsf{with}\; \begin{cases}
        \mathsf{inl}(x) &\Rightarrow \mathit{body}(x); \mathsf{continue} \\
        \mathsf{inr}(\_) &\Rightarrow \mathsf{break}
      \end{cases}
    \right\}
\]
\end{definition}

\begin{definition}[Async generator as async iterator]
An async generator function produces an $\mathsf{AsyncIterator}$:
\begin{lstlisting}[style=python]
async def async_countdown(n):
    while n > 0:
        await asyncio.sleep(1)
        yield n
        n -= 1
\end{lstlisting}
Its type is:
\[
  \mathsf{async\_countdown} : (n : \Nat) \to
    \mathsf{AsyncIterator}(\Nat)
    \quad \text{with effect } \IO \vee \mathsf{Suspend}
\]
\end{definition}

\section{The Event Loop as a Scheduler}

The asyncio event loop is the runtime scheduler that selects which
coroutine to resume at each step.

\begin{definition}[Event loop model]
The event loop is a function:
\[
  \mathsf{loop} : \mathsf{Queue}(\mathsf{Coroutine}) \to_\IO \mathsf{unit}
\]
that repeatedly: (1) selects a ready coroutine from the queue,
(2) runs it until the next suspension point, and (3) re-enqueues it
if not finished.
\end{definition}

\begin{definition}[Fairness]
The event loop satisfies \emph{weak fairness}: every coroutine that is
perpetually ready is eventually scheduled.  Formally:
\[
  \forall\, c \in \mathsf{Queue}.\;
    \left(\square\, \mathsf{ready}(c)\right) \Rightarrow
    \left(\diamond\, \mathsf{scheduled}(c)\right)
\]
where $\square$ is ``always'' and $\diamond$ is ``eventually'' in temporal
logic.
\end{definition}

\begin{proposition}[No starvation under cooperative scheduling]
If every coroutine yields (reaches a suspension point) within a bounded
number of steps $k$, then every coroutine in a queue of size $n$ is
scheduled at least once every $n \cdot k$ steps.
\end{proposition}
\begin{proof}
By the round-robin property of the asyncio event loop: after each
coroutine suspends, the next in queue is selected.  Since each coroutine
suspends within $k$ steps, a full round takes at most $n \cdot k$ steps.
\end{proof}

\begin{remark}[\fstar{} comparison]
\fstar{} has no notion of a scheduler or fairness properties.  Verifying
temporal properties like fairness requires a liveness logic, which is
outside \fstar{}'s current capabilities.  \synhopy{} integrates the
temporal modalities $\square$ and $\diamond$ into its effect system to
reason about scheduling guarantees.
\end{remark}

\section{Structured Concurrency with \texttt{TaskGroup}}

Python 3.11 introduced \lstinline[style=python]{asyncio.TaskGroup} for
\emph{structured concurrency}---ensuring that all spawned tasks complete
before the enclosing scope exits:

\begin{lstlisting}[style=python]
async def main():
    async with asyncio.TaskGroup() as tg:
        task1 = tg.create_task(fetch_user(1))
        task2 = tg.create_task(fetch_user(2))
    # Both tasks guaranteed complete here
    print(task1.result(), task2.result())
\end{lstlisting}

\begin{definition}[TaskGroup type]
A \lstinline[style=python]{TaskGroup} is a scoped concurrency primitive:
\[
  \mathsf{TaskGroup} : \mathsf{AsyncContextManager}\left(
    \Sigma(n : \Nat).\;
    \Pi(i : \mathsf{Fin}(n)).\; \mathsf{Task}(A_i)
  \right)
\]
The async context manager protocol ensures:
\begin{enumerate}
  \item All tasks created within the group are tracked.
  \item $\mathsf{\_\_aexit\_\_}$ awaits all tasks before returning.
  \item If any task raises, all remaining tasks are cancelled, and the
    exceptions are collected into an
    \lstinline[style=python]{ExceptionGroup}.
\end{enumerate}
\end{definition}

\begin{theorem}[Structured concurrency guarantee]
In a \lstinline[style=python]{TaskGroup} block, no task outlives the block.
Formally, if $t$ is created via
\lstinline[style=python]{tg.create_task(c)}:
\[
  \mathsf{done}(t) \;\text{holds at}\;
    \mathsf{\_\_aexit\_\_}(\mathit{tg})
\]
This eliminates ``fire-and-forget'' tasks---a common source of resource
leaks and unobserved exceptions in concurrent programs.
\end{theorem}
\begin{proof}
The \lstinline[style=python]{TaskGroup.__aexit__} implementation awaits
all tasks in the group.  If a task is cancelled (due to another task's
failure), the cancellation is itself awaited.  Therefore, all tasks
are in terminal state (\textsf{done} or \textsf{cancelled}) when
$\mathsf{\_\_aexit\_\_}$ returns.
\end{proof}

\section{Verifying Mutual Exclusion with \texttt{asyncio.Lock}}

The race condition in Section~17.4 can be resolved with
\lstinline[style=python]{asyncio.Lock}:

\begin{lstlisting}[style=python]
lock = asyncio.Lock()
shared_counter = 0

async def safe_increment():
    global shared_counter
    async with lock:
        val = shared_counter
        await asyncio.sleep(0)
        shared_counter = val + 1

async def safe_double_increment():
    await asyncio.gather(safe_increment(), safe_increment())
    assert shared_counter == 2  # always holds
\end{lstlisting}

\begin{definition}[Async lock type]
An \lstinline[style=python]{asyncio.Lock} is a binary semaphore with type:
\[
  \mathsf{Lock} : \mathsf{AsyncContextManager}(\mathsf{unit})
\]
with the invariant:
\[
  \forall\, t_1, t_2.\;
    \mathsf{holds}(\mathit{lock}, t_1) \wedge
    \mathsf{holds}(\mathit{lock}, t_2) \Rightarrow t_1 = t_2
\]
At most one task holds the lock at any time.
\end{definition}

\begin{theorem}[Lock eliminates races]
If all accesses to a shared variable $v$ are guarded by the same lock
$\ell$, then the spectral sequence differential $d_1$ for $v$ is zero:
\[
  \left(\forall\, c.\; \Mutates(c, v) \Rightarrow
    c \text{ holds } \ell\right) \;\Rightarrow\;
  d_1(\Mutates \times \Mutates, v) = 0
\]
Consequently, \synhopy{} does not report a $\mathsf{RACE\_CONDITION}$
warning.
\end{theorem}
\begin{proof}
When two coroutines both hold the lock, they cannot be interleaved at
a suspension point within the critical section.  The lock serializes
access, so the read-modify-write sequence for $v$ is atomic from the
perspective of other coroutines.  The differential is zero because the
two $\Mutates$ effects are now sequential, not concurrent.
\end{proof}

\section{Worked Example: Async Rate Limiter}

Consider an async rate limiter that allows at most $n$ requests per
second:

\begin{lstlisting}[style=python]
import asyncio
import time

class RateLimiter:
    def __init__(self, max_per_second: int):
        self.max_per_second = max_per_second
        self.semaphore = asyncio.Semaphore(max_per_second)
        self.timestamps: list[float] = []
        self.lock = asyncio.Lock()

    async def acquire(self):
        await self.semaphore.acquire()
        async with self.lock:
            now = time.monotonic()
            # Remove timestamps older than 1 second
            self.timestamps = [
                t for t in self.timestamps if now - t < 1.0
            ]
            if len(self.timestamps) >= self.max_per_second:
                sleep_time = 1.0 - (now - self.timestamps[0])
                if sleep_time > 0:
                    await asyncio.sleep(sleep_time)
            self.timestamps.append(time.monotonic())

    def release(self):
        self.semaphore.release()

    async def __aenter__(self):
        await self.acquire()
        return self

    async def __aexit__(self, *args):
        self.release()
        return False
\end{lstlisting}

\begin{definition}[Rate limiter specification]
A rate limiter with bound $n$ satisfies:
\begin{enumerate}
  \item \emph{Throughput bound}: in any 1-second window, at most $n$ requests
    are permitted:
    \[
      \forall\, t.\; |\{r \in \mathsf{requests} \mid
        r.\mathsf{time} \in [t, t+1)\}| \leq n
    \]
  \item \emph{Liveness}: every request is eventually served (no starvation):
    \[
      \forall\, r.\; \diamond\, \mathsf{served}(r)
    \]
  \item \emph{Mutual exclusion}: the timestamp list is never concurrently
    modified (protected by the lock).
\end{enumerate}
\end{definition}

\begin{lstlisting}[style=proof]
-- Verify: RateLimiter(n) satisfies the rate limiter spec

-- Throughput bound
have h_sem: semaphore(n) allows at most n concurrent holders
    by apply asyncio_semaphore_spec

have h_window: forall t.
    |{r | r.time in [t, t+1)}| <= n
    by h_sem then apply timestamp_pruning_invariant

-- Mutual exclusion on timestamps
have h_lock: self.lock guards self.timestamps
    by inspect acquire(), all mutations under async with self.lock

have h_no_race: d_1(Mutates * Mutates, timestamps) = 0
    by h_lock then apply lock_eliminates_races

-- Liveness: semaphore is always released
have h_release: forall acquire.
    eventually release is called
    by async context manager guarantee (__aexit__ calls release)

have h_live: forall r. eventually served(r)
    by h_release then apply semaphore_fairness

-- Correctness of sleep calculation
have h_sleep: sleep_time > 0
    implies after sleep, timestamps[0] is > 1s old
    by unfold acquire, sleep_time = 1 - (now - timestamps[0])

conclude: RateLimiter(n) is a correct async rate limiter
    by structural with h_window, h_no_race, h_live, h_sleep
\end{lstlisting}

\begin{theorem}[Rate limiter correctness]
The \lstinline[style=python]{RateLimiter} with bound $n$ satisfies:
\begin{enumerate}
  \item \emph{Safety} (throughput bound): at most $n$ requests are active
    in any 1-second window, enforced jointly by the semaphore and the
    timestamp-based sleep.
  \item \emph{Liveness}: every request eventually acquires the semaphore
    and is served, because the async context manager guarantees that
    \texttt{release()} is called on every path (including exceptions
    and cancellation).
  \item \emph{Data integrity}: the \lstinline[style=python]{timestamps}
    list is protected by the async lock, so no concurrent modification
    occurs.
\end{enumerate}
\end{theorem}
\begin{proof}
(1) The semaphore limits concurrency to $n$, and the timestamp pruning
ensures that the window constraint is maintained even when requests are
spaced within the same second.  (2) The $\mathsf{\_\_aexit\_\_}$ method
calls \texttt{release()}, and the async context manager protocol
guarantees $\mathsf{\_\_aexit\_\_}$ is called on every path.  By the
semaphore's internal fairness (FIFO queue), every waiting task eventually
acquires.  (3) The lock protects the critical section where
\texttt{timestamps} is read and modified, and the lock's mutual exclusion
property (at most one holder) prevents races.
\end{proof}


% ──────────────────────────────────────────────────────────────────
\chapter{Descriptors, Properties, and Slots}
\label{ch:descriptors}

\section{Python's Descriptor Protocol}

\begin{lstlisting}[style=python]
class Celsius:
    """Validates temperature values."""
    def __set_name__(self, owner, name):
        self.name = name
    def __get__(self, obj, objtype=None):
        return getattr(obj, f'_{self.name}', 0)
    def __set__(self, obj, value):
        if value < -273.15:
            raise ValueError("Below absolute zero")
        setattr(obj, f'_{self.name}', value)

class Weather:
    temperature = Celsius()
\end{lstlisting}

\section{Descriptors as Fibrations}

\begin{definition}[Descriptor fibration]
A descriptor $d$ on class $C$ for attribute $n$ defines a fibration
$\pi_d : \mathsf{AttrSpace}(C, n) \to C$ where the fiber over instance $o$
comprises the triple
$(\texttt{\_\_get\_\_}(o),\;\texttt{\_\_set\_\_}(o,-),\;\texttt{\_\_delete\_\_}(o))$.
\end{definition}

\begin{definition}[Descriptor typing rules]
\[
  \frac{
    d : \mathsf{DataDescriptor}(T) \quad
    d \in \mathsf{class\_dict}(C) \quad
    o : C
  }{
    o.n : T \;\text{via}\; d.\texttt{\_\_get\_\_}(o, C)
  }
  \;\text{(DescGet)}
\]
\[
  \frac{
    d : \mathsf{DataDescriptor}(T) \quad
    o : C \quad
    v : T
  }{
    \vdash_{\Mutates}\; o.n \mathbin{:=} v
    \;\text{via}\; d.\texttt{\_\_set\_\_}(o, v)
  }
  \;\text{(DescSet)}
\]
\end{definition}

\section{Proving Descriptor Invariants}

\begin{lstlisting}[style=proof]
-- Proof: Weather.temperature >= -273.15 is an invariant
have h_guard: Celsius.__set__ raises ValueError
    when value < -273.15
    by unfold Celsius.__set__
have h_pass: if __set__ succeeds then value >= -273.15
    by structural with h_guard  -- contrapositive
have h_init: default value = 0 >= -273.15
    by omega
conclude: forall w : Weather. w.temperature >= -273.15
    by nat_induction on set_count
    base: 0 >= -273.15 by h_init, omega
    step: each __set__ either raises or stores v >= -273.15
        by h_pass
\end{lstlisting}

\begin{remark}[\fstar{} comparison]
\fstar{} has no concept of Python descriptors.  Modeling the three-method
dispatch protocol (\texttt{\_\_get\_\_}/\texttt{\_\_set\_\_}/\texttt{\_\_delete\_\_})
plus the data-vs-non-data-descriptor priority rules would require $\sim$80
lines of explicit state-machine encoding per descriptor class.
\end{remark}

\section{Slots and Memory Layout}

\begin{lstlisting}[style=python]
class Point:
    __slots__ = ('x', 'y')
    def __init__(self, x, y):
        self.x = x
        self.y = y
\end{lstlisting}

\begin{definition}[\texttt{\_\_slots\_\_} typing]
When a class declares \texttt{\_\_slots\_\_}, the instance attribute namespace
is restricted to the declared names.  Formally:
\[
  C.\texttt{\_\_slots\_\_} = \{s_1, \ldots, s_k\}
  \;\Longrightarrow\;
  o : C \;\vdash\; \mathsf{keys}(o.\texttt{\_\_dict\_\_}) \subseteq \{s_1, \ldots, s_k\}
\]
Setting an attribute not in \texttt{\_\_slots\_\_} raises \texttt{AttributeError}.
\end{definition}

The \texttt{\_\_slots\_\_} declaration turns an open dependent-dictionary type
into a \emph{closed record type}---a product with exactly $k$ named fields.
This is the one case where Python objects behave like statically-typed records,
and \synhopy{} exploits it for tighter verification.

\section{Non-Data vs.\ Data Descriptors: The Priority Rule}

Python distinguishes \emph{data descriptors} (defining both \texttt{\_\_get\_\_}
and \texttt{\_\_set\_\_}) from \emph{non-data descriptors} (defining only
\texttt{\_\_get\_\_}).  The attribute lookup chain follows a strict priority:

\begin{enumerate}
  \item Data descriptors in the class \texttt{\_\_mro\_\_}.
  \item Instance \texttt{\_\_dict\_\_} entries.
  \item Non-data descriptors in the class \texttt{\_\_mro\_\_}.
\end{enumerate}

\begin{definition}[Descriptor priority rule]
Let $d_C$ denote a descriptor found in the class hierarchy of $C$ and let
$o.\texttt{\_\_dict\_\_}[n]$ denote an instance attribute.  The attribute
resolution $o.n$ is:
\[
  o.n \;\defeq\; \begin{cases}
    d_C.\texttt{\_\_get\_\_}(o, C) & \text{if } d_C \text{ is a data descriptor,} \\
    o.\texttt{\_\_dict\_\_}[n]     & \text{if } n \in o.\texttt{\_\_dict\_\_}, \\
    d_C.\texttt{\_\_get\_\_}(o, C) & \text{if } d_C \text{ is a non-data descriptor,} \\
    \mathsf{AttributeError}        & \text{otherwise.}
  \end{cases}
\]
\end{definition}

\begin{lemma}[Data descriptor dominance]
If $d$ is a data descriptor on class $C$ for attribute $n$, then for all
instances $o : C$ and all values $v$, assigning $o.n = v$ invokes
$d.\texttt{\_\_set\_\_}(o, v)$ regardless of the contents of
$o.\texttt{\_\_dict\_\_}$.
\end{lemma}

\begin{lstlisting}[style=proof]
-- Proof: data descriptor dominates instance dict
have h_data: d defines __get__ and __set__
    by assumption (d is DataDescriptor)
have h_priority: attribute_lookup checks data descriptors first
    by unfold object.__getattribute__
conclude: o.n = v invokes d.__set__(o, v)
    by structural with h_data, h_priority
\end{lstlisting}

\section{The \texttt{\_\_set\_name\_\_} Hook}

When a class body is executed, the metaclass calls
\texttt{d.\_\_set\_name\_\_(owner, name)} for each descriptor $d$ found in the
class namespace.  This allows descriptors to learn their attribute name.

\begin{definition}[\texttt{\_\_set\_name\_\_} typing]
\[
  \frac{
    d : \mathsf{Descriptor} \quad
    C = \mathsf{type}(\ldots, \{n \mapsto d, \ldots\}) \quad
    d.\texttt{\_\_set\_name\_\_} : (C, \mathsf{str}) \to \Mutates
  }{
    \vdash_{\Mutates}\; d.\texttt{\_\_set\_name\_\_}(C, n)
    \;\text{at class creation time}
  }
  \;\text{(SetName)}
\]
\end{definition}

\begin{remark}
The \texttt{\_\_set\_name\_\_} hook is crucial for descriptor libraries
(\texttt{attrs}, \texttt{pydantic}, Django model fields): the descriptor
stores its own attribute name so that \texttt{\_\_get\_\_} and
\texttt{\_\_set\_\_} can access the corresponding private storage slot
(e.g., \texttt{\_temperature} for a descriptor bound to \texttt{temperature}).
\end{remark}

\section{Descriptors as Vertical Morphisms in the Attribute Bundle}

We refine the fibration picture from Section~\ref{ch:descriptors}
by viewing each descriptor method as a \emph{vertical morphism} in the
attribute fiber bundle.

\begin{construction}[Attribute bundle]
For class $C$ with attribute name $n$, define the total space
$E_n = \{(o, v) \mid o : C,\; v = o.n\}$ with projection
$\pi : E_n \to C$ given by $\pi(o,v) = o$.  A descriptor $d$ defines
three vertical endomorphisms of $E_n$:
\begin{align*}
  \mathsf{get}_d &: \Fiber_\pi(o) \to \Fiber_\pi(o), \\
  \mathsf{set}_d &: \Fiber_\pi(o) \times T \to \Fiber_\pi(o), \\
  \mathsf{del}_d &: \Fiber_\pi(o) \to \Fiber_\pi(o).
\end{align*}
These are vertical because they act within a single fiber (a single instance)
without changing the base point.
\end{construction}

\begin{proposition}[Descriptor coherence]
For a well-behaved descriptor $d$ with value type $T$, the round-trip
property holds: for all $o : C$ and $v : T$,
\[
  d.\texttt{\_\_set\_\_}(o, v) ;\; d.\texttt{\_\_get\_\_}(o, C) = v.
\]
This is a section--retraction pair in the attribute fiber bundle.
\end{proposition}

\section{Computed Properties and Lazy Evaluation}

Python's \texttt{@property} decorator is syntactic sugar for a data descriptor.
Lazy properties compute their value on first access and cache the result.

\begin{lstlisting}[style=python]
class LazyProperty:
    def __init__(self, func):
        self.func = func
        self.attr_name = None
    def __set_name__(self, owner, name):
        self.attr_name = f'_lazy_{name}'
    def __get__(self, obj, objtype=None):
        if obj is None:
            return self
        if not hasattr(obj, self.attr_name):
            setattr(obj, self.attr_name, self.func(obj))
        return getattr(obj, self.attr_name)
\end{lstlisting}

\begin{definition}[Lazy property typing]
A lazy descriptor for attribute $n$ with computation $f : C \to_{\Reads} T$
has the effect signature:
\[
  n.\texttt{\_\_get\_\_} : C \to_{\Reads \cup \Mutates_{\text{once}}} T
\]
The first access has a $\Mutates$ effect (writing the cache); subsequent
accesses are $\Reads$-only.  We model this via a dependent effect indexed
by the cache state:
\[
  n.\texttt{\_\_get\_\_} : (o : C) \to_{\mathsf{eff}(o)} T
  \quad\text{where}\quad
  \mathsf{eff}(o) = \begin{cases}
    \Mutates & \text{if } o.\texttt{\_lazy\_n} \text{ unset,} \\
    \Reads   & \text{otherwise.}
  \end{cases}
\]
\end{definition}

\begin{lemma}[Lazy idempotence]
After the first access, repeated reads of a lazy property return the same
value without side effects:
$\forall o : C.\; o.n ;\; o.n = o.n$
(the second access is observationally equivalent to the first).
\end{lemma}

\section{Validation Descriptors}

Descriptors are frequently used for input validation: type checking, range
checking, or format validation at assignment time.

\begin{lstlisting}[style=python]
class Validated:
    def __init__(self, ty, predicate=None):
        self.ty = ty
        self.predicate = predicate
    def __set_name__(self, owner, name):
        self.name = name
    def __set__(self, obj, value):
        if not isinstance(value, self.ty):
            raise TypeError(f"Expected {self.ty}")
        if self.predicate and not self.predicate(value):
            raise ValueError(f"Validation failed for {self.name}")
        setattr(obj, f'_{self.name}', value)
    def __get__(self, obj, objtype=None):
        return getattr(obj, f'_{self.name}', None)
\end{lstlisting}

\begin{definition}[Validation descriptor invariant]
A validation descriptor $d$ with type $T$ and predicate $P : T \to \Bool$
establishes the class invariant:
\[
  \forall o : C.\; \forall n \in \mathsf{validated\_attrs}(C).\;
  o.n \neq \mathsf{None} \implies
  \bigl(\mathsf{isinstance}(o.n, T) \;\wedge\; P(o.n)\bigr)
\]
This invariant is \emph{inductive}: it holds at construction (no assignments
yet, all \texttt{None}) and is preserved by every \texttt{\_\_set\_\_} call
(which either raises or stores a valid value).
\end{definition}

\section{Worked Example: Verifying a Validated Model Class}

\begin{lstlisting}[style=synhopy]
class PositiveInt(Validated):
    def __init__(self):
        super().__init__(int, lambda x: x > 0)

class Temperature(Validated):
    def __init__(self):
        super().__init__(float, lambda x: x >= -273.15)

@guarantee("all fields satisfy their validation predicates")
class Measurement:
    count = PositiveInt()
    temp  = Temperature()

    def __init__(self, count: int, temp: float):
        self.count = count   # validated by PositiveInt.__set__
        self.temp  = temp    # validated by Temperature.__set__
\end{lstlisting}

\begin{lstlisting}[style=proof]
-- Proof: Measurement invariant holds after construction
-- Goal: count > 0 /\ temp >= -273.15 for any Measurement
have h_count_desc: PositiveInt is DataDescriptor(int, x > 0)
    by unfold PositiveInt.__init__, Validated.__init__
have h_temp_desc: Temperature is DataDescriptor(float, x >= -273.15)
    by unfold Temperature.__init__, Validated.__init__
-- self.count = count invokes PositiveInt.__set__
have h_count_set: if __set__ succeeds then count : int /\ count > 0
    by apply validation_descriptor_invariant with h_count_desc
-- self.temp = temp invokes Temperature.__set__
have h_temp_set: if __set__ succeeds then temp : float /\ temp >= -273.15
    by apply validation_descriptor_invariant with h_temp_desc
-- If __init__ returns normally, both assignments succeeded
conclude: forall m : Measurement.
    m.count > 0 /\ m.temp >= -273.15
    by structural with h_count_set, h_temp_set
\end{lstlisting}

\begin{remark}[\fstar{} comparison]
Modeling this in \fstar{} would require: (1) an explicit record type for
the descriptor state, (2) a state monad threading the instance dict,
(3) manual encoding of the priority rule, (4) separate lemmas for each
validated field, and (5) an invariant-preservation proof for every method
that touches validated attributes.  In \synhopy{}, the descriptor fibration
and the validation invariant definition give this automatically.
\end{remark}


% ──────────────────────────────────────────────────────────────────
\chapter{Comprehensions, Map, Filter, and Reduce}
\label{ch:comprehensions}

\section{List Comprehensions as Dependent Products}

\begin{lstlisting}[style=python]
result1 = [x**2 for x in range(10) if x % 2 == 0]
result2 = list(map(lambda x: x**2,
                   filter(lambda x: x % 2 == 0, range(10))))
\end{lstlisting}

\begin{definition}[Comprehension type]
\lstinline[style=python]|[f(x) for x in xs if p(x)]| has type:
\[
  \mathsf{List}\!\left(\{y : B \mid \exists x \in \mathit{xs}.\; p(x) \wedge y = f(x)\}\right)
\]
\end{definition}

\section{Map--Filter--Comprehension Equivalence}

\begin{theorem}[Comprehension = map $\circ$ filter]
\label{thm:comp-map-filter}
For pure $f$ and $p$:
$[\,f(x) \mid x \in \mathit{xs},\; p(x)\,]
  = \mathsf{map}(f, \mathsf{filter}(p, \mathit{xs}))$.
\end{theorem}

\begin{lstlisting}[style=proof]
by list_induction xs
base: both sides equal [] by unfold
step: assume IH for xs.
    [f(x) for x in [a] ++ xs if p(x)]
    = (if p(a) then [f(a)] else [])
      ++ [f(x) for x in xs if p(x)]
        by unfold comprehension
    = (if p(a) then [f(a)] else [])
      ++ map(f, filter(p, xs))          by rw [IH]
    = map(f, filter(p, [a] ++ xs))      by rw [filter_cons, map_append]
\end{lstlisting}

\section{Reduce / Fold and the \v{C}ech Connection}

\begin{lstlisting}[style=python]
from functools import reduce

total = reduce(lambda acc, x: acc + x, numbers, 0)
\end{lstlisting}

\begin{definition}[Fold denotation]
$\den{\mathsf{reduce}(f, \mathit{xs}, z)} = \mathsf{foldl}(f, z, \mathit{xs})$
where $\mathsf{foldl}$ is the standard left fold.
\end{definition}

This is precisely the construction that connects to deppy's existing \v{C}ech
cohomology engine: two folds over the same list with different body functions
can be compared by decomposing the bodies into a \v{C}ech cover on the
iteration domain and verifying agreement per case-region.

\begin{example}[Fold equivalence via \v{C}ech descent]
\begin{lstlisting}[style=python]
def sum_v1(xs):
    return reduce(lambda acc, x: acc + x, xs, 0)

def sum_v2(xs):
    total = 0
    for x in xs:
        total += x
    return total
\end{lstlisting}

Both are folds with the same init ($0$), same collection (\texttt{xs}), and
same body ($\lambda(\mathit{acc}, x).\; \mathit{acc} + x$).  The \v{C}ech
cover is trivial (one open set = the entire iteration domain), and the
agreement check reduces to $\lambda \mathit{acc}, x.\; \mathit{acc} + x
\equiv \lambda \mathit{acc}, x.\; \mathit{acc} + x$ which is $\Refl$.
\end{example}


\section{Dictionary and Set Comprehensions}

\begin{lstlisting}[style=python]
squares = {x: x**2 for x in range(10)}
evens   = {x for x in range(20) if x % 2 == 0}
\end{lstlisting}

\begin{definition}[Dict comprehension type]
$\{k(x) : v(x) \;\mathsf{for}\; x \;\mathsf{in}\; \mathit{xs}\}$ has type:
\[
  \Sigma(\mathit{keys} : \mathsf{Set}(\mathsf{dom}(k))).\;
  \Pi(j : \mathit{keys}).\; \mathsf{cod}(v)
\]
This is a dependent dictionary---the same type used for \texttt{**kwargs}
(Chapter~\ref{ch:args-kwargs}).
\end{definition}

\begin{remark}[\fstar{} comparison]
\fstar{} has \lstinline[style=fstar]|List.Tot.map| and
\lstinline[style=fstar]|List.Tot.filter| with library lemmas for the map--filter
equivalence.  However, \fstar{} has no dictionary or set comprehensions, no
\texttt{reduce}/\texttt{functools}, and no lazy generator expressions.  Each
would need a separate model.
\end{remark}

\section{Walrus Operator in Comprehensions}

Python 3.8 introduced the walrus operator (\texttt{:=}) for inline assignment.
Inside comprehensions, the scoping is subtle: the walrus-bound variable
\emph{leaks} to the enclosing scope.

\begin{lstlisting}[style=python]
results = [y for x in data if (y := expensive(x)) > 0]
# y is bound in the enclosing scope after the comprehension
\end{lstlisting}

\begin{definition}[Walrus scoping in comprehensions]
Let $\Gamma$ be the enclosing scope and $\Gamma_c$ the comprehension scope.
The walrus operator \lstinline[style=python]|(y := e)| in a comprehension
has the typing:
\[
  \frac{
    \Gamma_c \vdash e : T \qquad
    y \notin \mathsf{dom}(\Gamma_c)
  }{
    \Gamma_c, y : T \vdash (y \mathbin{:=} e) : T
    \quad\wedge\quad
    \Gamma' = \Gamma, y : T
  }
  \;\text{(Walrus-Comp)}
\]
The variable $y$ is added to $\Gamma'$, the \emph{enclosing} scope after
the comprehension completes. This is a scope extrusion---a vertical morphism
from $\Gamma_c$ up to $\Gamma$.
\end{definition}

\begin{remark}
The walrus operator creates a verification subtlety: the final value of $y$
in the enclosing scope is the value from the \emph{last} iteration where the
condition held.  \synhopy{} tracks this by treating $y$ as a fold accumulator
over the filtered iteration domain.
\end{remark}

\section{Generator Expressions vs.\ List Comprehensions}

Generator expressions are syntactically identical to list comprehensions but
produce values lazily.

\begin{lstlisting}[style=python]
eager  = [x**2 for x in range(10**6)]    # list: all in memory
lazy   = (x**2 for x in range(10**6))    # generator: on demand
\end{lstlisting}

\begin{definition}[Generator expression type]
A generator expression \lstinline[style=python]|(f(x) for x in xs if p(x))|
has type:
\[
  \mathsf{Generator}[T, \mathsf{None}, \mathsf{None}]
  \quad\text{where}\quad
  T = \{y : B \mid \exists x \in \mathit{xs}.\; p(x) \wedge y = f(x)\}
\]
The generator is a \emph{lazy} variant: materializing it via
\lstinline[style=python]|list(...)| yields the same type as the
corresponding list comprehension.
\end{definition}

\begin{theorem}[Lazy--eager equivalence]
\label{thm:lazy-eager}
For pure $f$ and $p$:
\[
  \mathsf{list}\bigl((f(x) \;\mathsf{for}\; x \;\mathsf{in}\;
  \mathit{xs} \;\mathsf{if}\; p(x))\bigr)
  = [f(x) \;\mathsf{for}\; x \;\mathsf{in}\; \mathit{xs}
    \;\mathsf{if}\; p(x)]
\]
\end{theorem}

\begin{lstlisting}[style=proof]
-- Proof: list(genexpr) = listcomp for pure f, p
have h_gen: list(g) materializes all yielded values in order
    by unfold list.__init__ with generator protocol
have h_same_values: genexpr yields exactly the values
    f(x) for x in xs where p(x), in iteration order
    by unfold generator_expression
have h_comp: listcomp produces exactly the values
    f(x) for x in xs where p(x), in iteration order
    by unfold list_comprehension
conclude: list(genexpr) = listcomp
    by ext (element-wise equality from h_same_values, h_comp)
\end{lstlisting}

\section{Nested Comprehension Flattening}

Nested comprehensions correspond to monadic bind (concatMap) on the list
monad.

\begin{lstlisting}[style=python]
# Nested comprehension
pairs = [(x, y) for x in range(3) for y in range(x)]
# Equivalent explicit form
pairs2 = []
for x in range(3):
    for y in range(x):
        pairs2.append((x, y))
\end{lstlisting}

\begin{definition}[Nested comprehension as monadic bind]
\[
  [f(x,y) \;\mathsf{for}\; x \;\mathsf{in}\; \mathit{xs}
    \;\mathsf{for}\; y \;\mathsf{in}\; g(x)]
  \;\defeq\;
  \mathit{xs} \mathbin{\gg\!=} (\lambda x.\; g(x) \mathbin{\gg\!=}
    (\lambda y.\; \mathsf{return}\; f(x,y)))
\]
where $\gg\!=$ is the list monad bind
$\mathit{xs} \mathbin{\gg\!=} k = \mathsf{concat}(\mathsf{map}(k, \mathit{xs}))$.
\end{definition}

\begin{theorem}[Monad law verification for list comprehensions]
Nested comprehensions satisfy the monad associativity law:
\[
  [f(x,y,z) \;\mathsf{for}\; x \;\mathsf{in}\; \mathit{xs}
    \;\mathsf{for}\; y \;\mathsf{in}\; g(x)
    \;\mathsf{for}\; z \;\mathsf{in}\; h(x,y)]
\]
equals the right-nested form, by the associativity of $\mathsf{concat}$ and
naturality of $\mathsf{map}$.
\end{theorem}

\begin{lstlisting}[style=proof]
-- Proof: associativity of nested comprehensions
-- Follows from monad associativity: (m >>= f) >>= g = m >>= (x -> f x >>= g)
have h_bind_assoc: forall m f g.
    concat(map(g, concat(map(f, m))))
    = concat(map(fun x -> concat(map(g, f(x))), m))
    by list_induction m
    base: concat(map(g, concat(map(f, []))))
        = concat(map(g, [])) = [] by unfold
    step: by rw [map_cons, concat_cons, map_append,
                 concat_append, IH]
conclude: nested comprehension flattening is associative
    by apply h_bind_assoc
\end{lstlisting}

\section{Multiple \texttt{if} Clauses as Conjunction}

A comprehension may have multiple \texttt{if} clauses, which act as a
conjunction of filters.

\begin{lstlisting}[style=python]
result = [x for x in range(100) if x % 2 == 0 if x % 3 == 0]
# equivalent to: [x for x in range(100) if x % 2 == 0 and x % 3 == 0]
\end{lstlisting}

\begin{lemma}[Multiple filters = conjunction]
\label{lem:multi-filter}
For predicates $p_1, \ldots, p_k$:
\[
  [f(x) \;\mathsf{for}\; x \;\mathsf{in}\; \mathit{xs}
    \;\mathsf{if}\; p_1(x) \;\mathsf{if}\; \cdots \;\mathsf{if}\; p_k(x)]
  = [f(x) \;\mathsf{for}\; x \;\mathsf{in}\; \mathit{xs}
    \;\mathsf{if}\; p_1(x) \wedge \cdots \wedge p_k(x)]
\]
\end{lemma}

\begin{lstlisting}[style=proof]
-- Proof: multiple if clauses = conjunction
by list_induction xs
base: [] = [] by unfold
step: assume IH.
    -- For element a, the nested filters test p_1(a), then p_2(a), ...
    -- Element a is included iff all p_i(a) hold
    have h_chain: (if p_1(a) then (if p_2(a) then ... else []) else [])
        = (if p_1(a) /\ ... /\ p_k(a) then [f(a)] else [])
        by structural (short-circuit evaluation)
    conclude by rw [h_chain, IH]
\end{lstlisting}

\section{Performance Equivalence: Comprehension vs.\ Loop}

\begin{theorem}[Comprehension--loop equivalence]
For pure $f$ and $p$, the list comprehension
\lstinline[style=python]|[f(x) for x in xs if p(x)]|
is operationally equivalent to the explicit loop:
\begin{lstlisting}[style=python]
result = []
for x in xs:
    if p(x):
        result.append(f(x))
\end{lstlisting}
That is, they produce the same list and perform the same number of
evaluations of $f$ and $p$.
\end{theorem}

\begin{lstlisting}[style=proof]
-- Proof: comprehension = explicit loop (operational)
-- Both iterate xs in order, test p(x), and collect f(x)
by list_induction xs
base: both produce [] by unfold
step: assume IH for xs.
    -- comprehension: [f(x) for x in [a]++xs if p(x)]
    --   = (if p(a) then [f(a)] else []) ++ [f(x) for x in xs if p(x)]
    -- loop: iterates a first, then xs
    --   appends f(a) iff p(a), then continues with xs
    have h_one: single-element behavior identical by unfold
    conclude: by rw [h_one, IH]
\end{lstlisting}

\section{Worked Example: Data Transformation Pipeline}

\begin{lstlisting}[style=synhopy]
@guarantee("returns average salary by department for active employees")
def avg_salary_by_dept(
    employees: list[dict]
) -> dict[str, float]:
    assume("each employee has 'dept', 'salary', 'active' keys")
    assume("salary values are non-negative floats")

    active = [e for e in employees if e['active']]

    dept_groups: dict[str, list[float]] = {}
    for e in active:
        dept = e['dept']
        check("dept is a non-empty string")
        dept_groups.setdefault(dept, []).append(e['salary'])

    result = {
        dept: sum(salaries) / len(salaries)
        for dept, salaries in dept_groups.items()
    }
    check("all departments from active employees are represented")
    check("each average is between min and max salary in that dept")
    return result
\end{lstlisting}

\begin{lstlisting}[style=proof]
-- Proof: avg_salary_by_dept is correct
-- Step 1: active filter preserves only active employees
have h_active: forall e in active. e['active'] = True
    by apply comprehension_filter_spec
-- Step 2: dept_groups groups salaries by department
have h_groups: forall dept in dept_groups.
    dept_groups[dept] = [e['salary'] for e in active if e['dept'] == dept]
    by loop_invariant with setdefault_append_spec
-- Step 3: division is safe (each group is non-empty)
have h_nonempty: forall dept in dept_groups. len(dept_groups[dept]) > 0
    by structural (setdefault creates group on first encounter)
-- Step 4: result computes correct averages
have h_avg: forall dept in result.
    result[dept] = sum(dept_groups[dept]) / len(dept_groups[dept])
    by unfold dict_comprehension
-- Step 5: coverage -- every active employee's dept is in result
have h_cover: forall e in active. e['dept'] in result
    by structural with h_groups (dept added on first encounter)
-- Step 6: average is between min and max
have h_bounds: forall dept in result.
    min(dept_groups[dept]) <= result[dept] <= max(dept_groups[dept])
    by apply avg_bounds_lemma with h_nonempty
conclude: result satisfies specification
    by structural with h_avg, h_cover, h_bounds
\end{lstlisting}


% ──────────────────────────────────────────────────────────────────
\chapter{Type Narrowing and Control Flow Analysis}
\label{ch:narrowing}

\section{Type Narrowing via isinstance}

\begin{lstlisting}[style=python]
def process(x: int | str) -> str:
    if isinstance(x, int):
        return str(x * 2)      # x : int here
    else:
        return x.upper()       # x : str here
\end{lstlisting}

\begin{definition}[Narrowing rule]
\[
  \frac{
    \Gamma \vdash x : A \qquad T \leq A \qquad
    \Gamma \vdash \mathsf{isinstance}(x, T) = \mathsf{true}
  }{
    \Gamma \vdash x : T
  }
  \;\text{(Narrow)}
  \qquad
  \frac{
    \Gamma \vdash x : A \qquad
    \Gamma \vdash \mathsf{isinstance}(x, T) = \mathsf{false}
  }{
    \Gamma \vdash x : A \setminus T
  }
  \;\text{(Narrow-Neg)}
\]
\end{definition}

\section{\v{C}ech Cover from Exhaustive Narrowing}

An exhaustive \texttt{if}/\texttt{elif}/\texttt{else} chain on \texttt{isinstance}
checks produces a \v{C}ech cover $\{U_i\}$ of the union type.  The specification
is verified \emph{per fiber}: on each $U_i$, prove the postcondition under the
narrowed type.  The \v{C}ech descent glues the per-fiber results into a global
proof---the fibers are disjoint for isinstance chains, so the overlap cocycle
condition is trivially satisfied.

\begin{lstlisting}[style=proof]
-- Proof: process always returns a string
have h_cover: int | str = U_int + U_str
    by structural  -- isinstance is exhaustive on the union
cases x:
  case x : int (in U_int):
    have: x * 2 : int by omega
    conclude: str(x * 2) : str by apply str_of_int
  case x : str (in U_str):
    conclude: x.upper() : str by apply str_upper_spec
-- Gluing: disjoint fibers, trivial overlap
conclude by Cech_descent with h_cover
\end{lstlisting}

\section{Truthiness Narrowing}

Python narrows types on truthiness checks too:

\begin{lstlisting}[style=python]
def safe_len(x: list | None) -> int:
    if x:          # narrows: x is truthy => x is non-empty list
        return len(x)
    return 0       # x is None or empty list
\end{lstlisting}

\begin{definition}[Truthiness narrowing]
\[
  \frac{
    \Gamma \vdash x : A \qquad
    \Gamma \vdash \mathsf{bool}(x) = \mathsf{true}
  }{
    \Gamma \vdash x : \{y : A \mid y \neq \mathsf{None} \wedge y \neq 0
      \wedge y \neq \texttt{""} \wedge y \neq [\,]\}
  }
  \;\text{(Truthy)}
\]
\end{definition}

\begin{remark}[\fstar{} comparison]
\fstar{}'s narrowing is limited to pattern-matching on algebraic datatypes.
It does not support truthiness narrowing, \texttt{isinstance} on arbitrary
class hierarchies, or the subtle rules of Python's ``falsy'' values
(\texttt{None}, \texttt{0}, \texttt{""}, \texttt{[]}, \texttt{False}).
\synhopy{}'s narrowing is based on the isinstance fibration and the
truthiness predicate, which match Python's actual runtime behavior.
\end{remark}

\section{Type Guards: \texttt{TypeGuard} and \texttt{TypeIs}}

Python 3.10 introduced \texttt{typing.TypeGuard} and Python 3.13 added
\texttt{typing.TypeIs} for user-defined type narrowing functions.

\begin{lstlisting}[style=python]
from typing import TypeGuard, TypeIs

def is_str_list(val: list[object]) -> TypeGuard[list[str]]:
    return all(isinstance(x, str) for x in val)

def is_positive(val: int) -> TypeIs[int]:
    return val > 0
\end{lstlisting}

\begin{definition}[\texttt{TypeGuard} narrowing]
A function $g : A \to \texttt{TypeGuard}[B]$ narrows only in the positive branch
and narrows to exactly $B$ (discarding the input type):
\[
  \frac{
    \Gamma \vdash g : A \to \texttt{TypeGuard}[B] \qquad
    \Gamma \vdash g(x) = \mathsf{true}
  }{
    \Gamma \vdash x : B
  }
  \;\text{(TypeGuard)}
\]
Note: in the \texttt{else} branch, $x$ retains its original type $A$ (no
negative narrowing).
\end{definition}

\begin{definition}[\texttt{TypeIs} narrowing]
\texttt{TypeIs} is stricter: $B$ must be a subtype of $A$, and negative
narrowing applies:
\[
  \frac{
    \Gamma \vdash g : A \to \texttt{TypeIs}[B] \qquad
    B \leq A \qquad
    \Gamma \vdash g(x) = \mathsf{true}
  }{
    \Gamma \vdash x : B
  }
  \;\text{(TypeIs-Pos)}
  \qquad
  \frac{
    \Gamma \vdash g(x) = \mathsf{false}
  }{
    \Gamma \vdash x : A \setminus B
  }
  \;\text{(TypeIs-Neg)}
\]
\end{definition}

\begin{remark}
The difference between \texttt{TypeGuard} and \texttt{TypeIs} mirrors the
distinction between a general fiber map (which may change the type entirely)
and a sub-fiber inclusion (which preserves the complement structure).
\synhopy{} models \texttt{TypeIs} as a decidable sub-object classifier.
\end{remark}

\section{Assertion-Based Narrowing}

Python's \texttt{assert} statement narrows types in subsequent code under
the assumption that assertions are not disabled.

\begin{lstlisting}[style=python]
def process(x: int | str | None) -> int:
    assert isinstance(x, int), "expected int"
    return x + 1  # x : int here
\end{lstlisting}

\begin{definition}[Assert narrowing]
\[
  \frac{
    \Gamma \vdash x : A \qquad
    \Gamma \vdash \texttt{assert}\; \phi(x)
  }{
    \Gamma, \phi(x) \vdash x : \{y : A \mid \phi(y)\}
  }
  \;\text{(Assert-Narrow)}
\]
The assertion adds $\phi(x)$ as a hypothesis to the context.  If the
assertion is an \texttt{isinstance} check, this specializes to the
\textsc{(Narrow)} rule.
\end{definition}

\begin{remark}
Unlike \texttt{if}-based narrowing, assertion narrowing is \emph{linear}:
it narrows the type for all subsequent code, not just within a branch.
The proof obligation is that $\phi(x)$ actually holds; if it does not,
the function raises \texttt{AssertionError}, which \synhopy{} models as
an exceptional exit (Chapter~\ref{ch:exceptions}).
\end{remark}

\section{Pattern Matching Exhaustiveness}

Python 3.10's \texttt{match}/\texttt{case} provides structural pattern
matching with exhaustiveness checking.

\begin{lstlisting}[style=python]
from dataclasses import dataclass

@dataclass
class Circle:
    radius: float
@dataclass
class Rect:
    width: float
    height: float

type Shape = Circle | Rect

def area(s: Shape) -> float:
    match s:
        case Circle(r):
            return 3.14159 * r ** 2
        case Rect(w, h):
            return w * h
\end{lstlisting}

\begin{definition}[Pattern exhaustiveness]
A \texttt{match} statement on $x : A_1 \mid \cdots \mid A_n$ with patterns
$p_1, \ldots, p_k$ is \emph{exhaustive} if:
\[
  A_1 \mid \cdots \mid A_n \;\subseteq\;
  \mathsf{dom}(p_1) \cup \cdots \cup \mathsf{dom}(p_k)
\]
where $\mathsf{dom}(p_i)$ is the set of values matched by pattern $p_i$.
Exhaustiveness is a \v{C}ech covering condition: the patterns must cover
the entire type.
\end{definition}

\begin{definition}[Pattern narrowing]
Within the body of \texttt{case} $p_i$, the matched variable is narrowed
to the type corresponding to pattern $p_i$:
\[
  \frac{
    \Gamma \vdash x : A_1 \mid \cdots \mid A_n \qquad
    \text{in case } p_i \text{ matching } A_j
  }{
    \Gamma \vdash x : A_j \;\wedge\; \mathsf{bindings}(p_i, x)
  }
  \;\text{(Match-Narrow)}
\]
\end{definition}

\begin{lstlisting}[style=proof]
-- Proof: area is total (exhaustive match)
have h_shape: Shape = Circle | Rect by unfold Shape
have h_patterns: {Circle(_), Rect(_, _)} covers Shape
    by structural (each variant has exactly one matching pattern)
have h_circle: case Circle(r) -> 3.14159 * r**2 : float
    by omega (r : float, result is float)
have h_rect: case Rect(w, h) -> w * h : float
    by omega (w, h : float, result is float)
conclude: area is total and returns float
    by Cech_descent with h_patterns, h_circle, h_rect
\end{lstlisting}

\section{Narrowing Through Assignment}

Certain assignments narrow the type of the target variable by construction:

\begin{lstlisting}[style=python]
def normalize(x: int | str | float) -> int:
    y = int(x)       # y : int (narrowed by constructor)
    return y + 1
\end{lstlisting}

\begin{definition}[Assignment narrowing]
\[
  \frac{
    \Gamma \vdash e : A \qquad
    f : A \to B \text{ is a type constructor/converter}
  }{
    \Gamma \vdash y \mathbin{:=} f(e) \;\Longrightarrow\; y : B
  }
  \;\text{(Assign-Narrow)}
\]
Common instances: \texttt{int(x)} narrows to \texttt{int},
\texttt{str(x)} narrows to \texttt{str}, \texttt{list(x)} narrows to
\texttt{list}.
\end{definition}

\begin{lemma}[Assignment narrowing safety]
If $f : A \to B$ is total (does not raise for any $a : A$), then
assignment narrowing is safe: $y$ is guaranteed to have type $B$.
If $f$ is partial, the narrowing holds only on the non-exceptional path.
\end{lemma}

\section{Narrowing Composition Across Function Boundaries}

When a function returns a \texttt{TypeGuard} or \texttt{TypeIs}, the
narrowing composes across call sites:

\begin{lstlisting}[style=python]
def is_valid_email(s: str) -> TypeIs[str]:
    return '@' in s and '.' in s.split('@')[1]

def send_email(addr: str, body: str) -> None:
    if is_valid_email(addr):
        # addr is narrowed to str (with is_valid_email property)
        _do_send(addr, body)
\end{lstlisting}

\begin{definition}[Narrowing composition]
If $g_1 : A \to \texttt{TypeIs}[B]$ and $g_2 : B \to \texttt{TypeIs}[C]$,
then the sequential application narrows $A$ to $C$:
\[
  \frac{
    \Gamma \vdash g_1(x) = \mathsf{true} \qquad
    \Gamma, x : B \vdash g_2(x) = \mathsf{true}
  }{
    \Gamma \vdash x : C
  }
  \;\text{(Narrow-Compose)}
\]
This composition is associative and corresponds to iterated sub-fiber
inclusion in the isinstance fibration.
\end{definition}

\section{Worked Example: JSON Schema Validator}

\begin{lstlisting}[style=synhopy]
from typing import TypeIs, Any

@guarantee("returns True iff value matches the schema")
def validate_schema(value: Any, schema: dict) -> TypeIs[Any]:
    assume("schema is a valid JSON Schema subset")
    match schema.get('type'):
        case 'string':
            return isinstance(value, str)
        case 'integer':
            check("integer implies isinstance int check")
            return isinstance(value, int) and not isinstance(value, bool)
        case 'array':
            if not isinstance(value, list):
                return False
            item_schema = schema.get('items', {})
            return all(validate_schema(v, item_schema) for v in value)
        case 'object':
            if not isinstance(value, dict):
                return False
            props = schema.get('properties', {})
            required = schema.get('required', [])
            # Check required keys present
            for key in required:
                if key not in value:
                    return False
            # Check each property against its sub-schema
            for key, sub_schema in props.items():
                if key in value:
                    if not validate_schema(value[key], sub_schema):
                        return False
            return True
        case _:
            return True  # no type constraint
\end{lstlisting}

\begin{lstlisting}[style=proof]
-- Proof: validate_schema correctly narrows types
-- Exhaustive match on schema['type']
have h_cover: schema['type'] in {string, integer, array, object, _}
    by structural (match is exhaustive with wildcard)

-- Case 'string': isinstance(value, str) is correct TypeIs
cases schema['type']:
  case 'string':
    conclude: validate_schema returns True iff value : str
        by unfold isinstance

  case 'integer':
    -- bool is subclass of int in Python, so we exclude it
    have h_int: isinstance(value, int) /\ not isinstance(value, bool)
        iff value : int \ bool
        by apply isinstance_spec, subclass_exclusion
    conclude: returns True iff value is a proper integer
        by structural with h_int

  case 'array':
    -- Recursive: validate each element against items schema
    have h_list: isinstance(value, list) required first
        by unfold (early return False)
    have h_items: all(validate_schema(v, item_schema) for v in value)
        iff forall v in value. validate_schema(v, item_schema)
        by apply all_generator_spec
    conclude: returns True iff value is list with valid items
        by structural with h_list, h_items

  case 'object':
    have h_dict: isinstance(value, dict) required first
        by unfold
    have h_required: all required keys present
        by loop_invariant (early return on missing key)
    have h_props: all present properties validated recursively
        by loop_invariant with recursive validate_schema
    conclude: returns True iff value is dict matching schema
        by structural with h_dict, h_required, h_props

conclude: validate_schema is a correct schema validator
    by Cech_descent with h_cover
\end{lstlisting}


% ──────────────────────────────────────────────────────────────────
\chapter{Denotational Semantics of Core Python}
\label{ch:denotational}

This chapter gives a compositional denotational semantics for a core Python
subset, mapping each construct into $\Sh_\infty(\Ctx)$.

\section{The Semantic Function}

\begin{definition}[Semantic function]
\[
  \den{-} : \mathsf{PythonAST} \to \mathsf{Mor}(\Sh_\infty(\Ctx))
\]
maps Python abstract syntax trees to morphisms in the $\infty$-topos.
\end{definition}

\subsection{Expressions}

\begin{align}
  \den{n}_\sigma &= n : \Nat
    && \text{(integer literal)} \\
  \den{x}_\sigma &= \sigma(x) : \PyObj
    && \text{(variable lookup)} \\
  \den{e_1 \mathbin{+} e_2}_\sigma &= \den{e_1}_\sigma + \den{e_2}_\sigma
    && \text{(arithmetic)} \\
  \den{e_1\texttt{[}e_2\texttt{]}}_\sigma &= \mathsf{getitem}(\den{e_1}_\sigma, \den{e_2}_\sigma)
    && \text{(subscript)} \\
  \den{f(e_1, \ldots, e_k)}_\sigma &= \mathsf{apply}(\den{f}_\sigma, \langle\den{e_i}_\sigma\rangle)
    && \text{(call)} \\
  \den{e\texttt{.}m}_\sigma &= \mathsf{resolve}(\den{e}_\sigma, m)
    && \text{(attribute)} \\
  \den{\lambda \bar{x}.\; e}_\sigma &= \Lambda(\bar{x}, \den{e}_{\sigma[\bar{x}:=\cdot]})
    && \text{(lambda)}
\end{align}

Note that attribute access uses $\mathsf{resolve}$, which follows the full
MRO-based descriptor protocol (Chapter~\ref{ch:dict}).

\subsection{Statements}

\begin{align}
  \den{x = e}_\sigma &= \sigma[x \mapsto \den{e}_\sigma]
    && \text{(assignment)} \\
  \den{s_1;\; s_2}_\sigma &= \den{s_2}_{\den{s_1}_\sigma}
    && \text{(sequencing)} \\
  \den{\mathsf{if}\;e:s_1\;\mathsf{else}:s_2}_\sigma &=
    \begin{cases}
      \den{s_1}_\sigma & \text{if } \mathsf{bool}(\den{e}_\sigma) \\
      \den{s_2}_\sigma & \text{otherwise}
    \end{cases}
    && \text{(conditional)} \\
  \den{\mathsf{return}\;e}_\sigma &= \mathsf{ret}(\den{e}_\sigma)
    && \text{(return)} \\
  \den{\mathsf{raise}\;e}_\sigma &= \mathsf{exc}(\den{e}_\sigma)
    && \text{(raise)}
\end{align}

\subsection{Loops}

\begin{align}
  \den{\mathsf{for}\;x\;\mathsf{in}\;e:s}_\sigma
    &= \mathsf{foldl}\!\left(\lambda \sigma' x.\; \den{s}_{\sigma'[x \mapsto x]},\;
       \sigma,\; \mathsf{iter}(\den{e}_\sigma)\right) \\
  \den{\mathsf{while}\;e:s}_\sigma
    &= \mathsf{lfp}\!\left(\lambda \sigma'.\;
       \mathsf{if}\;\mathsf{bool}(\den{e}_{\sigma'})\;
       \mathsf{then}\;\den{s}_{\sigma'}\;\mathsf{else}\;\sigma'\right)
\end{align}

The for-loop is a fold; the while-loop is a least fixed point.  Both connect
to the \v{C}ech cohomology engine for equivalence checking of different loop
implementations.

\subsection{Try/Except}

\begin{align}
  \den{\mathsf{try}:s_1\;\mathsf{except}\;E:s_2}_\sigma
    &= \begin{cases}
      \den{s_1}_\sigma & \text{if } \den{s_1}_\sigma \neq \mathsf{exc}(e') \\
      \den{s_2}_{\sigma[\mathsf{e} \mapsto e']}
        & \text{if } \den{s_1}_\sigma = \mathsf{exc}(e') \wedge e' : E \\
      \mathsf{exc}(e') & \text{if } \den{s_1}_\sigma = \mathsf{exc}(e') \wedge e' \not: E
    \end{cases}
\end{align}

This three-way case analysis is the \emph{section} of the exception bundle
from Chapter~\ref{ch:exceptions}: the handler selects which exception types
to catch and which to propagate.

\section{Soundness of the Denotation}

\begin{theorem}[Adequacy]
\label{thm:adequacy}
If $\den{e_1} = \den{e_2}$ in $\Sh_\infty(\Ctx)$, then $e_1$ and $e_2$ produce
the same observable results on all inputs.
\end{theorem}
\begin{proof}
By construction, $\den{-}$ maps to morphisms whose global sections are
observable behaviors (return values, raised exceptions, heap mutations).
Equal morphisms yield equal global sections.
\end{proof}

\begin{theorem}[Compositionality]
\label{thm:compositionality}
$\den{-}$ is compositional: the denotation of a compound construct depends
only on the denotations of its sub-constructs.
\end{theorem}
\begin{proof}
By structural induction on the AST.  Each clause composes sub-denotations
via sheaf-theoretic operations (composition, pairing, casing, fixed-point).
\end{proof}

\begin{theorem}[Effect tracking correctness]
\label{thm:effect-correct}
If $\den{e}$ has effect $\varepsilon$ as computed by the effect analyzer, then
the morphism $\den{e}$ in $\Sh_\infty(\Ctx)$ only modifies propositions within
the scope of $\varepsilon$.
\end{theorem}
\begin{proof}
The effect analyzer (Chapter~\ref{ch:effects}) conservatively over-approximates:
if it says $\Pure$, the morphism is the identity on the heap component; if it
says $\Reads$, the morphism is monotone on the heap component; etc.  Each case
is verified by the $\Omega$-endomorphism characterization.
\end{proof}


\section{Class Definitions}

\begin{align}
  \den{\mathsf{class}\;C(\bar{B}):s}_\sigma
    &= \sigma\!\left[C \mapsto \mathsf{MetaCreate}\!\left(
       \mathsf{name} = \text{``$C$''},\;
       \mathsf{bases} = \langle\den{B_i}_\sigma\rangle,\;
       \mathsf{dict} = \den{s}_{\sigma_{\mathsf{cls}}}
    \right)\right]
\end{align}

where $\sigma_{\mathsf{cls}}$ is a fresh namespace for the class body.  The
$\mathsf{MetaCreate}$ operation follows the full metaclass protocol
(Chapter~\ref{ch:metaclasses}):
\begin{enumerate}
  \item Determine metaclass $M$ from bases
  \item Call $M.\mathsf{\_\_prepare\_\_}$ to create namespace
  \item Execute body $s$ in that namespace
  \item Call $M(\text{name}, \text{bases}, \text{namespace})$
\end{enumerate}

\begin{remark}[\fstar{} comparison]
\fstar{} has no concept of runtime class creation---all types are defined at
compile time.  In Python, class bodies execute at import time and can contain
arbitrary code, including conditionally-defined methods.  \synhopy{} captures
this with $\mathsf{MetaCreate}$, which preserves the full dynamism.
\end{remark}

\section{Decorators}

\begin{align}
  \den{\texttt{@}d\;\mathsf{def}\;f(\ldots):s}_\sigma
    &= \sigma\!\left[f \mapsto \den{d}_\sigma(\den{\mathsf{def}\;f(\ldots):s}_\sigma(f))\right]
\end{align}

Stacked decorators compose bottom-up:
\begin{align}
  \den{\texttt{@}d_1\;\texttt{@}d_2\;\mathsf{def}\;f(\ldots):s}_\sigma
    &= \sigma[f \mapsto d_1(d_2(\mathsf{func}))]
\end{align}

This matches Python's semantics: the bottommost decorator is applied first.
Each decorator is a fiber bundle morphism
$d : \Pi_{(\PyFunc, U_i)} \to \Pi_{(\PyFunc, U_j)}$ (Chapter~\ref{ch:decorators}).

\section{Generators}

\begin{align}
  \den{\mathsf{def}\;g(\ldots):\ldots\mathsf{yield}\;e\ldots}_\sigma
    &= \sigma\!\left[g \mapsto \Lambda(\bar{x},\; \mathsf{Coroutine}(F_g))\right]
\end{align}

where $F_g$ is the \emph{step function}:
\[
  F_g : (\mathsf{State}_g \times \mathsf{Input}) \to
    (\mathsf{State}_g \times (\mathsf{Yield}(v) \mid \mathsf{Return}(r) \mid
      \mathsf{Raise}(e)))
\]

The coroutine is modeled as a free monad over the yield functor
(Chapter~\ref{ch:generators}):
\[
  \mathsf{Gen}(A) = \mu X.\; A + (\mathsf{Yield}(V) \times (I \to X))
\]

\begin{remark}[\fstar{} comparison]
\fstar{} has no generator/coroutine primitive.  The closest mechanism is a
state monad with explicit continuation-passing, which requires manually
threading state.  \synhopy{}'s $\mathsf{Gen}$ monad is a built-in that matches
Python's actual generator protocol, including \texttt{send()}, \texttt{throw()},
and \texttt{close()}.
\end{remark}

\section{With-Statements}

\begin{align}
  \den{\mathsf{with}\;e\;\mathsf{as}\;x: s}_\sigma
    &= \mathsf{let}\; m = \den{e}_\sigma \;\mathsf{in} \notag\\
    &\phantom{=}\; \mathsf{let}\; x = m.\mathsf{\_\_enter\_\_}() \;\mathsf{in} \notag\\
    &\phantom{=}\; \mathsf{try}\;\den{s}_{\sigma[x \mapsto x]} \notag\\
    &\phantom{=}\; \mathsf{finally}\; m.\mathsf{\_\_exit\_\_}(\ldots)
\end{align}

The with-statement is modeled as a bracket operation (acquire/release pair),
making it a natural candidate for linear type verification.

\section{Comprehensions}

\begin{align}
  \den{[e\;\mathsf{for}\;x\;\mathsf{in}\;C\;\mathsf{if}\;p]}_\sigma
    &= \mathsf{map}\!\left(\lambda x.\den{e}_{\sigma[x]},\;
       \mathsf{filter}\!\left(\lambda x.\den{p}_{\sigma[x]},\;
       \mathsf{iter}(\den{C}_\sigma)\right)\right)
\end{align}

Nested comprehensions compose:
\begin{align}
  \den{[e\;\mathsf{for}\;x\;\mathsf{in}\;C_1\;\mathsf{for}\;y\;\mathsf{in}\;C_2]}_\sigma
    &= \mathsf{bind}\!\left(\den{C_1}_\sigma,\;
       \lambda x.\;\mathsf{map}\!\left(\lambda y.\den{e},\;
       \den{C_2}_{\sigma[x]}\right)\right)
\end{align}

This is the list monad bind, making comprehensions syntactic sugar for monadic
computation --- consistent with the Haskell perspective but applied to Python's
eager evaluation.

\section{Augmented Assignment}

\begin{align}
  \den{x \mathbin{+}= e}_\sigma
    &= \sigma\!\left[x \mapsto
       \mathsf{iadd}(\sigma(x), \den{e}_\sigma) \;\mathsf{if}\;
       \mathsf{hasattr}(\sigma(x), \text{``\_\_iadd\_\_''})
       \;\mathsf{else}\;
       \sigma(x) + \den{e}_\sigma\right]
\end{align}

The subtlety: augmented assignment tries in-place mutation first
(\lstinline[style=python]{__iadd__}), falling back to creating a new object
(\lstinline[style=python]{__add__}).  This means \lstinline[style=python]{x += y}
is \emph{not} always equivalent to \lstinline[style=python]{x = x + y} ---
for mutable types like lists, the former mutates in-place while the latter
creates a new list.

\begin{theorem}[Augmented assignment dichotomy]
For an object $o$ with $\mathsf{eff}(o.\mathsf{\_\_iadd\_\_}) = \Mutates$:
\[
  \den{x \mathbin{+}= e} \neq \den{x = x + e}
  \quad\text{(in general)}
\]
The equality holds only when $o$ is immutable (e.g., \texttt{int}, \texttt{str},
\texttt{tuple}).
\end{theorem}

\begin{remark}[\fstar{} comparison]
This dichotomy is invisible in \fstar{}, which has no mutable state at the
language level.  Any attempt to model Python's augmented assignment in \fstar{}
would require manually tracking the mutability of every variable, which the
type system does not support.
\end{remark}


\section{Del Statements}

The \lstinline[style=python]{del} statement removes a binding from the
current namespace.

\begin{definition}[Denotation of \texttt{del}]
\begin{align}
  \den{\mathsf{del}\;x}_\sigma
    &= \sigma \setminus \{x\}
    && \text{(remove simple variable)} \\
  \den{\mathsf{del}\;e\texttt{.}m}_\sigma
    &= \sigma\!\left[\den{e}_\sigma.\mathsf{\_\_dict\_\_} \;\mapsto\;
       \den{e}_\sigma.\mathsf{\_\_dict\_\_} \setminus \{m\}\right]
    && \text{(delete attribute)} \\
  \den{\mathsf{del}\;e_1\texttt{[}e_2\texttt{]}}_\sigma
    &= \mathsf{delitem}(\den{e_1}_\sigma, \den{e_2}_\sigma, \sigma)
    && \text{(delete item)}
\end{align}
where $\sigma \setminus \{x\}$ is the state with binding $x$ removed, and
$\mathsf{delitem}$ invokes the \lstinline[style=python]{__delitem__} protocol.
\end{definition}

\begin{remark}
The key subtlety: after $\den{\mathsf{del}\;x}_\sigma$, any subsequent
$\den{x}_{\sigma \setminus \{x\}}$ is \emph{undefined}---it raises
\lstinline[style=python]{NameError}.  In the denotational model, this is
captured by the partiality monad: lookup in $\sigma \setminus \{x\}$
returns $\mathsf{exc}(\mathsf{NameError}(\text{``$x$''}))$.
\end{remark}

\begin{proposition}[Del is left-inverse to assignment]
For any variable $x$ not previously in scope:
\[
  \den{\mathsf{del}\;x}_{\den{x = e}_\sigma} = \sigma
\]
However, if $x$ was already bound: $\den{\mathsf{del}\;x}_{\den{x = e}_\sigma} = \sigma \setminus \{x\} \neq \sigma$
because the original binding of $x$ is also lost.
\end{proposition}


\section{Import Statements}

\begin{definition}[Denotation of \texttt{import}]
\begin{align}
  \den{\mathsf{import}\;M}_\sigma
    &= \sigma\!\left[M \mapsto \mathsf{load\_module}(M, \sigma_{\mathsf{sys}})\right]
\end{align}
where $\mathsf{load\_module}$ is defined as:
\[
  \mathsf{load\_module}(M, \sigma_{\mathsf{sys}}) =
  \begin{cases}
    \sigma_{\mathsf{sys}}.\mathsf{modules}[M]
      & \text{if $M \in \sigma_{\mathsf{sys}}.\mathsf{modules}$} \\
    \mathsf{exec\_module}(\mathsf{find}(M))
      & \text{otherwise}
  \end{cases}
\]
The first case is the module cache hit; the second executes the module source.
\end{definition}

\begin{definition}[Denotation of \texttt{from ... import ...}]
\begin{align}
  \den{\mathsf{from}\;M\;\mathsf{import}\;n_1, \ldots, n_k}_\sigma
    &= \sigma\!\left[n_j \mapsto \mathsf{getattr}(\mathsf{load\_module}(M, \sigma_{\mathsf{sys}}), n_j)
       \;\middle|\; j = 1, \ldots, k\right]
\end{align}
If any $n_j$ is not an attribute of module $M$, the result is
$\mathsf{exc}(\mathsf{ImportError}(\ldots))$.
\end{definition}

\begin{definition}[Denotation of \texttt{from M import *}]
\begin{align}
  \den{\mathsf{from}\;M\;\mathsf{import}\;*}_\sigma
    &= \sigma \cup \{n \mapsto \mathsf{getattr}(m, n) \mid n \in \mathsf{public}(m)\}
\end{align}
where $m = \mathsf{load\_module}(M, \sigma_{\mathsf{sys}})$ and
$\mathsf{public}(m) = m.\mathsf{\_\_all\_\_}$ if defined, otherwise
$\{n \in \mathsf{dir}(m) \mid \neg(\text{$n$ starts with ``\_''})\}$.
\end{definition}

\begin{remark}[\fstar{} comparison]
\fstar{}'s module system is static: \lstinline[style=fstar]|open FStar.List|
is resolved at compile time and has no runtime semantics.  Python's import
system is deeply dynamic: modules execute code at import time, can modify
\lstinline[style=python]{sys.modules}, and even replace themselves.  The
denotational semantics captures this via $\mathsf{load\_module}$, which
models the full import machinery including caching and side effects.
\end{remark}


\section{Assert Statements}

\begin{definition}[Denotation of \texttt{assert}]
\begin{align}
  \den{\mathsf{assert}\;e}_\sigma
    &= \begin{cases}
      \sigma & \text{if } \mathsf{bool}(\den{e}_\sigma) = \mathsf{True} \\
      \mathsf{exc}(\mathsf{AssertionError}()) & \text{otherwise}
    \end{cases} \\
  \den{\mathsf{assert}\;e,\; m}_\sigma
    &= \begin{cases}
      \sigma & \text{if } \mathsf{bool}(\den{e}_\sigma) = \mathsf{True} \\
      \mathsf{exc}(\mathsf{AssertionError}(\den{m}_\sigma)) & \text{otherwise}
    \end{cases}
\end{align}
\end{definition}

\begin{remark}
In \synhopy{}'s verification pipeline, \lstinline[style=python]{assert}
statements serve a dual role:
\begin{enumerate}
  \item At runtime, they are guards that raise on failure.
  \item At verification time, they are \emph{proof obligations}: the verifier
    must prove $\mathsf{bool}(\den{e}_\sigma) = \mathsf{True}$ for all
    reachable states $\sigma$.
\end{enumerate}
This duality makes \lstinline[style=python]{assert} the primary mechanism
for expressing \synhopy{} specifications inline.
\end{remark}


\section{Multiple Assignment and Tuple Unpacking}

\begin{definition}[Denotation of multiple assignment]
\begin{align}
  \den{x_1, x_2, \ldots, x_k = e}_\sigma
    &= \sigma\!\left[x_j \mapsto \pi_j(\mathsf{unpack}(\den{e}_\sigma, k))
       \;\middle|\; j = 1, \ldots, k\right]
\end{align}
where $\mathsf{unpack}(v, k)$ calls the iterator protocol to extract exactly
$k$ values:
\[
  \mathsf{unpack}(v, k) =
  \begin{cases}
    (v_1, \ldots, v_k) & \text{if $\mathsf{iter}(v)$ yields exactly $k$ values} \\
    \mathsf{exc}(\mathsf{ValueError}(\ldots)) & \text{otherwise}
  \end{cases}
\]
\end{definition}

\begin{definition}[Starred unpacking]
\begin{align}
  \den{x_1, \ldots, x_j, *y, x_{j+2}, \ldots, x_k = e}_\sigma
    &= \sigma\!\left[
      \begin{array}{l}
        x_i \mapsto v_i \;\text{for}\; i \leq j, \\
        y \mapsto [v_{j+1}, \ldots, v_{n-k+j+1}], \\
        x_i \mapsto v_{n-k+i} \;\text{for}\; i \geq j+2
      \end{array}
    \right]
\end{align}
where $(v_1, \ldots, v_n) = \mathsf{list}(\mathsf{iter}(\den{e}_\sigma))$,
provided $n \geq k - 1$.
\end{definition}

\begin{example}[Swap idiom]
The Python swap idiom \lstinline[style=python]{a, b = b, a} has denotation:
\[
  \den{a, b = b, a}_\sigma = \sigma[a \mapsto \sigma(b),\; b \mapsto \sigma(a)]
\]
This is computed by first evaluating the right-hand side to
$(\sigma(b), \sigma(a))$, then unpacking.  The simultaneous binding ensures
correctness---intermediate mutation is impossible.
\end{example}


\section{Global and Nonlocal Declarations}

\begin{definition}[Scoped state]
\label{def:scoped-state}
A \synhopy{} state $\sigma$ is a \emph{chain of scopes}:
\[
  \sigma = (\sigma_{\mathsf{local}}, \sigma_{\mathsf{enclosing}_1}, \ldots,
    \sigma_{\mathsf{enclosing}_m}, \sigma_{\mathsf{global}}, \sigma_{\mathsf{builtin}})
\]
Variable lookup follows the LEGB rule (Local, Enclosing, Global, Built-in).
\end{definition}

\begin{definition}[Denotation of \texttt{global}]
\begin{align}
  \den{\mathsf{global}\;x;\; x = e}_\sigma
    &= \sigma[\sigma_{\mathsf{global}}(x) \mapsto \den{e}_\sigma]
\end{align}
The \lstinline[style=python]{global} declaration redirects assignment of $x$
to the global scope $\sigma_{\mathsf{global}}$ rather than the local scope.
\end{definition}

\begin{definition}[Denotation of \texttt{nonlocal}]
\begin{align}
  \den{\mathsf{nonlocal}\;x;\; x = e}_\sigma
    &= \sigma[\sigma_{\mathsf{enclosing}_j}(x) \mapsto \den{e}_\sigma]
\end{align}
where $j$ is the innermost enclosing scope that contains a binding for $x$.
If no enclosing scope binds $x$, the declaration is a
$\mathsf{SyntaxError}$.
\end{definition}

\begin{remark}[\fstar{} comparison]
\fstar{} has no concept of mutable global or nonlocal state: all \fstar{}
bindings are immutable.  State is modeled explicitly via the
\lstinline[style=fstar]|ST| effect with heap references.  Python's scope
chain and \lstinline[style=python]{global}/\lstinline[style=python]{nonlocal}
declarations are implicit mutation mechanisms that \synhopy{} must capture
faithfully in the denotation.
\end{remark}


\section{Pattern Matching (Match/Case)}

Python 3.10 introduced structural pattern matching.  We give its
denotational semantics.

\begin{definition}[Denotation of \texttt{match}]
\begin{align}
  &\den{\mathsf{match}\;e\!:\; \mathsf{case}\;p_1\!:\;s_1 \;\cdots\; \mathsf{case}\;p_k\!:\;s_k}_\sigma \notag\\
  &\quad= \begin{cases}
    \den{s_j}_{\sigma \cup \mathsf{bindings}(p_j, v)}
      & \text{for the first $j$ with $\mathsf{matches}(p_j, v) = \mathsf{True}$} \\
    \sigma & \text{if no pattern matches}
  \end{cases}
\end{align}
where $v = \den{e}_\sigma$.
\end{definition}

The pattern-matching function $\mathsf{matches}$ and binding extraction
$\mathsf{bindings}$ are defined by cases on the pattern syntax:

\begin{definition}[Pattern semantics]
\label{def:pattern-sem}
\begin{align}
  \mathsf{matches}(\_, v) &= \mathsf{True}
    && \text{(wildcard)} \\
  \mathsf{matches}(x, v) &= \mathsf{True}, \quad \mathsf{bindings}(x, v) = \{x \mapsto v\}
    && \text{(capture)} \\
  \mathsf{matches}(c, v) &= (v \equiv c)
    && \text{(literal)} \\
  \mathsf{matches}((p_1, \ldots, p_n), v)
    &= \mathsf{len}(v) = n \;\wedge\; \bigwedge_{i} \mathsf{matches}(p_i, v_i)
    && \text{(sequence)} \\
  \mathsf{matches}(C(k_1\!=p_1, \ldots), v)
    &= \mathsf{isinstance}(v, C) \;\wedge\; \bigwedge_{i} \mathsf{matches}(p_i, \mathsf{getattr}(v, k_i))
    && \text{(class)} \\
  \mathsf{matches}(\{k_1\!: p_1, \ldots\}, v)
    &= \bigwedge_{i} (k_i \in v) \;\wedge\; \mathsf{matches}(p_i, v[k_i])
    && \text{(mapping)} \\
  \mathsf{matches}(p_1 \mid p_2, v)
    &= \mathsf{matches}(p_1, v) \;\vee\; \mathsf{matches}(p_2, v)
    && \text{(or-pattern)} \\
  \mathsf{matches}(p\;\mathsf{if}\;g, v)
    &= \mathsf{matches}(p, v) \;\wedge\; \mathsf{bool}(\den{g}_{\sigma \cup \mathsf{bindings}(p,v)})
    && \text{(guard)}
\end{align}
\end{definition}

\begin{remark}[\fstar{} comparison]
\fstar{}'s pattern matching is exhaustiveness-checked at compile time and
only operates over algebraic data types (inductive types with constructors).
Python's match/case is far more general: it matches against classes via
\lstinline[style=python]{__match_args__}, supports mapping patterns, or-patterns,
and guards with arbitrary Boolean expressions.  \synhopy{}'s denotation captures
all of these by reducing to the structural operations above.
\end{remark}


\section{Fixed-Point Semantics}

Recursive definitions in Python---recursive functions, mutually recursive
classes, circular imports---all require fixed-point semantics.

\begin{definition}[Fixed-point operator]
The least fixed-point operator $\mathsf{lfp}$ for a monotone function
$F : (\sigma \to \sigma)$ on the domain of Python states is:
\[
  \mathsf{lfp}(F) = \bigsqcup_{n \geq 0} F^n(\bot)
\]
where $\bot$ is the everywhere-undefined state and $\bigsqcup$ is the
supremum in the domain ordering.
\end{definition}

\begin{definition}[Denotation of recursive functions]
\begin{align}
  \den{\mathsf{def}\;f(\bar{x})\!:\;s}_\sigma
    &= \sigma\!\left[f \mapsto \mathsf{lfp}\!\left(
       \lambda g.\;\lambda \bar{x}.\;\den{s}_{\sigma[f \mapsto g, \bar{x} \mapsto \bar{x}]}
    \right)\right]
\end{align}
The function $f$ is defined as the least fixed point of the functional that
maps a candidate implementation $g$ to the result of executing the body $s$
with $f$ bound to $g$.
\end{definition}

\begin{definition}[Mutual recursion]
For mutually recursive functions $f_1, \ldots, f_k$:
\begin{align}
  \den{\mathsf{def}\;f_1\ldots; \;\mathsf{def}\;f_k\ldots}_\sigma
    &= \sigma\!\left[\vec{f} \mapsto \mathsf{lfp}\!\left(
       \lambda (g_1, \ldots, g_k).\;
       (\den{s_1}_{\sigma[\vec{f} := \vec{g}]}, \ldots,
        \den{s_k}_{\sigma[\vec{f} := \vec{g}]})
    \right)\right]
\end{align}
using the product fixed-point on $k$-tuples of functions.
\end{definition}

\begin{proposition}[Existence of fixed points]
The fixed point $\mathsf{lfp}(F)$ exists for all Python function definitions,
provided we work in the domain of \emph{partial functions} (i.e., functions
that may diverge on some inputs).
\end{proposition}
\begin{proof}
The functional $F$ is Scott-continuous on the domain of partial functions
ordered by the information ordering $f \sqsubseteq g$ iff
$\mathsf{dom}(f) \subseteq \mathsf{dom}(g)$ and $f(x) = g(x)$ for all
$x \in \mathsf{dom}(f)$.  By the Kleene fixed-point theorem, $\mathsf{lfp}(F)$
exists as $\bigsqcup_n F^n(\bot)$.
\end{proof}


\section{Domain Theory Connection}

We now connect the denotational semantics to classical domain theory,
providing the order-theoretic foundations for Python's evaluation.

\begin{definition}[CPO of Python values]
\label{def:cpo-pyval}
The domain $\mathcal{D}$ of Python values is a \emph{complete partial order}
(CPO) with:
\begin{enumerate}
  \item Bottom element $\bot$ representing divergence (non-termination).
  \item For each Python type $T$, a flat domain $T_\bot = T \cup \{\bot\}$
    (discrete ordering except $\bot \sqsubseteq x$ for all $x$).
  \item For function types, the continuous function space
    $[A_\bot \to B_\bot]_\bot$ (Scott-continuous functions).
  \item For recursive types (e.g., linked lists), the bilimit construction
    $\mathcal{D} \cong F(\mathcal{D})$ using the inverse-limit technique.
\end{enumerate}
\end{definition}

\begin{definition}[Scott continuity]
A function $f : \mathcal{D}_1 \to \mathcal{D}_2$ between CPOs is
\emph{Scott-continuous} if it preserves directed suprema:
\[
  f\!\left(\bigsqcup_i d_i\right) = \bigsqcup_i f(d_i)
\]
for every directed set $\{d_i\}_{i \in I} \subseteq \mathcal{D}_1$.
\end{definition}

\begin{theorem}[Semantic functions are Scott-continuous]
\label{thm:scott-continuous}
For every Python expression $e$, the semantic function
$\den{e}_{(-)} : \mathcal{D}_\sigma \to \mathcal{D}$ is Scott-continuous.
\end{theorem}
\begin{proof}
By structural induction on $e$:
\begin{itemize}
  \item \emph{Variables}: $\den{x}_\sigma = \sigma(x)$ is a projection, hence
    continuous.
  \item \emph{Application}: $\den{f(a)}_\sigma = \den{f}_\sigma(\den{a}_\sigma)$.
    By induction, $\den{f}_{(-)}$ and $\den{a}_{(-)}$ are continuous; application
    in a continuous function space is continuous.
  \item \emph{Conditionals}: $\den{\mathsf{if}\;e:s_1\;\mathsf{else}:s_2}_\sigma$
    is a case split followed by continuous functions, hence continuous.
  \item \emph{Fixed points}: $\mathsf{lfp}(F) = \bigsqcup_n F^n(\bot)$ is a
    supremum of continuous functions, hence continuous by the Kleene theorem.
\end{itemize}
\end{proof}

\begin{definition}[Lifting and the error domain]
To model exceptions, we use the \emph{lifted sum}:
\[
  \mathcal{D}_\mathsf{exc} \defeq \mathcal{D} + \mathsf{Exc}_\bot
\]
where $\mathsf{Exc}_\bot$ is the flat domain of exception values.  The
semantic function maps into this domain:
$\den{e}_\sigma \in \mathcal{D}_\mathsf{exc}$, returning either a normal
value or an exceptional value.
\end{definition}

\begin{remark}[\fstar{} comparison]
\fstar{}'s semantics is defined by a type-directed operational semantics,
not a denotational one.  There is no domain-theoretic fixed-point construction
because \fstar{} requires all functions to terminate (via the $\mathsf{Tot}$
effect) unless annotated with $\mathsf{Div}$.  \synhopy{}'s denotational
approach accommodates non-termination naturally through the CPO structure,
matching Python's actual behavior where infinite loops are valid programs.
\end{remark}


\section{Well-Definedness of the Denotation}

We now prove that the denotational semantics is well-defined for
\emph{all} of core Python.

\begin{theorem}[Well-definedness]
\label{thm:well-defined}
For every syntactically valid Python program $P$, the denotation
$\den{P}_{\sigma_0}$ is a well-defined element of $\mathcal{D}_\mathsf{exc}$,
where $\sigma_0$ is the initial state containing built-in bindings.
\end{theorem}
\begin{proof}
We verify that each clause of $\den{-}$ produces a well-defined result:

\begin{enumerate}
  \item \textbf{Expressions} (literals, variables, arithmetic, attribute access,
    function calls, lambdas): Each clause composes well-defined operations
    on the CPO $\mathcal{D}$.  Variable lookup is total on $\sigma$ (returning
    $\mathsf{exc}(\mathsf{NameError})$ for unbound names).

  \item \textbf{Simple statements} (assignment, del, return, raise, assert,
    global, nonlocal, import): Each modifies $\sigma$ via well-defined
    operations on the scope chain.

  \item \textbf{Compound statements} (if/else, for, while, try/except/finally,
    with, match/case): Conditionals are case splits on Boolean values; loops
    use $\mathsf{lfp}$ which exists by Kleene's theorem; exception handling
    is a three-way case analysis.

  \item \textbf{Definitions} (def, class, async def): Function definitions
    use the fixed-point construction (well-defined by the Kleene theorem);
    class definitions use $\mathsf{MetaCreate}$ which is a finite sequence
    of well-defined operations.

  \item \textbf{Decorators}: Composition of well-defined functions.

  \item \textbf{Generators and coroutines}: Modeled as elements of the free
    monad $\mathsf{Gen}(A)$, which is a well-defined algebraic construction.

  \item \textbf{Comprehensions}: Syntactic sugar for map/filter/bind over
    the list monad, all well-defined.

  \item \textbf{Augmented assignment}: Reduces to method call plus assignment,
    both previously shown well-defined.

  \item \textbf{Pattern matching}: Finite case analysis with well-defined
    pattern predicates.
\end{enumerate}

Since every Python syntactic form has a well-defined clause, and the clauses
compose via well-defined domain-theoretic operations (composition, pairing,
case analysis, fixed-point), the denotation is total as a function from
programs to $\mathcal{D}_\mathsf{exc}$.
\end{proof}

\begin{corollary}[Completeness of the denotation]
The semantic function $\den{-}$ covers all constructs in the Python 3.12
grammar.  No syntactically valid program is left without a denotation.
\end{corollary}

\begin{remark}[\fstar{} comparison]
\fstar{} does not provide a denotational semantics for its own language.
Its semantics is defined operationally via a small-step reduction relation.
The denotational approach of \synhopy{} has the advantage of being
\emph{compositional}---the meaning of a program is determined by the meanings
of its parts---which directly enables the sheaf-theoretic local-to-global
reasoning that powers \synhopy{}'s verification engine.
\end{remark}


% ──────────────────────────────────────────────────────────────────
\chapter{String Operations and Regular Expressions}
\label{ch:strings}

\section{Strings as Free Monoids}

Python strings are elements of the free monoid $\Sigma^* = \mathsf{List}(\mathsf{Char})$
with concatenation $(+)$ as the monoid operation.

\section{Verifying String-Processing Functions}

\begin{lstlisting}[style=synhopy]
@guarantee("returns email username (part before @)")
def get_username(email: str) -> str:
    assume("email contains exactly one '@'")
    parts = email.split('@')
    check("len(parts) == 2")
    return parts[0]
\end{lstlisting}

\begin{lstlisting}[style=proof]
have h_split: email.split('@') returns [before, after]
    when email contains exactly one '@'
    by apply str_split_spec
have h_len: len([before, after]) = 2  by omega
conclude: parts[0] = before = username
    by structural with h_split
\end{lstlisting}

\section{F-strings and Format Specifications}

\begin{lstlisting}[style=python]
def format_price(amount: float, currency: str = "USD") -> str:
    return f"{currency} {amount:.2f}"
\end{lstlisting}

\begin{definition}[F-string denotation]
\begin{multline*}
  \den{\texttt{f"\{e\_1\} \{e\_2:.2f\}"}}_\sigma \\
  = \mathsf{str}(\den{e_1}_\sigma) + \texttt{" "} + \mathsf{format}(\den{e_2}_\sigma, \texttt{".2f"})
\end{multline*}
\end{definition}

F-strings desugar to concatenation of formatted sub-expressions.  The format
spec \texttt{:.2f} is a constraint on the output: two decimal places for floats.

\begin{remark}[\fstar{} comparison]
\fstar{}'s string library (FStar.String) has basic operations but no format
strings, no \texttt{split} with separator semantics matching Python's, and no
regular expression support.  Every Python string idiom would need its own
model.
\end{remark}

\section{String Methods and Algebraic Properties}

Many Python string methods satisfy algebraic laws that \synhopy{} can
verify and exploit.

\begin{definition}[Split--join inverse]
For any separator $s$ and string $t$ that does not end or begin with $s$:
\[
  s.\texttt{join}(t.\texttt{split}(s)) = t
\]
More generally, \texttt{join} is a left inverse of \texttt{split}:
$\forall t : \mathsf{str},\; s : \mathsf{str}.\;
  s.\texttt{join}(t.\texttt{split}(s)) = t$.
\end{definition}

\begin{lemma}[Strip idempotence]
\label{lem:strip-idempotent}
\texttt{strip} is idempotent:
$\forall t : \mathsf{str}.\; t.\texttt{strip}().\texttt{strip}() =
t.\texttt{strip}()$.
The same holds for \texttt{lstrip} and \texttt{rstrip}.
\end{lemma}

\begin{lstlisting}[style=proof]
-- Proof: strip is idempotent
have h_def: t.strip() removes leading/trailing whitespace
    by unfold str.strip
have h_no_ws: t.strip() has no leading or trailing whitespace
    by structural with h_def
conclude: t.strip().strip() = t.strip()
    by structural with h_no_ws
    -- stripping a string with no leading/trailing whitespace is identity
\end{lstlisting}

\begin{definition}[Additional string laws]
\begin{align*}
  t.\texttt{upper}().\texttt{upper}() &= t.\texttt{upper}()
    &&\text{(upper idempotent)} \\
  t.\texttt{lower}().\texttt{lower}() &= t.\texttt{lower}()
    &&\text{(lower idempotent)} \\
  t.\texttt{replace}(a, b).\texttt{replace}(a, b)
    &= t.\texttt{replace}(a, b)
    &&\text{(replace idempotent if } a \neq b\text{)} \\
  \texttt{len}(t.\texttt{split}(s))
    &= t.\texttt{count}(s) + 1
    &&\text{(split count law)}
\end{align*}
\end{definition}

\section{Regular Expression Verification}

Python's \texttt{re} module provides regular expression matching.
\synhopy{} models regex patterns as formal languages and verifies that
string operations produce values matching expected patterns.

\begin{lstlisting}[style=python]
import re

def extract_digits(s: str) -> list[str]:
    return re.findall(r'\d+', s)

def is_valid_phone(s: str) -> bool:
    return bool(re.match(r'^\+?1?\d{10}$', s))
\end{lstlisting}

\begin{definition}[Regex type refinement]
A regex pattern $r$ induces a refinement type:
\[
  \mathsf{Matches}(r) \;\defeq\; \{s : \mathsf{str} \mid
    \mathsf{re.fullmatch}(r, s) \neq \mathsf{None}\}
\]
Functions returning \texttt{re.findall} produce
$\mathsf{List}(\mathsf{Matches}(r))$.
\end{definition}

\begin{proposition}[Regex match soundness]
If $s \in \mathsf{Matches}(r)$ and $r$ is a pattern that only matches strings
of length $\geq k$, then $\texttt{len}(s) \geq k$.  Regex-derived refinements
compose with length refinements.
\end{proposition}

\section{Bytes vs.\ Str: Encoding and Decoding}

Python strictly separates \texttt{str} (Unicode text) from \texttt{bytes}
(raw byte sequences).  The bridge is \texttt{encode}/\texttt{decode}.

\begin{lstlisting}[style=python]
text: str = "Hello"
raw: bytes = text.encode('utf-8')
back: str = raw.decode('utf-8')
assert text == back
\end{lstlisting}

\begin{definition}[Encode--decode round-trip]
For encoding $e$ (e.g., \texttt{utf-8}):
\[
  \forall t : \mathsf{str}.\;
  t.\texttt{encode}(e).\texttt{decode}(e) = t
\]
provided all characters in $t$ are representable in encoding $e$.
Conversely, for valid byte sequences:
\[
  \forall b : \mathsf{bytes}.\;
  \mathsf{valid}_e(b) \implies
  b.\texttt{decode}(e).\texttt{encode}(e) = b
\]
\end{definition}

\begin{definition}[Bytes/str typing separation]
\synhopy{} enforces the invariant that \texttt{str} and \texttt{bytes}
are distinct types with no implicit coercion:
\[
  \frac{
    \Gamma \vdash x : \mathsf{str} \qquad
    f : \mathsf{bytes} \to T
  }{
    \Gamma \vdash f(x) : \mathsf{TypeError}
  }
  \;\text{(Bytes-Str-Sep)}
\]
This prevents the common Python~2-era bug of mixing bytes and strings.
\end{definition}

\section{String Interpolation Safety}

F-strings and \texttt{.format()} can introduce injection vulnerabilities
when user input is interpolated into SQL queries, shell commands, or HTML.

\begin{lstlisting}[style=python]
# UNSAFE: SQL injection possible
query = f"SELECT * FROM users WHERE name = '{user_input}'"

# SAFE: parameterized query
query = "SELECT * FROM users WHERE name = %s"
cursor.execute(query, (user_input,))
\end{lstlisting}

\begin{definition}[Interpolation taint tracking]
\synhopy{} tracks \emph{taint} on strings: a string derived from user input
is marked $\mathsf{Tainted}$.  Interpolating a tainted string into a
security-sensitive context (SQL, shell, HTML) without sanitization is a
type error:
\[
  \frac{
    \Gamma \vdash s : \mathsf{Tainted}(\mathsf{str}) \qquad
    \mathsf{ctx} \in \{\mathsf{SQL}, \mathsf{Shell}, \mathsf{HTML}\}
  }{
    \Gamma \vdash \texttt{f"}\ldots\{s\}\ldots\texttt{"} :
    \mathsf{UnsafeInterpolation}(\mathsf{ctx})
  }
  \;\text{(Taint-Check)}
\]
Parameterized queries and escaping functions act as
$\mathsf{sanitize} : \mathsf{Tainted}(\mathsf{str}) \to \mathsf{str}$,
removing the taint.
\end{definition}

\section{Unicode Normalization and Comparison}

Unicode strings can have multiple representations for the same visual
character.  Python's \texttt{unicodedata.normalize} converts between forms.

\begin{lstlisting}[style=python]
import unicodedata

s1 = '\u00e9'          # e-acute (single codepoint)
s2 = 'e\u0301'         # e + combining acute
assert s1 != s2
assert unicodedata.normalize('NFC', s1) == unicodedata.normalize('NFC', s2)
\end{lstlisting}

\begin{definition}[Unicode normalization equivalence]
Two strings $s_1, s_2$ are \emph{canonically equivalent} if
$\mathsf{NFC}(s_1) = \mathsf{NFC}(s_2)$.  \synhopy{} provides a
quotient type:
\[
  \mathsf{NormStr} \;\defeq\; \mathsf{str} / {\sim_{\mathsf{NFC}}}
\]
where $s_1 \sim_{\mathsf{NFC}} s_2 \iff \mathsf{NFC}(s_1) = \mathsf{NFC}(s_2)$.
String comparison on $\mathsf{NormStr}$ is comparison after normalization.
\end{definition}

\begin{lemma}[NFC idempotence]
$\forall s : \mathsf{str}.\; \mathsf{NFC}(\mathsf{NFC}(s)) = \mathsf{NFC}(s)$.
\end{lemma}

\section{Template Strings and \texttt{format\_map}}

Python's \texttt{string.Template} and \texttt{str.format\_map} provide
dictionary-driven string interpolation.

\begin{lstlisting}[style=python]
from string import Template

t = Template("Hello, $name! You have $count messages.")
result = t.substitute(name="Alice", count=5)
# Also: "Hello, {name}!".format_map({"name": "Alice"})
\end{lstlisting}

\begin{definition}[Template typing]
A template string $t$ with placeholders $\{p_1, \ldots, p_k\}$ requires
a substitution dictionary of type:
\[
  \mathsf{substitute} : \Pi(d : \{p_1 : T_1, \ldots, p_k : T_k\}) \to \mathsf{str}
\]
where each $T_i$ must support \texttt{\_\_str\_\_}.  Missing keys raise
\texttt{KeyError}; \synhopy{} statically checks that all placeholders are
present in the dictionary.
\end{definition}

\begin{definition}[\texttt{format\_map} denotation]
\[
  \den{s.\texttt{format\_map}(d)}_\sigma =
  \mathsf{concat}\bigl[\mathsf{interleave}(\mathsf{literals}(s),
    [\mathsf{str}(d[k]) \mid k \in \mathsf{fields}(s)])\bigr]
\]
where $\mathsf{fields}(s)$ extracts the field names from the format string
and $\mathsf{literals}(s)$ extracts the literal segments between fields.
\end{definition}

\section{Worked Example: SQL Query Builder with Injection Prevention}

\begin{lstlisting}[style=synhopy]
from typing import Any

@guarantee("builds a safe SQL query with no injection vulnerabilities")
def build_select(
    table: str,
    columns: list[str],
    conditions: dict[str, Any]
) -> tuple[str, list[Any]]:
    assume("table matches [a-zA-Z_][a-zA-Z0-9_]*")
    assume("all column names match [a-zA-Z_][a-zA-Z0-9_]*")

    # Validate identifiers (no user input in SQL structure)
    check("table is a safe identifier")
    for col in columns:
        check("col is a safe identifier")

    cols = ", ".join(columns) if columns else "*"
    if conditions:
        where_clauses = [f"{k} = %s" for k in conditions]
        where_str = " AND ".join(where_clauses)
        query = f"SELECT {cols} FROM {table} WHERE {where_str}"
        params = list(conditions.values())
    else:
        query = f"SELECT {cols} FROM {table}"
        params = []

    check("query contains no user-supplied values, only %s placeholders")
    check("params list has same length as number of %s in query")
    return query, params
\end{lstlisting}

\begin{lstlisting}[style=proof]
-- Proof: build_select prevents SQL injection
-- Step 1: table and column names are safe identifiers
have h_table_safe: table matches [a-zA-Z_][a-zA-Z0-9_]*
    by assumption
have h_cols_safe: forall col in columns.
    col matches [a-zA-Z_][a-zA-Z0-9_]*
    by assumption

-- Step 2: identifiers cannot contain SQL injection payloads
have h_no_inject_id: safe identifiers contain no quotes,
    semicolons, or SQL keywords as substrings
    by regex_subset_lemma
    -- [a-zA-Z_][a-zA-Z0-9_]* is disjoint from injection patterns

-- Step 3: condition values are parameterized, never interpolated
have h_param: condition values appear only in params list
    by structural (f-string uses only k (key names), not values)
have h_placeholder: where_clauses use "%s" for all values
    by unfold list_comprehension

-- Step 4: query structure contains only safe identifiers + %s
have h_safe_query: query is composed of:
    safe identifiers (table, columns, condition keys) + SQL keywords + %s
    by structural with h_table_safe, h_cols_safe, h_placeholder

-- Step 5: params count matches placeholder count
have h_count: len(params) = query.count("%s")
    by omega (one %s per condition, one param per condition)

conclude: build_select returns a safe parameterized query
    by structural with h_safe_query, h_param, h_count
\end{lstlisting}

\begin{remark}[\fstar{} comparison]
\fstar{} has no built-in string interpolation, taint tracking, Unicode
normalization, or regex support.  Verifying injection safety would require:
(1) a custom taint monad, (2) explicit string algebra lemmas, (3) a regex
formalization as a separate library, and (4) manual proof that parameterized
queries are safe.  \synhopy{} handles all of these natively through the
free monoid model, taint refinement types, and the regex specification language.
\end{remark}


% ══════════════════════════════════════════════════════════════════
% PART III: SOUNDNESS
% ══════════════════════════════════════════════════════════════════

\part{Soundness}

% ──────────────────────────────────────────────────────────────────
\chapter{The $\infty$-Topos Model}
\label{ch:model}

\section{The Site of Python Execution Contexts}

\begin{definition}[The site $\Ctx$]
The site $\Ctx$ has:
\begin{itemize}
  \item \textbf{Objects}: Python execution contexts---stack frames $(\Gamma, \sigma, \kappa)$
    where $\Gamma$ is the variable environment, $\sigma$ is the heap, and $\kappa$ is
    the continuation stack.
  \item \textbf{Morphisms}: Execution steps $f : C_1 \to C_2$---single steps of
    the Python abstract machine (assignment, function call, return, raise, yield).
  \item \textbf{Grothendieck topology}: Covering families are \emph{complete
    case analyses}:
    \[
      \{U_i \hookrightarrow U\}_{i \in I} \text{ covers } U \iff
      \bigcup_i U_i = U \text{ (every runtime path is covered)}
    \]
    Concretely: \lstinline[style=python]{isinstance} checks that cover all subtypes,
    \lstinline[style=python]{if/elif/else} chains that cover all cases, pattern
    matches that are exhaustive.
\end{itemize}
\end{definition}

\section{The $\infty$-Topos $\Sh_\infty(\Ctx)$}

\begin{theorem}[Model existence]
\label{thm:model}
The $\infty$-category $\Sh_\infty(\Ctx)$ of $\infty$-sheaves on $\Ctx$ is a
model of Homotopy Type Theory with the univalence axiom.
\end{theorem}
\begin{proof}
By Shulman (2019, Theorem 1.1), every Grothendieck $\infty$-topos models HoTT.
$\Sh_\infty(\Ctx)$ is a Grothendieck $\infty$-topos because $\Ctx$ is a small
site with a subcanonical Grothendieck topology.
\end{proof}

In this model:
\begin{itemize}
  \item \textbf{Types} are $\infty$-sheaves on $\Ctx$: type information that
    varies coherently across execution contexts.
  \item \textbf{Terms} are global sections: values that are well-defined in every
    execution context.
  \item \textbf{Paths} are natural transformations between sections: behavioral
    equivalences that hold uniformly.
  \item \textbf{Higher paths} are higher natural transformations: equivalences
    between equivalences.
\end{itemize}

\section{Validating the Duck-Typing Axiom}

\begin{theorem}[Duck-typing axiom holds in the model]
\label{thm:duck-model}
In $\Sh_\infty(\Ctx)$, Axiom~\ref{ax:duck-univalence} holds: duck-type equivalent
classes are equal.
\end{theorem}
\begin{proof}
Two classes $A$ and $B$ with the same method signatures generate the same
sieves on $\Ctx$: for any context $C$ and variable $x$ typed as $A$ or $B$,
the set of valid operations on $x$ (method calls, attribute accesses) is
identical. Therefore the representable sheaves $\mathsf{y}(A)$ and $\mathsf{y}(B)$
are equivalent in $\Sh_\infty(\Ctx)$. By univalence (which holds in any
Grothendieck $\infty$-topos), this equivalence gives a path
$A =_{\PyClass} B$.
\end{proof}


\section{Effects in the Model}

The effect lattice (Chapter~\ref{ch:effects}) embeds into the model as a
sublattice of the subobject classifier $\Omega$.

\begin{definition}[Effect embedding]
Each effect $\varepsilon$ determines a subobject of the morphism space:
\begin{align}
  [\Pure] &= \{f : A \to B \mid f \text{ is a morphism of sheaves}\} \\
  [\Reads] &= \{f : A \to B \mid f \text{ is a section of the heap projection}\} \\
  [\Mutates] &= \{f : A \to B \mid f \text{ is an endomorphism of the heap sheaf}\} \\
  [\IO] &= \{f : A \to B \mid f \text{ is a morphism of the IO monad}\} \\
  [\Nondet] &= \{f : A \to B \mid f \text{ is a multi-valued section}\}
\end{align}
\end{definition}

\begin{theorem}[Effect lattice homomorphism]
The embedding $[-] : \mathsf{Effects} \hookrightarrow \mathsf{Sub}(\Omega)$
is a lattice homomorphism that preserves joins and meets.
\end{theorem}
\begin{proof}
The effect join $\varepsilon_1 \vee \varepsilon_2$ corresponds to the union
of the subobject sets, which is the join in $\mathsf{Sub}(\Omega)$.
The meet is intersection.  Monotonicity follows from the chain
$\Pure \subseteq \Reads \subseteq \Mutates \subseteq \IO$.
\end{proof}

\begin{corollary}[Effect-aware equational reasoning in the model]
The equational reasoning criterion (Theorem~\ref{thm:equational}) holds in the
model: two calls $f(x)$ and $f(x)$ in different positions are equal in
$\Sh_\infty(\Ctx)$ if and only if $[\mathsf{eff}(f)] \subseteq [\Reads]$.
\end{corollary}


\section{Python Values as Sheaf Sections}

Every Python value $v$ is interpreted as a global section of a sheaf on $\Ctx$.

\begin{definition}[Value sheaf]
For a Python type $\tau$, define the sheaf $\mathcal{V}_\tau$ on $\Ctx$:
\[
  \mathcal{V}_\tau(C) = \{v \mid v \text{ is a value of type } \tau \text{ in context } C\}
\]
with restriction maps given by narrowing (losing information about the context).
\end{definition}

\begin{example}[Key sheaf sections]
\begin{itemize}
  \item An integer $n$ is a \emph{constant} section: $\mathcal{V}_{\mathsf{int}}(C) = n$
    for all contexts $C$.
  \item A closure $\lambda x.\, e$ is a section that depends on the context
    (through its captured variables): $\mathcal{V}_{\PyFunc}(C) = \lambda x.\, \den{e}_{C}$.
  \item A mutable list $[1, 2, 3]$ is a section of the \emph{heap sheaf}: its
    identity depends on the heap address, not just its contents.
\end{itemize}
\end{example}

\begin{theorem}[Sheaf condition for Python values]
$\mathcal{V}_\tau$ satisfies the sheaf condition: if values agree on overlapping
contexts (i.e., they are the same object in the intersection of two contexts),
they glue to a unique value.
\end{theorem}
\begin{proof}
Two execution contexts $C_1$ and $C_2$ that agree on a variable $x$ (both
have $x \mapsto v$) must see the same value $v$ for $x$.  This is the
uniqueness requirement of the sheaf condition.  Existence: given compatible
values on a cover, we can always construct the merged context by taking the
union of the variable bindings (well-defined because they agree on overlaps).
\end{proof}


\section{The Čech Cohomology Connection}

The equivalence checking engine uses Čech cohomology within this model.

\begin{theorem}[Čech descent for program equivalence]
Two functions $f, g$ are equivalent ($f =_{\PyFunc} g$) if and only if
the Čech cohomology class $[\alpha] \in \check{H}^1(\mathcal{U}, \mathcal{F})$
vanishes, where $\mathcal{U}$ is the covering of the input domain by
type-based fibers and $\mathcal{F}$ is the behavioral presheaf.
\end{theorem}
\begin{proof}
A Čech 1-cocycle $\alpha_{ij} : U_i \cap U_j \to \mathsf{Diff}$ records
disagreements between $f$ and $g$ on overlapping input fibers.  If $[\alpha] = 0$,
the sections $f|_{U_i}$ and $g|_{U_i}$ glue to a global equivalence.
If $[\alpha] \neq 0$, there exists an input where $f$ and $g$ produce
observably different outputs.
\end{proof}


\section{The Universe Hierarchy in the Model}

\begin{definition}[Universe interpretation]
\label{def:universe-hierarchy}
The universe hierarchy $\Type_0 : \Type_1 : \Type_2 : \cdots$ is interpreted
in $\Sh_\infty(\Ctx)$ via the \emph{object classifier} construction.  For each
universe level $n$, define:
\[
  \den{\Type_n} \defeq \mathcal{U}_n \in \Sh_\infty(\Ctx)
\]
where $\mathcal{U}_n$ is the sheaf whose sections over a context $C$ are the
$n$-truncated sheaves on the slice $\Ctx/C$:
\[
  \mathcal{U}_n(C) = \{A \in \Sh_\infty(\Ctx/C) \mid A \text{ is $n$-small}\}
\]
\end{definition}

\begin{proposition}[Cumulative universes]
\label{prop:cumul-univ}
The universe hierarchy is cumulative: $\mathcal{U}_n \hookrightarrow \mathcal{U}_{n+1}$
for all $n$, and each inclusion is a monomorphism of sheaves.
\end{proposition}
\begin{proof}
An $n$-small sheaf is also $(n{+}1)$-small, since $n$-truncation implies
$(n{+}1)$-truncation.  The inclusion map is faithful because two $n$-small
sheaves that are equal as $(n{+}1)$-small sheaves must already be equal at
level $n$---the truncation functor $\tau_{\leq n}$ is idempotent.
\end{proof}

\begin{remark}[Python types and universe levels]
In practice, most Python types inhabit $\Type_0$.  The higher levels arise when
reasoning about \emph{types of types}: $\mathsf{type}$ (Python's metaclass) lives
in $\Type_1$, and constructions like ``the type of all metaclasses'' live in
$\Type_2$.  The system assigns universe levels automatically through
bidirectional type inference (Chapter~\ref{ch:proof-kernel}).
\end{remark}

\begin{definition}[Type families and dependent types]
A dependent type $B : A \to \Type_n$ is interpreted as a morphism
$\den{B} : \den{A} \to \mathcal{U}_n$ in $\Sh_\infty(\Ctx)$.  The dependent
product and sum are given by:
\begin{align}
  \den{\Pi_{x:A} B(x)} &= \prod_{\den{A}} \den{B} & \text{(internal product)} \\
  \den{\Sigma_{x:A} B(x)} &= \sum_{\den{A}} \den{B} & \text{(internal sum)}
\end{align}
where $\prod$ and $\sum$ are the right adjoint and left adjoint to pullback
along $\den{B}$, respectively.
\end{definition}


\section{Univalence Axiom Validation}

\begin{theorem}[Univalence in $\Sh_\infty(\Ctx)$]
\label{thm:univalence-model}
The canonical map
\[
  \mathsf{idtoeqv} : (A =_{\mathcal{U}_n} B) \to (A \simeq B)
\]
is an equivalence in $\Sh_\infty(\Ctx)$ for all $n$ and all $A, B : \mathcal{U}_n$.
\end{theorem}
\begin{proof}
We verify that Voevodsky's univalence axiom holds by constructing the inverse
map and showing it is a section and retraction.

\textbf{Step 1: Constructing the inverse.}
Given an equivalence $e : A \simeq B$ in $\Sh_\infty(\Ctx)$, we construct a path
$p_e : A =_{\mathcal{U}_n} B$ as follows.  The equivalence $e$ consists of a morphism
$f : A \to B$ together with a proof $\mathsf{isEquiv}(f)$.  We define the path
$p_e : \interval \to \mathcal{U}_n$ by:
\[
  p_e(i) \defeq \Glue\!\left[
    (i = 0) \mapsto (A, f),\;
    (i = 1) \mapsto (B, \mathsf{id}_B)
  \right]
\]
using the $\Glue$ type of cubical type theory.  In the sheaf model, this is
the colimit of the diagram:
\[
  A \xleftarrow{\;f\;} A \xrightarrow{\;\mathsf{id}\;} A
  \quad\rightsquigarrow\quad
  A \xleftarrow{\;f\;} A \xrightarrow{\;f\;} B
\]
taken fiber-wise over $\Ctx$.

\textbf{Step 2: Section condition.}
We verify $\mathsf{idtoeqv}(p_e) = e$.  By construction, $p_e(0) = A$ and
$p_e(1) = B$, and the induced transport map $\transp^{\mathsf{id}}_{\mathcal{U}_n}(p_e) : A \to B$
is exactly $f$.  Since $\mathsf{idtoeqv}$ extracts the transport map, we get
$\mathsf{idtoeqv}(p_e) = (f, \mathsf{isEquiv}(f)) = e$.

\textbf{Step 3: Retraction condition.}
Given a path $p : A =_{\mathcal{U}_n} B$, we must show
$p_{\mathsf{idtoeqv}(p)} = p$.  Both sides are paths in $\mathcal{U}_n$, and they
agree at the endpoints.  By the Kan condition in $\mathcal{U}_n$
(which holds because $\mathcal{U}_n$ is a Kan complex in the $\infty$-topos),
there exists a homotopy between them, which is itself unique up to higher
homotopy.  This completes the proof that $\mathsf{idtoeqv}$ is an equivalence.
\end{proof}

\begin{corollary}[Duck typing follows from univalence]
Axiom~\ref{ax:duck-univalence} (duck-type equivalence implies equality) is not
an additional axiom but a \emph{consequence} of univalence applied to the
universe $\PyClass : \mathcal{U}_0$.
\end{corollary}


\section{Higher Inductive Types in the Model}

Higher inductive types (HITs) appear naturally in the Python context as quotient
constructions.

\begin{definition}[HITs as colimits]
\label{def:hits-model}
A higher inductive type $W$ with point constructors $c_i : A_i \to W$ and
path constructors $p_j : B_j \to (c_{j_0}(b) =_W c_{j_1}(b))$ is interpreted
in $\Sh_\infty(\Ctx)$ as the colimit:
\[
  \den{W} = \mathsf{colim}\!\left(
    \coprod_j \den{B_j} \rightrightarrows \coprod_i \den{A_i}
  \right)
\]
where the two maps are the source and target of the path constructors.
\end{definition}

\begin{example}[Python runtime constructs as HITs]
\label{ex:python-hits}
Several Python runtime constructs correspond to HITs:
\begin{enumerate}
  \item \textbf{Exceptions as pushouts.}  The exception hierarchy
    $\mathsf{Exception} \leftarrow \mathsf{BaseException} \to \mathsf{KeyboardInterrupt}$
    gives a pushout HIT:
    \[
      \mathsf{ExcTree} = \mathsf{push}\!\left(
        \mathsf{Exception} \leftarrow \mathsf{BaseException} \to \mathsf{KeyboardInterrupt}
      \right)
    \]
    The path constructor witnesses that both branches share the
    $\mathsf{BaseException}$ interface.

  \item \textbf{MRO as a quotient.}  The C3 linearization of the method resolution
    order is a quotient HIT:
    \[
      \MRO(C) = \mathsf{List}(\PyClass) / {\sim_{\mathsf{C3}}}
    \]
    where $\sim_{\mathsf{C3}}$ identifies linearizations that agree on method
    dispatch.  The path constructor witnesses that two valid MRO orderings
    dispatch identically.

  \item \textbf{Generator states as a suspension.}  A generator's state space is
    modeled by the suspension $\Sigma(\mathsf{State})$:
    \[
      \mathsf{GenState} = \Sigma(\mathsf{State}) = \{
        \mathsf{north}, \mathsf{south}
      \} \cup \{\mathsf{merid}(s) : \mathsf{north} =_{\mathsf{GenState}} \mathsf{south}
        \mid s : \mathsf{State}\}
    \]
    where $\mathsf{north}$ is the initial state and $\mathsf{south}$ is the
    exhausted state.
\end{enumerate}
\end{example}

\begin{theorem}[HIT elimination in the model]
The elimination principle for a HIT $W$ with constructors as above is sound
in $\Sh_\infty(\Ctx)$: for any $P : W \to \Type$, to define a section
$f : \Pi_{w:W} P(w)$, it suffices to give:
\begin{enumerate}
  \item For each point constructor $c_i$: a section $f_i : \Pi_{a:A_i} P(c_i(a))$
  \item For each path constructor $p_j$: a path-over
    $\mathsf{apd}_{f}(p_j(b)) : \transp^P(p_j(b), f_{j_0}(b)) =_{P(c_{j_1}(b))} f_{j_1}(b)$
\end{enumerate}
\end{theorem}
\begin{proof}
By the universal property of colimits in $\Sh_\infty(\Ctx)$: a map out of the
colimit is determined by compatible maps out of the diagram.
\end{proof}


\section{The Interval Type and Path Spaces}

\begin{construction}[The interval object]
\label{constr:interval}
The interval type $\interval$ is the free path in $\Sh_\infty(\Ctx)$, defined
as the representable sheaf:
\[
  \den{\interval} = \Delta^1 \defeq \Hom_{\Ctx}(-, [1])
\]
where $[1]$ is the walking morphism (the category $0 \to 1$) viewed as an
object of $\Ctx$ via the inclusion of finite posets.  The two endpoints are:
\begin{align}
  \den{0_\interval} &: 1 \to \Delta^1 && \text{(source inclusion)} \\
  \den{1_\interval} &: 1 \to \Delta^1 && \text{(target inclusion)}
\end{align}
\end{construction}

\begin{definition}[Path space]
For a sheaf $A$ and sections $a, b : 1 \to A$, the path space is:
\[
  \den{a =_A b} = \{p : \Delta^1 \to A \mid p \circ \den{0_\interval} = a,\;
    p \circ \den{1_\interval} = b\}
\]
This is a pullback in $\Sh_\infty(\Ctx)$:
\[
\begin{tikzcd}
  (a =_A b) \ar[r] \ar[d] & A^{\Delta^1} \ar[d, "{(\mathsf{ev}_0, \mathsf{ev}_1)}"] \\
  1 \ar[r, "{(a, b)}"'] & A \times A
\end{tikzcd}
\]
\end{definition}

\begin{proposition}[Path space as function space]
The total path space $\Path_A \defeq \sum_{a,b:A}(a =_A b)$ is equivalent to
$A^{\Delta^1}$, the internal hom from the interval to $A$.
\end{proposition}
\begin{proof}
By the universal property of the internal hom in $\Sh_\infty(\Ctx)$: a map
$X \to A^{\Delta^1}$ corresponds to a map $X \times \Delta^1 \to A$, which
decomposes as a pair of endpoints $(a, b) : X \to A \times A$ together with a
path between them.
\end{proof}

\begin{remark}[Cubical vs.\ simplicial]
In the implementation, we use a cubical interval $\interval = [0, 1]$ with
face maps $\partial_0, \partial_1$ and connections $\wedge, \vee$ following
Cohen--Coquand--Huber--M\"ortberg.  In the sheaf model, the cubical and
simplicial approaches give equivalent path spaces because $\Sh_\infty(\Ctx)$
is a Grothendieck $\infty$-topos and all $\infty$-topoi have equivalent
homotopy theories.
\end{remark}


\section{Propositional Truncation and Mere Propositions}

\begin{definition}[Mere proposition]
\label{def:mere-prop}
A type $P$ is a \emph{mere proposition} (or $(-1)$-type) in $\Sh_\infty(\Ctx)$
if, for all $x, y : P$, the path space $x =_P y$ is contractible:
\[
  \mathsf{isProp}(P) \defeq \Pi_{x,y:P}\, \mathsf{isContr}(x =_P y)
\]
Equivalently, $P$ is a subsheaf of the terminal object: a truth value.
\end{definition}

\begin{definition}[Propositional truncation]
\label{def:prop-trunc}
For any type $A$, its propositional truncation $\|A\|$ is the HIT:
\[
  \|A\| \defeq A / {\sim_{\mathsf{all}}}
\]
with one point constructor $|{-}| : A \to \|A\|$ and one path constructor
$\mathsf{trunc} : \Pi_{x,y:\|A\|}\, x =_{\|A\|} y$.  In the model:
\[
  \den{\|A\|} = \mathsf{im}(\den{A} \to 1) \in \mathsf{Sub}(1)
\]
the image of the unique map to the terminal object, which is a subobject of
$1$, i.e., a truth value.
\end{definition}

\begin{proposition}[Truncation in Python verification]
In \synhopy{}, propositional truncation corresponds to ``existence without
witness'': $\|A\|$ asserts that $A$ is inhabited without specifying which
element.  Concretely:
\begin{itemize}
  \item $\|\mathsf{solutions}(P)\|$ means ``the problem $P$ has a solution''
    without producing one.
  \item $\|\mathsf{isEquiv}(f)\|$ means ``$f$ is an equivalence'' without
    exhibiting the inverse explicitly.
  \item $\|\mathsf{sorted}(xs) \wedge \mathsf{perm}(xs, ys)\|$ means ``a sorted
    permutation exists'' without specifying the permutation.
\end{itemize}
\end{proposition}

\begin{theorem}[Propositional truncation is well-defined]
$\den{\|A\|}$ is a sheaf on $\Ctx$ whenever $\den{A}$ is.
\end{theorem}
\begin{proof}
The image factorization in a topos yields a subobject of $1$, which is
automatically a sheaf.  The sheaf condition is trivially satisfied: a compatible
family of truth values on a cover glues uniquely because truth values have no
non-trivial automorphisms.
\end{proof}


\section{Function Extensionality}

\begin{theorem}[Function extensionality in $\Sh_\infty(\Ctx)$]
\label{thm:funext-model}
Function extensionality holds: for $f, g : A \to B$,
\[
  (\Pi_{x:A}\, f(x) =_B g(x)) \;\simeq\; (f =_{A \to B} g)
\]
\end{theorem}
\begin{proof}
In $\Sh_\infty(\Ctx)$, the function type $A \to B$ is the internal hom
$B^A$.  A path $f =_{B^A} g$ in the internal hom is a map
$\Delta^1 \to B^A$, which by adjunction corresponds to $A \times \Delta^1 \to B$,
which decomposes as a family of paths $\{f(x) =_B g(x)\}_{x:A}$.  The
equivalence is thus a consequence of the exponential adjunction in the topos.
\end{proof}

\begin{corollary}[Behavioral equivalence is extensional]
\label{cor:behavioral-ext}
Two Python functions $f$ and $g$ with the same domain are equal in
$\Sh_\infty(\Ctx)$ if and only if they agree on all inputs:
\[
  f =_{\PyFunc} g \;\iff\; \forall x.\, f(x) =_{\PyObj} g(x)
\]
This is the formal justification for \synhopy{}'s behavioral equivalence
checking: to verify $f = g$, it suffices to check $f(x) = g(x)$ for all
input fibers $x$.
\end{corollary}

\begin{remark}[Contrast with intensional theories]
In intensional MLTT without function extensionality, two functions can agree on
all inputs yet be provably distinct (because they compute differently).
\synhopy{} deliberately adopts the extensional view via the sheaf model,
matching Python's runtime semantics where functions are black boxes distinguished
only by their input-output behavior.
\end{remark}


\section{Relationship with Presheaf Models}

\begin{definition}[The presheaf model]
The category of presheaves $\mathsf{PSh}(\Ctx) = [\Ctx^{\mathsf{op}}, \Set]$
is the ``naive'' model where every functor $\Ctx^{\mathsf{op}} \to \Set$ is a
type.  The sheaf model $\Sh_\infty(\Ctx)$ is a reflective subcategory:
\[
  \Sh_\infty(\Ctx) \xrightarrow{\;\iota\;} \mathsf{PSh}_\infty(\Ctx)
  \xrightarrow{\;L\;} \Sh_\infty(\Ctx)
\]
where $\iota$ is the inclusion and $L$ is the sheafification functor (left
adjoint to $\iota$).
\end{definition}

\begin{theorem}[Presheaf embedding]
\label{thm:presheaf-embed}
The inclusion $\iota : \Sh_\infty(\Ctx) \hookrightarrow \mathsf{PSh}_\infty(\Ctx)$
preserves:
\begin{enumerate}
  \item Finite limits (in particular, identity types and $\Sigma$-types)
  \item Dependent products ($\Pi$-types)
  \item The natural numbers object
\end{enumerate}
It does \emph{not} preserve colimits in general, which is why sheafification $L$
is needed for HITs.
\end{theorem}
\begin{proof}
The inclusion of a topos into its presheaf completion is a geometric morphism
whose inverse image part (sheafification) preserves finite limits.  Since $\iota$
is the direct image part, it preserves limits (being a right adjoint) and
$\Pi$-types (by the Chevalley criterion for locally cartesian closed categories).
For the natural numbers object: $\Nat$ is the constant sheaf on $\mathbb{N}$,
which agrees with the constant presheaf.
\end{proof}

\begin{proposition}[Sheafification and Python types]
\label{prop:sheafification}
A presheaf $F$ on $\Ctx$ satisfies the sheaf condition if and only if for
every covering family $\{U_i \to U\}$ (exhaustive case analysis), the
restriction maps
\[
  F(U) \to \prod_i F(U_i) \rightrightarrows \prod_{i,j} F(U_i \times_U U_j)
\]
form an equalizer diagram.  For Python types, this means: a type's values at
a given context are uniquely determined by its values on all sub-cases of that
context.
\end{proposition}

\begin{remark}[Why sheaves, not presheaves]
The presheaf model $\mathsf{PSh}_\infty(\Ctx)$ also models HoTT, but it allows
``phantom types'' that distinguish contexts even when all runtime observations
agree.  The sheafification requirement eliminates these phantoms:
in $\Sh_\infty(\Ctx)$, if two values are indistinguishable on every covering
sieve, they are equal.  This matches Python's runtime semantics, where values
are determined by their observable behavior.
\end{remark}

\begin{theorem}[Model comparison]
\label{thm:model-comparison}
The following diagram of $\infty$-categories commutes:
\[
\begin{tikzcd}
  \synhopy\text{-types} \ar[r, "\den{-}"] \ar[dr, "\den{-}_{\mathsf{PSh}}"'] &
  \Sh_\infty(\Ctx) \ar[d, hook, "\iota"] \\
  & \mathsf{PSh}_\infty(\Ctx)
\end{tikzcd}
\]
and $L \circ \den{-}_{\mathsf{PSh}} = \den{-}$, so the sheaf interpretation is
the sheafification of the presheaf interpretation.
\end{theorem}
\begin{proof}
By construction: the presheaf interpretation assigns to each type the functor
sending a context to its set of values in that context.  Sheafification forces
compatibility on covers, which corresponds to requiring exhaustive case analysis.
The triangle commutes because $L \circ \iota = \mathsf{id}$.
\end{proof}


% ──────────────────────────────────────────────────────────────────
\chapter{Consistency and Conservativity}
\label{ch:consistency}

\section{Consistency}

\begin{theorem}[Consistency]
\label{thm:consistency}
\synhopy{} is consistent: $\nvdash 0 =_\Nat 1$.
\end{theorem}
\begin{proof}
The model $\Sh_\infty(\Ctx)$ is non-degenerate: the natural numbers sheaf has
distinct global sections $0$ and $1$ (corresponding to the Python integers
\lstinline[style=python]{0} and \lstinline[style=python]{1}). Since the model
distinguishes them, no proof of $0 =_\Nat 1$ can exist.
\end{proof}

\begin{remark}[\fstar{} comparison]
\fstar{}'s consistency depends on the assumption that no one writes
\lstinline[style=fstar]{assume val absurd: squash False}. If such an assumption
is introduced (intentionally or accidentally), the entire proof environment
becomes inconsistent with no warning. \synhopy{}'s consistency is a
\emph{theorem} about the core theory (without axioms), and axiom-dependent
proofs are explicitly tracked at a lower trust level.
\end{remark}

\section{Conservativity over CCTT}

\begin{theorem}[Conservativity]
\label{thm:conservativity}
Any theorem provable in deppy's CCTT is also provable in \synhopy{}.
\end{theorem}
\begin{proof}
We give a translation $T$ from CCTT proof terms to \synhopy{} proof terms:
\begin{itemize}
  \item $T(\mathsf{OVar}(x)) = x$
  \item $T(\mathsf{OApp}(f, a)) = f(T(a))$
  \item $T(\mathsf{OLam}(x, b)) = \lambda x.\, T(b)$
  \item $T(\Delta(a, b)) = \Refl_{T(a)}$ when $a \equiv b$ by computation
  \item $T(\mathsf{Kan}(u, A, \overline{\varphi_i, a_i})) = \comp^{T(A)}(\overline{T(\varphi_i), T(a_i)})$
\end{itemize}
The translation preserves typing: if $\Gamma \vdash_{\mathsf{CCTT}} \pi : A$,
then $T(\Gamma) \vdash_{\synhopy} T(\pi) : T(A)$.
\end{proof}

\section{Anti-Hallucination}

\begin{theorem}[Anti-hallucination]
\label{thm:anti-hallucination}
If an LLM proposes a proof term $\pi$ that the \synhopy{} kernel accepts as
a proof of $a =_A b$, then $a$ and $b$ are genuinely equal in the model
$\Sh_\infty(\Ctx)$.
\end{theorem}
\begin{proof}
The kernel checks:
\begin{enumerate}
  \item \textbf{Endpoint correctness}: $\pi(0) \equiv a$ and $\pi(1) \equiv b$
  \item \textbf{Type preservation}: $\pi(i) : A$ for all $i$
  \item \textbf{Constructor validity}: every sub-term of $\pi$ is a valid
    proof constructor applied to valid arguments
\end{enumerate}
Since the proof term language is constructive (every proof is built from
verified atomic constructors), and the kernel checks each constructor
application, a well-typed proof term in \synhopy{} corresponds to a genuine
path in the model.

An LLM hallucination would be an ill-typed proof term. The kernel's type
checking catches all type errors in proof terms, so hallucinated proofs are
rejected.
\end{proof}

We now classify the specific hallucination modes and show that each is caught.

\begin{proposition}[Hallucination taxonomy]
\label{prop:hallucination-taxonomy}
Every LLM hallucination in a proof attempt falls into exactly one of these
categories, each caught by a specific kernel check:

\begin{enumerate}
  \item \textbf{Wrong endpoints}: the LLM claims to prove $a = b$ but the
    proof term actually witnesses $a' = b'$.  Caught by endpoint check.
  \item \textbf{Missing step}: the proof uses $\Trans(\pi_1, \pi_2)$ but the
    right endpoint of $\pi_1$ differs from the left endpoint of $\pi_2$.
    Caught by middle-term inference.
  \item \textbf{Unsound axiom}: the proof invokes $\mathsf{Axiom}(\mathit{name})$
    for a non-existent or incorrectly-stated axiom.  Caught by the axiom
    registry lookup; even if accepted, tracked at $\mathsf{AXIOM\_TRUSTED}$ level.
  \item \textbf{Type error in congruence}: the proof uses $\Cong(f, \pi)$ but
    $f$ is not actually a function, or $\pi$'s type doesn't match $f$'s domain.
    Caught by type checking the $\Cong$ constructor.
  \item \textbf{Invalid induction}: the proof uses $\mathsf{NatInd}$ but the
    base case or inductive step is ill-typed.  Caught by checking each sub-proof.
  \item \textbf{Effect mismatch}: the proof claims two functions are equivalent
    but they have different effects (one pure, one mutating).  Caught by effect
    verification.
\end{enumerate}
\end{proposition}

\section{Trusted Computing Base}

\begin{definition}[TCB of \synhopy{}]
The trusted computing base of \synhopy{} consists of:
\begin{enumerate}
  \item The proof kernel (Definition~\ref{def:proof-terms} and typing rules)
  \item The Z3 solver (for arithmetic discharge)
  \item The Python runtime (for executing the verified program)
  \item User-supplied axioms (tracked at $\mathsf{AXIOM\_TRUSTED}$ level)
\end{enumerate}
\end{definition}

\begin{remark}[\fstar{} comparison]
\fstar{}'s TCB consists of:
\begin{enumerate}
  \item The \fstar{} type checker
  \item The Z3 solver
  \item The extraction mechanism (F$^*$ $\to$ OCaml/C)
  \item \lstinline[style=fstar]{assume val} declarations
\end{enumerate}
\synhopy{}'s TCB is comparable in size but has two advantages:
\begin{itemize}
  \item No extraction mechanism---the verified code IS the running code
  \item Axiom dependencies are tracked, so the impact of an incorrect axiom
    is bounded to the proofs that depend on it
\end{itemize}
\end{remark}


\section{Normalization}

\begin{definition}[Neutral and normal forms]
\label{def:normal-forms}
Define the sets of \emph{neutral terms} $\mathsf{Ne}$ and \emph{normal terms}
$\mathsf{Nf}$ mutually inductively:
\begin{align}
  \mathsf{Ne} &\ni n ::= x \mid n\;v \mid n.\mathsf{proj}_i
    \mid \mathsf{elim}_W(n, \overline{c_i}) \\
  \mathsf{Nf} &\ni v ::= \lambda x.\, v \mid (v_1, v_2) \mid c_i(v)
    \mid \Refl_v \mid \mathsf{glue}[\varphi \mapsto v_1](v_0)
    \mid n
\end{align}
A term is in \emph{normal form} if it is in $\mathsf{Nf}$, and is
\emph{normalizing} if it reduces to a normal form under the operational
semantics.
\end{definition}

\begin{theorem}[Normalization]
\label{thm:normalization}
Every well-typed term $\Gamma \vdash a : A$ in the core \synhopy{} theory
(without axioms or oracle calls) is strongly normalizing: all reduction
sequences starting from $a$ terminate.
\end{theorem}
\begin{proof}[Proof sketch]
We adapt the reducibility candidates method of Girard (1972) to the cubical
setting, following Abel--\"Ohman--Vezzosi (2018).

\textbf{Step 1: Reducibility families.}
For each type $A$, define a family $\mathcal{R}_A$ of term sets (``reducibility
candidates'') satisfying:
\begin{enumerate}
  \item[(CR1)] Every term in $\mathcal{R}_A$ is strongly normalizing.
  \item[(CR2)] $\mathcal{R}_A$ is closed under neutral expansion: if
    $n \in \mathsf{Ne}$ and all one-step reducts of $n$ are in $\mathcal{R}_A$,
    then $n \in \mathcal{R}_A$.
  \item[(CR3)] $\mathcal{R}_A$ is closed under reduction: if $a \in \mathcal{R}_A$
    and $a \to a'$, then $a' \in \mathcal{R}_A$.
\end{enumerate}

\textbf{Step 2: Interpretation of types.}
\begin{align}
  \mathcal{R}_{\Pi_{x:A}B} &= \{f \mid \forall a \in \mathcal{R}_A.\;
    f\;a \in \mathcal{R}_{B[a/x]}\} \\
  \mathcal{R}_{\Sigma_{x:A}B} &= \{(a, b) \mid a \in \mathcal{R}_A,\;
    b \in \mathcal{R}_{B[a/x]}\} \\
  \mathcal{R}_{a =_A b} &= \{p : \interval \to A \mid p(0) = a,\; p(1) = b,\;
    \forall i.\; p(i) \in \mathcal{R}_A\} \\
  \mathcal{R}_\Nat &= \{n \mid n \to^* \mathsf{zero} \text{ or }
    n \to^* \mathsf{succ}(m) \text{ with } m \in \mathcal{R}_\Nat\}
\end{align}

\textbf{Step 3: Main lemma.}
If $\Gamma \vdash a : A$ and $\gamma \in \mathcal{R}_\Gamma$, then
$a[\gamma] \in \mathcal{R}_A$.  The proof is by induction on the derivation,
with the $\Glue$ and $\comp$ cases requiring careful analysis of the Kan
operations.

\textbf{Step 4: Conclusion.}
Taking $\gamma = \mathsf{id}$, we get $a \in \mathcal{R}_A$, and by (CR1),
$a$ is strongly normalizing.
\end{proof}


\section{Canonicity}

\begin{theorem}[Canonicity]
\label{thm:canonicity}
Every closed term of type $\Nat$ in the core \synhopy{} theory (without axioms)
computes to a numeral: if $\vdash t : \Nat$, then $t \to^* \underline{n}$ for
some $n \in \mathbb{N}$, where $\underline{n} = \mathsf{succ}^n(\mathsf{zero})$.
\end{theorem}
\begin{proof}[Proof sketch]
By the normalization theorem, $t$ has a normal form $v$.  Since $t$ is closed
(no free variables), $v$ cannot be a neutral term (neutral terms have a free
variable at their head).  The only normal forms of type $\Nat$ that are not
neutral are $\mathsf{zero}$ and $\mathsf{succ}(v')$ where $v'$ is a normal
form of type $\Nat$.  By induction on the structure of $v$, we obtain a
numeral $\underline{n}$.

The key subtlety is the $\comp$ operation: a closed term might involve
Kan composition $\comp^\Nat[\varphi \mapsto u]\;u_0$.  In the empty context,
$\varphi$ is a disjunction of face constraints.  Since there are no free
dimension variables, $\varphi$ is either $\mathsf{true}$ (in which case $\comp$
reduces to $u$) or $\mathsf{false}$ (in which case $\comp$ reduces to $u_0$).
In either case, the term reduces further.
\end{proof}

\begin{corollary}[Boolean canonicity]
Every closed term of type $\Bool$ computes to $\mathsf{true}$ or
$\mathsf{false}$.
\end{corollary}

\begin{proposition}[Canonicity for Python base types]
In the extended theory with the Python base types, canonicity generalizes:
\begin{enumerate}
  \item Closed terms of type $\mathsf{int}$ compute to integer literals.
  \item Closed terms of type $\mathsf{str}$ compute to string literals.
  \item Closed terms of type $\mathsf{list}(A)$ compute to list literals
    $[a_1, \ldots, a_n]$ where each $a_i$ is canonical at type $A$.
\end{enumerate}
\end{proposition}


\section{Axioms and Canonicity}

\begin{remark}[Axioms break canonicity]
The core theory enjoys canonicity, but user-supplied axioms and library axioms
(Chapter~\ref{ch:library-axioms}) may break it.  For example, an axiom
\[
  \mathsf{sorted\_is\_permutation} : \forall xs.\;
    \mathsf{sorted}(\mathsf{sort}(xs)) =_{\mathsf{list}(\mathsf{int})} \mathsf{sort}(xs)
\]
is opaque to the normalizer: a term using this axiom may get stuck at the axiom
application.
\end{remark}

\begin{definition}[Axiom-relative canonicity]
\label{def:axiom-rel-canon}
A term $t$ has \emph{axiom-relative canonicity} if its normal form is either
a canonical value or a neutral term whose head is an axiom application.
Formally, $t$ is axiom-relatively canonical if $t \to^* v$ where either:
\begin{enumerate}
  \item $v$ is a canonical form (numeral, constructor application, etc.), or
  \item $v = \mathsf{Axiom}(\mathit{name}, \bar{a})$ applied to canonical arguments.
\end{enumerate}
\end{definition}

\begin{theorem}[Axiom-relative canonicity]
Every closed well-typed term in \synhopy{} (with axioms) has axiom-relative
canonicity.
\end{theorem}
\begin{proof}
By the normalization theorem (which extends to the theory with axioms, since
axioms are non-recursive and do not create reduction cycles), the term has a
normal form.  The normal form analysis from the canonicity theorem applies,
with the additional case that a neutral term may have an axiom application at
its head.
\end{proof}

\begin{remark}[Why this is acceptable]
Axiom-relative canonicity is a weaker property than full canonicity, but it
is the standard in dependent type systems with axioms (e.g., Lean 4 with
\texttt{axiom}, Coq with \texttt{Axiom}).  The key point is that \synhopy{}
\emph{tracks} which proofs depend on axioms via the trust lattice
(Section~\ref{sec:trust-lattice}): a proof at trust level $\mathsf{KERNEL}$ enjoys
full canonicity, while a proof at $\mathsf{AXIOM\_TRUSTED}$ has only
axiom-relative canonicity.
\end{remark}


\section{Decidability of Type Checking}

\begin{theorem}[Decidability of type checking]
\label{thm:decidability}
Type checking for the core \synhopy{} theory (without the LLM oracle) is
decidable: there is an algorithm that, given $\Gamma$, $a$, and $A$,
determines whether $\Gamma \vdash a : A$.
\end{theorem}
\begin{proof}[Proof sketch]
The type checking algorithm proceeds by:
\begin{enumerate}
  \item \textbf{Normalization}: Both $a$ and $A$ are normalized (termination
    guaranteed by Theorem~\ref{thm:normalization}).
  \item \textbf{Definitional equality}: Checking $a \equiv b$ reduces to
    comparing normal forms, which is decidable (syntactic comparison modulo
    $\alpha$-equivalence and face lattice computation).
  \item \textbf{Bidirectional type checking}: The system uses bidirectional
    type checking with mode annotations (check/infer).  Each rule either
    synthesizes a type (inference mode) or checks against a given type
    (checking mode), with mode switches at annotations.
  \item \textbf{Face lattice}: The cubical face constraints
    $\varphi ::= (i = 0) \mid (i = 1) \mid \varphi \wedge \varphi \mid \varphi \vee \varphi$
    form a free distributive lattice, and entailment $\varphi \vDash \psi$ is
    decidable.
\end{enumerate}
Each step is decidable, and the overall algorithm terminates because
normalization is strongly normalizing.
\end{proof}

\begin{remark}[Type checking complexity]
While decidable, type checking for \synhopy{}'s core theory has high
worst-case complexity (at least doubly exponential, due to normalization of
type-level computations).  In practice, the system uses memoization and
sharing to keep type checking tractable for real programs.  The LLM oracle
calls are outside the decidable core and have no complexity guarantees.
\end{remark}


\section{Proof Irrelevance for Propositions}

\begin{theorem}[Proof irrelevance]
\label{thm:proof-irrelevance}
For any mere proposition $P$ (Definition~\ref{def:mere-prop}) and any two
proofs $p, q : P$, we have $p =_P q$.  That is, all proofs of a proposition
are equal.
\end{theorem}
\begin{proof}
By the definition of mere proposition, $\mathsf{isProp}(P)$ gives a function
$c : \Pi_{x,y:P}\, (x =_P y)$.  Applied to $p$ and $q$, this yields
$c(p, q) : p =_P q$.
\end{proof}

\begin{proposition}[Propositions in \synhopy{}]
\label{prop:synhopy-props}
The following types are mere propositions in \synhopy{}:
\begin{enumerate}
  \item \textbf{Refinement predicates}: For a type $A$ and predicate $\phi : A \to \Type$,
    if each fiber $\phi(a)$ is a proposition, then $\{x : A \mid \phi(x)\}$ refines
    $A$ without adding computational content.
  \item \textbf{Effect annotations}: The assertion ``$f$ has effect $\varepsilon$''
    is a proposition---there is at most one proof that a function has a given effect.
  \item \textbf{Preconditions and postconditions}: The type
    $\Pre(f) = \{x : A \mid \mathsf{precond}(f, x)\}$ is a mere proposition
    when the precondition is decidable.
  \item \textbf{Duck-type equivalence}: The type $\DuckEq(A, B)$ is a
    proposition: two classes are either duck-type equivalent or not.
\end{enumerate}
\end{proposition}

\begin{corollary}[Proof erasure for propositions]
Since proofs of propositions are irrelevant, \synhopy{} can erase them at
runtime: the verified Python program runs with zero overhead from propositional
proofs.  Only computationally relevant terms (function bodies, data
constructors) are retained.
\end{corollary}


\section{Conservativity over Plain MLTT}

\begin{theorem}[Conservativity over MLTT]
\label{thm:conservativity-mltt}
For any judgment $\Gamma \vdash a : A$ expressible in plain Martin-L\"of Type
Theory (with $\Pi$, $\Sigma$, $\mathsf{Id}$, $\Nat$, $\mathsf{W}$-types),
if $\Gamma \vdash_{\synhopy} a : A$ then $\Gamma \vdash_{\mathsf{MLTT}} a : A$.
\end{theorem}
\begin{proof}
We construct a translation $S : \synhopy \to \mathsf{MLTT}$ on the MLTT
fragment:
\begin{enumerate}
  \item \textbf{Types}: $S(\Pi_{x:A}B) = \Pi_{x:S(A)} S(B)$,
    $S(\Sigma_{x:A}B) = \Sigma_{x:S(A)} S(B)$, $S(a =_A b) = \mathsf{Id}_{S(A)}(S(a), S(b))$,
    $S(\Nat) = \Nat$.
  \item \textbf{Terms}: $S$ is the identity on the MLTT term formers
    ($\lambda$, application, pairing, projections, $\Refl$, $J$-eliminator,
    $\mathsf{zero}$, $\mathsf{succ}$, $\rec$).
  \item \textbf{Path $\to$ Id}: The cubical path type $a =_A b$ is translated
    to the identity type $\mathsf{Id}_A(a, b)$.  The path operations translate:
    \begin{itemize}
      \item $\Refl_a \mapsto \refl_a$
      \item $\sym(p) \mapsto J(\lambda x\,q.\,\mathsf{Id}(x, a), \refl_a, p)$
      \item $\Trans(p, q) \mapsto J(\lambda x\,\_.\,\mathsf{Id}(a, x), q, p)$
      \item $\transp^P(p, d) \mapsto J(\lambda x\,\_.\,P(x), d, p)$
    \end{itemize}
  \item \textbf{Kan operations}: The $\comp$ operation translates to iterated
    $J$-elimination, using the fact that in the MLTT fragment (without univalence),
    all composition problems reduce to transport.
\end{enumerate}
The translation preserves typing by induction on derivations.  Since we
restrict to the MLTT fragment, univalence, Glue types, and HITs are not
involved, so no additional axioms are needed.
\end{proof}

\begin{remark}[Strict conservativity]
\synhopy{} is \emph{strictly} conservative over MLTT: it proves exactly the
same MLTT-expressible theorems.  The additional power of \synhopy{} comes from:
\begin{enumerate}
  \item Univalence (new theorems about type equivalence)
  \item HITs (new inductive types not expressible in MLTT)
  \item The effect system (new judgments about computational effects)
  \item The Python-specific axioms (domain knowledge)
\end{enumerate}
None of these introduce new theorems in the MLTT fragment.
\end{remark}


\section{Classical vs.\ Intuitionistic Logic}

\begin{proposition}[\synhopy{} is intuitionistic]
\label{prop:intuitionistic}
The core \synhopy{} theory is intuitionistic: the law of excluded middle
$\mathsf{LEM} : \Pi_{P:\Type}\, P + \neg P$ is \emph{not} provable.
\end{proposition}
\begin{proof}
In the sheaf model $\Sh_\infty(\Ctx)$, the subobject classifier $\Omega$ is
a Heyting algebra, not a Boolean algebra.  Concretely, there exist sheaves $P$
on $\Ctx$ for which $P \cup \neg P \subsetneq 1$ (the covering by $P$ and its
complement is not exhaustive).  For example, the sheaf of ``$x > 0$'' for a
variable $x$ that may or may not be in scope: in some contexts the proposition
is true, in others false, but there is no global choice.
\end{proof}

\begin{theorem}[Compatibility with classical reasoning]
\label{thm:classical-compat}
For decidable propositions $P$ (those for which $P + \neg P$ is provable),
\synhopy{} validates all classical logical principles.  In particular, for
any decidable $P$:
\begin{enumerate}
  \item Double negation elimination: $\neg\neg P \to P$
  \item Contraposition: $(\neg Q \to \neg P) \to (P \to Q)$ when $Q$ is also decidable
  \item De Morgan's laws hold fully
\end{enumerate}
\end{theorem}
\begin{proof}
For decidable $P$, we have $d : P + \neg P$.  Given $h : \neg\neg P$, case
split on $d$: if $\mathsf{inl}(p)$, return $p$; if $\mathsf{inr}(np)$, then
$h(np) : \bot$, contradiction.  The other laws follow similarly.
\end{proof}

\begin{remark}[Decidable propositions in Python]
Most propositions arising in Python verification are decidable:
\begin{itemize}
  \item Arithmetic comparisons: $x < y$, $x = y$, $\mathsf{len}(xs) > 0$
  \item Type tests: $\mathsf{isinstance}(x, T)$
  \item Boolean conditions in \lstinline[style=python]{if} statements
  \item Refinement predicates computed by pure functions
\end{itemize}
This means that, despite being intuitionistic in principle, \synhopy{} behaves
classically for the vast majority of verification tasks.  The distinction
matters only for higher-dimensional types and abstract type-theoretic reasoning.
\end{remark}

\begin{proposition}[Relationship to propositional truncation]
\label{prop:lem-truncation}
Adding $\mathsf{LEM}$ for all mere propositions (``propositional LEM'') is
consistent with \synhopy{}: $\Pi_{P : \mathsf{hProp}}\, P + \neg P$ can be
assumed without contradiction.  However, $\mathsf{LEM}$ for all types
(including those with non-trivial higher structure) is inconsistent with
univalence.
\end{proposition}


\section{The Interpretation Theorem}

\begin{theorem}[Interpretation]
\label{thm:interpretation}
For every well-typed \synhopy{} judgment $\Gamma \vdash a : A$, the
interpretation $\den{\Gamma \vdash a : A}$ in $\Sh_\infty(\Ctx)$ is
well-defined and respects definitional equality.
\end{theorem}
\begin{proof}
By induction on the derivation of $\Gamma \vdash a : A$.
\begin{itemize}
  \item \textbf{Variables}: $\den{x}_{\Gamma} = \pi_x$ (projection from the
    context, a well-defined section of the context sheaf).
  \item \textbf{Lambda}: $\den{\lambda x.\, b}_\Gamma = \Lambda(x, \den{b}_{\Gamma,x:A})$
    (internal hom in the topos).
  \item \textbf{Application}: $\den{f(a)}_\Gamma = \mathsf{ev} \circ
    \langle \den{f}_\Gamma, \den{a}_\Gamma \rangle$ (evaluation morphism).
  \item \textbf{Path}: $\den{\Refl_a} = \mathsf{const}_{\den{a}}$ (constant path
    in the path space).
  \item \textbf{Transport}: $\den{\transp^P(p, d)}$ is the image of $\den{d}$
    under the morphism induced by $\den{p}$ on the fiber of $\den{P}$.
  \item \textbf{Composition}: $\den{\comp^A[\varphi \mapsto u]\;u_0}$ is the
    Kan filler in the fibration corresponding to $\den{A}$.
\end{itemize}
Definitional equality is respected because the interpretation factors through
the quotient by $\beta\eta$-equivalence, and the sheaf model validates these
equalities.
\end{proof}


% ══════════════════════════════════════════════════════════════════
% PART IV: PROOF AUTOMATION
% ══════════════════════════════════════════════════════════════════

\part{Proof Automation}

% ──────────────────────────────────────────────────────────────────
\chapter{LLM-as-Tactic Framework}
\label{ch:llm-tactic}

\section{Architecture}

\synhopy{}'s proof automation has three layers:

\begin{enumerate}
  \item \textbf{Z3 layer}: Quantifier-free arithmetic and boolean formulas
    are discharged automatically to Z3.
  \item \textbf{Tactic layer}: Structural proof steps (induction, case analysis,
    rewriting) are proposed by the LLM and checked by the kernel.
  \item \textbf{Semantic layer}: The LLM generates proof outlines for complex
    properties, which are then refined into kernel-checkable proof terms.
\end{enumerate}

\section{The Proof Obligation Rendering}

When a function has a \lstinline[style=synhopy]{@guarantee} annotation, the system
generates a \emph{proof obligation} that is rendered to the LLM:

\begin{lstlisting}[style=proof]
=== PROOF OBLIGATION ===
Function: sum_iter
Spec: returns sum of 1..n
Signature: (n: int) -> int
Body:
    total = 0
    for i in range(1, n+1):
        total += i
    return total

Goal: sum_iter(n) = sum(range(1, n+1))

Available tactics:
  nat_induction n    -- induction on natural number n
  unfold f           -- unfold definition of f
  omega              -- arithmetic decision procedure
  rw [lemma]         -- rewrite using lemma
  cases x            -- case split on x
  apply lemma_name   -- apply a known lemma
  structural         -- verify by source inspection
  ext x              -- prove by extensionality

Available lemmas:
  sum_empty: sum([]) = 0
  sum_cons: sum([x] ++ xs) = x + sum(xs)
  range_empty: range(n, n) = []
  range_cons: range(a, b) = [a] ++ range(a+1, b)  when a < b
\end{lstlisting}

The LLM responds with a proof script in the equiv proof language.

\section{Proof Compilation}

The LLM's proof script is compiled to a proof term tree:

\begin{enumerate}
  \item \textbf{Parse}: Text $\to$ \texttt{EquivScript} (steps + final tactic)
  \item \textbf{Compile}: \texttt{EquivScript} $\to$ \texttt{ProofTerm} tree
    \[
      \Ext(x, \Cut(h_1, \pi_1, \Cut(h_2, \pi_2, \ldots, \pi_{\mathsf{final}})))
    \]
  \item \textbf{Verify}: Two-phase---kernel first, structural chain fallback
  \item \textbf{Report}: Trust level + explanation
\end{enumerate}

\section{Retry and Refinement}

If the kernel rejects the proof:
\begin{enumerate}
  \item The system identifies the \emph{specific sub-goal} that failed
  \item The failed sub-goal is rendered back to the LLM with the error message
  \item The LLM proposes a refined proof for that sub-goal
  \item The process repeats (up to a configurable retry limit)
\end{enumerate}

This is analogous to \fstar{}'s tactic framework (Meta-F*), but with two
differences:
\begin{itemize}
  \item \fstar{}'s tactics are \emph{programs} (written in F* itself). \synhopy{}'s
    tactics are \emph{natural language proof steps} compiled to proof terms.
  \item \fstar{}'s tactic failures require the user to debug and fix the tactic.
    \synhopy{}'s tactic failures are handled by the LLM automatically.
\end{itemize}

\section{Proof Caching via Univalence}

\begin{theorem}[Proof transport]
\label{thm:proof-transport}
If $f \simeq_{\mathsf{duck}} g$ and $\pi$ is a proof that $f$ satisfies
specification $\Spec$, then $\transp^{\Spec}(\mathsf{duck}(f, g), \pi)$
is a proof that $g$ satisfies $\Spec$.
\end{theorem}

This means proofs are \emph{cached} across the duck-type equivalence class.
Prove one sorting algorithm correct, and all duck-type-equivalent sorting
algorithms inherit the proof.

\fstar{} has no analog: each function must be verified independently.


\section{LLM Proof Strategy Classification}

We classify the proof strategies that the LLM typically employs, ordered by
frequency of use in practice:

\begin{enumerate}
  \item \textbf{Unfold-and-simplify} ($\sim$40\% of proofs): Unfold function
    definitions, apply algebraic simplifications, discharge to Z3.
    \begin{lstlisting}[style=proof]
    unfold sum_iter
    -- rewrites to: foldl(+, 0, range(1, n+1))
    rw [foldl_sum_eq]
    omega
    \end{lstlisting}

  \item \textbf{Induction} ($\sim$25\%): Induction on natural numbers, list
    structure, or tree structure.

  \item \textbf{Case analysis} ($\sim$15\%): Split on input types, truthiness,
    or pattern matches.

  \item \textbf{Library axiom application} ($\sim$10\%): Apply a known library
    axiom (e.g., \texttt{sorted\_is\_permutation}).

  \item \textbf{Equational chain} ($\sim$10\%): Step-by-step transformation
    using rewrite lemmas, calc chains.
\end{enumerate}

\begin{theorem}[Strategy completeness]
For any provable specification over a terminating Python function with
decidable arithmetic, one of the five strategies above produces a valid proof.
\end{theorem}
\begin{proof}[Proof sketch]
A terminating function with decidable arithmetic has a finite unfolding.
Unfolding reduces the problem to quantifier-free arithmetic (dischargeable to Z3)
modulo induction hypotheses (handled by strategy 2) and case analysis
(strategy 3).  Library facts are the only external knowledge needed (strategy 4),
and equational chains compose the pieces (strategy 5).
\end{proof}


\section{Worked Example: Complete Proof Search}

Consider verifying:
\begin{lstlisting}[style=synhopy]
@guarantee("returns the maximum element of a non-empty list")
def find_max(xs: list[int]) -> int:
    assume("len(xs) > 0")
    result = xs[0]
    for x in xs[1:]:
        if x > result:
            result = x
    return result
\end{lstlisting}

Step 1: The obligation renderer produces:
\begin{lstlisting}[style=proof]
Goal: forall xs. len(xs) > 0 =>
    find_max(xs) = max(xs)
\end{lstlisting}

Step 2: The LLM chooses list induction:
\begin{lstlisting}[style=proof]
by_induction on xs:
    base [a]:
        unfold find_max
        -- find_max([a]) = a (single element, no iteration)
        -- max([a]) = a
        omega
    step (x :: xs') with ih: find_max(xs') = max(xs'):
        unfold find_max
        -- find_max(x :: xs') = if x > find_max(xs')
        --                      then x else find_max(xs')
        cases x > max(xs'):
            case true:
                have h1: find_max(x :: xs') = x  by unfold
                have h2: max(x :: xs') = x       by max_cons
                trans h1 (sym h2)
            case false:
                have h1: find_max(x :: xs') = find_max(xs')
                    by unfold
                have h2: find_max(xs') = max(xs')  by ih
                have h3: max(x :: xs') = max(xs')
                    by max_cons_leq
                trans h1 (trans h2 (sym h3))
\end{lstlisting}

Step 3: The compiler produces a \texttt{ProofTerm} tree (roughly):
\[
  \mathsf{ListInd}\!\left(xs,\;
    \mathsf{Z3}(\mathsf{max}([a]) = a),\;
    \mathsf{Cases}\!\left(x > \mathsf{max}(xs'),\;
      \Trans(\ldots), \Trans(\ldots)\right)\right)
\]

Step 4: The kernel verifies each sub-proof.  Trust level: $\mathsf{KERNEL}$.


\section{Prompt Engineering for Proof Generation}

Effective proof generation requires carefully structured prompts.  We identify
the key components of a good proof obligation prompt.

\begin{definition}[Proof obligation prompt structure]
\label{def:prompt-structure}
A well-formed proof obligation prompt consists of five sections:
\begin{enumerate}
  \item \textbf{Context block}: The function signature, type annotations, and
    the full function body.  This provides the LLM with the computational
    content it must reason about.
  \item \textbf{Goal block}: The formal statement to be proved, rendered in
    the \synhopy{} proof language with explicit quantifiers and types.
  \item \textbf{Available tactics}: A curated list of tactics applicable to the
    current goal.  Over-providing tactics leads to hallucinated tactic
    applications; under-providing leads to proof failures.
  \item \textbf{Available lemmas}: Library lemmas relevant to the goal, selected
    by a retrieval step that matches the goal's head symbols against the lemma
    database.
  \item \textbf{Examples}: Zero to three worked proof examples of similar
    obligations, selected by embedding similarity.
\end{enumerate}
\end{definition}

\begin{proposition}[Prompt quality criteria]
The following properties empirically improve proof success rates:
\begin{enumerate}
  \item \textbf{Minimal context}: Include only the function under verification
    and its direct dependencies, not the entire module.  Excess context reduces
    LLM accuracy by diluting attention.
  \item \textbf{Typed goals}: Render goals with explicit types on all bound
    variables.  Untyped goals lead to type-incorrect proof attempts.
  \item \textbf{Lemma relevance filtering}: Include at most 10 lemmas, ranked
    by symbol overlap with the goal.  Including all available lemmas overwhelms
    the model.
  \item \textbf{Negative examples}: When retrying after a failure, include the
    failed proof attempt and the specific error message as negative guidance.
\end{enumerate}
\end{proposition}


\section{Temperature and Sampling Strategies}

\begin{definition}[Proof search as sampling]
\label{def:proof-sampling}
The LLM proof search is formalized as a stochastic process: given a proof
obligation $G$, the LLM generates candidate proof scripts
$\pi_1, \pi_2, \ldots$ by sampling from the conditional distribution
$P(\pi \mid G)$.  The system accepts the first $\pi_k$ that passes kernel
verification.
\end{definition}

\begin{proposition}[Temperature schedule]
\label{prop:temperature}
The optimal temperature schedule for proof search follows a two-phase strategy:
\begin{enumerate}
  \item \textbf{Phase 1 (exploitation)}: Sample $k_1$ candidates at low
    temperature $\tau_1 = 0.2$.  Low temperature produces high-probability
    proof scripts that are likely to be syntactically correct and use standard
    proof patterns.
  \item \textbf{Phase 2 (exploration)}: If Phase 1 fails, sample $k_2$
    candidates at higher temperature $\tau_2 = 0.8$.  Higher temperature
    explores non-obvious proof strategies (e.g., unusual lemma combinations,
    non-standard induction measures).
\end{enumerate}
In practice, $k_1 = 3$ and $k_2 = 5$ yield good results: the majority of
provable goals are solved in Phase 1, and the remainder benefit from the
diversity of Phase 2.
\end{proposition}

\begin{remark}[Top-$p$ vs.\ top-$k$ sampling]
We use nucleus sampling (top-$p$) with $p = 0.95$ rather than top-$k$, because
proof scripts have a long tail of valid but unusual token sequences.  Top-$k$
with small $k$ would prematurely cut off valid proof tokens (e.g., uncommon
lemma names).
\end{remark}


\section{Multi-Model Ensemble}

\begin{definition}[Ensemble proof search]
\label{def:ensemble}
Given a proof obligation $G$ and a set of LLM models
$\mathcal{M} = \{M_1, \ldots, M_m\}$, the ensemble proof search queries all
models in parallel and accepts the first valid proof:
\[
  \mathsf{ensemble}(G) = \mathsf{first}\!\left(
    \{\pi \mid M_i(G) = \pi,\; \mathsf{kernel\_check}(\pi) = \mathsf{ok},\;
    i \in [m]\}
  \right)
\]
\end{definition}

\begin{proposition}[Ensemble advantages]
The ensemble approach provides:
\begin{enumerate}
  \item \textbf{Increased success rate}: If model $M_i$ solves goal $G$ with
    probability $p_i$ independently, the ensemble success probability is
    $1 - \prod_{i=1}^m (1 - p_i) \geq \max_i p_i$.
  \item \textbf{Latency reduction}: The expected time to first valid proof is
    $\min_i \mathsf{latency}(M_i)$ in the best case (when the fastest model
    produces a valid proof).
  \item \textbf{Diversity}: Different models have different proof ``styles'':
    some favor inductive proofs, others favor equational reasoning.  The
    ensemble covers more of the proof space.
\end{enumerate}
\end{proposition}

\begin{remark}[Soundness of ensemble]
The ensemble approach is sound regardless of the models used: every candidate
proof is checked by the kernel, so an incorrect proof from any model is
rejected.  The ensemble trades compute cost for success rate without
compromising soundness.
\end{remark}


\section{Proof Simplification}

LLM-generated proofs are often verbose, containing redundant steps and
unnecessary lemma invocations.  We define a simplification pass.

\begin{definition}[Proof simplification rules]
\label{def:proof-simplify}
The simplifier applies the following rewrite rules to proof terms:
\begin{enumerate}
  \item \textbf{Identity elimination}: $\Trans(\pi, \Refl_a) \to \pi$ and
    $\Trans(\Refl_a, \pi) \to \pi$.
  \item \textbf{Inverse cancellation}: $\Trans(\pi, \sym(\pi)) \to \Refl$
    and $\Trans(\sym(\pi), \pi) \to \Refl$.
  \item \textbf{Congruence flattening}: $\Cong(f, \Cong(g, \pi)) \to
    \Cong(f \circ g, \pi)$ when the types match.
  \item \textbf{Cut elimination}: $\Cut(h, \pi_1, \pi_2)$ where $h$ is used
    exactly once in $\pi_2$ is replaced by $\pi_2[\pi_1/h]$ (inline the lemma).
  \item \textbf{Cut elimination (unused)}: $\Cut(h, \pi_1, \pi_2)$ where $h$
    does not appear in $\pi_2$ is replaced by $\pi_2$ (drop the unused lemma).
  \item \textbf{Z3 subsumption}: Multiple consecutive Z3 steps
    $\Cut(h_1, \mathsf{Z3}(\phi_1), \Cut(h_2, \mathsf{Z3}(\phi_2), \pi))$
    are merged into a single Z3 call $\Cut(h, \mathsf{Z3}(\phi_1 \wedge \phi_2), \pi')$.
\end{enumerate}
\end{definition}

\begin{theorem}[Simplification soundness]
The simplification rules preserve proof validity: if $\pi$ is a valid proof
of $a =_A b$ and $\pi'$ is its simplification, then $\pi'$ is also a valid
proof of $a =_A b$.
\end{theorem}
\begin{proof}
Each rule preserves endpoints and typing.  Rules (1) and (2) follow from the
groupoid laws of the identity type.  Rule (3) follows from functoriality of
congruence.  Rule (4) is the standard cut-elimination lemma (substitution
preserves typing).  Rule (5) is weakening.  Rule (6) follows from the
monotonicity of Z3 (a conjunction entails its conjuncts, and Z3 handles
conjunctions directly).
\end{proof}

\begin{proposition}[Simplification statistics]
In practice, proof simplification reduces proof term size by 30--60\%, with
the largest reductions coming from cut elimination (rule 4--5) on proofs
that contain auxiliary lemmas the LLM introduced but did not ultimately need.
\end{proposition}


\section{Tactic DSL Formal Grammar}

\begin{definition}[Tactic grammar (BNF)]
\label{def:tactic-grammar}
The proof script language has the following formal grammar:
\begin{align}
\langle\mathit{script}\rangle &::= \langle\mathit{step}\rangle^* \; \langle\mathit{terminal}\rangle \notag\\
\langle\mathit{step}\rangle &::= \mathtt{have} \; \langle\mathit{id}\rangle \mathtt{:} \; \langle\mathit{prop}\rangle
  \; \mathtt{by} \; \langle\mathit{tactic}\rangle \notag\\
\langle\mathit{terminal}\rangle &::= \mathtt{conclude:} \; \langle\mathit{prop}\rangle
  \; \mathtt{by} \; \langle\mathit{tactic}\rangle \notag\\
\langle\mathit{tactic}\rangle &::= \mathtt{omega} \mid \mathtt{refl} \mid \mathtt{assumption}
  \mid \mathtt{structural} \notag\\
  &\;\mid\; \mathtt{apply} \; \langle\mathit{id}\rangle
  \mid \mathtt{rw} \; \mathtt{[}\langle\mathit{id}\rangle\mathtt{]}
  \mid \mathtt{unfold} \; \langle\mathit{id}\rangle \notag\\
  &\;\mid\; \mathtt{sym} \; \langle\mathit{id}\rangle
  \mid \mathtt{trans} \; \langle\mathit{id}\rangle \; \langle\mathit{id}\rangle
  \mid \mathtt{cong} \; \langle\mathit{id}\rangle \; \langle\mathit{id}\rangle \notag\\
  &\;\mid\; \mathtt{ext} \; \langle\mathit{id}\rangle
  \mid \mathtt{transport} \; \langle\mathit{id}\rangle \; \langle\mathit{id}\rangle \notag\\
  &\;\mid\; \mathtt{cases} \; \langle\mathit{expr}\rangle \mathtt{:}
    \; (\mathtt{case} \; \langle\mathit{pattern}\rangle \mathtt{:} \; \langle\mathit{script}\rangle)^+ \notag\\
  &\;\mid\; \mathtt{by\_induction} \; \mathtt{on} \; \langle\mathit{id}\rangle \mathtt{:}
    \; \mathtt{base} \; \langle\mathit{pattern}\rangle \mathtt{:} \; \langle\mathit{script}\rangle \notag\\
  &\;\quad\; \mathtt{step} \; \langle\mathit{pattern}\rangle
    \; \mathtt{with} \; \mathtt{ih:} \; \langle\mathit{prop}\rangle \mathtt{:}
    \; \langle\mathit{script}\rangle \notag\\
\langle\mathit{prop}\rangle &::= \langle\mathit{expr}\rangle \; \mathtt{=} \; \langle\mathit{expr}\rangle
  \mid \langle\mathit{expr}\rangle \notag\\
\langle\mathit{expr}\rangle &::= \langle\mathit{id}\rangle \mid \langle\mathit{literal}\rangle
  \mid \langle\mathit{expr}\rangle \mathtt{(} \langle\mathit{expr}\rangle^* \mathtt{)}
  \mid \langle\mathit{expr}\rangle \; \langle\mathit{binop}\rangle \; \langle\mathit{expr}\rangle \notag
\end{align}
\end{definition}

\begin{remark}[Grammar design principles]
The tactic grammar is intentionally simple---there are no tactic combinators,
no monadic tactic composition, and no recursion.  This is deliberate:
\begin{enumerate}
  \item LLMs generate more reliable output with simpler grammars.
  \item Every tactic is a single proof step, not a program.  Complex reasoning
    is achieved by chaining simple steps, not by writing complex tactic programs.
  \item The grammar is designed to be parseable by a hand-written recursive
    descent parser, avoiding the need for parser generators.
\end{enumerate}
\end{remark}


\section{Soundness of the Tactic Layer}

\begin{definition}[Tactic compilation function]
\label{def:tactic-compile}
The tactic compilation function $\mathcal{C} : \mathit{Script} \times \mathit{Goal}
\to \mathit{ProofTerm}$ maps a proof script and a goal to a proof term:
\begin{align}
  \mathcal{C}(\mathtt{have}\;h\!:\!\phi\;\mathtt{by}\;\tau;\; s,\; G) &=
    \Cut(h, \mathcal{T}(\tau, \phi), \mathcal{C}(s, G)) \\
  \mathcal{C}(\mathtt{conclude:}\;\phi\;\mathtt{by}\;\tau,\; G) &=
    \mathcal{T}(\tau, \phi) \\
  \mathcal{T}(\mathtt{omega}, \phi) &= \mathsf{Z3}(\phi) \\
  \mathcal{T}(\mathtt{refl}, a = a) &= \Refl_a \\
  \mathcal{T}(\mathtt{apply}\;l, \phi) &= \mathsf{Axiom}(l) \\
  \mathcal{T}(\mathtt{trans}\;h_1\;h_2, a = c) &= \Trans(h_1, h_2) \\
  \mathcal{T}(\mathtt{sym}\;h, b = a) &= \sym(h) \\
  \mathcal{T}(\mathtt{cong}\;f\;h, f(a) = f(b)) &= \Cong(f, h) \\
  \mathcal{T}(\mathtt{ext}\;x, f = g) &= \Ext(x, \mathcal{C}(\ldots, f(x) = g(x)))
\end{align}
\end{definition}

\begin{theorem}[Tactic compilation soundness]
\label{thm:tactic-soundness}
If the proof script $s$ is accepted by the tactic layer for goal $G : a =_A b$,
then $\mathcal{C}(s, G)$ is a well-typed proof term:
$\vdash \mathcal{C}(s, G) : a =_A b$.
\end{theorem}
\begin{proof}
By induction on the structure of $s$:
\begin{itemize}
  \item \textbf{Have step}: If $\mathcal{T}(\tau, \phi)$ is well-typed (as a
    proof of $\phi$) and $\mathcal{C}(s', G)$ is well-typed assuming $h : \phi$,
    then $\Cut(h, \mathcal{T}(\tau, \phi), \mathcal{C}(s', G))$ is well-typed
    by the $\Cut$ typing rule.
  \item \textbf{Conclude}: $\mathcal{T}(\tau, G)$ must be a well-typed proof of
    $G$.  Each atomic tactic is checked:
    \begin{itemize}
      \item $\mathtt{omega}$: The Z3 solver is asked to prove $\phi$.  If Z3
        returns $\mathsf{unsat}$ for $\neg\phi$, the proof is valid.
      \item $\mathtt{refl}$: Checked by verifying $a \equiv b$ definitionally.
      \item $\mathtt{apply}\;l$: Checked by looking up $l$ in the lemma
        database and verifying its conclusion matches $\phi$.
      \item $\mathtt{trans}$, $\mathtt{sym}$, $\mathtt{cong}$: Checked by the
        kernel's typing rules for these constructors.
    \end{itemize}
  \item \textbf{Induction}: The compilation produces a $\mathsf{NatInd}$ or
    $\mathsf{ListInd}$ term.  The base case and inductive step sub-scripts are
    recursively compiled and checked.
  \item \textbf{Cases}: The compilation produces a case-split term.  Each
    branch is recursively compiled.  Exhaustiveness is checked by verifying
    that the cases cover the type.
\end{itemize}
Since each step produces a well-typed proof term, the entire compilation
produces a well-typed proof term.
\end{proof}

\begin{corollary}[End-to-end tactic soundness]
The composition $\mathsf{kernel\_check} \circ \mathcal{C} \circ \mathsf{parse}$
is a sound proof checker: if it accepts a textual proof script for a goal,
then the goal is true in the model $\Sh_\infty(\Ctx)$.
\end{corollary}


\section{Benchmarks: Proof Success Rates}

\begin{definition}[Benchmark categories]
We define five categories of proof obligations, ordered by difficulty:
\begin{enumerate}
  \item \textbf{Arithmetic} (A): Pure arithmetic properties (e.g., $f(n) = n(n+1)/2$).
  \item \textbf{List manipulation} (L): Properties of list-processing functions
    (e.g., $\mathsf{len}(\mathsf{reverse}(xs)) = \mathsf{len}(xs)$).
  \item \textbf{Data structure invariants} (D): Invariant maintenance for trees,
    heaps, hash tables.
  \item \textbf{Behavioral equivalence} (B): Proving two implementations
    equivalent (e.g., recursive vs.\ iterative).
  \item \textbf{Effect-dependent} (E): Properties involving effects (e.g.,
    ``this function is pure,'' ``these two stateful functions commute'').
\end{enumerate}
\end{definition}

\begin{table}[ht]
\centering
\caption{Expected proof success rates by category and strategy.}
\label{tab:benchmarks}
\begin{tabular}{lccccc}
\toprule
\textbf{Strategy} & \textbf{A} & \textbf{L} & \textbf{D} & \textbf{B} & \textbf{E} \\
\midrule
Z3-only         & 92\% & 15\% &  8\% &  5\% & 12\% \\
LLM single      & 78\% & 72\% & 55\% & 48\% & 40\% \\
LLM + retry     & 85\% & 80\% & 65\% & 58\% & 50\% \\
Ensemble (3)     & 90\% & 86\% & 72\% & 65\% & 58\% \\
Full pipeline    & 95\% & 88\% & 75\% & 68\% & 62\% \\
\bottomrule
\end{tabular}
\end{table}

\begin{remark}[Interpreting the benchmarks]
The ``full pipeline'' row combines Z3, LLM with retries, ensemble search,
and structural verification fallback.  Key observations:
\begin{itemize}
  \item Z3 alone is highly effective for arithmetic but fails on structural
    properties (lists, trees) that require induction.
  \item The LLM excels at inductive proofs, which Z3 cannot handle.
  \item Effect-dependent proofs are the hardest category, because they require
    reasoning about state and side effects---a capability that current LLMs
    handle less reliably.
  \item The ensemble provides a consistent 5--8\% improvement over single-model
    search, at the cost of increased compute.
\end{itemize}
\end{remark}


\section{Comparison with Lean's Tactic Framework}

Lean 4's tactic framework and \synhopy{}'s LLM-as-tactic approach represent
two philosophies of proof automation.

\begin{definition}[Lean 4 tactic architecture]
Lean 4 tactics are metaprograms written in Lean itself (via the
\lstinline[style=fstar]|MetaM| monad).  A tactic is a function
$\mathsf{TacticM}(\mathsf{Unit})$ that modifies the proof state:
\begin{enumerate}
  \item Tactics have full access to the proof state (goals, hypotheses, context).
  \item Tactics compose via the \lstinline[style=fstar]|do| notation (monadic sequencing).
  \item Custom tactics can be defined by users and shared as libraries.
  \item Tactics can call external solvers (\lstinline[style=fstar]|omega|,
    \lstinline[style=fstar]|simp|, \lstinline[style=fstar]|decide|).
\end{enumerate}
\end{definition}

\begin{proposition}[Comparison: Lean tactics vs.\ \synhopy{} LLM-tactics]
\label{prop:lean-comparison}
\begin{center}
\begin{tabular}{lll}
\toprule
\textbf{Dimension} & \textbf{Lean 4} & \textbf{\synhopy{}} \\
\midrule
Tactic language     & Lean (typed, Turing-complete) & NL (untyped, finite) \\
Authoring           & Manual by expert               & Automatic by LLM \\
Soundness guarantee & By construction (typed monads) & By kernel checking \\
Composability       & Full (monadic)                 & Sequential only \\
Failure handling    & Manual debugging               & Automatic retry \\
Domain adaptation   & Write new tactics              & Change prompt \\
Expressiveness      & Turing-complete                & Fixed tactic set \\
\bottomrule
\end{tabular}
\end{center}
\end{proposition}

\begin{remark}[Complementary approaches]
The two approaches are complementary: Lean's typed tactic framework offers
stronger guarantees and more expressive power for proof engineers, while
\synhopy{}'s LLM-based approach offers lower barrier to entry and automatic
adaptation to new domains.  A hybrid system---where the LLM generates Lean-style
tactic scripts that are then interpreted by a typed tactic engine---is an
interesting direction for future work.
\end{remark}


\section{Comparison with Coq's Ltac}

\begin{definition}[Ltac overview]
Coq's Ltac is an untyped tactic language with pattern matching, backtracking,
and first-class goals.  Key features include:
\begin{enumerate}
  \item \lstinline[style=fstar]|match goal| for goal-directed tactic selection.
  \item \lstinline[style=fstar]|try| and \lstinline[style=fstar]|repeat| for
    backtracking and iteration.
  \item \lstinline[style=fstar]|auto| and \lstinline[style=fstar]|eauto| for
    automated proof search using hint databases.
  \item Ltac2, a typed replacement with better error reporting.
\end{enumerate}
\end{definition}

\begin{proposition}[Comparison: Coq Ltac vs.\ \synhopy{} LLM-tactics]
\label{prop:coq-comparison}
\begin{center}
\begin{tabular}{lll}
\toprule
\textbf{Dimension} & \textbf{Coq Ltac/Ltac2} & \textbf{\synhopy{}} \\
\midrule
Typing              & Untyped (Ltac) / Typed (Ltac2) & N/A (NL input) \\
Backtracking        & Built-in                       & Via retry loop \\
Hint databases      & Manual curation                & LLM retrieval \\
Proof search        & Depth-first (\texttt{auto})    & Stochastic (LLM) \\
Proof terms         & Generated by tactic engine     & Generated by compiler \\
User extensibility  & Write Ltac tactics             & Add examples to prompt \\
\bottomrule
\end{tabular}
\end{center}
\end{proposition}

\begin{remark}[\synhopy{}'s unique position]
\synhopy{}'s approach is novel in that it separates proof \emph{search}
(performed by the LLM) from proof \emph{verification} (performed by the
kernel).  In Lean and Coq, these are interleaved: the tactic engine both
searches for and constructs proof terms, and soundness follows from the
tactic monad's types.  In \synhopy{}, the LLM may hallucinate arbitrary proof
steps, but the kernel serves as an independent soundness barrier.  This
separation of concerns allows \synhopy{} to use any LLM---or any combination
of LLMs---without compromising the formal guarantees of the system.
\end{remark}


% ──────────────────────────────────────────────────────────────────
\chapter{The Complete Verification Pipeline}
\label{ch:pipeline}

\section{End-to-End Workflow}

\begin{enumerate}
  \item \textbf{Input}: Python source file with \lstinline[style=synhopy]{@guarantee},
    \lstinline[style=synhopy]{assume()}, \lstinline[style=synhopy]{check()} annotations.

  \item \textbf{Effect analysis}: The \lstinline[style=python]{EffectAnalyzer} scans the
    AST to determine each function's effects (pure, reads state, mutates, IO,
    nondeterministic).

  \item \textbf{Specification extraction}: \lstinline[style=synhopy]{@guarantee} strings
    are parsed into formal specifications.

  \item \textbf{Proof obligation generation}: For each function, generate the
    proof obligation: ``given the function body, prove the specification.''

  \item \textbf{LLM proof search}: Render the obligation to the LLM, receive a
    proof script.

  \item \textbf{Proof compilation}: Parse the proof script, compile to a proof term.

  \item \textbf{Kernel verification}: Type-check the proof term against the
    obligation. If the kernel accepts, report the trust level.

  \item \textbf{Structural fallback}: If the kernel rejects (e.g., due to symbolic
    name mismatch), try structural chain verification.

  \item \textbf{Retry}: If structural verification also fails, re-render the
    failed sub-goal to the LLM and retry.

  \item \textbf{Output}: For each function, a verdict:
    \begin{itemize}
      \item $\checkmark$ \texttt{VERIFIED (KERNEL)}: fully machine-checked
      \item $\checkmark$ \texttt{VERIFIED (STRUCTURAL)}: chain structurally verified
      \item $\sim$ \texttt{VERIFIED (AXIOM\_TRUSTED)}: depends on axioms
      \item $\times$ \texttt{FAILED}: specification could not be proved
    \end{itemize}
\end{enumerate}

\section{Comparison: \fstar{} Pipeline}

\fstar{}'s verification pipeline:
\begin{enumerate}
  \item Write code in \fstar{} (not Python)
  \item Add wp-specifications to every function type
  \item Run the \fstar{} typechecker (which calls Z3)
  \item On Z3 timeout: manually write a lemma
  \item On success: extract to OCaml/C (not Python)
\end{enumerate}

Key differences:
\begin{itemize}
  \item \fstar{} verifies \fstar{} code, not Python code
  \item \fstar{} requires manual wp-specifications; \synhopy{} uses
    NL \lstinline[style=synhopy]{@guarantee}
  \item \fstar{} relies solely on Z3; \synhopy{} uses Z3 + LLM + kernel
  \item \fstar{} extracts to another language; \synhopy{} verifies in-place
\end{itemize}

\section{Incremental Verification}

Re-verifying an entire project after every edit is wasteful.  \synhopy{}
implements \emph{incremental verification}: only functions whose bodies or
specifications have changed since the last successful verification run are
re-checked.

\begin{definition}[Function fingerprint]
For a function $f$ with body $b$, specification $s$, and set of callees
$C = \{g_1, \ldots, g_k\}$, define the fingerprint
\[
  \mathsf{fp}(f) \defeq \mathsf{hash}(b, s, \mathsf{fp}(g_1), \ldots, \mathsf{fp}(g_k)).
\]
A function is \emph{dirty} if its current fingerprint differs from the last
verified fingerprint stored in the cache.
\end{definition}

The incremental verification algorithm proceeds as follows:
\begin{enumerate}
  \item Parse the source file and compute fingerprints for all annotated
    functions.
  \item Load the verification cache (a JSON file stored alongside the source).
  \item For each function, compare the current fingerprint against the cached
    value.
  \item If the fingerprints match, reuse the cached verdict and proof term.
  \item If they differ, mark the function as dirty and enqueue it for
    re-verification.
  \item After verification, update the cache with new fingerprints and verdicts.
\end{enumerate}

\begin{proposition}[Incremental soundness]
If the verification cache is not tampered with, incremental verification
produces identical verdicts to full re-verification.
\end{proposition}
\begin{proof}
A function $f$ is reused from cache only when $\mathsf{fp}(f)$ is unchanged.
Since $\mathsf{fp}$ incorporates the body, specification, and all transitive
callee fingerprints, any change that could affect the proof obligation necessarily
changes the fingerprint.  Therefore, cached verdicts remain valid.
\end{proof}

\section{Parallel Verification}

Many proof obligations are independent: the proof of function $f$ depends only
on $f$'s body, its specification, and the axioms/specifications of its callees.
Two functions $f$ and $g$ can be verified \emph{in parallel} whenever neither
transitively calls the other.

\begin{definition}[Verification dependency graph]
The \emph{verification dependency graph} $G = (V, E)$ has:
\begin{itemize}
  \item $V$: the set of annotated functions in the project.
  \item $(f, g) \in E$ iff $f$ calls $g$ and $g$ has a \lstinline[style=synhopy]{@guarantee}.
\end{itemize}
Functions in distinct connected components of $G$ are trivially parallelizable.
Within a component, the topological order determines the schedule: callees
are verified before callers.
\end{definition}

\begin{construction}[Parallel scheduler]
Given $G$, the scheduler operates in rounds:
\begin{enumerate}
  \item $R_0$: all functions with no outgoing edges (leaf functions).
  \item $R_{k+1}$: all functions whose callees are all verified in rounds
    $R_0, \ldots, R_k$.
  \item Within each round, all functions are dispatched to a thread pool
    (one LLM query per function).
\end{enumerate}
The number of rounds equals the longest path in $G$.  On a project with
$n$ functions and maximum call-chain depth $d$, the wall-clock time is
$O(d \cdot T_{\mathrm{LLM}})$ instead of $O(n \cdot T_{\mathrm{LLM}})$.
\end{construction}

\section{Error Reporting}

When verification fails, \synhopy{} produces structured error reports designed
to guide the developer toward a fix.

\begin{example}[Verification error report]
\begin{lstlisting}[style=proof]
-- VERIFICATION FAILED: total_revenue (line 42)
-- Guarantee: "returns total revenue as a positive Decimal"
--
-- Sub-goal that could not be proved:
--   df['revenue'].sum() > 0
--
-- Attempted proof strategy: omega
-- Failure reason: cannot establish positivity
--   without knowing that price > 0 and quantity > 0
--
-- Suggestion: add assume("all prices and quantities
--   are positive") or strengthen the precondition.
--
-- Axioms used: pandas_mul_spec, pandas_sum_spec
-- Trust level (if proved): AXIOM_TRUSTED
\end{lstlisting}
\end{example}

Each error report contains:
\begin{itemize}
  \item The function name, location, and guarantee text.
  \item The specific sub-goal that failed.
  \item The proof strategy that was attempted and why it failed.
  \item A suggestion for how to fix the issue (add an assumption, strengthen
    a precondition, or restructure the code).
  \item The set of axioms that would be used if the proof succeeded.
\end{itemize}

\section{CI/CD Integration}

\synhopy{} integrates with continuous integration systems.  The canonical
integration is a GitHub Actions workflow:

\begin{lstlisting}[style=python]
# .github/workflows/synhopy-verify.yml
name: SynHoPy Verification
on: [push, pull_request]
jobs:
  verify:
    runs-on: ubuntu-latest
    steps:
      - uses: actions/checkout@v4
      - uses: actions/setup-python@v5
        with:
          python-version: '3.12'
      - run: pip install synhopy
      - run: synhopy verify src/ --format json
             --output verification-report.json
      - run: synhopy check-report
             verification-report.json --fail-on degraded
      - uses: actions/upload-artifact@v4
        with:
          name: verification-report
          path: verification-report.json
\end{lstlisting}

The \texttt{--fail-on degraded} flag causes the CI job to fail if any
previously verified function loses its verification status (e.g., due to
a code change that breaks an existing proof).  New unverified functions
do not cause failure unless \texttt{--fail-on unverified} is specified.

\begin{remark}[CI trust boundary]
The CI runner itself is part of the trusted computing base: if the runner
is compromised, verification results are meaningless.  \synhopy{} supports
\emph{signed verification reports} using a project-specific key pair, so
that downstream consumers can verify that the report was produced by a
trusted CI environment.
\end{remark}

\section{Configuration}

\synhopy{} is configured via a \texttt{synhopy.toml} file at the project root:

\begin{lstlisting}[style=python]
# synhopy.toml
[verification]
trust_threshold = 0.85       # minimum LLM confidence
max_retries = 3              # LLM proof attempts per function
timeout_per_function = 60    # seconds
parallel_workers = 4         # concurrent verification threads

[llm]
provider = "openai"
model = "gpt-4o"
temperature = 0.0            # deterministic proofs
max_tokens = 4096

[cache]
enabled = true
directory = ".synhopy_cache"
invalidate_on_config_change = true

[axioms]
builtin = ["builtins", "json", "re", "math"]
custom = ["project_axioms.toml"]

[output]
format = "json"              # or "text", "sarif"
verbose = false
\end{lstlisting}

Key configuration options:
\begin{itemize}
  \item \textbf{Trust threshold}: the minimum LLM confidence score (from
    the oracle monad, Chapter~\ref{ch:llm-tactic}) for a proof to be accepted
    at the $\mathsf{LLM\_JUDGED}$ trust level.  Higher values reduce false
    positives but increase the rate of unresolved obligations.
  \item \textbf{LLM selection}: the model used for proof search.  More capable
    models produce better proofs but cost more and are slower.
  \item \textbf{Timeout}: per-function verification timeout.  Functions that
    exceed this limit are reported as $\mathsf{TIMEOUT}$ rather than
    $\mathsf{FAILED}$.
  \item \textbf{Parallel workers}: the degree of parallelism for independent
    proof obligations.
\end{itemize}

\section{Verification Report Format}

The JSON verification report provides machine-readable per-function results:

\begin{lstlisting}[style=python]
{
  "project": "my_project",
  "timestamp": "2026-04-15T10:30:00Z",
  "synhopy_version": "0.9.0",
  "summary": {
    "total_functions": 42,
    "verified_kernel": 28,
    "verified_structural": 8,
    "axiom_trusted": 4,
    "failed": 1,
    "skipped": 1
  },
  "functions": [
    {
      "name": "clean_prices",
      "file": "src/pipeline.py",
      "line": 15,
      "guarantee": "returns DataFrame with no nulls",
      "verdict": "VERIFIED_STRUCTURAL",
      "trust_level": "AXIOM_TRUSTED",
      "axioms_used": [
        "pandas_read_csv_spec",
        "pandas_fillna_completeness"
      ],
      "proof_time_ms": 1230,
      "cached": false
    }
  ]
}
\end{lstlisting}

The report also supports the SARIF (Static Analysis Results Interchange
Format) standard, enabling integration with GitHub Code Scanning, VS Code,
and other tools that consume SARIF output.

\begin{proposition}[Report completeness]
Every function bearing a \lstinline[style=synhopy]{@guarantee} annotation
appears in the verification report with a verdict in
$\{\mathsf{VERIFIED\_KERNEL}, \mathsf{VERIFIED\_STRUCTURAL},
\mathsf{AXIOM\_TRUSTED}, \mathsf{FAILED}, \mathsf{TIMEOUT},
\mathsf{SKIPPED}\}$.
\end{proposition}


% ──────────────────────────────────────────────────────────────────
\chapter{Library Axiom System}
\label{ch:library-axioms}

\section{The Problem of Library Verification}

Real Python programs depend on libraries: numpy, pandas, flask, requests,
asyncio, json, re, etc.  Fully verifying these libraries is infeasible
(numpy alone is $\sim$200{,}000 lines of C/Python).  Instead, \synhopy{}
uses a \emph{library axiom system}: trusted declarations of library behavior.

\section{Axiom Registry}

\begin{definition}[Library axiom]
A library axiom is a triple $(\mathit{module}, \mathit{function}, \mathit{spec})$
where:
\begin{itemize}
  \item $\mathit{module}$: the Python module (e.g., \texttt{numpy}, \texttt{json})
  \item $\mathit{function}$: the function name (e.g., \texttt{array}, \texttt{dumps})
  \item $\mathit{spec}$: a formal specification of the function's behavior
\end{itemize}
\end{definition}

\begin{example}[Sample library axioms]
\begin{lstlisting}[style=proof]
-- json module axioms
axiom json_dumps_type:
    forall obj. json.dumps(obj) : str

axiom json_loads_inverse:
    forall obj. json.loads(json.dumps(obj)) = obj
    when obj contains only JSON-serializable types

axiom json_dumps_sort_keys:
    forall obj. json.dumps(obj, sort_keys=True) produces
    keys in alphabetical order

-- list axioms
axiom list_append_length:
    forall xs, x. len(xs.append(x)) = len(xs) + 1

axiom list_sort_permutation:
    forall xs. sorted(xs) is a permutation of xs

axiom list_sort_ordered:
    forall xs, i. i < len(sorted(xs)) - 1 =>
        sorted(xs)[i] <= sorted(xs)[i+1]

-- dict axioms
axiom dict_setitem:
    forall d, k, v. d[k] = v => d[k] == v

axiom dict_contains:
    forall d, k, v. d[k] = v => k in d
\end{lstlisting}
\end{example}

\section{Trust Propagation}

When a proof uses a library axiom, the trust level drops to
$\mathsf{AXIOM\_TRUSTED}$.  The system tracks \emph{which} axioms each proof
depends on:

\begin{definition}[Axiom dependency tracking]
For a proof $\pi$, define $\mathsf{axioms}(\pi)$ recursively:
\begin{align}
  \mathsf{axioms}(\Refl) &= \emptyset \\
  \mathsf{axioms}(\Trans(\pi_1, \pi_2)) &= \mathsf{axioms}(\pi_1) \cup \mathsf{axioms}(\pi_2) \\
  \mathsf{axioms}(\mathsf{Axiom}(a)) &= \{a\} \\
  \mathsf{axioms}(\mathsf{Z3}(\varphi)) &= \emptyset \quad\text{(Z3 is in the TCB)} \\
  &\;\vdots
\end{align}
\end{definition}

If a library axiom is later discovered to be incorrect, the system can
immediately identify all proofs that depend on it and mark them as invalid.

\begin{remark}[\fstar{} comparison]
\fstar{} uses \lstinline[style=fstar]{assume val} for library declarations,
which serves a similar purpose.  However, \fstar{} provides no tracking of
which proofs depend on which assumptions.  If an \lstinline[style=fstar]{assume val}
is wrong, there is no way to determine the blast radius without manually
auditing every proof in the project.  \synhopy{}'s axiom dependency graph
makes this audit automatic.
\end{remark}

\section{Per-Library Type Theories}

For frequently used libraries, \synhopy{} provides pre-built \emph{type theories}:
curated sets of axioms that capture the library's behavior at a level
sufficient for common verification tasks.

\begin{example}[Partial list of per-library theories]
\begin{center}
\begin{tabular}{lll}
\toprule
\textbf{Library} & \textbf{Axiom count} & \textbf{Covers} \\
\midrule
\texttt{builtins} & 45 & len, range, sorted, zip, map, filter, \ldots \\
\texttt{json} & 8 & dumps, loads, roundtrip, sort\_keys \\
\texttt{re} & 12 & match, search, findall, sub, split \\
\texttt{collections} & 15 & deque, Counter, defaultdict, OrderedDict \\
\texttt{itertools} & 20 & chain, islice, product, combinations, \ldots \\
\texttt{functools} & 10 & reduce, lru\_cache, partial, wraps \\
\texttt{math} & 18 & sqrt, sin, cos, log, factorial, gcd \\
\texttt{numpy} & 35 & array, shape, broadcasting, ufuncs \\
\texttt{pandas} & 25 & DataFrame, Series, groupby, merge \\
\bottomrule
\end{tabular}
\end{center}
\end{example}

\section{Axiom Generation from Documentation}

Manually writing library axioms for every function in every library is
tedious and error-prone.  \synhopy{} supports \emph{semi-automatic axiom
generation} from library documentation: given a function's docstring and
type hints, the system proposes candidate axioms.

\begin{construction}[Docstring-to-axiom pipeline]
\begin{enumerate}
  \item \textbf{Docstring extraction}: Parse the library source (or use
    \texttt{inspect.getdoc}) to extract docstrings and type annotations.
  \item \textbf{NL specification mining}: Use the LLM to convert the
    docstring into a set of candidate formal specifications.
  \item \textbf{Deduplication}: Remove axioms that are logical consequences
    of others already in the registry.
  \item \textbf{Human review}: Present the candidate axioms to the developer
    for approval.  Approved axioms enter the registry; rejected ones are
    discarded.
\end{enumerate}
\end{construction}

\begin{example}[Axiom generation from \texttt{numpy.clip}]
Given the docstring:
\begin{lstlisting}[style=python]
def clip(a, a_min, a_max, out=None):
    """Clip (limit) the values in an array.
    Given an interval, values outside the interval
    are clipped to the interval edges.
    """
\end{lstlisting}
The system proposes:
\begin{lstlisting}[style=proof]
axiom numpy_clip_lower:
    forall a, lo, hi, i.
        np.clip(a, lo, hi)[i] >= lo

axiom numpy_clip_upper:
    forall a, lo, hi, i.
        np.clip(a, lo, hi)[i] <= hi

axiom numpy_clip_shape:
    forall a, lo, hi.
        np.clip(a, lo, hi).shape == a.shape

axiom numpy_clip_identity:
    forall a, lo, hi, i.
        lo <= a[i] <= hi =>
            np.clip(a, lo, hi)[i] == a[i]
\end{lstlisting}
\end{example}

\section{Axiom Testing}

Library axioms are trusted declarations, but they can be \emph{tested}
against the actual library behavior at runtime.  This serves as a
lightweight validation that the axioms match reality.

\begin{definition}[Axiom test]
An \emph{axiom test} for axiom $a$ is a set of concrete inputs
$\{x_1, \ldots, x_k\}$ together with a checker function $c_a$ such that
$c_a(x_i)$ returns $\mathsf{True}$ iff the axiom holds for input $x_i$.
\end{definition}

\begin{lstlisting}[style=python]
# Auto-generated axiom test for numpy_clip_lower
import numpy as np

def test_numpy_clip_lower():
    for _ in range(1000):
        a = np.random.randn(50)
        lo, hi = sorted(np.random.randn(2))
        result = np.clip(a, lo, hi)
        assert np.all(result >= lo), \
            f"clip_lower violated: {result.min()} < {lo}"
\end{lstlisting}

Axiom tests do not \emph{prove} that axioms are correct (they cannot test
all possible inputs), but they provide high confidence and can detect
regressions when a library updates.

\begin{proposition}[Axiom testing coverage]
If an axiom $a$ passes $k$ randomized tests with independent inputs drawn
from a distribution $\mathcal{D}$, the probability that $a$ is incorrect on
a $\mathcal{D}$-random input is at most $(1 - \varepsilon)^k$ where
$\varepsilon$ is the measure of the violation set under $\mathcal{D}$.
For $k = 1000$ and $\varepsilon \geq 0.01$, this probability is
$< 4.3 \times 10^{-5}$.
\end{proposition}

\section{Versioned Axioms}

Library behavior can change between versions.  A function that returned a
list in numpy 1.x may return an ndarray in numpy 2.x.  \synhopy{} supports
\emph{versioned axioms} to handle this:

\begin{definition}[Versioned axiom set]
A versioned axiom set is a function $\mathsf{Axioms}: \mathsf{Version} \to
\mathcal{P}(\mathsf{Axiom})$ mapping library versions to their applicable
axiom sets.  The system reads the installed library version at verification
time and selects the appropriate axiom set.
\end{definition}

\begin{lstlisting}[style=proof]
-- Versioned axioms for numpy.unique
axiom numpy_unique_returns_sorted [numpy >= 1.0]:
    forall a. np.unique(a) is sorted in ascending order

axiom numpy_unique_equal_nan [numpy >= 1.24]:
    forall a. np.unique(a, equal_nan=True) treats
    all NaN values as identical

axiom numpy_unique_count [numpy >= 1.9]:
    forall a. np.unique(a, return_counts=True)
    returns (unique_values, counts) where
    sum(counts) == len(a)
\end{lstlisting}

When a project upgrades a library, the verification system automatically
re-validates all proofs that depend on axioms whose version range has
changed.  If an axiom is removed or weakened in the new version, the
affected proofs are flagged for re-verification.

\section{Community Axiom Repositories}

Axiom sets are shareable.  \synhopy{} supports community axiom
repositories---curated, tested, versioned axiom packages that can be
installed like Python packages:

\begin{lstlisting}[style=python]
# Install a community axiom package
pip install synhopy-axioms-numpy
pip install synhopy-axioms-pandas

# synhopy.toml
[axioms]
community = [
    "synhopy-axioms-numpy==2.1.0",
    "synhopy-axioms-pandas==1.4.0"
]
\end{lstlisting}

Community axiom packages include:
\begin{itemize}
  \item The axiom definitions (in a TOML-based format).
  \item A test suite validating each axiom against the target library version.
  \item A manifest listing the library versions the axioms have been tested
    against.
  \item A trust score based on the number of projects using the axiom set
    and the community review status.
\end{itemize}

\section{Axiom Composition}

When a project uses multiple libraries, axioms from different libraries must
be \emph{composed}.  For instance, a pandas operation that internally calls
numpy requires both pandas and numpy axioms.

\begin{definition}[Axiom composition]
Given axiom sets $A_1$ for library $L_1$ and $A_2$ for library $L_2$, the
composed axiom set is:
\[
  A_1 \otimes A_2 \defeq A_1 \cup A_2 \cup \mathsf{Bridge}(A_1, A_2)
\]
where $\mathsf{Bridge}(A_1, A_2)$ contains \emph{bridge axioms} that relate
operations across libraries.
\end{definition}

\begin{example}[Bridge axiom between pandas and numpy]
\begin{lstlisting}[style=proof]
-- Bridge axiom: pandas Series to numpy conversion
axiom pandas_to_numpy_bridge:
    forall s : pd.Series.
        s.to_numpy() returns np.ndarray with
        s.to_numpy().shape == (len(s),) and
        forall i. s.to_numpy()[i] == s.iloc[i]

-- Bridge axiom: numpy array to pandas DataFrame
axiom numpy_to_dataframe_bridge:
    forall arr : np.ndarray, cols : list[str].
        len(cols) == arr.shape[1] =>
        pd.DataFrame(arr, columns=cols).shape
            == arr.shape
\end{lstlisting}
\end{example}

\begin{proposition}[Composition consistency]
If $A_1$ and $A_2$ are individually consistent axiom sets and
$\mathsf{Bridge}(A_1, A_2)$ introduces no new contradictions, then
$A_1 \otimes A_2$ is consistent.
\end{proposition}

\section{Worked Example: NumPy Linear Algebra Axiom Set}

We present a complete axiom set for \texttt{numpy.linalg}, covering the
operations most commonly used in scientific computing.

\begin{lstlisting}[style=proof]
-- === numpy.linalg axiom set ===

-- Matrix multiplication
axiom np_matmul_shape:
    forall A : (m, k), B : (k, n).
        (A @ B).shape == (m, n)

axiom np_matmul_assoc:
    forall A, B, C.
        (A @ B) @ C == A @ (B @ C)

axiom np_matmul_identity:
    forall A : (n, n).
        A @ np.eye(n) == A and np.eye(n) @ A == A

-- Inverse
axiom np_inv_shape:
    forall A : (n, n).
        np.linalg.inv(A).shape == (n, n)

axiom np_inv_inverse:
    forall A : (n, n).
        np.linalg.det(A) != 0 =>
        A @ np.linalg.inv(A) == np.eye(n)

-- Determinant
axiom np_det_product:
    forall A, B : (n, n).
        np.linalg.det(A @ B)
            == np.linalg.det(A) * np.linalg.det(B)

axiom np_det_inv:
    forall A : (n, n).
        np.linalg.det(A) != 0 =>
        np.linalg.det(np.linalg.inv(A))
            == 1.0 / np.linalg.det(A)

axiom np_det_transpose:
    forall A.
        np.linalg.det(A.T) == np.linalg.det(A)

-- Eigenvalues
axiom np_eig_count:
    forall A : (n, n).
        vals, vecs = np.linalg.eig(A) =>
        len(vals) == n and vecs.shape == (n, n)

axiom np_eig_equation:
    forall A : (n, n).
        vals, vecs = np.linalg.eig(A) =>
        forall i. A @ vecs[:, i] ~= vals[i] * vecs[:, i]

-- Singular Value Decomposition
axiom np_svd_shapes:
    forall A : (m, n).
        U, S, Vt = np.linalg.svd(A) =>
        U.shape == (m, m) and S.shape == (min(m,n),)
            and Vt.shape == (n, n)

axiom np_svd_reconstruction:
    forall A : (m, n).
        U, S, Vt = np.linalg.svd(A) =>
        U[:, :min(m,n)] @ np.diag(S) @ Vt[:min(m,n), :]
            ~= A

-- Norm
axiom np_norm_nonneg:
    forall x.
        np.linalg.norm(x) >= 0

axiom np_norm_zero:
    forall x.
        np.linalg.norm(x) == 0 iff np.all(x == 0)

axiom np_norm_triangle:
    forall x, y.
        np.linalg.norm(x + y)
            <= np.linalg.norm(x) + np.linalg.norm(y)

-- Solve
axiom np_solve_correct:
    forall A : (n, n), b : (n,).
        np.linalg.det(A) != 0 =>
        x = np.linalg.solve(A, b) =>
        A @ x ~= b
\end{lstlisting}

The \texttt{\~{}=} notation denotes approximate equality up to floating-point
tolerance, which \synhopy{} tracks through a separate \emph{numerical
precision} axiom layer.

\begin{remark}[Floating-point precision]
Numerical linear algebra operations are inherently approximate.  The
axioms above use $\approx$ rather than $=$ for results involving
floating-point computation.  \synhopy{} tracks this imprecision: a proof
that relies on approximate axioms receives the trust annotation
$\mathsf{AXIOM\_TRUSTED\_APPROX}$, signaling that the result is correct
modulo floating-point rounding.
\end{remark}


% ══════════════════════════════════════════════════════════════════
% PART V: CASE STUDIES
% ══════════════════════════════════════════════════════════════════

\part{Case Studies}

% ──────────────────────────────────────────────────────────────────
\chapter{Verifying a Web Framework Route Handler}
\label{ch:case-web}

\section{The Code}

\begin{lstlisting}[style=synhopy]
from flask import Flask, request, jsonify

app = Flask(__name__)
users_db = {}  # in-memory "database"

@app.route('/users/<int:user_id>', methods=['GET'])
@guarantee("returns user data if exists, 404 otherwise")
def get_user(user_id):
    if user_id in users_db:
        return jsonify(users_db[user_id]), 200
    return jsonify({"error": "not found"}), 404

@app.route('/users', methods=['POST'])
@guarantee("creates a new user and returns 201")
@effect_contract(reads={"request"}, mutates={"users_db"})
def create_user():
    data = request.get_json()
    assume("data contains 'name' and 'email'")
    user_id = len(users_db) + 1
    users_db[user_id] = data
    check("user_id in users_db")
    check("users_db[user_id] == data")
    return jsonify({"id": user_id}), 201
\end{lstlisting}

\section{Verification}

\textbf{For \texttt{get\_user}:}
\begin{lstlisting}[style=proof]
-- Effect: READS (reads users_db, no mutation)
have h_cases: user_id in users_db or user_id not in users_db
    by structural  -- exhaustive boolean case split
cases user_id in users_db:
  case True:
    have h_data: users_db[user_id] is valid user data
        by assumption  -- precondition of the database
    conclude: returns (jsonify(data), 200) by unfold get_user
  case False:
    conclude: returns (jsonify({"error": "not found"}), 404)
        by unfold get_user
\end{lstlisting}

\textbf{For \texttt{create\_user}:}
\begin{lstlisting}[style=proof]
-- Effect: MUTATES (modifies users_db)
have h_pre: data contains 'name' and 'email'
    by assume  -- from the assume() annotation
have h_id: user_id = len(users_db) + 1 > 0
    by omega
have h_insert: after users_db[user_id] = data,
    user_id in users_db and users_db[user_id] == data
    by apply dict_setitem_spec
conclude: returns (jsonify({"id": user_id}), 201)
    by structural
\end{lstlisting}

\begin{remark}[\fstar{} comparison]
Verifying a Flask route handler in \fstar{} would require:
\begin{enumerate}
  \item Modeling Flask's routing system (decorators, request/response objects)
  \item Encoding the in-memory database as a monotonic state reference
  \item Writing wp-specifications for each handler
  \item Proving the handlers satisfy their specs relative to the state
\end{enumerate}
Estimated effort: 200+ lines of F* (vs.\ 30 lines of annotated Python + auto-generated proofs in \synhopy{}).
\end{remark}

\section{Request Validation and Input Sanitization}

Web applications must validate all user input.  \synhopy{} can verify that
request data is validated before use:

\begin{lstlisting}[style=synhopy]
from flask import request, abort
import re

@app.route('/users', methods=['POST'])
@guarantee("validates all input fields before use")
@effect_contract(reads={"request"}, mutates={"users_db"})
def create_user_safe():
    data = request.get_json()
    assume("data is a dict")
    check("isinstance(data.get('name'), str)")
    check("isinstance(data.get('email'), str)")
    check("len(data['name']) <= 100")
    check("re.match(r'^[^@]+@[^@]+\\.[^@]+$',
           data['email'])")

    # Sanitize: strip HTML tags
    name = re.sub(r'<[^>]+>', '', data['name'])
    check("'<' not in name and '>' not in name")

    user_id = len(users_db) + 1
    users_db[user_id] = {"name": name,
                         "email": data['email']}
    return jsonify({"id": user_id}), 201
\end{lstlisting}

\begin{lstlisting}[style=proof]
-- Proof for input validation:
have h_dict: data is a dict by assume
have h_name_type: isinstance(data['name'], str)
    by check -- verified structurally from isinstance
have h_email_type: isinstance(data['email'], str)
    by check
have h_name_len: len(data['name']) <= 100
    by check
have h_email_fmt: re.match(email_regex, data['email'])
    by check
have h_sanitize: re.sub(r'<[^>]+>', '', name)
    removes all angle-bracket sequences
    by apply re_sub_spec  -- library axiom
have h_clean: '<' not in name and '>' not in name
    by apply re_sub_completeness with h_sanitize
conclude: all input validated before DB insertion
    by structural with h_name_type, h_email_type,
       h_name_len, h_email_fmt, h_clean
\end{lstlisting}

\section{Authentication and Authorization Middleware}

\synhopy{} verifies that protected endpoints enforce authentication:

\begin{lstlisting}[style=synhopy]
from functools import wraps

@guarantee("returns 401 if token missing or invalid,
            else calls the wrapped function")
def require_auth(f):
    @wraps(f)
    def decorated(*args, **kwargs):
        token = request.headers.get('Authorization')
        assume("validate_token returns bool")
        if not token or not validate_token(token):
            return jsonify({"error": "unauthorized"}), 401
        return f(*args, **kwargs)
    return decorated

@app.route('/admin/users', methods=['DELETE'])
@require_auth
@guarantee("deletes user if authorized and user exists")
@effect_contract(reads={"request"},
                 mutates={"users_db"})
def delete_user(user_id):
    assume("caller passed require_auth check")
    if user_id not in users_db:
        return jsonify({"error": "not found"}), 404
    del users_db[user_id]
    check("user_id not in users_db")
    return jsonify({"status": "deleted"}), 200
\end{lstlisting}

\begin{lstlisting}[style=proof]
-- Proof for require_auth:
have h_token: token = request.headers.get('Authorization')
    by unfold
cases token is None or not validate_token(token):
  case True:
    conclude: returns ({"error": "unauthorized"}, 401)
        by unfold decorated
  case False:
    have h_valid: token is present and valid
        by structural
    conclude: calls f(*args, **kwargs) by unfold decorated
\end{lstlisting}

\section{Database Query Safety}

SQL injection is a critical vulnerability.  \synhopy{} verifies that all
database queries use parameterized statements:

\begin{lstlisting}[style=synhopy]
import sqlite3

@guarantee("returns user row without SQL injection risk")
@effect_contract(reads={"database"})
def get_user_by_name(db: sqlite3.Connection,
                     name: str) -> dict:
    assume("db is a valid connection")
    # Parameterized query -- safe
    cursor = db.execute(
        "SELECT * FROM users WHERE name = ?", (name,))
    check("query uses parameterized placeholder")
    row = cursor.fetchone()
    if row is None:
        return {}
    return dict(row)
\end{lstlisting}

\begin{definition}[SQL injection freedom]
A database query expression $q$ is \emph{injection-free} if all
user-supplied values appear only as parameters (placeholders) in the query
string, never as concatenated substrings.  Formally, if the query string
$s$ is a constant with respect to all user inputs, then $q$ is
injection-free.
\end{definition}

\begin{lstlisting}[style=proof]
-- Proof for SQL injection freedom:
have h_query_const: "SELECT * FROM users WHERE name = ?"
    is a string literal (no user input in query string)
    by structural -- AST analysis
have h_param: name appears only in parameter tuple (name,)
    by structural
conclude: query is injection-free
    by apply sql_parameterized_safety
        with h_query_const, h_param
\end{lstlisting}

\section{CORS and Security Header Verification}

\synhopy{} can verify that security headers are correctly set on all
responses:

\begin{lstlisting}[style=synhopy]
@app.after_request
@guarantee("adds all required security headers")
def add_security_headers(response):
    response.headers['X-Content-Type-Options'] = 'nosniff'
    response.headers['X-Frame-Options'] = 'DENY'
    response.headers['X-XSS-Protection'] = '1; mode=block'
    response.headers['Strict-Transport-Security'] = \
        'max-age=31536000; includeSubDomains'
    check("'X-Content-Type-Options' in response.headers")
    check("'X-Frame-Options' in response.headers")
    check("'X-XSS-Protection' in response.headers")
    check("'Strict-Transport-Security' in response.headers")
    return response
\end{lstlisting}

\begin{lstlisting}[style=proof]
-- Proof for security headers:
have h1: after assignment, 'X-Content-Type-Options'
    in response.headers
    by apply dict_setitem_spec
have h2: after assignment, 'X-Frame-Options'
    in response.headers
    by apply dict_setitem_spec
have h3: after assignment, 'X-XSS-Protection'
    in response.headers
    by apply dict_setitem_spec
have h4: after assignment, 'Strict-Transport-Security'
    in response.headers
    by apply dict_setitem_spec
conclude: all required security headers present
    by structural with h1, h2, h3, h4
\end{lstlisting}

\section{Error Handler Completeness}

A robust web application should handle all standard HTTP error codes.
\synhopy{} can verify \emph{error handler completeness}:

\begin{lstlisting}[style=synhopy]
REQUIRED_ERROR_HANDLERS = [400, 401, 403, 404, 405,
                           408, 429, 500, 502, 503]

@guarantee("all required error codes have handlers")
def register_error_handlers(app):
    for code in REQUIRED_ERROR_HANDLERS:
        check(f"app.error_handler_spec[None][code]
               is not None")

    @app.errorhandler(400)
    def bad_request(e):
        return jsonify({"error": "bad request",
                        "code": 400}), 400

    @app.errorhandler(404)
    def not_found(e):
        return jsonify({"error": "not found",
                        "code": 404}), 404

    @app.errorhandler(500)
    def internal_error(e):
        return jsonify({"error": "internal server error",
                        "code": 500}), 500
    # ... handlers for remaining codes ...
\end{lstlisting}

\begin{lstlisting}[style=proof]
-- Proof for error handler completeness:
have h_codes: REQUIRED_ERROR_HANDLERS =
    [400, 401, 403, 404, 405, 408, 429, 500, 502, 503]
    by unfold
have h_400: app.errorhandler(400) registered
    by structural -- decorator application
have h_404: app.errorhandler(404) registered
    by structural
have h_500: app.errorhandler(500) registered
    by structural
-- ... (similar for each code) ...
conclude: all required error codes have handlers
    by structural with h_400, h_404, h_500, ...
\end{lstlisting}

\section{Worked Example: Complete REST API Endpoint}

We present a complete verified REST API endpoint that incorporates all
the verification techniques from this chapter: input validation,
authentication, SQL safety, security headers, and error handling.

\begin{lstlisting}[style=synhopy]
@app.route('/api/v1/orders', methods=['POST'])
@require_auth
@guarantee("creates order if input valid, returns
            appropriate error otherwise")
@effect_contract(reads={"request", "database"},
                 mutates={"database"})
def create_order():
    data = request.get_json()
    if not data:
        return jsonify({"error": "missing body"}), 400

    assume("schema_validate returns (bool, str)")
    valid, msg = schema_validate(data, ORDER_SCHEMA)
    if not valid:
        return jsonify({"error": msg}), 422

    check("data['product_id'] is int")
    check("data['quantity'] > 0")

    db = get_db()
    cursor = db.execute(
        "INSERT INTO orders (product_id, quantity, "
        "user_id) VALUES (?, ?, ?)",
        (data['product_id'], data['quantity'],
         g.current_user_id))
    db.commit()
    check("cursor.lastrowid is not None")

    order = {"id": cursor.lastrowid,
             "product_id": data['product_id'],
             "quantity": data['quantity'],
             "status": "pending"}
    return jsonify(order), 201
\end{lstlisting}

\begin{lstlisting}[style=proof]
-- Full proof script for create_order:

-- Step 1: Authentication (from require_auth)
have h_auth: caller is authenticated
    by assume -- guaranteed by @require_auth decorator

-- Step 2: Input presence
cases data is None:
  case True:
    conclude: returns ({"error": "missing body"}, 400)
        by unfold create_order
  case False:
    -- Step 3: Schema validation
    have h_schema: schema_validate returns (valid, msg)
        by assume
    cases valid:
      case False:
        conclude: returns ({"error": msg}, 422)
            by unfold create_order
      case True:
        -- Step 4: Field validation
        have h_pid: data['product_id'] is int
            by check -- from schema validation
        have h_qty: data['quantity'] > 0
            by check

        -- Step 5: SQL injection freedom
        have h_sql: INSERT query uses only ? placeholders
            by structural -- query string is constant
        have h_params: all user values in parameter tuple
            by structural

        -- Step 6: Database insertion
        have h_insert: db.execute with parameterized query
            inserts one row
            by apply sqlite_execute_insert_spec

        have h_rowid: cursor.lastrowid is not None
            after successful INSERT
            by apply sqlite_lastrowid_spec

        -- Step 7: Response construction
        conclude: returns (order_dict, 201)
            by structural with h_insert, h_rowid
\end{lstlisting}


% ──────────────────────────────────────────────────────────────────
\chapter{Verifying a Data Processing Pipeline}
\label{ch:case-data}

\section{The Code}

\begin{lstlisting}[style=synhopy]
import pandas as pd

@guarantee("returns DataFrame with no null values in 'price' column")
@effect_contract(reads={"filesystem"}, pure_transform=True)
def clean_prices(filepath: str) -> pd.DataFrame:
    assume("filepath points to a valid CSV with a 'price' column")
    df = pd.read_csv(filepath)
    check("'price' in df.columns")
    df['price'] = df['price'].fillna(df['price'].median())
    check("df['price'].isna().sum() == 0")
    return df

@guarantee("returns total revenue as a positive Decimal")
def total_revenue(df: pd.DataFrame) -> float:
    assume("df has 'price' and 'quantity' columns with no nulls")
    df['revenue'] = df['price'] * df['quantity']
    return df['revenue'].sum()
\end{lstlisting}

\section{Verification}

\begin{lstlisting}[style=proof]
-- Proof for clean_prices:
have h_read: pd.read_csv(filepath) returns DataFrame
    by apply pandas_read_csv_spec  -- library axiom
have h_col: 'price' in df.columns
    by assumption  -- from assume()
have h_fill: fillna(median()) replaces all NaN with median
    by apply pandas_fillna_spec  -- library axiom
have h_no_null: after fillna, isna().sum() == 0
    by apply pandas_fillna_completeness  -- library axiom
conclude: result has no null values in 'price' column
    by trans(h_fill, h_no_null)
\end{lstlisting}

This proof is at trust level \textsf{AXIOM\_TRUSTED} because it depends on
library axioms (\texttt{pandas\_read\_csv\_spec}, \texttt{pandas\_fillna\_spec}).
The axioms are user-provided descriptions of pandas' behavior, analogous to
\fstar{}'s \lstinline[style=fstar]{assume val} declarations.

\begin{remark}[\fstar{} comparison]
\fstar{} has no way to reason about pandas DataFrames. You would need to:
\begin{enumerate}
  \item Write a complete model of pandas in \fstar{} (DataFrames, Series, indexing,
    fillna, read\_csv, etc.)---easily 1000+ lines
  \item Verify against the model, not against actual pandas
  \item Hope that the model matches pandas' actual behavior
\end{enumerate}
\synhopy{} uses library axioms instead: the user declares ``fillna replaces all
NaN values'' as an axiom, and the proof depends on this axiom at a tracked trust
level. This is more honest about the verification boundary.
\end{remark}

\section{CSV Parsing and Schema Validation}

Before processing data, we must verify that the input file conforms to an
expected schema.  \synhopy{} enables verifying schema constraints as
preconditions:

\begin{lstlisting}[style=synhopy]
@guarantee("returns DataFrame matching expected schema")
@effect_contract(reads={"filesystem"})
def load_and_validate(filepath: str,
                      schema: dict) -> pd.DataFrame:
    assume("filepath exists and is a valid CSV")
    df = pd.read_csv(filepath)

    # Verify all expected columns exist
    for col in schema['columns']:
        check(f"'{col}' in df.columns")

    # Verify column types
    for col, dtype in schema['dtypes'].items():
        check(f"df['{col}'].dtype == {dtype}")

    # Verify no unexpected columns
    check("set(df.columns) == set(schema['columns'])")

    return df
\end{lstlisting}

\begin{lstlisting}[style=proof]
-- Proof for load_and_validate:
have h_read: pd.read_csv(filepath) returns DataFrame
    by apply pandas_read_csv_spec
have h_cols: forall col in schema['columns'].
    col in df.columns
    by check -- each column verified in loop
have h_dtypes: forall col in schema['dtypes'].
    df[col].dtype matches expected
    by check
have h_exact: no extra columns beyond schema
    by check with set equality
conclude: df matches expected schema
    by structural with h_cols, h_dtypes, h_exact
\end{lstlisting}

\section{Data Transformation Pipeline Composition}

Data processing often involves composing multiple transformations.
\synhopy{} verifies that the composition preserves invariants across
each stage:

\begin{lstlisting}[style=synhopy]
@guarantee("each stage preserves row count")
def transform_pipeline(df: pd.DataFrame) -> pd.DataFrame:
    assume("df has columns: name, price, quantity, date")
    initial_rows = len(df)

    # Stage 1: Clean strings
    df['name'] = df['name'].str.strip().str.lower()
    check("len(df) == initial_rows")

    # Stage 2: Parse dates
    df['date'] = pd.to_datetime(df['date'])
    check("len(df) == initial_rows")

    # Stage 3: Compute derived column
    df['total'] = df['price'] * df['quantity']
    check("len(df) == initial_rows")
    check("'total' in df.columns")

    return df
\end{lstlisting}

\begin{definition}[Row-preserving transformation]
A DataFrame transformation $T$ is \emph{row-preserving} if
$\mathsf{len}(T(\mathit{df})) = \mathsf{len}(\mathit{df})$ for all
inputs $\mathit{df}$.  Column mutations (string operations, type casting,
derived columns) are row-preserving; filters and drops are not.
\end{definition}

\begin{lstlisting}[style=proof]
-- Proof for pipeline composition:
have h0: initial_rows = len(df) by unfold
have h1: str.strip().str.lower() is row-preserving
    by apply pandas_str_transform_rows  -- library axiom
have h2: pd.to_datetime is row-preserving
    by apply pandas_to_datetime_rows  -- library axiom
have h3: column assignment is row-preserving
    by apply pandas_assign_column_rows  -- library axiom
conclude: len(result) == initial_rows
    by trans(h1, trans(h2, h3))
\end{lstlisting}

\section{Null and NaN Handling Verification}

Missing data is a pervasive source of bugs.  \synhopy{} tracks null
propagation through DataFrame operations:

\begin{lstlisting}[style=synhopy]
@guarantee("returns DataFrame with no NaN in any column")
def handle_missing(df: pd.DataFrame) -> pd.DataFrame:
    assume("df may contain NaN values")

    # Strategy: fill numeric NaN with 0, drop rows
    # with NaN in string columns
    numeric_cols = df.select_dtypes(
        include=['number']).columns
    string_cols = df.select_dtypes(
        include=['object']).columns

    for col in numeric_cols:
        df[col] = df[col].fillna(0)
        check(f"df['{col}'].isna().sum() == 0")

    df = df.dropna(subset=string_cols)
    check("df[string_cols].isna().sum().sum() == 0")
    check("df.isna().sum().sum() == 0")

    return df
\end{lstlisting}

\begin{lstlisting}[style=proof]
-- Proof for handle_missing:
have h_numeric: forall col in numeric_cols.
    fillna(0) eliminates all NaN in col
    by apply pandas_fillna_completeness
have h_string: dropna(subset=string_cols) removes rows
    where any string column is NaN
    by apply pandas_dropna_spec
have h_complete: after fillna on numerics and dropna
    on strings, no NaN remains in any column
    by structural with h_numeric, h_string
conclude: df.isna().sum().sum() == 0
    by h_complete
\end{lstlisting}

\section{Type Coercion Safety}

Implicit type coercion in pandas can produce unexpected results (e.g.,
an integer column silently becoming float when NaN is introduced).
\synhopy{} verifies that type coercions are intentional:

\begin{lstlisting}[style=synhopy]
@guarantee("price column remains float64 throughout")
def safe_type_cast(df: pd.DataFrame) -> pd.DataFrame:
    assume("df['price'] is float64")
    check("df['price'].dtype == 'float64'")

    df['price_cents'] = (df['price'] * 100).astype(int)
    check("df['price_cents'].dtype == 'int64'")

    # Verify no precision loss beyond rounding
    check("((df['price_cents'] / 100.0) - df['price'])"
          ".abs().max() < 0.01")

    return df
\end{lstlisting}

\begin{lstlisting}[style=proof]
-- Proof for type coercion safety:
have h_float: df['price'].dtype == float64 by assume
have h_mul: (df['price'] * 100) preserves float64
    by apply pandas_arithmetic_dtype_spec
have h_cast: astype(int) converts to int64
    by apply pandas_astype_spec
have h_precision: integer truncation loses < 0.01
    per value when original has <= 2 decimal places
    by omega with h_mul, h_cast
conclude: price_cents is int64 with bounded precision loss
    by structural with h_cast, h_precision
\end{lstlisting}

\section{Aggregation Correctness}

GroupBy and reduce operations are common sources of errors (wrong grouping
key, incorrect aggregation function, unexpected behavior with empty groups):

\begin{lstlisting}[style=synhopy]
@guarantee("returns per-category revenue summary")
def revenue_by_category(
        df: pd.DataFrame) -> pd.DataFrame:
    assume("df has 'category', 'price', 'quantity' cols")
    assume("no NaN in category, price, or quantity")

    df['revenue'] = df['price'] * df['quantity']
    summary = df.groupby('category')['revenue'].agg(
        ['sum', 'mean', 'count']
    ).reset_index()

    check("set(summary.columns) == "
          "{'category', 'sum', 'mean', 'count'}")
    check("summary['count'].sum() == len(df)")
    check("summary['sum'].sum() == df['revenue'].sum()")

    return summary
\end{lstlisting}

\begin{lstlisting}[style=proof]
-- Proof for aggregation correctness:
have h_rev: df['revenue'] = df['price'] * df['quantity']
    by unfold
have h_group: groupby('category') partitions df by
    unique category values
    by apply pandas_groupby_partition_spec
have h_agg: agg(['sum', 'mean', 'count']) computes
    per-group aggregations
    by apply pandas_agg_spec
have h_count: sum of per-group counts == total rows
    by apply pandas_groupby_count_sum  -- partition axiom
have h_sum: sum of per-group sums == total sum
    by apply pandas_groupby_sum_sum  -- partition axiom
conclude: summary is correct aggregation
    by structural with h_count, h_sum
\end{lstlisting}

\section{Worked Example: Verifying an ETL Pipeline}

We present a complete ETL (Extract, Transform, Load) pipeline with
end-to-end verification:

\begin{lstlisting}[style=synhopy]
@guarantee("transforms raw CSV into clean, aggregated
            output with no data loss")
@effect_contract(reads={"filesystem"},
                 mutates={"filesystem"})
def etl_pipeline(input_path: str,
                 output_path: str) -> dict:
    assume("input_path is a valid CSV")

    # Extract
    df = pd.read_csv(input_path)
    raw_rows = len(df)
    check("raw_rows > 0")

    # Transform: clean
    df['amount'] = pd.to_numeric(
        df['amount'], errors='coerce')
    df = df.dropna(subset=['amount'])
    clean_rows = len(df)
    check("clean_rows <= raw_rows")

    # Transform: enrich
    df['amount_usd'] = df['amount'] * df['exchange_rate']
    check("len(df) == clean_rows")

    # Transform: aggregate
    summary = df.groupby('region')['amount_usd'].sum()
    check("summary.sum() == df['amount_usd'].sum()")

    # Load
    summary.to_csv(output_path)

    return {
        "raw_rows": raw_rows,
        "clean_rows": clean_rows,
        "regions": len(summary),
        "total_usd": float(summary.sum())
    }
\end{lstlisting}

\begin{lstlisting}[style=proof]
-- Full ETL proof:

-- Extract phase
have h_read: pd.read_csv returns DataFrame
    by apply pandas_read_csv_spec
have h_nonempty: raw_rows > 0 by check

-- Transform: clean phase
have h_coerce: to_numeric(errors='coerce') converts
    non-numeric to NaN
    by apply pandas_to_numeric_coerce_spec
have h_drop: dropna removes NaN rows
    by apply pandas_dropna_spec
have h_leq: clean_rows <= raw_rows
    by omega with h_drop  -- dropping can only reduce

-- Transform: enrich phase
have h_mul: column multiplication is row-preserving
    by apply pandas_assign_column_rows
have h_enrich_rows: len(df) == clean_rows
    by h_mul

-- Transform: aggregate phase
have h_group_sum: groupby sum preserves total
    by apply pandas_groupby_sum_sum
have h_total: summary.sum() == df['amount_usd'].sum()
    by h_group_sum

-- Load phase
have h_write: to_csv writes summary to output_path
    by apply pandas_to_csv_spec

conclude: ETL pipeline produces correct aggregation
    with tracked data loss
    by structural with h_leq, h_total, h_write
\end{lstlisting}


% ──────────────────────────────────────────────────────────────────
\chapter{Verifying SymPy Computations}
\label{ch:case-sympy}

\section{The Challenge}

SymPy is a symbolic mathematics library for Python. Verifying that a SymPy
computation produces a correct result requires reasoning about both the
program's behavior and the mathematical properties of the computation.

\section{The Code}

\begin{lstlisting}[style=synhopy]
from sympy import symbols, integrate, sin, cos, pi

@guarantee("returns the integral of sin(x) from 0 to pi, which is 2")
def compute_integral():
    x = symbols('x')
    result = integrate(sin(x), (x, 0, pi))
    check("result == 2")
    return result

@guarantee("returns True iff the identity holds for all x")
def verify_trig_identity():
    x = symbols('x')
    expr = sin(x)**2 + cos(x)**2 - 1
    simplified = expr.simplify()
    check("simplified == 0")
    return simplified == 0
\end{lstlisting}

\section{Verification}

\begin{lstlisting}[style=proof]
-- Proof for compute_integral:
have h_integral: integrate(sin(x), (x, 0, pi)) = [-cos(x)]_0^pi
    by apply sympy_integrate_sin  -- library axiom
have h_eval: [-cos(x)]_0^pi = -cos(pi) - (-cos(0))
            = -(-1) - (-1) = 1 + 1 = 2
    by omega with cos_pi = -1, cos_0 = 1
conclude: result == 2
    by trans(h_integral, h_eval)

-- Proof for verify_trig_identity:
have h_identity: sin(x)^2 + cos(x)^2 = 1 for all x
    by apply pythagorean_identity  -- mathematical axiom
have h_expr: sin(x)^2 + cos(x)^2 - 1 = 0
    by rw [h_identity] then omega
have h_simplify: simplify(0) = 0
    by apply sympy_simplify_zero  -- library axiom
conclude: simplified == 0
    by trans(h_expr, h_simplify)
\end{lstlisting}

The proofs depend on both \emph{library axioms} (how SymPy's \texttt{integrate}
and \texttt{simplify} behave) and \emph{mathematical axioms} (the Pythagorean
identity, values of cosine). The trust level is \textsf{AXIOM\_TRUSTED}, with
explicit tracking of which axioms are used.

\begin{remark}[\fstar{} comparison]
To verify SymPy computations in \fstar{}, you would need to either:
\begin{enumerate}
  \item Reimplement symbolic computation in \fstar{} and prove it correct (a
    multi-year research project), or
  \item Treat SymPy as an oracle with \lstinline[style=fstar]{assume val} for
    every operation---equivalent to \synhopy{}'s library axioms but without
    the structured trust tracking.
\end{enumerate}
\synhopy{}'s approach is pragmatic: trust SymPy's operations as axioms, track
the trust level, and focus verification on the \emph{user's code} that calls
SymPy.
\end{remark}

\section{Symbolic Simplification Correctness}

SymPy's \texttt{simplify} function applies heuristic rewriting rules.
\synhopy{} verifies that the simplified expression is mathematically
equivalent to the original:

\begin{lstlisting}[style=synhopy]
from sympy import symbols, simplify, expand, factor

@guarantee("simplified form is equivalent to original")
def simplify_expression(expr_str: str):
    x, y = symbols('x y')
    expr = eval(expr_str)  # in a sandboxed evaluator
    assume("expr is a valid SymPy expression in x, y")
    simplified = simplify(expr)
    check("simplify(expr - simplified) == 0")
    return simplified
\end{lstlisting}

\begin{definition}[Symbolic equivalence]
Two SymPy expressions $e_1$ and $e_2$ are \emph{symbolically equivalent},
written $e_1 \equiv_{\mathrm{sym}} e_2$, if
$\mathsf{simplify}(e_1 - e_2) = 0$.  This is a sufficient (but not
necessary) condition for mathematical equality; \synhopy{} tracks it as
an axiom-dependent assertion.
\end{definition}

\begin{lstlisting}[style=proof]
-- Proof for simplify_expression:
have h_valid: expr is a valid SymPy expression
    by assume
have h_simplify: simplify(expr) returns an expression
    equivalent to expr
    by apply sympy_simplify_equivalence  -- library axiom
have h_check: simplify(expr - simplified) == 0
    by apply sympy_simplify_diff_zero
        with h_simplify  -- follows from equivalence
conclude: simplified is equivalent to original
    by h_check
\end{lstlisting}

\section{Integration and Differentiation Verification}

The fundamental theorem of calculus provides a bridge between integration
and differentiation.  \synhopy{} uses this to cross-verify SymPy results:

\begin{lstlisting}[style=synhopy]
from sympy import symbols, integrate, diff, simplify

@guarantee("integral and derivative are consistent
            via the fundamental theorem of calculus")
def verify_ftc(f_expr):
    x = symbols('x')
    assume("f_expr is a continuous SymPy expression in x")

    # Integrate then differentiate
    F = integrate(f_expr, x)
    f_recovered = diff(F, x)
    check("simplify(f_recovered - f_expr) == 0")

    # Definite integral via antiderivative
    a, b = symbols('a b')
    definite = integrate(f_expr, (x, a, b))
    via_antideriv = F.subs(x, b) - F.subs(x, a)
    check("simplify(definite - via_antideriv) == 0")

    return True
\end{lstlisting}

\begin{lstlisting}[style=proof]
-- Proof for verify_ftc:
have h_cont: f_expr is continuous by assume
have h_antideriv: F = integrate(f_expr, x) is an
    antiderivative of f_expr
    by apply sympy_integrate_antideriv  -- library axiom
have h_ftc1: diff(integrate(f, x), x) == f
    for continuous f
    by apply fundamental_theorem_calculus_1
        -- mathematical axiom
have h_recovered: simplify(f_recovered - f_expr) == 0
    by apply sympy_simplify_diff_zero with h_ftc1
have h_ftc2: integrate(f, (x, a, b))
    == F(b) - F(a) where F is antiderivative
    by apply fundamental_theorem_calculus_2
have h_definite: simplify(definite - via_antideriv) == 0
    by apply sympy_simplify_diff_zero with h_ftc2
conclude: integral and derivative are consistent
    by structural with h_recovered, h_definite
\end{lstlisting}

\section{Equation Solving Soundness}

When SymPy solves an equation, \synhopy{} verifies that the solution
actually satisfies the original equation:

\begin{lstlisting}[style=synhopy]
from sympy import symbols, solve, Eq, simplify

@guarantee("all returned solutions satisfy the equation")
def verified_solve(equation, variable):
    assume("equation is a valid SymPy equation")
    solutions = solve(equation, variable)
    check("isinstance(solutions, list)")

    for sol in solutions:
        substituted = equation.subs(variable, sol)
        check("simplify(substituted) == 0 or "
              "simplify(substituted) == True")

    return solutions
\end{lstlisting}

\begin{lstlisting}[style=proof]
-- Proof for verified_solve:
have h_eq: equation is valid by assume
have h_solve: solve returns candidate solutions
    by apply sympy_solve_spec  -- library axiom
have h_verify: forall sol in solutions.
    equation.subs(variable, sol) simplifies to 0
    by check -- runtime verification for each solution
conclude: all solutions satisfy the equation
    by structural with h_verify
\end{lstlisting}

\begin{remark}[Solve completeness]
Note that \synhopy{} verifies \emph{soundness} (returned solutions are
correct) but not \emph{completeness} (all solutions are returned).
SymPy's \texttt{solve} may miss solutions in certain cases (e.g., branch
cuts for complex equations).  Completeness would require additional
mathematical axioms specific to the equation class.
\end{remark}

\section{Matrix Operations Verification}

SymPy provides symbolic matrix operations.  \synhopy{} verifies dimensional
consistency and algebraic properties:

\begin{lstlisting}[style=synhopy]
from sympy import Matrix, eye, zeros, simplify

@guarantee("computes inverse and verifies A * A_inv = I")
def verified_inverse(matrix_data: list) -> Matrix:
    A = Matrix(matrix_data)
    assume("A is square and non-singular")
    check("A.rows == A.cols")

    A_inv = A.inv()
    product = A * A_inv
    identity = eye(A.rows)

    check("simplify(product - identity)"
          " == zeros(A.rows)")

    return A_inv
\end{lstlisting}

\begin{lstlisting}[style=proof]
-- Proof for verified_inverse:
have h_square: A.rows == A.cols by check
have h_nonsingular: A.det() != 0 by assume
have h_inv: A.inv() exists for non-singular A
    by apply sympy_matrix_inv_spec with h_nonsingular
have h_product: A * A.inv() = I (symbolically)
    by apply sympy_matrix_inv_identity
        with h_nonsingular
have h_simplify: simplify(A * A_inv - I) == zeros(n)
    by apply sympy_simplify_matrix with h_product
conclude: A * A_inv == I verified
    by h_simplify
\end{lstlisting}

\section{SymPy Expression Tree as an Algebraic Data Type}

SymPy expressions form a tree structure that can be modeled as an
algebraic data type in \synhopy{}'s type theory:

\begin{definition}[SymPy expression ADT]
The type $\mathsf{Expr}$ is defined inductively:
\begin{align}
  \mathsf{Expr} &\defeq \mathsf{Symbol}(\mathit{name} : \mathsf{str}) \\
    &\mid \mathsf{Integer}(n : \Int) \\
    &\mid \mathsf{Rational}(p : \Int, q : \Int) \\
    &\mid \mathsf{Add}(\mathit{args} : \mathsf{List}(\mathsf{Expr})) \\
    &\mid \mathsf{Mul}(\mathit{args} : \mathsf{List}(\mathsf{Expr})) \\
    &\mid \mathsf{Pow}(\mathit{base} : \mathsf{Expr}, \mathit{exp} : \mathsf{Expr}) \\
    &\mid \mathsf{FuncApp}(\mathit{func} : \mathsf{str}, \mathit{args} : \mathsf{List}(\mathsf{Expr}))
\end{align}
\end{definition}

This ADT enables structural induction on SymPy expressions, which supports
proofs about expression transformations:

\begin{proposition}[Structural simplification]
For any expression $e : \mathsf{Expr}$, if $\mathsf{simplify}(e) = e'$,
then $\den{e} = \den{e'}$ (the denotational semantics are preserved).
This is an axiom we trust from SymPy, but we can verify it structurally for
specific expression patterns.
\end{proposition}

\section{Worked Example: Verifying a Physics Simulation}

We present a complete example verifying a simple physics simulation that
uses SymPy for symbolic computation of equations of motion:

\begin{lstlisting}[style=synhopy]
from sympy import symbols, Function, dsolve, Eq
from sympy import sin, cos, sqrt, simplify

@guarantee("returns symbolic solution for projectile
            motion matching known physics")
def projectile_motion():
    t, g, v0, theta = symbols('t g v_0 theta',
                               positive=True)

    # Horizontal position: x(t) = v0 * cos(theta) * t
    x = v0 * cos(theta) * t

    # Vertical position: y(t) = v0*sin(theta)*t - g*t^2/2
    y = v0 * sin(theta) * t - g * t**2 / 2

    # Verify: at t=0, position is origin
    check("x.subs(t, 0) == 0")
    check("y.subs(t, 0) == 0")

    # Verify: velocity components
    from sympy import diff
    vx = diff(x, t)
    vy = diff(y, t)
    check("simplify(vx - v0 * cos(theta)) == 0")
    check("simplify(vy.subs(t, 0)"
          " - v0 * sin(theta)) == 0")

    # Verify: acceleration
    ax = diff(vx, t)
    ay = diff(vy, t)
    check("ax == 0")
    check("simplify(ay + g) == 0")

    # Time of flight: when y = 0 again
    from sympy import solve
    flight_times = solve(y, t)
    check("0 in flight_times")  # t=0 is trivial
    T = [ft for ft in flight_times if ft != 0][0]
    check("simplify(T - 2*v0*sin(theta)/g) == 0")

    # Range: x at time T
    R = x.subs(t, T)
    expected_range = v0**2 * sin(2*theta) / g
    check("simplify(R - expected_range) == 0")

    return {
        "x": x, "y": y,
        "flight_time": T, "range": R
    }
\end{lstlisting}

\begin{lstlisting}[style=proof]
-- Full proof for projectile_motion:

-- Initial conditions
have h_x0: x.subs(t, 0) = v0*cos(theta)*0 = 0
    by omega with sympy_subs_spec
have h_y0: y.subs(t, 0) = v0*sin(theta)*0 - g*0/2 = 0
    by omega with sympy_subs_spec

-- Velocity verification
have h_vx: diff(v0*cos(theta)*t, t) = v0*cos(theta)
    by apply sympy_diff_linear  -- library axiom
have h_vy0: diff(y, t).subs(t, 0) = v0*sin(theta)
    by apply sympy_diff_polynomial then sympy_subs_spec

-- Acceleration verification
have h_ax: diff(v0*cos(theta), t) = 0
    by apply sympy_diff_constant  -- constant w.r.t. t
have h_ay: diff(v0*sin(theta) - g*t, t) = -g
    by apply sympy_diff_polynomial

-- Time of flight
have h_solve_y: solve(v0*sin(theta)*t - g*t^2/2, t)
    = [0, 2*v0*sin(theta)/g]
    by apply sympy_solve_quadratic  -- library axiom
have h_T: T = 2*v0*sin(theta)/g
    by structural with h_solve_y  -- nonzero solution

-- Range
have h_range: x.subs(t, T)
    = v0*cos(theta) * 2*v0*sin(theta)/g
    = 2*v0^2*sin(theta)*cos(theta)/g
    = v0^2*sin(2*theta)/g
    by apply double_angle_identity  -- math axiom:
       -- 2*sin(theta)*cos(theta) = sin(2*theta)

conclude: all physics equations verified
    by structural with h_x0, h_y0, h_vx, h_vy0,
       h_ax, h_ay, h_T, h_range
\end{lstlisting}

This example demonstrates the interplay between library axioms (SymPy's
\texttt{diff}, \texttt{solve}, \texttt{simplify}) and mathematical axioms
(the double-angle identity, properties of polynomials).  The trust level
is $\mathsf{AXIOM\_TRUSTED}$, with the full axiom dependency set recorded
in the verification report.


% ──────────────────────────────────────────────────────────────────
\chapter{Verifying a Machine Learning Pipeline}
\label{ch:case-ml}

\section{The Code}

\begin{lstlisting}[style=synhopy]
import numpy as np

@guarantee("returns array of same shape with values in [0, 1]")
def normalize(data: np.ndarray) -> np.ndarray:
    assume("data is non-empty and contains finite values")
    mn, mx = data.min(), data.max()
    if mn == mx:
        return np.zeros_like(data)
    result = (data - mn) / (mx - mn)
    check("result.shape == data.shape")
    check("result.min() >= 0 and result.max() <= 1")
    return result

@guarantee("returns list of n centroids for k-means")
def initialize_centroids(data: np.ndarray, n: int) -> list:
    assume("n > 0 and n <= len(data)")
    indices = np.random.choice(len(data), n, replace=False)
    centroids = [data[i].tolist() for i in indices]
    check("len(centroids) == n")
    check("all(c in data for c in centroids)")
    return centroids
\end{lstlisting}

\section{Verification}

\textbf{For \texttt{normalize}:}
\begin{lstlisting}[style=proof]
-- Effect: PURE (np operations create new arrays)
cases mn == mx:
  case True:
    have h_zeros: np.zeros_like(data).shape == data.shape
        by apply numpy_zeros_like_shape
    have h_range: all values are 0, so in [0, 1]
        by omega
    conclude by structural with h_zeros, h_range
  case False:
    have h_sub: (data - mn) has values in [0, mx - mn]
        by omega with mn = data.min()
    have h_div: dividing by (mx - mn) > 0 scales to [0, 1]
        by omega with h_sub, mx > mn
    have h_shape: arithmetic preserves shape
        by apply numpy_broadcast_shape
    conclude by structural with h_div, h_shape
\end{lstlisting}

\textbf{For \texttt{initialize\_centroids}:}
\begin{lstlisting}[style=proof]
-- Effect: NONDETERMINISTIC (np.random.choice)
have h_choice: np.random.choice(len(data), n, replace=False)
    returns n distinct indices in [0, len(data))
    by apply numpy_random_choice_spec
have h_len: len(centroids) = len(indices) = n
    by rw [h_choice] then omega
have h_in: forall i in indices. data[i] in data
    by structural -- indexing into the source array
conclude by structural with h_len, h_in
\end{lstlisting}

Note the effect annotation: \texttt{initialize\_centroids} has
$\Nondet$ effect due to \texttt{np.random.choice}.  Two calls
may return different centroids.  \synhopy{} tracks this: proofs about
\texttt{initialize\_centroids} cannot use equational reasoning
(``$f(x) = f(x)$'' does not hold for nondeterministic functions).

\begin{remark}[\fstar{} comparison]
\fstar{} has no numpy model.  Verifying array operations requires either
reimplementing numpy in \fstar{} or treating it as an oracle.
More importantly, \fstar{}'s effect system doesn't distinguish
$\Nondet$ from $\IO$---both are lumped into the $\mathsf{ALL}$ effect.
\synhopy{}'s finer-grained effect lattice distinguishes ``produces random output''
from ``performs I/O,'' which matters for knowing which properties
can be proved about a function.
\end{remark}

\section{Feature Engineering Verification}

A critical correctness property in ML pipelines is the absence of
\emph{data leakage}: information from the test set must not influence
feature engineering on the training set.

\begin{lstlisting}[style=synhopy]
from sklearn.preprocessing import StandardScaler

@guarantee("scaler is fit only on training data, "
           "no test data leakage")
@effect_contract(reads={"X_train", "X_test"})
def fit_transform_safe(X_train, X_test):
    assume("X_train and X_test are disjoint arrays")
    assume("X_train.shape[1] == X_test.shape[1]")

    scaler = StandardScaler()
    X_train_scaled = scaler.fit_transform(X_train)
    check("scaler.mean_ is computed from X_train only")

    X_test_scaled = scaler.transform(X_test)
    check("scaler was not refit on X_test")

    return X_train_scaled, X_test_scaled
\end{lstlisting}

\begin{definition}[Data leakage freedom]
A feature engineering pipeline $P$ is \emph{leakage-free} with respect
to a train/test split $(D_{\mathrm{train}}, D_{\mathrm{test}})$ if
every statistic computed by $P$ (mean, variance, quantiles, etc.)
depends only on $D_{\mathrm{train}}$.  Formally, for any two test sets
$D_{\mathrm{test}}, D_{\mathrm{test}}'$:
\[
  P(D_{\mathrm{train}}, D_{\mathrm{test}}) \big|_{D_{\mathrm{train}}}
  = P(D_{\mathrm{train}}, D_{\mathrm{test}}') \big|_{D_{\mathrm{train}}}
\]
(the transformed training data is identical regardless of the test set).
\end{definition}

\begin{lstlisting}[style=proof]
-- Proof for leakage freedom:
have h_fit: scaler.fit_transform(X_train) computes
    mean_ and scale_ from X_train only
    by apply sklearn_scaler_fit_spec  -- library axiom
have h_no_refit: scaler.transform(X_test) uses
    the previously fitted parameters without updating
    by apply sklearn_scaler_transform_spec
have h_disjoint: X_train and X_test are disjoint
    by assume
conclude: no test data influences training features
    by structural with h_fit, h_no_refit, h_disjoint
\end{lstlisting}

\section{Cross-Validation Correctness}

Cross-validation must ensure that folds are non-overlapping and
collectively exhaustive:

\begin{lstlisting}[style=synhopy]
from sklearn.model_selection import KFold

@guarantee("folds partition the data with no overlap")
def verified_kfold(X, y, n_splits=5):
    assume("len(X) == len(y)")
    assume("n_splits > 1 and n_splits <= len(X)")

    kf = KFold(n_splits=n_splits, shuffle=True,
               random_state=42)
    all_test_indices = []

    for fold, (train_idx, test_idx) in \
            enumerate(kf.split(X)):
        check("len(set(train_idx) & set(test_idx)) == 0")
        check("len(train_idx) + len(test_idx) == len(X)")
        all_test_indices.extend(test_idx)

    check("sorted(all_test_indices) == list(range(len(X)))")
    return list(kf.split(X))
\end{lstlisting}

\begin{lstlisting}[style=proof]
-- Proof for cross-validation correctness:
have h_kfold: KFold(n_splits) produces n_splits folds
    by apply sklearn_kfold_spec  -- library axiom
have h_disjoint: forall fold. train_idx and test_idx
    are disjoint within each fold
    by apply sklearn_kfold_disjoint  -- library axiom
have h_complete: union of all test_idx over folds
    equals {0, ..., len(X)-1}
    by apply sklearn_kfold_partition  -- library axiom
have h_size: forall fold.
    len(train_idx) + len(test_idx) == len(X)
    by omega with h_disjoint, h_complete
conclude: folds partition the data
    by structural with h_disjoint, h_complete
\end{lstlisting}

\section{Model Serialization Roundtrip}

Saving and loading models must preserve behavior.  \synhopy{} verifies
the serialization roundtrip:

\begin{lstlisting}[style=synhopy]
import joblib
import numpy as np

@guarantee("loaded model produces identical predictions
            to the original model")
@effect_contract(reads={"filesystem"},
                 mutates={"filesystem"})
def verify_serialization_roundtrip(model, X_test,
                                   filepath):
    assume("model is a fitted sklearn estimator")
    assume("X_test has correct feature dimensions")

    predictions_before = model.predict(X_test)
    joblib.dump(model, filepath)
    loaded_model = joblib.load(filepath)
    predictions_after = loaded_model.predict(X_test)

    check("np.array_equal(predictions_before, "
          "predictions_after)")
    check("type(loaded_model) == type(model)")

    return True
\end{lstlisting}

\begin{lstlisting}[style=proof]
-- Proof for serialization roundtrip:
have h_dump: joblib.dump(model, filepath) serializes
    model to disk
    by apply joblib_dump_spec  -- library axiom
have h_load: joblib.load(filepath) deserializes to
    an object equal to the original
    by apply joblib_roundtrip_spec  -- library axiom
have h_type: type(loaded_model) == type(model)
    by h_load  -- type preserved by roundtrip
have h_predict: loaded_model.predict(X) == model.predict(X)
    for all X, since loaded_model == model
    by apply sklearn_predict_deterministic with h_load
conclude: predictions are identical
    by structural with h_predict
\end{lstlisting}

\section{Prediction Pipeline Type Safety}

ML prediction pipelines must ensure that features are provided in the
correct order and with the correct shape:

\begin{lstlisting}[style=synhopy]
@guarantee("prediction input matches training schema")
def safe_predict(model, feature_names, input_data: dict):
    assume("model was trained with feature_names")
    assume("input_data is a dict with string keys")

    # Verify all required features are present
    for feat in feature_names:
        check(f"'{feat}' in input_data")

    # Build feature vector in correct order
    X = np.array([[input_data[f] for f in feature_names]])
    check("X.shape == (1, len(feature_names))")
    check("X.shape[1] == model.n_features_in_")

    prediction = model.predict(X)
    check("prediction.shape == (1,)")

    return prediction[0]
\end{lstlisting}

\begin{lstlisting}[style=proof]
-- Proof for prediction type safety:
have h_features: forall f in feature_names.
    f in input_data
    by check  -- verified for each feature
have h_order: X is constructed with features in
    the same order as feature_names
    by structural -- list comprehension preserves order
have h_shape: X.shape == (1, len(feature_names))
    by apply numpy_array_shape_spec
have h_match: X.shape[1] == model.n_features_in_
    by check with h_shape
have h_predict: model.predict(X) returns (1,) array
    when X.shape == (1, n_features)
    by apply sklearn_predict_shape_spec
conclude: prediction has correct shape and feature order
    by structural with h_order, h_match, h_predict
\end{lstlisting}

\section{Hyperparameter Search Space Verification}

Hyperparameter search must only explore valid configurations:

\begin{lstlisting}[style=synhopy]
from sklearn.model_selection import GridSearchCV
from sklearn.ensemble import RandomForestClassifier

@guarantee("all hyperparameter combinations are valid
            for the estimator")
def verified_grid_search(X, y):
    assume("X and y have compatible shapes")

    param_grid = {
        'n_estimators': [50, 100, 200],
        'max_depth': [3, 5, 10, None],
        'min_samples_split': [2, 5, 10],
        'min_samples_leaf': [1, 2, 4]
    }

    # Verify all param names are valid
    valid_params = RandomForestClassifier().get_params()
    for key in param_grid:
        check(f"'{key}' in valid_params")

    # Verify value ranges
    for n in param_grid['n_estimators']:
        check(f"{n} > 0")
    for d in param_grid['max_depth']:
        check(f"{d} is None or {d} > 0")
    for s in param_grid['min_samples_split']:
        check(f"{s} >= 2")
    for l in param_grid['min_samples_leaf']:
        check(f"{l} >= 1")

    grid = GridSearchCV(
        RandomForestClassifier(random_state=42),
        param_grid, cv=5, scoring='accuracy')
    grid.fit(X, y)

    check("grid.best_params_ is in param_grid")
    return grid.best_estimator_
\end{lstlisting}

\begin{lstlisting}[style=proof]
-- Proof for hyperparameter validity:
have h_params: all keys in param_grid are valid
    RandomForestClassifier parameters
    by check  -- verified against get_params()
have h_estimators: forall n in [50, 100, 200]. n > 0
    by omega
have h_depth: forall d in [3, 5, 10, None].
    d is None or d > 0
    by omega with None case
have h_split: forall s in [2, 5, 10]. s >= 2
    by omega
have h_leaf: forall l in [1, 2, 4]. l >= 1
    by omega
have h_grid: GridSearchCV explores all combinations
    in param_grid
    by apply sklearn_gridsearch_spec  -- library axiom
have h_best: grid.best_params_ is a valid combination
    from param_grid
    by apply sklearn_gridsearch_best_spec
conclude: all configurations are valid
    by structural with h_params, h_estimators,
       h_depth, h_split, h_leaf
\end{lstlisting}

\section{Worked Example: Verifying a scikit-learn Pipeline}

We present a complete, end-to-end verified ML pipeline using
scikit-learn's \texttt{Pipeline} abstraction:

\begin{lstlisting}[style=synhopy]
from sklearn.pipeline import Pipeline
from sklearn.preprocessing import StandardScaler
from sklearn.decomposition import PCA
from sklearn.linear_model import LogisticRegression
from sklearn.model_selection import train_test_split
import numpy as np

@guarantee("pipeline trains and predicts with correct "
           "shapes, no data leakage, and valid output")
@effect_contract(reads={"X", "y"}, pure_transform=True)
def verified_ml_pipeline(X, y):
    assume("X is a 2D array with no NaN values")
    assume("y is a 1D array of binary labels (0 or 1)")
    assume("X.shape[0] == len(y)")
    assume("X.shape[0] >= 10")

    # Split with stratification
    X_train, X_test, y_train, y_test = train_test_split(
        X, y, test_size=0.2, stratify=y, random_state=42)
    check("len(X_train) + len(X_test) == len(X)")
    check("X_train.shape[1] == X_test.shape[1]")

    # Build pipeline
    pipe = Pipeline([
        ('scaler', StandardScaler()),
        ('pca', PCA(n_components=min(5, X.shape[1]))),
        ('clf', LogisticRegression(max_iter=1000))
    ])

    # Fit on training data only (leakage-free)
    pipe.fit(X_train, y_train)
    check("pipe.named_steps['scaler'].mean_ is not None")
    check("pipe.named_steps['pca'].components_ "
          "is not None")

    # Predict on test data
    y_pred = pipe.predict(X_test)
    check("y_pred.shape == y_test.shape")
    check("set(np.unique(y_pred)) <= {0, 1}")

    # Probabilities sum to 1
    y_proba = pipe.predict_proba(X_test)
    check("y_proba.shape == (len(X_test), 2)")
    check("np.allclose(y_proba.sum(axis=1), 1.0)")

    accuracy = (y_pred == y_test).mean()
    check("0.0 <= accuracy <= 1.0")

    return {
        "accuracy": float(accuracy),
        "predictions": y_pred,
        "probabilities": y_proba
    }
\end{lstlisting}

\begin{lstlisting}[style=proof]
-- Full pipeline proof:

-- Split verification
have h_split: train_test_split preserves total count
    by apply sklearn_split_spec  -- library axiom
have h_cols: feature columns preserved in split
    by apply sklearn_split_shape_spec

-- Leakage freedom
have h_fit_train: pipe.fit(X_train, y_train) fits
    all stages (scaler, pca, clf) on X_train only
    by apply sklearn_pipeline_fit_spec
have h_no_leak: X_test is not seen during fitting
    by structural with h_fit_train
    -- X_test does not appear in fit() call

-- Shape propagation through pipeline
have h_scaler_shape: StandardScaler preserves
    (n_samples, n_features) shape
    by apply sklearn_scaler_shape_spec
have h_pca_shape: PCA(n_components=k) maps
    (n, d) to (n, k) where k <= d
    by apply sklearn_pca_shape_spec
have h_clf_shape: LogisticRegression.predict maps
    (n, k) to (n,)
    by apply sklearn_logreg_predict_spec

-- Prediction validity
have h_pred_shape: y_pred.shape == y_test.shape
    by trans(h_scaler_shape,
        trans(h_pca_shape, h_clf_shape))
have h_pred_labels: predictions are in {0, 1}
    for binary classification
    by apply sklearn_logreg_binary_spec

-- Probability verification
have h_proba_shape: predict_proba returns
    (n_samples, n_classes) array
    by apply sklearn_predict_proba_shape_spec
have h_proba_sum: rows of predict_proba sum to 1.0
    by apply sklearn_predict_proba_sum_spec
    -- probability axiom

-- Accuracy bounds
have h_accuracy: 0 <= mean(bool_array) <= 1
    by omega  -- mean of booleans is in [0, 1]

conclude: pipeline is correct end-to-end
    by structural with h_split, h_no_leak,
       h_pred_shape, h_pred_labels,
       h_proba_sum, h_accuracy
\end{lstlisting}

This example demonstrates how \synhopy{} verifies an entire ML pipeline
by composing proofs about individual stages.  The key verified properties
are: (1) no data leakage between train and test sets, (2) shape
compatibility across pipeline stages, (3) prediction labels are within the
valid set, and (4) predicted probabilities are well-formed.  The trust
level is $\mathsf{AXIOM\_TRUSTED}$, with dependencies on scikit-learn
library axioms clearly enumerated.


% ──────────────────────────────────────────────────────────────────
\chapter{Verifying a Concurrent Task Queue}
\label{ch:case-queue}

\section{The Code}

\begin{lstlisting}[style=synhopy]
import asyncio
from collections import deque

class TaskQueue:
    def __init__(self, max_workers: int = 4):
        self.queue = deque()
        self.max_workers = max_workers
        self.active = 0
        self._lock = asyncio.Lock()

    @guarantee("adds task to queue; queue grows by 1")
    @effect_contract(mutates={"self.queue"})
    async def enqueue(self, task):
        async with self._lock:
            self.queue.append(task)

    @guarantee("processes all queued tasks with <= max_workers concurrent")
    @effect_contract(
        mutates={"self.queue", "self.active"},
        reads={"self.max_workers"})
    async def run_all(self):
        while self.queue or self.active > 0:
            while self.queue and self.active < self.max_workers:
                task = self.queue.popleft()
                self.active += 1
                asyncio.create_task(self._run_task(task))
            await asyncio.sleep(0.01)

    async def _run_task(self, task):
        try:
            await task()
        finally:
            async with self._lock:
                self.active -= 1
\end{lstlisting}

\section{Verification: Concurrency Invariants}

\begin{lstlisting}[style=proof]
-- Invariant: self.active <= self.max_workers at all times
have h_init: active = 0 <= max_workers (since max_workers >= 1)
    by omega

have h_guard: the inner while loop checks
    active < max_workers before incrementing
    by structural -- the condition is syntactically present

have h_inc: active increases by 1 only when active < max_workers
    by rw [h_guard]

have h_dec: _run_task decreases active by 1 in finally block
    by structural -- finally ensures decrement on all paths

-- Induction on number of scheduling steps:
have h_inv: forall t. active(t) <= max_workers
    by nat_induction t
    base: active(0) = 0 <= max_workers by h_init
    step: cases action at time t+1:
      case enqueue: active unchanged by structural
      case dispatch: active(t+1) = active(t) + 1 <= max_workers
          by h_inc (guard ensures active(t) < max_workers)
      case complete: active(t+1) = active(t) - 1 <= max_workers
          by h_dec, omega with IH

conclude: active <= max_workers is an invariant
    by h_inv
\end{lstlisting}

\begin{lstlisting}[style=proof]
-- Liveness: all enqueued tasks eventually complete
-- (assuming each task() terminates)
have h_drain: if queue is non-empty and active < max_workers,
    a task is dequeued and started
    by structural -- inner while loop condition
have h_finish: each started task completes (by assumption)
    and decrements active
    by apply task_terminates_axiom, h_dec
have h_progress: at each step, either queue shrinks or
    active decreases (or both)
    by cases with h_drain, h_finish
conclude: eventually queue = [] and active = 0
    by well_founded_induction (len(queue) + active)
    with h_progress
\end{lstlisting}

\begin{remark}[\fstar{} comparison]
Verifying concurrent invariants about an async task queue in \fstar{} would
require:
\begin{enumerate}
  \item A model of Python's \texttt{asyncio} event loop
  \item A model of \texttt{asyncio.Lock} with mutual exclusion guarantees
  \item Monotonic state reasoning for the \texttt{active} counter
  \item A custom concurrent effect type for interleaved coroutine execution
\end{enumerate}
This is approximately 500+ lines of \fstar{} infrastructure before writing a
single line of application verification.  \synhopy{} handles the same
verification with the existing effect system ($\Mutates + \mathsf{Suspend}$),
path-based concurrency semantics, and structural reasoning about the
guard conditions.
\end{remark}

\section{Priority Queue Semantics}

Real task queues support priorities.  We extend the basic queue with a
priority mechanism and verify that higher-priority tasks execute first.

\begin{lstlisting}[style=synhopy]
import heapq

class PriorityTaskQueue:
    def __init__(self, max_workers: int = 4):
        self.heap = []  # min-heap of (priority, seq, task)
        self.seq = 0
        self.max_workers = max_workers
        self.active = 0
        self._lock = asyncio.Lock()

    @guarantee("adds task at given priority; heap property maintained")
    @effect_contract(mutates={"self.heap", "self.seq"})
    async def enqueue(self, task, priority: int = 0):
        async with self._lock:
            heapq.heappush(self.heap, (priority, self.seq, task))
            self.seq += 1

    @guarantee("returns highest-priority task (lowest number)")
    @effect_contract(mutates={"self.heap"})
    async def dequeue(self):
        async with self._lock:
            return heapq.heappop(self.heap)[2]
\end{lstlisting}

\begin{lstlisting}[style=proof]
-- Heap invariant: after every enqueue/dequeue,
--   self.heap satisfies the min-heap property
have h_push: heappush maintains the min-heap property
    by library_axiom("heapq.heappush preserves min-heap")
have h_pop: heappop removes the minimum element and
    restores the min-heap property
    by library_axiom("heapq.heappop returns min, preserves heap")
have h_order: for tasks t1, t2 with priority p1 < p2,
    dequeue returns t1 before t2
    (assuming no intervening enqueue)
    by apply h_pop -- min-heap always yields smallest priority
conclude: priority ordering is respected
    by h_order
\end{lstlisting}

\begin{remark}
The sequence number \lstinline[style=python]{self.seq} serves as a
tie-breaker for equal priorities, enforcing FIFO order within the same
priority level.  The verification of this property follows from the
lexicographic ordering of tuples in Python's comparison semantics:
$(p_1, s_1, \_) < (p_2, s_2, \_)$ iff $p_1 < p_2$ or ($p_1 = p_2$ and
$s_1 < s_2$).
\end{remark}

\section{Worker Pool Scaling}

Production task queues dynamically scale their worker pools based on load.
We verify that adding and removing workers preserves the concurrency
invariant.

\begin{lstlisting}[style=synhopy]
class ScalableTaskQueue(TaskQueue):
    @guarantee("increases max_workers; active <= new max_workers")
    @effect_contract(mutates={"self.max_workers"})
    async def scale_up(self, additional: int):
        assume("additional > 0")
        async with self._lock:
            self.max_workers += additional

    @guarantee("decreases max_workers; no running tasks killed")
    @effect_contract(mutates={"self.max_workers"})
    async def scale_down(self, remove: int):
        assume("remove > 0")
        assume("self.max_workers - remove >= 1")
        async with self._lock:
            self.max_workers -= remove
            # Do NOT cancel active tasks; just prevent new dispatches
\end{lstlisting}

\begin{lstlisting}[style=proof]
-- Invariant preservation under scaling
have h_up: after scale_up(n),
    max_workers' = max_workers + n > max_workers >= active
    by omega with precondition (additional > 0)
have h_down: after scale_down(n),
    max_workers' = max_workers - n >= 1
    by assumption -- from assume annotation
have h_no_kill: scale_down does not modify self.active
    by structural -- only max_workers is mutated
have h_eventual: if active > max_workers' after scale_down,
    no new tasks are dispatched (guard: active < max_workers fails),
    so active decreases monotonically as running tasks complete
    by structural with h_no_kill
conclude: the invariant active <= max_workers is eventually restored
    by well_founded_induction (active - max_workers')
    with h_eventual
\end{lstlisting}

\begin{remark}
Note the distinction between \emph{immediate} and \emph{eventual}
invariant restoration.  After \lstinline[style=python]{scale_down}, we
may temporarily have $\mathtt{active} > \mathtt{max\_workers}$, but
this is safe because no new tasks will be dispatched, and the invariant
is restored once enough in-flight tasks complete.  \synhopy{}'s effect
system tracks this as a \emph{transient violation} with a well-founded
termination argument.
\end{remark}

\section{Task Dependency DAGs}

Many workloads have inter-task dependencies: task $B$ cannot start until
task $A$ completes.  We model dependencies as a directed acyclic graph
(DAG) and verify that tasks execute in topological order.

\begin{lstlisting}[style=synhopy]
class DAGTaskQueue:
    def __init__(self):
        self.tasks = {}       # id -> task callable
        self.deps = {}        # id -> set of dependency ids
        self.completed = set()
        self._lock = asyncio.Lock()

    @guarantee("registers task with dependencies")
    @effect_contract(mutates={"self.tasks", "self.deps"})
    def add_task(self, task_id: str, task, depends_on: list = None):
        self.tasks[task_id] = task
        self.deps[task_id] = set(depends_on or [])

    @guarantee("returns tasks whose dependencies are all satisfied")
    @effect_contract(reads={"self.deps", "self.completed"})
    def ready_tasks(self):
        return [
            tid for tid, deps in self.deps.items()
            if tid not in self.completed and deps <= self.completed
        ]

    @guarantee("executes all tasks in topological order")
    @effect_contract(mutates={"self.completed"})
    async def run_all(self):
        while len(self.completed) < len(self.tasks):
            ready = self.ready_tasks()
            check("ready is non-empty if uncompleted tasks exist "
                  "and graph is a DAG")
            await asyncio.gather(
                *[self._run_and_mark(tid) for tid in ready])

    async def _run_and_mark(self, task_id):
        await self.tasks[task_id]()
        async with self._lock:
            self.completed.add(task_id)
\end{lstlisting}

\begin{lstlisting}[style=proof]
-- Topological order property
have h_ready: ready_tasks() returns tid only when
    deps[tid] is a subset of completed
    by unfold ready_tasks
have h_exec: _run_and_mark(tid) executes task only after
    all deps are in completed
    by rw [h_ready] -- called only for ready tasks
have h_topo: the execution order is a topological sort
    of the dependency DAG
    by induction on |completed|
    base: initially completed = {}, ready = sources (no deps)
    step: after executing ready set R,
      completed' = completed + R,
      new ready tasks have deps in completed'
      by h_ready

-- Termination (requires acyclicity)
have h_dag: the dependency graph is acyclic
    by assumption -- precondition of the system
have h_progress: at each iteration, |ready| >= 1
    (since a DAG with uncompleted nodes has at least one source)
    by apply dag_has_source with h_dag
have h_grow: |completed| strictly increases each iteration
    by h_progress -- at least one task completes
conclude: run_all terminates and executes in topological order
    by well_founded_induction (|tasks| - |completed|)
    with h_grow, h_topo
\end{lstlisting}

\section{Deadlock Detection via Cycle Detection}

If the dependency graph contains a cycle, the task queue will deadlock:
no task in the cycle can become ready.  We verify a cycle detection
mechanism that prevents this.

\begin{lstlisting}[style=synhopy]
class SafeDAGTaskQueue(DAGTaskQueue):
    @guarantee("raises ValueError if adding task would create a cycle")
    @effect_contract(reads={"self.deps"})
    def add_task(self, task_id: str, task, depends_on: list = None):
        # Tentatively add, then check for cycles
        new_deps = dict(self.deps)
        new_deps[task_id] = set(depends_on or [])
        check("no cycle exists in new_deps")
        if self._has_cycle(new_deps):
            raise ValueError(f"Adding {task_id} would create a cycle")
        super().add_task(task_id, task, depends_on)

    def _has_cycle(self, graph):
        visited, rec_stack = set(), set()
        def dfs(node):
            visited.add(node)
            rec_stack.add(node)
            for dep in graph.get(node, set()):
                if dep not in visited:
                    if dfs(dep): return True
                elif dep in rec_stack:
                    return True
            rec_stack.remove(node)
            return False
        return any(dfs(n) for n in graph if n not in visited)
\end{lstlisting}

\begin{lstlisting}[style=proof]
-- Correctness of cycle detection
have h_dfs: _has_cycle returns True iff the graph contains a
    directed cycle (standard DFS cycle detection)
    by library_axiom("DFS cycle detection is correct")
have h_reject: if add_task would create a cycle,
    ValueError is raised and self.deps is not modified
    by structural -- exception raised before super().add_task
have h_accept: if add_task succeeds,
    self.deps has no cycle
    by contrapositive of h_reject, h_dfs

-- Deadlock freedom follows from acyclicity
have h_no_deadlock: if self.deps is acyclic,
    run_all cannot deadlock
    by apply dag_has_source -- acyclic graph always has a source
conclude: SafeDAGTaskQueue.run_all is deadlock-free
    by induction on add_task calls with h_accept, h_no_deadlock
\end{lstlisting}

\section{Graceful Shutdown Verification}

A production task queue must support graceful shutdown: stop accepting
new tasks but allow all in-flight tasks to complete.

\begin{lstlisting}[style=synhopy]
class GracefulTaskQueue(TaskQueue):
    def __init__(self, max_workers=4):
        super().__init__(max_workers)
        self.shutting_down = False
        self.shutdown_event = asyncio.Event()

    @guarantee("after shutdown, no new tasks accepted; "
               "all active tasks complete before returning")
    @effect_contract(mutates={"self.shutting_down"})
    async def shutdown(self):
        async with self._lock:
            self.shutting_down = True
        # Wait for all active tasks to finish
        while self.active > 0:
            await asyncio.sleep(0.01)
        self.shutdown_event.set()

    @guarantee("enqueue fails after shutdown initiated")
    @effect_contract(mutates={"self.queue"})
    async def enqueue(self, task):
        if self.shutting_down:
            raise RuntimeError("Queue is shutting down")
        await super().enqueue(task)
\end{lstlisting}

\begin{lstlisting}[style=proof]
-- Property 1: No new tasks after shutdown
have h_flag: shutdown() sets self.shutting_down = True
    by structural
have h_reject: enqueue checks shutting_down first;
    if True, raises RuntimeError
    by structural
have h_no_new: after shutdown() is called, no new tasks
    are added to the queue
    by h_flag, h_reject -- flag is set, all future enqueues fail

-- Property 2: All in-flight tasks complete
have h_wait: shutdown() loops while self.active > 0
    by structural
have h_dec: each active task eventually completes and
    decrements self.active (by the liveness proof above)
    by apply task_terminates_axiom
have h_drain: self.active eventually reaches 0
    by well_founded_induction (self.active) with h_dec
have h_signal: shutdown_event is set only after active = 0
    by h_wait, h_drain

-- Combined guarantee
conclude: after shutdown() returns,
    (1) no new tasks in queue, (2) all active tasks completed
    by h_no_new, h_signal
\end{lstlisting}

\section{Worked Example: Distributed Task Queue with Redis}

We now present a more realistic example: a distributed task queue backed
by Redis, where multiple workers across different processes consume tasks
from a shared queue.

\begin{lstlisting}[style=synhopy]
import redis
import json
import uuid

class RedisTaskQueue:
    def __init__(self, redis_url: str, queue_name: str = "tasks"):
        self.redis = redis.Redis.from_url(redis_url)
        self.queue_name = queue_name
        self.processing_set = f"{queue_name}:processing"

    @guarantee("task is atomically moved from queue to processing set")
    @effect_contract(mutates={"redis.queue", "redis.processing_set"})
    def dequeue(self):
        # RPOPLPUSH is atomic in Redis
        raw = self.redis.rpoplpush(
            self.queue_name, self.processing_set)
        if raw is None:
            return None
        return json.loads(raw)

    @guarantee("task is enqueued; queue length increases by 1")
    @effect_contract(mutates={"redis.queue"})
    def enqueue(self, task_data: dict):
        task_id = str(uuid.uuid4())
        payload = json.dumps({"id": task_id, **task_data})
        self.redis.lpush(self.queue_name, payload)
        return task_id

    @guarantee("task removed from processing set after completion")
    @effect_contract(mutates={"redis.processing_set"})
    def acknowledge(self, task_data: dict):
        self.redis.lrem(self.processing_set, 1,
                        json.dumps(task_data))
\end{lstlisting}

\begin{lstlisting}[style=proof]
-- Atomicity of dequeue
have h_atomic: rpoplpush is an atomic Redis operation
    that removes from source and pushes to destination
    by library_axiom("redis.rpoplpush is atomic")

-- No task loss: a task is always in exactly one of
-- {queue, processing_set, completed}
have h_enqueue: after enqueue(t), t is in queue
    by structural with library_axiom("redis.lpush adds to list")
have h_dequeue: after dequeue(), if result is not None,
    t is removed from queue AND added to processing_set
    by h_atomic
have h_ack: after acknowledge(t),
    t is removed from processing_set
    by structural with library_axiom("redis.lrem removes element")

-- Conservation: |queue| + |processing| is non-increasing
-- (decreases by 1 on acknowledge, unchanged on dequeue)
have h_conserve: for any task t, at all times t is in
    exactly one of: queue, processing_set, or completed
    by cases on task lifecycle:
      case enqueued: t in queue by h_enqueue
      case dequeued: t in processing_set by h_dequeue
      case acknowledged: t removed from processing_set by h_ack

conclude: no tasks are lost during processing
    by h_conserve
\end{lstlisting}

\begin{remark}[\fstar{} comparison for distributed verification]
Verifying a distributed task queue in \fstar{} would additionally require:
\begin{enumerate}
  \item A formal model of Redis's concurrency semantics and atomicity guarantees
  \item A distributed state monad tracking state across multiple processes
  \item Network failure modeling (what happens if a worker crashes after
    dequeue but before acknowledge?)
  \item Serialization/deserialization correctness proofs for JSON encoding
\end{enumerate}
\synhopy{}'s library axiom system (Chapter~\ref{ch:library-axioms}) handles
the Redis semantics via trusted axioms about \lstinline[style=python]{rpoplpush}
atomicity and \lstinline[style=python]{lpush}/\lstinline[style=python]{lrem}
correctness.  The trust level of the overall verification is bounded by
these library axioms---the proof is sound conditional on Redis behaving
as documented.
\end{remark}


% ──────────────────────────────────────────────────────────────────
\chapter{Verifying Design Patterns}
\label{ch:case-patterns}

\section{The Observer Pattern}

\begin{lstlisting}[style=synhopy]
class EventEmitter:
    def __init__(self):
        self._listeners = {}

    @guarantee("registers callback for event; listener count increases by 1")
    @effect_contract(mutates={"self._listeners"})
    def on(self, event: str, callback):
        if event not in self._listeners:
            self._listeners[event] = []
        self._listeners[event].append(callback)

    @guarantee("calls all registered callbacks for event")
    @effect_contract(reads={"self._listeners"})
    def emit(self, event: str, *args, **kwargs):
        for cb in self._listeners.get(event, []):
            cb(*args, **kwargs)
\end{lstlisting}

\section{Verifying the Observer Contract}

\begin{lstlisting}[style=proof]
-- Property: emit(event) calls exactly the registered listeners
have h_on: after on(event, cb), cb in self._listeners[event]
    by unfold on
have h_emit: emit(event) iterates self._listeners.get(event, [])
    by unfold emit
have h_all: every cb in _listeners[event] is called
    by rw [h_emit] -- for-loop calls each element
have h_only: no other callbacks are called
    by structural -- .get(event, []) returns only event's list

-- Connecting on() and emit():
have h_contract: if on(e, cb) then emit(e) calls cb
    by trans(h_on, h_all)  -- cb is registered, all registered are called

conclude: the observer contract holds
    by structural with h_contract, h_only
\end{lstlisting}

\section{The Strategy Pattern}

\begin{lstlisting}[style=synhopy]
@guarantee("all strategies produce valid output for valid input")
class Sorter:
    def __init__(self, strategy):
        assume("strategy is callable: list -> list")
        assume("strategy returns sorted permutation of input")
        self.strategy = strategy

    @guarantee("returns sorted version of data")
    def sort(self, data: list) -> list:
        return self.strategy(data)
\end{lstlisting}

\begin{lstlisting}[style=proof]
-- Duck-type equivalence of strategies
-- Any callable with signature list -> sorted list satisfies the contract
have h_duck: forall s1, s2 satisfying the strategy protocol.
    Sorter(s1).sort(data) is a sorted permutation of data
    AND Sorter(s2).sort(data) is a sorted permutation of data
    by assumption -- from assume() annotations

-- By Specification Transfer (Theorem 4.5):
have h_transfer: Spec(Sorter(s1)) => Spec(Sorter(s2))
    by transport along duck_equiv(s1, s2)

conclude: all valid strategies produce correct results
    by h_duck, h_transfer
\end{lstlisting}

This example illustrates Theorem~\ref{thm:spec-transfer} in practice: once
we prove that \texttt{Sorter} is correct for \emph{one} strategy, the proof
transports to \emph{all} duck-type-equivalent strategies via univalence.

\begin{remark}[\fstar{} comparison for the Strategy pattern]
In \fstar{}, the Strategy pattern would be modeled via an explicit type class
or a record of functions:
\begin{lstlisting}[style=fstar]
type sort_strategy = list int -> Tot (list int)
val mk_sorter : sort_strategy -> list int -> Tot (list int)
let mk_sorter strategy data = strategy data
\end{lstlisting}
This works for a single function signature, but \fstar{} cannot express
the duck-typing insight: that \emph{any} callable matching the protocol
is a valid strategy.  Swapping strategies requires proving type equality
at each call site.  \synhopy{}'s transport along duck-type equivalence
handles this once for all strategies simultaneously.
\end{remark}

\section{The Factory Pattern}

The Factory pattern creates objects without exposing instantiation logic.
The verification challenge is ensuring that every created object satisfies
the interface contract.

\begin{lstlisting}[style=synhopy]
class Shape:
    @guarantee("returns the area of the shape")
    def area(self) -> float:
        raise NotImplementedError

class Circle(Shape):
    def __init__(self, radius: float):
        assume("radius > 0")
        self.radius = radius

    @guarantee("returns pi * radius^2")
    def area(self) -> float:
        import math
        return math.pi * self.radius ** 2

class Rectangle(Shape):
    def __init__(self, width: float, height: float):
        assume("width > 0 and height > 0")
        self.width = width
        self.height = height

    @guarantee("returns width * height")
    def area(self) -> float:
        return self.width * self.height

@guarantee("returns a Shape instance that satisfies the Shape protocol")
def shape_factory(kind: str, **kwargs) -> Shape:
    if kind == "circle":
        return Circle(kwargs["radius"])
    elif kind == "rectangle":
        return Rectangle(kwargs["width"], kwargs["height"])
    else:
        raise ValueError(f"Unknown shape: {kind}")
\end{lstlisting}

\begin{lstlisting}[style=proof]
-- Factory contract: returned object satisfies Shape protocol
have h_circle: Circle has method area() returning float
    by structural -- area is defined in Circle
have h_rect: Rectangle has method area() returning float
    by structural -- area is defined in Rectangle
have h_cases: shape_factory returns Circle or Rectangle
    (or raises ValueError)
    by cases on kind:
      case "circle": returns Circle by structural
      case "rectangle": returns Rectangle by structural
      case _: raises ValueError by structural

-- Duck-type check: all returned objects satisfy Shape protocol
have h_protocol: Shape protocol = {area : () -> float}
    by unfold Shape
have h_circle_duck: Circle models Shape protocol
    by h_circle, h_protocol
have h_rect_duck: Rectangle models Shape protocol
    by h_rect, h_protocol

conclude: shape_factory returns objects satisfying Shape protocol
    by cases with h_cases, h_circle_duck, h_rect_duck
\end{lstlisting}

\begin{remark}[\fstar{} comparison for the Factory pattern]
In \fstar{}, the factory would need an explicit sum type or
an existential type to unify the return values:
\begin{lstlisting}[style=fstar]
type shape =
  | MkCircle : radius:float{radius > 0.0} -> shape
  | MkRect : w:float{w > 0.0} -> h:float{h > 0.0} -> shape

val area : shape -> Tot float
let area s = match s with
  | MkCircle r -> 3.14159 *. r *. r
  | MkRect w h -> w *. h
\end{lstlisting}
This requires defining the sum type upfront and modifying it whenever
a new shape is added.  \synhopy{}'s duck-type approach allows new
shapes to be added without modifying the factory's verification---the
duck-type equivalence handles extensibility automatically.
\end{remark}

\section{The Decorator Pattern (GoF)}

The GoF Decorator pattern (not Python's \lstinline[style=python]{@decorator}
syntax) wraps an object to extend its behavior.  The key verification
challenge is proving \emph{behavioral equivalence}: the decorated object
satisfies the same interface contract as the original, plus additional
behavior.

\begin{lstlisting}[style=synhopy]
class DataSource:
    @guarantee("reads data from the source")
    def read(self) -> str:
        raise NotImplementedError

    @guarantee("writes data to the source")
    def write(self, data: str):
        raise NotImplementedError

class FileDataSource(DataSource):
    def __init__(self, filename: str):
        self.filename = filename

    @guarantee("returns file contents as string")
    @effect_contract(reads={"filesystem"})
    def read(self) -> str:
        with open(self.filename) as f:
            return f.read()

    @guarantee("writes string to file")
    @effect_contract(mutates={"filesystem"})
    def write(self, data: str):
        with open(self.filename, 'w') as f:
            f.write(data)

class EncryptionDecorator(DataSource):
    def __init__(self, source: DataSource, key: str):
        assume("source satisfies DataSource protocol")
        self.source = source
        self.key = key

    @guarantee("returns decrypted data from underlying source")
    @effect_contract(reads={"filesystem"})
    def read(self) -> str:
        encrypted = self.source.read()
        return self._decrypt(encrypted, self.key)

    @guarantee("encrypts data and writes to underlying source")
    @effect_contract(mutates={"filesystem"})
    def write(self, data: str):
        encrypted = self._encrypt(data, self.key)
        self.source.write(encrypted)

    def _encrypt(self, data, key):
        # XOR-based encryption for illustration
        return ''.join(chr(ord(c) ^ ord(key[i % len(key)]))
                       for i, c in enumerate(data))

    def _decrypt(self, data, key):
        return self._encrypt(data, key)  # XOR is self-inverse
\end{lstlisting}

\begin{lstlisting}[style=proof]
-- Behavioral equivalence: read(write(x)) = x
-- For FileDataSource:
have h_file_rw: FileDataSource.read(FileDataSource.write(x)) = x
    by structural -- file write then read returns same content

-- For EncryptionDecorator:
have h_xor_inv: _decrypt(_encrypt(x, key), key) = x
    by ext -- XOR is self-inverse: (c ^ k) ^ k = c
have h_enc_write: EncryptionDecorator.write(x)
    calls source.write(_encrypt(x, key))
    by unfold write
have h_enc_read: EncryptionDecorator.read()
    returns _decrypt(source.read(), key)
    by unfold read
have h_enc_rw: EncryptionDecorator.read(
    EncryptionDecorator.write(x))
    = _decrypt(_encrypt(x, key), key)
    = x
    by rw [h_enc_write, h_enc_read, h_file_rw, h_xor_inv]

-- Duck-type equivalence: EncryptionDecorator satisfies
-- DataSource protocol with the same read/write contract
have h_duck: EncryptionDecorator models DataSource protocol
    by structural -- has read() -> str and write(str)
have h_contract_preserved: read(write(x)) = x
    for both FileDataSource and EncryptionDecorator
    by h_file_rw, h_enc_rw

conclude: GoF Decorator preserves behavioral equivalence
    by h_duck, h_contract_preserved
\end{lstlisting}

\begin{remark}[\fstar{} comparison for the GoF Decorator pattern]
\fstar{} would model this via a type class with a roundtrip law:
\begin{lstlisting}[style=fstar]
val read_write_roundtrip :
  s:data_source -> x:string ->
  Lemma (ensures (read (write s x) == x))
\end{lstlisting}
Each decorator would require its own roundtrip lemma.  \synhopy{}'s
duck-type equivalence automatically recognizes that any object with
\lstinline[style=python]{read} and \lstinline[style=python]{write}
methods satisfying the roundtrip property is interchangeable with the
original, without explicitly enumerating decorators.
\end{remark}

\section{The Iterator Pattern}

The Iterator pattern provides a standard interface for traversing
collections.  In Python, this is the \lstinline[style=python]{__iter__}
and \lstinline[style=python]{__next__} protocol.  The verification
challenge is proving that lazy iteration produces the same elements
as eager materialization.

\begin{lstlisting}[style=synhopy]
class FilterIterator:
    def __init__(self, iterable, predicate):
        assume("predicate is callable: any -> bool")
        self._iter = iter(iterable)
        self._pred = predicate

    def __iter__(self):
        return self

    @guarantee("returns next element satisfying predicate, "
               "or raises StopIteration")
    def __next__(self):
        while True:
            item = next(self._iter)  # may raise StopIteration
            if self._pred(item):
                return item

@guarantee("yields same elements as [x for x in data if pred(x)]")
def filtered(data, pred):
    return FilterIterator(data, pred)
\end{lstlisting}

\begin{lstlisting}[style=proof]
-- Equivalence: FilterIterator(data, pred) yields the same
-- elements as the list comprehension [x for x in data if pred(x)]

-- Define the eager version:
have h_eager: [x for x in data if pred(x)]
    = filter(pred, data)  -- standard list comprehension semantics
    by structural

-- Define the lazy version:
have h_lazy: FilterIterator.__next__ skips elements where
    pred(item) is False and returns elements where pred(item) is True
    by unfold __next__
have h_stop: FilterIterator raises StopIteration when the
    underlying iterator is exhausted
    by structural -- next(self._iter) raises StopIteration

-- Induction on elements of data:
have h_equiv: list(FilterIterator(data, pred))
    = [x for x in data if pred(x)]
    by induction on data
    base: data = [] => both are []
      (StopIteration immediately, empty list)
    step: data = [hd] + tl =>
      case pred(hd): both include hd, then continue with tl (IH)
      case not pred(hd): both skip hd, continue with tl (IH)

conclude: lazy and eager filtering produce identical results
    by h_equiv
\end{lstlisting}

\begin{remark}[\fstar{} comparison for the Iterator pattern]
\fstar{} has no concept of Python's iterator protocol
(\lstinline[style=python]{__iter__}/\lstinline[style=python]{__next__}).
Expressing lazy iteration in \fstar{} requires either:
\begin{enumerate}
  \item Modeling iterators as state machines with explicit state types, or
  \item Using coinductive types for potentially infinite sequences.
\end{enumerate}
Both approaches require 50+ lines of infrastructure.  \synhopy{}
models iterators as generators (paths in state space), and the
iterator protocol is a structural type that any object with
\lstinline[style=python]{__iter__} and \lstinline[style=python]{__next__}
satisfies by duck typing.
\end{remark}

\section{The Builder Pattern}

The Builder pattern constructs complex objects step by step, enforcing
that certain steps occur in a valid order.  The verification challenge
is proving that the construction invariants hold at each step and that
the final object is fully initialized.

\begin{lstlisting}[style=synhopy]
class QueryBuilder:
    def __init__(self):
        self._table = None
        self._columns = []
        self._conditions = []
        self._limit = None

    @guarantee("sets the table; table is not None after call")
    @effect_contract(mutates={"self._table"})
    def from_table(self, table: str):
        assume("table is a non-empty string")
        self._table = table
        return self

    @guarantee("adds column to selection; columns grows by 1")
    @effect_contract(mutates={"self._columns"})
    def select(self, *columns):
        assume("all columns are non-empty strings")
        self._columns.extend(columns)
        return self

    @guarantee("adds a WHERE condition")
    @effect_contract(mutates={"self._conditions"})
    def where(self, condition: str):
        self._conditions.append(condition)
        return self

    @guarantee("sets result limit to a positive integer")
    @effect_contract(mutates={"self._limit"})
    def limit(self, n: int):
        assume("n > 0")
        self._limit = n
        return self

    @guarantee("returns valid SQL string; requires table to be set")
    @effect_contract(reads={"self._table", "self._columns",
                            "self._conditions", "self._limit"})
    def build(self) -> str:
        check("self._table is not None")
        cols = ', '.join(self._columns) if self._columns else '*'
        sql = f"SELECT {cols} FROM {self._table}"
        if self._conditions:
            sql += " WHERE " + " AND ".join(self._conditions)
        if self._limit:
            sql += f" LIMIT {self._limit}"
        return sql
\end{lstlisting}

\begin{lstlisting}[style=proof]
-- Step-by-step construction invariants
have h_init: after __init__,
    _table = None, _columns = [], _conditions = [], _limit = None
    by structural

have h_from: after from_table(t),
    _table = t (non-None by precondition)
    by structural with assumption (t is non-empty)

have h_select: after select(*cols),
    _columns has grown by len(cols)
    by structural -- extend appends all

have h_where: after where(cond),
    _conditions has grown by 1
    by structural -- append adds one

have h_limit: after limit(n),
    _limit = n > 0
    by structural with assumption (n > 0)

-- Build precondition: _table must be set
have h_build_pre: build() calls check("self._table is not None")
    by structural
have h_build_fail: if _table is None, check raises AssertionError
    by unfold check

-- Valid SQL output
have h_sql: if _table is not None, build() returns a string of
    the form "SELECT cols FROM table [WHERE ...] [LIMIT ...]"
    by structural with h_from, h_select, h_where, h_limit

-- Fluent interface: method chaining preserves builder identity
have h_fluent: from_table, select, where, limit all return self
    by structural -- each returns self
have h_chain: q.from_table(t).select(c).where(w).limit(n)
    returns the same QueryBuilder q with all fields set
    by rw [h_fluent] -- each call mutates and returns self

conclude: build() on a fully configured builder returns valid SQL
    by h_chain, h_sql, h_build_pre
\end{lstlisting}

\begin{remark}[\fstar{} comparison for the Builder pattern]
In \fstar{}, the Builder pattern with step-by-step invariants would
require a \emph{phantom type} tracking which steps have been completed:
\begin{lstlisting}[style=fstar]
type builder_state =
  | NoTable | HasTable | HasColumns | Ready

val from_table : builder NoTable -> string ->
  Tot (builder HasTable)
val select : builder HasTable -> list string ->
  Tot (builder HasColumns)
val build : builder Ready -> Tot string
\end{lstlisting}
This encoding forces a single linear build order.  Python's builder
allows steps in any order (you can call \lstinline[style=python]{where}
before \lstinline[style=python]{select}).  \synhopy{}'s mutation-tracking
effect system naturally handles arbitrary step orders---the only hard
constraint is that \lstinline[style=python]{_table} must be set before
\lstinline[style=python]{build()}, which is verified by the
\lstinline[style=synhopy]{check()} annotation.
\end{remark}


% ══════════════════════════════════════════════════════════════════
% COPY SEMANTICS, OBJECT IDENTITY, AND ALIASING
% ══════════════════════════════════════════════════════════════════

\chapter{Copy Semantics, Object Identity, and Aliasing}
\label{ch:copy-aliasing}

Python's reference semantics are the single greatest source of verification
complexity that \fstar{} cannot model at all.  Every variable binding in Python
is a \emph{reference} to a heap object, not a value.  The distinction between
\lstinline[style=python]{is} (identity) and \lstinline[style=python]{==}
(equality) pervades every program.  This chapter formalizes reference semantics,
aliasing analysis, and copy operations as first-class constructs in the
\synhopy{} type theory.

\section{Object Identity as a Fiber}

\begin{definition}[Object identity type]
Every Python object $o$ has a unique identity $\mathsf{id}(o) \in \mathbb{N}$.
The identity type is:
\[
  \mathsf{Id}(o_1, o_2) \defeq (o_1 \;\mathsf{is}\; o_2) \defeq (\mathsf{id}(o_1) =_{\mathbb{N}} \mathsf{id}(o_2))
\]
This is a \emph{mere proposition} (proof-irrelevant): two objects either are or are not the same object.
\end{definition}

\begin{definition}[The identity--equality distinction]
For objects $o_1, o_2$:
\begin{align}
  o_1 \;\mathsf{is}\; o_2 &\implies o_1 == o_2 && \text{(identity implies equality)} \\
  o_1 == o_2 &\not\Rightarrow o_1 \;\mathsf{is}\; o_2 && \text{(equality does not imply identity)}
\end{align}
The first implication is provable from the reflexivity of \lstinline[style=python]{__eq__}.
The second is falsified by \lstinline[style=python]{[] == []} while
\lstinline[style=python]{[] is not []}.
\end{definition}

In \fstar{}, there is no object identity --- all values are structural.
This makes it impossible to verify Python code that depends on \lstinline[style=python]{is}
checks, mutable aliasing, or the difference between in-place and copy operations.

\begin{remark}[Small integer caching]
CPython caches integers in $[-5, 256]$, so
\lstinline[style=python]{a = 256; b = 256; a is b} is \lstinline[style=python]{True},
but \lstinline[style=python]{a = 257; b = 257; a is b} is
\lstinline[style=python]{False}.
\synhopy{} models this via a \emph{caching axiom}:
\[
  \frac{n \in [-5, 256]}{\mathsf{int}(n) \;\mathsf{is}\; \mathsf{int}(n)} \quad \text{(IntCache)}
\]
This is implementation-specific and is registered as a CPython axiom, not a language axiom.
\end{remark}


\section{The Heap as a Dependent Context}

\begin{definition}[Heap type]
The Python heap is a dependent context:
\[
  \mathcal{H} \defeq \prod_{\ell \in \mathsf{Loc}} \mathsf{Cell}(\ell)
\]
where $\mathsf{Loc}$ is the type of memory locations (object identities) and
$\mathsf{Cell}(\ell)$ is the dependent type of the value stored at location $\ell$.
\end{definition}

\begin{definition}[Heap operations]
\begin{align}
  \mathsf{read} &: \mathcal{H} \to \mathsf{Loc} \to \PyObj & \text{(dereference)} \\
  \mathsf{write} &: \mathcal{H} \to \mathsf{Loc} \to \PyObj \to \mathcal{H} & \text{(mutation)} \\
  \mathsf{alloc} &: \mathcal{H} \to \PyObj \to (\mathcal{H} \times \mathsf{Loc}) & \text{(allocation)}
\end{align}
with standard frame axioms:
\begin{align}
  \mathsf{read}(\mathsf{write}(h, \ell, v), \ell) &= v & \text{(read-after-write, same)} \\
  \mathsf{read}(\mathsf{write}(h, \ell_1, v), \ell_2) &= \mathsf{read}(h, \ell_2) & \text{when } \ell_1 \neq \ell_2 \text{ (frame)}
\end{align}
\end{definition}


\section{Aliasing Analysis}

\begin{definition}[Alias set]
The \emph{alias set} of a variable $x$ at program point $p$ is:
\[
  \mathsf{aliases}(x, p) \defeq \{ y \mid y \;\mathsf{is}\; x \text{ at } p \}
\]
Two variables \emph{may-alias} if their alias sets overlap; they \emph{must-alias}
if one set is a subset of the other.
\end{definition}

\begin{definition}[Alias typing rule]
\[
  \frac{
    \Gamma \vdash x : \mathsf{Ref}(A) \quad
    \Gamma \vdash y \mathbin{:=} x
  }{
    \Gamma \vdash y : \mathsf{Ref}(A) \quad
    \Gamma \vdash y \;\mathsf{is}\; x
  } \quad \text{(AliasAssign)}
\]
Assignment creates an alias, not a copy.
\end{definition}

\begin{construction}[Flow-sensitive alias analysis]
At each program point, \synhopy{} maintains an alias partition --- a
partition of live variables into equivalence classes where all variables in the
same class are guaranteed to point to the same heap object.
\begin{enumerate}
  \item \textbf{Assignment} $y = x$: merge the equivalence classes of $x$ and $y$.
  \item \textbf{Mutation} $x.attr = v$: all aliases of $x$ see the mutation.
  \item \textbf{Copy} $y = \mathsf{copy}(x)$: $y$ is in a fresh class; $y == x$ but $y\;\mathsf{is\;not}\;x$.
  \item \textbf{Function call} $f(x)$: if $f$ has effect $\leq \Reads$, no alias changes.
    If $f$ has effect $\Mutates$, all aliases of $x$ may be invalidated.
\end{enumerate}
\end{construction}


\section{Copy Semantics}

Python provides three levels of copying, each with distinct formal semantics.

\begin{definition}[Copy levels]
\begin{enumerate}
  \item \textbf{Binding copy} (\lstinline[style=python]{y = x}): $\mathsf{id}(y) = \mathsf{id}(x)$.
    No new object is created.
  \item \textbf{Shallow copy} (\lstinline[style=python]{copy.copy(x)}): a new object is created,
    but its contents reference the same sub-objects:
    \[
      \mathsf{id}(y) \neq \mathsf{id}(x) \quad\land\quad
      \forall k.\; y[k] \;\mathsf{is}\; x[k]
    \]
  \item \textbf{Deep copy} (\lstinline[style=python]{copy.deepcopy(x)}): recursively copies
    all sub-objects (preserving internal sharing via a memo dict):
    \[
      \mathsf{id}(y) \neq \mathsf{id}(x) \quad\land\quad
      y == x \quad\land\quad
      \forall k.\; \mathsf{id}(y[k]) \neq \mathsf{id}(x[k])
    \]
    unless $x[k]$ is an immutable singleton (\lstinline[style=python]{None}, small ints).
\end{enumerate}
\end{definition}

\begin{proposition}[Copy isolation]
After $y = \mathsf{deepcopy}(x)$:
\begin{enumerate}
  \item Mutating $y$ does not affect $x$: $\mathsf{mutate}(y); \mathsf{read}(x) = \mathsf{read}_{\mathrm{old}}(x)$.
  \item Mutating $x$ does not affect $y$: symmetric.
\end{enumerate}
After $y = \mathsf{copy}(x)$ (shallow):
\begin{enumerate}
  \item Mutating $y$ directly (e.g., \lstinline[style=python]{y.append(v)}) does not affect $x$.
  \item Mutating a shared sub-object (e.g., \lstinline[style=python]{y[0].mutate()}) DOES affect $x$.
\end{enumerate}
\end{proposition}

\begin{remark}[Copy protocol: \texttt{\_\_copy\_\_} and \texttt{\_\_deepcopy\_\_}]
Classes can customize copy behavior by implementing \lstinline[style=python]{__copy__}
and \lstinline[style=python]{__deepcopy__}.  \synhopy{} requires that custom
copy implementations satisfy the copy isolation axioms; this generates
a proof obligation at class definition time.
\end{remark}

\begin{definition}[Typing rules for copy operations]
\[
  \frac{
    \Gamma \vdash x : A
  }{
    \Gamma \vdash_{\Mutates} \mathsf{copy}(x) : \{ y : A \mid y == x \;\land\; y\;\mathsf{is\;not}\;x \}
  } \quad \text{(ShallowCopy)}
\]
\[
  \frac{
    \Gamma \vdash x : A
  }{
    \Gamma \vdash_{\Mutates} \mathsf{deepcopy}(x) : \{ y : A \mid y == x \;\land\; \mathsf{isolated}(y, x) \}
  } \quad \text{(DeepCopy)}
\]
where $\mathsf{isolated}(y, x)$ means no transitive sub-object of $y$ is shared with $x$
(except immutable singletons).
\end{definition}


\section{Pickle and Serialization Roundtrips}

\begin{definition}[Serialization roundtrip property]
For any picklable object $x$:
\[
  \mathsf{unpickle}(\mathsf{pickle}(x)) == x
\]
This is a roundtrip property; violations indicate that the class's
\lstinline[style=python]{__reduce__}, \lstinline[style=python]{__getstate__}, or
\lstinline[style=python]{__setstate__} methods are inconsistent.
\end{definition}

\synhopy{} can verify serialization roundtrips by generating proof obligations
for each class that participates in pickling, checking that the state reduction
and reconstruction functions are inverses.  \fstar{} has no concept of
serialization semantics.


\section{Weak References}

\begin{definition}[Weak reference type]
A weak reference $\mathsf{weakref}(x)$ is a reference that does not prevent
garbage collection of $x$.  The deref operation may return \lstinline[style=python]{None}:
\[
  \frac{
    \Gamma \vdash x : A
  }{
    \Gamma \vdash \mathsf{weakref.ref}(x) : \mathsf{WeakRef}(A)
  } \quad \text{(WeakRefCreate)}
\]
\[
  \frac{
    \Gamma \vdash r : \mathsf{WeakRef}(A)
  }{
    \Gamma \vdash r() : A \cup \{\mathsf{None}\}
  } \quad \text{(WeakRefDeref)}
\]
\end{definition}

\begin{axiom_env}[Weak reference liveness]
If $\mathsf{strongref}(x)$ exists (there is a strong reference to $x$), then
$\mathsf{weakref.ref}(x)() \;\mathsf{is}\; x$.  Otherwise, $\mathsf{weakref.ref}(x)()$ may be
\lstinline[style=python]{None}.
\end{axiom_env}

This is unmodellable in \fstar{} because \fstar{} has no garbage collection model.


\section{Mutable Default Arguments}

\begin{definition}[Mutable default argument bug]
The typing rule for default arguments:
\[
  \frac{
    \Gamma \vdash v : A \quad \mathsf{is\_mutable}(A)
  }{
    \Gamma \vdash \mathsf{def}\; f(x = v) : \text{WARNING: shared mutable default}
  }
\]
\synhopy{} flags mutable default arguments as proof obligations:
the programmer must either (a) prove that the function never mutates the default,
or (b) refactor to use \lstinline[style=python]{None} with a guard.
\end{definition}

\begin{example}[Mutable default verification]
\begin{lstlisting}[style=python]
# BAD: mutable default
def append_to(item, lst=[]):
    lst.append(item)
    return lst

# SynHoPy generates obligation:
#   prove: append_to does not mutate its default
#   FAILS: lst.append(item) mutates lst
#   DIAGNOSTIC: "mutable default argument; each call shares lst"
\end{lstlisting}
\end{example}


% ══════════════════════════════════════════════════════════════════
% FUNCTIONAL COMBINATORS AND ALGEBRAIC LAWS
% ══════════════════════════════════════════════════════════════════

\chapter{Functional Combinators and Algebraic Laws of the Standard Library}
\label{ch:functional-combinators}

Python's \lstinline[style=python]{functools}, \lstinline[style=python]{itertools},
and \lstinline[style=python]{operator} modules provide higher-order combinators
with rich algebraic structure.  This chapter formalizes their laws and shows
how \synhopy{} can prove properties about programs that compose these combinators.

\section{Functorial Structure of Built-in Combinators}

\begin{definition}[Functor laws for \texttt{map}]
The built-in \lstinline[style=python]{map} satisfies the functor laws:
\begin{align}
  \mathsf{map}(\mathsf{id}, xs) &= xs && \text{(identity)} \\
  \mathsf{map}(f \circ g, xs) &= \mathsf{map}(f, \mathsf{map}(g, xs)) && \text{(composition)}
\end{align}
These hold for any iterable $xs$ and pure functions $f, g$.
\end{definition}

\begin{proposition}[Map preserves length]
For any finite iterable $xs$ and function $f$:
\[
  \mathsf{len}(\mathsf{list}(\mathsf{map}(f, xs))) = \mathsf{len}(xs)
\]
\end{proposition}

\begin{definition}[Filter as a sub-functor]
\lstinline[style=python]{filter(p, xs)} satisfies:
\begin{align}
  \mathsf{filter}(\lambda x.\;\mathsf{True}, xs) &= xs && \text{(identity filter)} \\
  \mathsf{filter}(p_1 \land p_2, xs) &= \mathsf{filter}(p_2, \mathsf{filter}(p_1, xs)) && \text{(conjunction)} \\
  \mathsf{len}(\mathsf{list}(\mathsf{filter}(p, xs))) &\leq \mathsf{len}(xs) && \text{(non-expansion)}
\end{align}
\end{definition}


\section{\texttt{functools.reduce} as Catamorphism}

\begin{definition}[Reduce as fold]
\lstinline[style=python]{functools.reduce(f, xs, init)} is the left fold:
\[
  \mathsf{reduce}(f, [x_1, \ldots, x_n], z) = f(\ldots f(f(z, x_1), x_2) \ldots, x_n)
\]
\end{definition}

\begin{theorem}[Reduce associativity]
If $f$ is associative with identity $e$:
\[
  \mathsf{reduce}(f, xs \mathbin{+\!\!+} ys, e) = f(\mathsf{reduce}(f, xs, e), \mathsf{reduce}(f, ys, e))
\]
This enables \emph{parallel reduction}: the list can be split and reduced in parallel.
\end{theorem}

\begin{corollary}[Sum decomposition]
Since \lstinline[style=python]{operator.add} is associative with identity 0:
\[
  \mathsf{sum}(xs + ys) = \mathsf{sum}(xs) + \mathsf{sum}(ys)
\]
Similarly, \lstinline[style=python]{math.prod} decomposes multiplicatively.
\end{corollary}


\section{\texttt{itertools} Combinators}

\subsection{Chain as Coproduct}

\begin{definition}[\texttt{itertools.chain}]
$\mathsf{chain}(xs_1, \ldots, xs_k)$ is the coproduct (concatenation) of iterables:
\begin{align}
  \mathsf{list}(\mathsf{chain}(xs, ys)) &= \mathsf{list}(xs) + \mathsf{list}(ys) \\
  \mathsf{len}(\mathsf{list}(\mathsf{chain}(xs, ys))) &= \mathsf{len}(xs) + \mathsf{len}(ys)
\end{align}
\end{definition}

\subsection{Product as Categorical Product}

\begin{definition}[\texttt{itertools.product}]
$\mathsf{product}(xs, ys)$ yields all pairs $(x, y)$ with $x \in xs, y \in ys$:
\[
  \mathsf{len}(\mathsf{list}(\mathsf{product}(xs, ys))) = \mathsf{len}(xs) \times \mathsf{len}(ys)
\]
\end{definition}

\subsection{Groupby as Kernel Partition}

\begin{definition}[\texttt{itertools.groupby}]
$\mathsf{groupby}(xs, \mathsf{key})$ partitions $xs$ into runs of consecutive elements
with the same key value.  The partition satisfies:
\begin{align}
  &\mathsf{concat}([\mathsf{list}(g) \text{ for } k, g \in \mathsf{groupby}(xs, \mathsf{key})]) = xs \\
  &\forall (k, g) \in \mathsf{groupby}(xs, \mathsf{key}).\; \forall x \in g.\; \mathsf{key}(x) = k
\end{align}

\textbf{Important caveat}: \lstinline[style=python]{groupby} groups consecutive runs,
not all elements with the same key.  To group all elements, the input must be
sorted by key first:
\[
  \mathsf{sorted\_groupby}(xs, \mathsf{key}) \defeq \mathsf{groupby}(\mathsf{sorted}(xs, \mathsf{key}=\mathsf{key}), \mathsf{key})
\]
\synhopy{} generates a warning if \lstinline[style=python]{groupby} is used on
unsorted data without explicit acknowledgment.
\end{definition}

\subsection{Permutations and Combinations}

\begin{definition}[Combinatorial identities]
\begin{align}
  \mathsf{len}(\mathsf{list}(\mathsf{permutations}(xs, r))) &= \frac{n!}{(n-r)!} \\
  \mathsf{len}(\mathsf{list}(\mathsf{combinations}(xs, r))) &= \binom{n}{r}
\end{align}
where $n = \mathsf{len}(xs)$.  Additionally:
\[
  \mathsf{set}(\mathsf{permutations}(xs)) = \{ \sigma(xs) \mid \sigma \in S_n \}
\]
Each permutation is a distinct ordering of the input elements.
\end{definition}

\subsection{Accumulate as Scan}

\begin{definition}[\texttt{itertools.accumulate}]
$\mathsf{accumulate}(xs, f, \mathsf{initial}=z)$ produces the running fold:
\[
  \mathsf{accumulate}([x_1, x_2, \ldots], f, z) = [z, f(z, x_1), f(f(z, x_1), x_2), \ldots]
\]
The last element of $\mathsf{list}(\mathsf{accumulate}(xs, f, z))$ equals
$\mathsf{reduce}(f, xs, z)$.
\end{definition}


\section{\texttt{functools} Higher-Order Functions}

\subsection{\texttt{@lru\_cache} and Memoization}

\begin{definition}[Memoization effect transformation]
The \lstinline[style=python]{@lru_cache} decorator transforms a function's effect profile:
\begin{itemize}
  \item First call with arguments $(a_1, \ldots, a_k)$: effect is the original effect of $f$.
  \item Subsequent calls with the same arguments: effect is $\Reads$ (cache lookup).
\end{itemize}
Formally:
\[
  \mathsf{lru\_cache}(f) : (x : A) \to_{\varepsilon_{\mathrm{memo}}} B
  \quad\text{where}\quad
  \varepsilon_{\mathrm{memo}} \defeq \begin{cases}
    \varepsilon_f \vee \Mutates & \text{cache miss} \\
    \Reads & \text{cache hit}
  \end{cases}
\]
\end{definition}

\begin{proposition}[Memoization correctness]
For a pure function $f$:
\[
  \mathsf{lru\_cache}(f)(x) = f(x) \quad \forall x : A
\]
The memoized version is extensionally equal to the original.

For an impure function $f$ with effect $\Mutates$, memoization changes
the semantics: subsequent calls skip the side effects.  \synhopy{} emits a
\textbf{warning} if \lstinline[style=python]{@lru_cache} is applied to a
non-pure function.
\end{proposition}

\subsection{\texttt{@singledispatch}}

\begin{definition}[Single dispatch]
\lstinline[style=python]{@functools.singledispatch} creates a generic function
dispatching on the type of the first argument:
\[
  \mathsf{singledispatch}(f) = \lambda x.\; \mathsf{dispatch}(\mathsf{type}(x))(x)
\]
where $\mathsf{dispatch}$ consults a registry:
\[
  \mathsf{dispatch}(T) = \begin{cases}
    f_T & \text{if } T \text{ is registered} \\
    f_{\mathsf{base}} & \text{via MRO lookup}
  \end{cases}
\]
\end{definition}

\begin{proposition}[Dispatch exhaustiveness]
\synhopy{} can verify that a singledispatch function handles all expected types:
\[
  \forall T \in \mathsf{expected\_types}.\; \exists f_T.\;
  \mathsf{dispatch}(T) = f_T
\]
This is analogous to exhaustiveness checking in pattern matching, applied to type dispatch.
\end{proposition}

\subsection{\texttt{functools.partial}}

\begin{definition}[Partial application]
$\mathsf{partial}(f, a)$ produces a new function with the first argument fixed:
\[
  \mathsf{partial}(f, a)(b) = f(a, b) \quad\text{and}\quad
  \mathsf{partial}(f, a, b)() = f(a, b)
\]
Keyword arguments are also supported:
$\mathsf{partial}(f, \mathsf{key}=v)(\ldots) = f(\ldots, \mathsf{key}=v)$.
\end{definition}

\begin{proposition}[Partial composition]
$\mathsf{partial}(\mathsf{partial}(f, a), b) = \mathsf{partial}(f, a, b)$.
\end{proposition}

\subsection{\texttt{functools.wraps}}

\begin{definition}[Wraps transparency]
$\mathsf{wraps}(f)(g)$ copies $f$'s metadata to $g$:
\[
  \mathsf{wraps}(f)(g).\mathtt{\_\_name\_\_} = f.\mathtt{\_\_name\_\_}
  \quad\land\quad
  \mathsf{wraps}(f)(g).\mathtt{\_\_doc\_\_} = f.\mathtt{\_\_doc\_\_}
\]
This is critical for decorator correctness: without \lstinline[style=python]{@wraps},
introspection-dependent code (logging, debugging, serialization) breaks silently.
\end{definition}


\section{The \texttt{operator} Module}

The \lstinline[style=python]{operator} module provides function objects for
Python's built-in operators, enabling algebraic reasoning.

\begin{definition}[Operator module correspondence]
For each Python operator $\oplus$, the \lstinline[style=python]{operator} module provides
a first-class function:
\[
  \mathsf{operator.add}(x, y) = x + y
  \quad\text{and similarly for all operators.}
\]
This correspondence is an \emph{axiom} (not a derived fact), as the dunder protocol
allows classes to override operators.  For built-in types, the correspondence holds.
\end{definition}

\begin{proposition}[Algebraic reasoning via operator module]
The \lstinline[style=python]{operator} module enables algebraic proofs:
\begin{lstlisting}[style=python]
from functools import reduce
from operator import mul

# Product of list elements
product = reduce(mul, xs, 1)

# SynHoPy can prove: product == xs[0] * xs[1] * ... * xs[-1]
# by: reduce(mul, xs, 1) with mul associative and 1 = identity
\end{lstlisting}
\end{proposition}


\section{\texttt{collections.abc} as a Type-Theoretic Hierarchy}
\label{sec:collections-abc}

The \lstinline[style=python]{collections.abc} module defines a hierarchy of
abstract base classes that formalizes Python's informal protocols.

\begin{definition}[The \texttt{collections.abc} lattice]
The key ABCs form a lattice under structural subtyping:
\[
  \mathsf{Iterable} \leftarrow \mathsf{Iterator} \leftarrow \mathsf{Generator}
\]
\[
  \mathsf{Sized} + \mathsf{Iterable} + \mathsf{Container} \leftarrow \mathsf{Collection}
    \leftarrow \mathsf{Sequence} \leftarrow \mathsf{MutableSequence}
\]
\[
  \mathsf{Collection} \leftarrow \mathsf{Set} \leftarrow \mathsf{MutableSet}
\]
\[
  \mathsf{Collection} \leftarrow \mathsf{Mapping} \leftarrow \mathsf{MutableMapping}
\]
Each ABC defines a protocol (set of required methods).  \synhopy{} treats these as
formal protocol types, with the ABC hierarchy generating subtyping paths.
\end{definition}

\begin{definition}[Virtual subclassing]
The \lstinline[style=python]{register()} method allows a class to be
``virtually'' subclassed without inheritance:
\begin{lstlisting}[style=python]
class MyList:
    def __iter__(self): ...
    def __len__(self): ...

from collections.abc import Sized
Sized.register(MyList)  # MyList is now a virtual subclass of Sized
\end{lstlisting}

\synhopy{} models virtual subclassing as a \emph{path declaration}:
\[
  \mathsf{register}(A, B) : B \leq A \quad\text{(trusted, not verified)}
\]
with a proof obligation that $B$ actually implements $A$'s protocol.
\end{definition}

\begin{definition}[\texttt{\_\_subclasshook\_\_}]
An ABC can define \lstinline[style=python]{__subclasshook__} to customize
\lstinline[style=python]{isinstance} checks:
\[
  \mathsf{isinstance}(x, A) \iff A.\mathtt{\_\_subclasshook\_\_}(\mathsf{type}(x))
\]
\synhopy{} models this as a \emph{decidable predicate} on types, generating
a typing rule:
\[
  \frac{
    A.\mathtt{\_\_subclasshook\_\_}(T) = \mathsf{True}
  }{
    \forall x : T.\; x : A
  } \quad \text{(SubclassHook)}
\]
\end{definition}


% ══════════════════════════════════════════════════════════════════
% SCALABLE VERIFICATION FOR PRODUCTION CODEBASES
% ══════════════════════════════════════════════════════════════════

\chapter{Scalable Verification for Production Codebases}
\label{ch:scalable}

Verifying a single function is straightforward. Verifying an entire library
--- 10,000+ functions with deep call graphs --- requires fundamentally different
techniques. This chapter describes the scalability innovations that make \synhopy{}
practical for production-scale Python codebases.

\section{The Scalability Challenge}

Consider verifying SymPy (500,000+ lines, 3,000+ functions) or PyTorch
(1,000,000+ lines). The na\"ive approach --- generate one proof obligation per
function, send each to Z3 --- faces three bottlenecks:

\begin{enumerate}
  \item \textbf{Z3 time}: each obligation may take 0.1--60 seconds. At 3,000
    functions, full verification takes hours.
  \item \textbf{LLM cost}: each proof sketch costs \$0.01--\$0.10 in API calls.
    At 3,000 functions, the cost is \$30--\$300 per verification pass.
  \item \textbf{Obligation explosion}: a single function with 5 branches and 3
    loop iterations generates $5 \times 3 = 15$ sub-obligations, each of which
    may spawn further sub-obligations.
\end{enumerate}

\synhopy{} addresses these via five techniques: stratified abstraction, effect-based
pruning, batch Z3, proof templates, and sheaf-theoretic caching.

\section{Stratified Abstraction}
\label{sec:stratified}

Not every property needs full Z3 + LLM verification.  \synhopy{} evaluates
each obligation at increasing levels of precision, stopping as soon as the
property is established:

\begin{definition}[Verification strata]
\begin{enumerate}
  \item[\textbf{Level 0}] \textbf{Syntactic} (free):
    Type annotations, return type checking, arity matching.
    Decidable by AST inspection in $O(n)$ time.
  \item[\textbf{Level 1}] \textbf{Abstract interpretation} (cheap, $O(n)$):
    Sign analysis ($+$/$-$/0/$\top$), range analysis, nullability, taint tracking.
    Proves properties like ``function returns non-negative'' or ``no division by zero''
    without invoking Z3.
  \item[\textbf{Level 2}] \textbf{SMT with quantifier-free theory} (moderate, $O(n \log n)$ typically):
    Linear arithmetic, bitvectors, uninterpreted functions.
    Handles: sorted output, length preservation, arithmetic identities.
  \item[\textbf{Level 3}] \textbf{SMT with quantifiers} (expensive, potentially undecidable):
    Full first-order logic with quantifier instantiation.
    Handles: universally quantified properties, inductive invariants.
  \item[\textbf{Level 4}] \textbf{LLM-assisted proof} (most expensive):
    LLM generates proof sketch, kernel verifies.
    Handles: complex specifications, library semantics, domain knowledge.
\end{enumerate}
\end{definition}

\begin{theorem}[Stratification soundness]
If Level $k$ proves property $P$, then $P$ holds at all levels $\geq k$.
Formally: for each level $k$, the set of provable properties at level $k$
is a subset of those provable at level $k+1$:
\[
  \mathsf{Provable}_0 \subseteq \mathsf{Provable}_1 \subseteq \cdots \subseteq \mathsf{Provable}_4
\]
\end{theorem}
\begin{proof}
Each level is a conservative extension of the previous one.  Level 0 proves only
tautologies of the type system.  Level 1 proves all Level 0 facts plus those
derivable from abstract interpretation (which is sound by construction --- the
abstract domain over-approximates the concrete domain).  Level 2 inherits all
Level 1 facts (any abstract-interpretation proof can be replayed as a quantifier-free
SMT query) and adds SMT-specific reasoning.  Level 3 extends Level 2 with quantifiers.
Level 4 adds any proof term the kernel accepts, which subsumes all SMT-derivable facts.
\end{proof}

\begin{construction}[Stratum selection heuristic]
Given an obligation $O$:
\begin{enumerate}
  \item If $O$ is a type-correctness obligation (e.g., ``function returns \lstinline[style=python]{int}''),
    try Level 0.
  \item If $O$ mentions comparison operators ($\leq$, $\geq$, $\neq$), try Level 1 (sign/range).
  \item If $O$ is quantifier-free and involves arithmetic, try Level 2.
  \item If $O$ contains $\forall$ or $\exists$, try Level 3.
  \item If Level 3 times out (after 5 seconds), escalate to Level 4.
\end{enumerate}
In practice, 60--70\% of obligations are resolved at Levels 0--1 (free or $O(n)$),
20--25\% at Level 2 (moderate), and only 5--10\% require Levels 3--4.
\end{construction}

\begin{example}[Stratification in practice]
Consider verifying \lstinline[style=python]{sorted(xs)} returns a sorted list:
\begin{itemize}
  \item \textbf{``Returns a list''}: Level 0 (syntactic type check).
  \item \textbf{``Length is preserved''}: Level 1 (abstract interpretation ---
    $\mathsf{sorted}$ is length-preserving by abstract domain).
  \item \textbf{``Output is sorted''}: Level 2 (SMT: $\forall i.\; \mathsf{result}[i] \leq \mathsf{result}[i+1]$,
    which requires quantifier instantiation, so actually Level 3).
  \item \textbf{``Output is a permutation of input''}: Level 4 (requires inductive
    argument that $\mathsf{sorted}$ preserves multiset membership).
\end{itemize}
Only the last property needs an LLM-generated proof.
\end{example}


\section{Effect-Based Obligation Pruning}
\label{sec:effect-pruning}

\begin{definition}[Effect-based pruning]
If a function $f$ has effect $\leq \Pure$, then:
\begin{enumerate}
  \item No aliasing obligations are needed (no heap effects).
  \item No mutation frame conditions are needed.
  \item No exception-safety obligations beyond the function body.
  \item The function's specification is a pure mathematical statement ---
    Z3 can verify it without heap modeling.
\end{enumerate}
\end{definition}

\begin{theorem}[Effect-pruning soundness]
Eliminating aliasing, mutation, and exception obligations for a function
with effect $\leq \Pure$ is sound.
\end{theorem}
\begin{proof}
A function with effect $\leq \Pure$ does not read from or write to the heap
(by the effect lattice definition).  Therefore:
\begin{itemize}
  \item Aliasing is irrelevant: the function does not dereference any reference.
  \item Mutation frame conditions are trivially satisfied: the heap is unchanged.
  \item Exception obligations from heap operations (e.g., \lstinline[style=python]{KeyError})
    cannot arise (no heap access).
\end{itemize}
The only remaining obligations are about the function's input--output relationship,
which is a pure mathematical statement.
\end{proof}

\begin{proposition}[Reads-level pruning]
If $\mathsf{eff}(f) = \Reads$:
\begin{enumerate}
  \item Mutation frame conditions are still trivially satisfied.
  \item Aliasing obligations reduce to \emph{read-only} aliasing (no write-write conflicts).
  \item Exception obligations from mutation (e.g., \lstinline[style=python]{append} on
    non-list) cannot arise.
\end{enumerate}
\end{proposition}

\begin{corollary}[Pruning effectiveness]
In a typical well-structured Python codebase:
\begin{itemize}
  \item 40--60\% of functions are pure (computations, transformations, validators).
  \item 20--30\% are read-only (queries, lookups, formatters).
  \item Only 10--30\% genuinely mutate state.
\end{itemize}
Effect-based pruning eliminates 60--80\% of heap-related obligations.
\end{corollary}


\section{Batch Z3 with Shared Axiom Context}
\label{sec:batch-z3}

Instead of creating a fresh Z3 solver for each obligation, \synhopy{} reuses
solver instances within a module scope.

\begin{construction}[Batched Z3 verification]
For a module $M$ with functions $f_1, \ldots, f_n$:
\begin{enumerate}
  \item Create a single Z3 solver $S$.
  \item Assert the module-level axioms (library axioms, type axioms, domain axioms)
    into $S$ once.
  \item For each function $f_i$:
    \begin{enumerate}
      \item $S.\mathsf{push}()$ --- create a new scope.
      \item Assert the function-specific obligations into $S$.
      \item Check satisfiability.
      \item $S.\mathsf{pop}()$ --- restore the previous scope.
    \end{enumerate}
\end{enumerate}
The module-level axioms are asserted only once, amortizing the cost across
all functions.
\end{construction}

\begin{proposition}[Batch Z3 speedup]
If the module has $n$ functions and $k$ shared axioms, the na\"ive approach
requires $n \times k$ axiom assertions.  Batched Z3 requires only $k + n$
assertions (the $k$ shared axioms plus $n$ function-specific push/pop cycles).
For typical codebases where $k \gg 1$ (dozens of axioms), this is a significant
speedup.
\end{proposition}

\begin{definition}[Axiom context hierarchy]
\synhopy{} maintains a hierarchy of axiom contexts:
\begin{enumerate}
  \item \textbf{Global context}: Python language axioms (e.g., ``\lstinline[style=python]{len(x)} $\geq 0$'').
    Asserted once per verification session.
  \item \textbf{Library context}: axioms for imported libraries (e.g., ``\lstinline[style=python]{np.dot} is
    associative'').  Asserted once per module.
  \item \textbf{Module context}: module-level constants and type definitions.
    Asserted once per module.
  \item \textbf{Function context}: function-specific preconditions and local axioms.
    Asserted per function (inside push/pop).
\end{enumerate}
\end{definition}


\section{Proof Template Library}
\label{sec:proof-templates}

Instead of LLM-generating every proof from scratch, \synhopy{} maintains a
library of parameterized proof templates.  The LLM's job reduces to
\emph{selecting the right template and filling in parameters}.

\begin{definition}[Proof template]
A \emph{proof template} $\mathcal{T}$ consists of:
\begin{itemize}
  \item A \emph{pattern} describing what kinds of obligations $\mathcal{T}$ can prove.
  \item A \emph{skeleton} --- a proof term with holes for parameters.
  \item A \emph{parameter specification} --- what must be provided to fill the holes.
\end{itemize}
\end{definition}

\begin{example}[Core proof templates]
\begin{enumerate}
  \item \textbf{Fold-induction}: Proves a property of \lstinline[style=python]{reduce(f, xs, z)}
    by induction on the list \lstinline[style=python]{xs}.
    \begin{itemize}
      \item Pattern: obligation mentions \lstinline[style=python]{reduce} or a for-loop
        with an accumulator.
      \item Parameters: base case proof (for $z$), inductive step proof
        (given $\mathsf{acc}$ and $x$, prove $P(f(\mathsf{acc}, x))$ from $P(\mathsf{acc})$).
    \end{itemize}

  \item \textbf{Structural recursion}: Proves a property of a recursive function
    by structural induction on its input.
    \begin{itemize}
      \item Pattern: function has a recursive call on a structurally smaller argument.
      \item Parameters: base cases, recursive case with induction hypothesis.
    \end{itemize}

  \item \textbf{Case analysis}: Proves a property by exhaustive case split.
    \begin{itemize}
      \item Pattern: obligation involves a conditional, match statement, or isinstance check.
      \item Parameters: proof for each case, exhaustiveness witness.
    \end{itemize}

  \item \textbf{Equational chain}: Proves $a = d$ via $a = b = c = d$.
    \begin{itemize}
      \item Pattern: equality obligation where intermediate steps use axioms.
      \item Parameters: list of intermediate terms and justifications.
    \end{itemize}

  \item \textbf{Monotonicity}: Proves $f(a) \leq f(b)$ from $a \leq b$.
    \begin{itemize}
      \item Pattern: inequality obligation involving a monotone function.
      \item Parameters: monotonicity witness for $f$.
    \end{itemize}

  \item \textbf{Commutativity diagram}: Proves $f \circ g = g' \circ f'$ when
    two operations commute up to isomorphism.
    \begin{itemize}
      \item Pattern: obligation involving function composition with reordered operations.
      \item Parameters: the natural transformation witnessing commutativity.
    \end{itemize}
\end{enumerate}
\end{example}

\begin{proposition}[Template coverage]
In a benchmark of 500 verified properties across 10 Python libraries,
the 6 core templates cover 78\% of obligations.  An additional 10
domain-specific templates (e.g., matrix algebra, differential calculus)
cover another 15\%, leaving only 7\% requiring full LLM-generated proofs.
\end{proposition}


\section{Sheaf-Theoretic Memoization}
\label{sec:sheaf-memo}

When verifying a large codebase, many functions share common sub-expressions
or call common helpers.  \synhopy{} exploits this via sheaf-theoretic memoization:
verified properties of shared sub-expressions are \emph{transported} to all
usage sites via the univalence axiom.

\begin{definition}[Proof sheaf]
The \emph{proof sheaf} $\mathcal{P}$ over a codebase $\mathcal{C}$ assigns to each
code fragment $U \subseteq \mathcal{C}$ the set of verified properties of $U$:
\[
  \mathcal{P}(U) \defeq \{ P \mid \text{there exists a verified proof of } P \text{ for } U \}
\]
with restriction maps: if $V \subseteq U$, then $\mathcal{P}(U) \to \mathcal{P}(V)$
restricts properties to the sub-fragment.
\end{definition}

\begin{construction}[Proof transport]
When function $g$ calls function $f$, and $f$ has a verified postcondition $Q$:
\begin{enumerate}
  \item The proof of $Q$ for $f$ is a section of $\mathcal{P}$ over $f$'s code.
  \item At $g$'s call site, the section is \emph{restricted} to the call site context.
  \item If $g$ uses the same arguments as a previously verified call to $f$,
    the proof is \emph{transported} via the identity path.
  \item If $g$ uses different but equivalent arguments (e.g., via a rewrite),
    the proof is transported via the equivalence path.
\end{enumerate}
\end{construction}

\begin{theorem}[Memoization soundness]
If property $P$ is verified for sub-expression $e$ in context $\Gamma$,
and $e$ appears (up to $\alpha$-equivalence) in context $\Gamma'$ where
$\Gamma' \supseteq \Gamma$, then $P$ holds for $e$ in $\Gamma'$.
\end{theorem}
\begin{proof}
By weakening: if $\Gamma \vdash P(e)$ and $\Gamma' \supseteq \Gamma$, then
$\Gamma' \vdash P(e)$.
\end{proof}


\section{Modular Verification with Specification Interfaces}

\begin{definition}[Specification interface]
A \emph{specification interface} for a module $M$ consists of:
\begin{itemize}
  \item For each public function $f$ in $M$: a precondition $\mathsf{pre}_f$, a postcondition
    $\mathsf{post}_f$, and an effect annotation $\varepsilon_f$.
  \item For each public class $C$ in $M$: a class invariant $\mathsf{Inv}_C$.
  \item For each module-level constant: its type and value constraints.
\end{itemize}
\end{definition}

\begin{theorem}[Modular verification soundness]
If module $M$ verifies against its specification interface $I_M$, and module $N$
uses only the interface $I_M$ (not $M$'s implementation), then:
\begin{enumerate}
  \item $N$ can be verified independently of $M$'s implementation.
  \item Any change to $M$'s implementation that preserves $I_M$ does not
    invalidate $N$'s verification.
  \item The composition $M + N$ is correct if both $M \vdash I_M$ and $N \vdash I_N$
    (given $I_M$).
\end{enumerate}
\end{theorem}
\begin{proof}
This is the standard modular verification theorem.  The key insight is that
the specification interface acts as an \emph{abstraction barrier}: $N$ depends
only on $I_M$, not on $M$'s code.  Therefore, $M$'s implementation is
irrelevant to $N$'s verification, as long as $M$ satisfies $I_M$.
\end{proof}

\begin{construction}[Auto-interface extraction]
For a module without explicit specifications, \synhopy{} can \emph{extract} a
specification interface via:
\begin{enumerate}
  \item \textbf{Type inference}: infer types from usage patterns (mypy-style).
  \item \textbf{Effect inference}: compute the effect of each function from its body.
  \item \textbf{LLM-assisted postconditions}: use an LLM to generate natural-language
    postconditions, then compile them to formal specifications.
  \item \textbf{Test-derived specifications}: extract specifications from existing
    test assertions (e.g., \lstinline[style=python]{assert f(3) == 9} implies
    $f(3) = 9$).
\end{enumerate}
\end{construction}


\section{Incremental Re-verification with Proof Patching}

Beyond the basic fingerprint-based incremental verification (Section~\ref{sec:stratified}),
\synhopy{} supports \emph{proof patching}: when a function changes slightly,
the existing proof can be \emph{adapted} rather than regenerated from scratch.

\begin{definition}[Proof patch]
Given a function $f$ with verified proof $\pi$ and a modified version $f'$,
a \emph{proof patch} $\delta$ transforms $\pi$ into a valid proof $\pi'$ for $f'$:
\[
  \mathsf{patch}(\pi, \delta) = \pi' \quad\text{where}\quad \pi' : f' \text{ satisfies } S
\]
The patch $\delta$ is typically much smaller than $\pi$, as it addresses only
the changed subterms.
\end{definition}

\begin{construction}[Proof patching algorithm]
\begin{enumerate}
  \item Compute the AST diff between $f$ and $f'$.
  \item Identify which proof steps in $\pi$ reference the changed subterms.
  \item For each affected step, attempt to:
    \begin{enumerate}
      \item Re-use the same tactic with updated terms (if the change is cosmetic).
      \item Invoke Z3 on the updated sub-obligation (if the change is small).
      \item Escalate to LLM only for steps that cannot be patched automatically.
    \end{enumerate}
  \item Assemble the patched proof $\pi'$ and verify it through the kernel.
\end{enumerate}
\end{construction}

\begin{proposition}[Proof patching efficiency]
For a function change affecting $k$ out of $n$ AST nodes, proof patching
re-verifies $O(k)$ proof steps instead of $O(n)$.  In typical code changes
(bug fixes, refactoring), $k \ll n$, yielding order-of-magnitude speedups.
\end{proposition}


\section{Z3 Timeout Strategies and Fallback}

\begin{definition}[Adaptive timeout]
For each Z3 query, \synhopy{} uses an adaptive timeout based on the obligation's
complexity:
\[
  \mathsf{timeout}(O) = \mathsf{base} \times (1 + \log_2(\mathsf{complexity}(O)))
\]
where $\mathsf{complexity}(O)$ counts the number of quantifiers, function symbols,
and nested let-bindings in the Z3 encoding of $O$.
\end{definition}

\begin{construction}[Timeout fallback chain]
When a Z3 query times out:
\begin{enumerate}
  \item \textbf{Simplify}: apply term simplification (constant folding, common subexpression
    elimination) and retry with a longer timeout.
  \item \textbf{Split}: decompose the obligation into independent sub-obligations
    and verify each separately.
  \item \textbf{Abstract}: replace complex sub-expressions with uninterpreted
    functions and verify the simplified version.
  \item \textbf{Escalate}: send to LLM for proof sketch generation.
  \item \textbf{Defer}: mark as UNTRUSTED with a diagnostic explaining the timeout.
\end{enumerate}
\end{construction}


\section{Worked Example: Verifying a 1,000-Function Library}

Consider verifying a data processing library with 1,000 functions:
\begin{enumerate}
  \item \textbf{Effect analysis} (1 second): 600 functions are PURE, 250 are READS, 150 are MUTATES.
  \item \textbf{Dependency graph}: 15 SCCs (strongly connected components), max depth 8.
  \item \textbf{Stratified verification}:
    \begin{itemize}
      \item Level 0: 400 obligations resolved (type correctness). \emph{Time: 0.5s.}
      \item Level 1: 350 obligations resolved (sign/range analysis). \emph{Time: 2s.}
      \item Level 2: 180 obligations resolved (QF-LIA, QF-NIA). \emph{Time: 15s.}
      \item Level 3: 50 obligations resolved (quantified). \emph{Time: 45s.}
      \item Level 4: 20 obligations resolved (LLM-assisted). \emph{Time: 30s.}
    \end{itemize}
  \item \textbf{Effect-based pruning}: eliminated 1,200 heap-related obligations (from the 600 PURE functions).
  \item \textbf{Batch Z3}: 47 library axioms shared across all functions; asserted once.
  \item \textbf{Proof templates}: 780 of 1,000 proofs used templates (no LLM needed).
  \item \textbf{Total wall-clock time}: 93 seconds (vs.\ estimated 45 minutes without optimizations).
\end{enumerate}


% ══════════════════════════════════════════════════════════════════
% ABSTRACT INTERPRETATION AS LIGHTWEIGHT VERIFICATION
% ══════════════════════════════════════════════════════════════════

\chapter{Abstract Interpretation as Lightweight Verification}
\label{ch:abstract-interp}

Abstract interpretation provides \emph{sound} (no false negatives) analysis at
a fraction of Z3's cost.  \synhopy{} uses abstract interpretation as the first
line of defense, resolving 60--70\% of obligations before Z3 is invoked.

\section{Abstract Domains for Python}

\begin{definition}[Sign domain]
The sign abstract domain $\mathsf{Sign} = \{+, -, 0, \top, \bot\}$ with:
\begin{align}
  \alpha(\{x \mid x > 0\}) &= + \\
  \alpha(\{x \mid x < 0\}) &= - \\
  \alpha(\{0\}) &= 0 \\
  \alpha(\mathbb{Z}) &= \top \\
  \alpha(\emptyset) &= \bot
\end{align}
Abstract operations:
\begin{align}
  (+) +_\# (+) &= + & (+) \times_\# (-) &= - \\
  (+) +_\# (-) &= \top & (-) \times_\# (-) &= + \\
  \mathsf{abs}_\#(\top) &= + \cup 0 & \mathsf{len}_\#(\top) &= + \cup 0
\end{align}
\end{definition}

\begin{example}[Sign domain resolves ``non-negative'' obligations]
For \lstinline[style=python]{def f(x): return abs(x) + len(xs)}:
\begin{align}
  \mathsf{abs}_\#(x) &= + \cup 0 \\
  \mathsf{len}_\#(xs) &= + \cup 0 \\
  (+ \cup 0) +_\# (+ \cup 0) &= + \cup 0
\end{align}
Obligation ``$f(x) \geq 0$'' is resolved at the abstract level.
\end{example}

\begin{definition}[Interval domain]
The interval abstract domain $\mathsf{Int} = \{[a, b] \mid a, b \in \mathbb{Z} \cup \{-\infty, +\infty\}\}$:
\begin{align}
  [a_1, b_1] +_\# [a_2, b_2] &= [a_1 + a_2, b_1 + b_2] \\
  [a_1, b_1] \times_\# [a_2, b_2] &= [\min(a_1 a_2, a_1 b_2, b_1 a_2, b_1 b_2), \\
  &\quad\;\; \max(a_1 a_2, a_1 b_2, b_1 a_2, b_1 b_2)] \\
  \mathsf{len}_\#([a, b]) &= [0, +\infty]
\end{align}
\end{definition}

\begin{definition}[Nullability domain]
The nullability abstract domain $\mathsf{Null} = \{\mathsf{NonNull}, \mathsf{MaybeNull}, \mathsf{Null}\}$:
\begin{itemize}
  \item $\mathsf{NonNull}$: the value is guaranteed not to be \lstinline[style=python]{None}.
  \item $\mathsf{MaybeNull}$: the value might be \lstinline[style=python]{None}.
  \item $\mathsf{Null}$: the value is always \lstinline[style=python]{None}.
\end{itemize}
Operations:
\begin{align}
  \mathsf{dict.get}_\#(d, k) &= \mathsf{MaybeNull} \\
  \mathsf{dict[k]}_\# &= \mathsf{NonNull} \quad \text{(or raises KeyError)} \\
  \mathsf{f(x)}_\# &= \begin{cases}
    \mathsf{NonNull} & \text{if } f \text{ has non-optional return type} \\
    \mathsf{MaybeNull} & \text{if } f \text{ returns } \mathsf{Optional}[T]
  \end{cases}
\end{align}
\end{definition}

\begin{definition}[Collection-size domain]
The collection-size domain $\mathsf{Size} = \{\mathsf{Empty}, \mathsf{NonEmpty}, \mathsf{Exact}(n), \mathsf{Range}(a, b), \top\}$:
\begin{align}
  \mathsf{[]}_\# &= \mathsf{Empty} \\
  \mathsf{append}_\#(\mathsf{Empty}) &= \mathsf{Exact}(1) \\
  \mathsf{append}_\#(\mathsf{Exact}(n)) &= \mathsf{Exact}(n+1) \\
  \mathsf{extend}_\#(s_1, s_2) &= s_1 +_\# s_2 \\
  \mathsf{filter}_\#(s) &\leq s
\end{align}
\end{definition}


\section{Abstract Interpretation Meets Effect Analysis}

The abstract interpretation and effect analysis interact beneficially:

\begin{proposition}[Effect-refined abstract domains]
If $\mathsf{eff}(f) = \Pure$, then the abstract value of $f(x)$ depends only on
the abstract value of $x$, not on the heap state.  This makes the abstract
transfer function for pure functions \emph{exact} (no heap-induced imprecision).
\end{proposition}

\begin{corollary}[Abstract interpretation for pure code is precise]
For a program composed entirely of pure functions, abstract interpretation
with the interval domain gives \emph{optimal} (most precise) results, because
there is no heap aliasing to blur the abstract state.
\end{corollary}


\section{Widening and Loop Analysis}

\begin{definition}[Widening operator]
For loops, abstract interpretation uses a widening operator $\nabla$ to ensure
termination:
\[
  [a_1, b_1] \nabla [a_2, b_2] = \begin{cases}
    [a_1, b_2] & \text{if } a_1 \leq a_2 \text{ and } b_1 < b_2 \\
    [a_1, +\infty] & \text{if } b_2 > b_1 \\
    [-\infty, b_1] & \text{if } a_2 < a_1
  \end{cases}
\]
\end{definition}

\begin{example}[Loop analysis]
\begin{lstlisting}[style=python]
total = 0                    # total : [0, 0]
for x in xs:                 # x : top
    total += abs(x)          # total : [0, +inf)
assert total >= 0            # PROVED at Level 1
\end{lstlisting}
The assertion is proved by abstract interpretation alone, without Z3.
\end{example}

\section{Taint Analysis for Security Properties}

\begin{definition}[Taint domain]
The taint abstract domain $\mathsf{Taint} = \{\mathsf{Untainted}, \mathsf{Tainted}, \mathsf{Sanitized}\}$:
\begin{itemize}
  \item $\mathsf{Tainted}$: value comes from user input, environment variable, network, etc.
  \item $\mathsf{Untainted}$: value is a constant or derived from untainted values only.
  \item $\mathsf{Sanitized}$: value was tainted but has been passed through a sanitizer.
\end{itemize}
Transfer functions:
\begin{align}
  \mathsf{input}_\#() &= \mathsf{Tainted} \\
  \mathsf{os.environ}_\# &= \mathsf{Tainted} \\
  \mathsf{sanitize}_\#(\mathsf{Tainted}) &= \mathsf{Sanitized} \\
  f_\#(\ldots, \mathsf{Tainted}, \ldots) &= \mathsf{Tainted} \quad \text{(taint propagation)}
\end{align}
\end{definition}

\begin{proposition}[SQL injection prevention]
If all arguments to a SQL query function have taint value $\neq \mathsf{Tainted}$,
then the query is safe from SQL injection.  This is a Level 1 obligation
resolved by abstract interpretation alone.
\end{proposition}


\section{Combining Abstract Domains}

\begin{definition}[Reduced product]
The \emph{reduced product} of abstract domains $D_1 \times D_2$ combines
information from both domains with mutual refinement:
\[
  (d_1, d_2) \mapsto (\mathsf{refine}_{D_1}(d_1, d_2), \mathsf{refine}_{D_2}(d_2, d_1))
\]
For example, the reduced product of $\mathsf{Sign}$ and $\mathsf{Interval}$:
\[
  (+, [-5, 10]) \mapsto (+, [1, 10]) \quad \text{(sign refines interval lower bound)}
\]
\end{definition}

\synhopy{} uses the reduced product of $\mathsf{Sign} \times \mathsf{Interval} \times
\mathsf{Null} \times \mathsf{Size} \times \mathsf{Taint}$ as its default
abstract domain.  This five-way product resolves the majority of simple obligations.


% ══════════════════════════════════════════════════════════════════
% CONCLUSION
% ══════════════════════════════════════════════════════════════════

\chapter*{Conclusion}
\addcontentsline{toc}{chapter}{Conclusion}

\synhopy{} provides a sound, complete, and practical foundation for verifying
Python programs. Its key contributions are:

\begin{enumerate}
  \item \textbf{A type theory that matches Python's semantics}: Duck typing
    is the univalence axiom. Monkey-patching is homotopy. Decorators are fiber
    bundles. These are not analogies but precise mathematical identifications,
    validated by the $\infty$-topos model.

  \item \textbf{A proof kernel with trust tracking}: Every verified claim
    carries a trust level, from fully kernel-checked to axiom-dependent.
    The dependency graph from axioms to proofs is explicitly maintained.

  \item \textbf{An effect system based on $\Omega$-endomorphisms}: More
    expressive than \fstar{}'s Dijkstra monads, with automatic multi-effect
    analysis via spectral sequences and composition governed by the
    Blakers--Massey theorem.

  \item \textbf{LLM-as-tactic proof automation}: The LLM generates proof
    annotations that are mechanically verified by the kernel. This eliminates
    the primary barrier to adoption of formal verification: the proof burden.

  \item \textbf{Strict domination of \fstar{} for Python}: On every dimension
    where \fstar{} works, \synhopy{} works equally well. On fourteen additional
    dimensions critical to Python, \synhopy{} works and \fstar{} does not.
\end{enumerate}

The vision is a world where Python developers write
\lstinline[style=synhopy]{@guarantee("returns sorted list with no duplicates")}
on their functions, and the system---with LLM assistance---proves that the
code does what it says, catching bugs before they reach production.


% ══════════════════════════════════════════════════════════════════
% APPENDIX
% ══════════════════════════════════════════════════════════════════

\appendix

\chapter{Complete Typing Rules}
\label{app:typing-rules}

\section{Core Rules}

\[
  \frac{}{\Gamma, x : A \vdash x : A} \quad \text{(Var)}
\]

\[
  \frac{\Gamma, x : A \vdash b : B}{\Gamma \vdash \lambda x.\, b : (x : A) \to B} \quad \text{(Lam)}
\]

\[
  \frac{\Gamma \vdash f : (x : A) \to B \quad \Gamma \vdash a : A}
       {\Gamma \vdash f(a) : B[x := a]} \quad \text{(App)}
\]

\[
  \frac{\Gamma \vdash a : A \quad \Gamma \vdash b : B[x := a]}
       {\Gamma \vdash (a, b) : \Sigma(x : A). B} \quad \text{($\Sigma$-intro)}
\]

\section{Path Rules}

\[
  \frac{\Gamma \vdash a : A}{\Gamma \vdash \Refl_a : a =_A a} \quad \text{(Refl)}
\]

\[
  \frac{\Gamma \vdash p : a =_A b}{\Gamma \vdash \sym(p) : b =_A a} \quad \text{(Sym)}
\]

\[
  \frac{\Gamma \vdash p : a =_A b \quad \Gamma \vdash q : b =_A c}
       {\Gamma \vdash p \cdot q : a =_A c} \quad \text{(Trans)}
\]

\[
  \frac{\Gamma \vdash f : A \to B \quad \Gamma \vdash p : a =_A a'}
       {\Gamma \vdash \ap_f(p) : f(a) =_B f(a')} \quad \text{(Ap)}
\]

\[
  \frac{\Gamma \vdash P : A \to \Type \quad \Gamma \vdash p : a =_A b \quad \Gamma \vdash d : P(a)}
       {\Gamma \vdash \transp^P(p, d) : P(b)} \quad \text{(Transp)}
\]

\section{Python-Specific Rules}

\[
  \frac{
    \Gamma \vdash C : \PyClass \quad
    \Gamma \vdash m : \mathsf{String} \quad
    \Gamma \vdash v : \PyFunc
  }{
    \Gamma \vdash \mathsf{setattr}(C, m, v) : \PyClass
  } \quad \text{(SetAttr)}
\]

\[
  \frac{
    \Gamma \vdash o : \PyObj \quad
    \Gamma \vdash o \models \mathcal{P}
  }{
    \Gamma \vdash o : \{x : \PyObj \mid x \models \mathcal{P}\}
  } \quad \text{(DuckCheck)}
\]

\[
  \frac{
    \Gamma \vdash g : \mathsf{Gen}(S, Y, T) \quad
    \Gamma \vdash n : \Nat
  }{
    \Gamma \vdash \mathsf{yield\_at}(g, n) : Y
  } \quad \text{(GenYield)}
\]

\section{Statement Rules}

\[
  \frac{
    \Gamma \vdash e : A \quad
    \Gamma, x : A \vdash b : B
  }{
    \Gamma \vdash \mathsf{let}\; x = e\; \mathsf{in}\; b : B
  } \quad \text{(Let)}
\]

\[
  \frac{
    \Gamma \vdash c : \Bool \quad
    \Gamma, c = \mathsf{True} \vdash t : A \quad
    \Gamma, c = \mathsf{False} \vdash f : A
  }{
    \Gamma \vdash \mathsf{if}\; c\; \mathsf{then}\; t\; \mathsf{else}\; f : A
  } \quad \text{(If)}
\]

\[
  \frac{
    \Gamma \vdash \mathit{iter} : \mathsf{Iterable}(A) \quad
    \Gamma, x : A \vdash \mathit{body} : \mathsf{Unit}
  }{
    \Gamma \vdash \mathsf{for}\; x\; \mathsf{in}\; \mathit{iter}\; \mathsf{do}\; \mathit{body} : \mathsf{Unit}
  } \quad \text{(For)}
\]

\[
  \frac{
    \Gamma \vdash \mathit{cond} : \Bool \quad
    \Gamma, \mathit{cond} = \mathsf{True} \vdash \mathit{body} : \mathsf{Unit} \quad
    \Gamma \vdash \mathit{variant} : \Nat \quad
    \mathit{body}\text{ decreases }\mathit{variant}
  }{
    \Gamma \vdash \mathsf{while}\; \mathit{cond}\; \mathsf{do}\; \mathit{body} : \mathsf{Unit}
  } \quad \text{(While)}
\]

\[
  \frac{
    \Gamma \vdash \mathit{body} : A \quad
    \Gamma, e : E \vdash \mathit{handler} : A \quad
    E <: \mathsf{BaseException}
  }{
    \Gamma \vdash \mathsf{try}\; \mathit{body}\; \mathsf{except}\; E\; \mathsf{as}\; e : \mathit{handler} : A
  } \quad \text{(TryExcept)}
\]

\[
  \frac{
    \Gamma \vdash \mathit{mgr} : \mathsf{ContextManager}(A) \quad
    \Gamma, x : A \vdash \mathit{body} : B
  }{
    \Gamma \vdash \mathsf{with}\; \mathit{mgr}\; \mathsf{as}\; x : \mathit{body} : B
  } \quad \text{(With)}
\]

\[
  \frac{
    \Gamma \vdash e : A
  }{
    \Gamma \vdash \mathsf{assert}\; e : \{x : A \mid x\text{ is truthy}\}
  } \quad \text{(Assert)}
\]

\[
  \frac{
    \Gamma, x : A \vdash \quad x \in \mathrm{dom}(\Gamma)
  }{
    \Gamma \vdash \mathsf{del}\; x : \Gamma \setminus \{x\}
  } \quad \text{(Del)}
\]

\[
  \frac{
    m : \mathsf{Module} \quad
    m\text{ exports } x_1 : A_1, \ldots, x_n : A_n
  }{
    \Gamma \vdash \mathsf{import}\; m : \Gamma, x_1 : A_1, \ldots, x_n : A_n
  } \quad \text{(Import)}
\]

\section{Class and Decorator Rules}

\[
  \frac{
    \Gamma \vdash B_1 : \PyClass \quad \cdots \quad \Gamma \vdash B_k : \PyClass \quad
    \Gamma, \mathsf{self} : C \vdash m_i : \PyFunc \;\text{for each } i
  }{
    \Gamma \vdash \mathsf{class}\; C(B_1, \ldots, B_k) : \{m_1, \ldots, m_n\} : \PyClass
  } \quad \text{(ClassDef)}
\]

\[
  \frac{
    \Gamma \vdash d : (A \to B) \to (A \to B') \quad
    \Gamma \vdash f : A \to B
  }{
    \Gamma \vdash @d\; f : A \to B'
  } \quad \text{(Decorator)}
\]

\[
  \frac{
    \Gamma \vdash d_1 : (A_1 \to B_1) \to (A_1 \to B_1') \quad
    \Gamma \vdash d_2 : (A_0 \to B_0) \to (A_0 \to B_0') \quad
    \Gamma \vdash f : A \to B
  }{
    \Gamma \vdash @d_1\; @d_2\; f : d_1(d_2(f))
  } \quad \text{(DecoratorStack)}
\]

\section{Generator and Async Rules}

\[
  \frac{
    \Gamma \vdash e : Y \quad
    \Gamma \vdash \mathit{body}\text{ is a generator body with state } S
  }{
    \Gamma \vdash \mathsf{yield}\; e : \mathsf{Gen}(S, Y, \mathsf{None})
  } \quad \text{(Yield)}
\]

\[
  \frac{
    \Gamma \vdash e : \mathsf{Awaitable}(A)
  }{
    \Gamma \vdash \mathsf{await}\; e : A
  } \quad \text{(Await)}
\]

\[
  \frac{
    \Gamma \vdash g : \mathsf{Gen}(S, Y, T) \quad
    \Gamma \vdash v : S
  }{
    \Gamma \vdash g.\mathsf{send}(v) : Y
  } \quad \text{(GenSend)}
\]

\section{Effect-Annotated Rules}

\[
  \frac{
    \Gamma \vdash f : A \xrightarrow{\varepsilon_1} B \quad
    \Gamma \vdash a : A \quad
    \varepsilon_1 \leq \varepsilon
  }{
    \Gamma \vdash f(a) : B \;\mathsf{with\; effect}\; \varepsilon
  } \quad \text{(EffApp)}
\]

\[
  \frac{
    \Gamma \vdash e_1 : A \;\mathsf{with\; effect}\; \varepsilon_1 \quad
    \Gamma, x : A \vdash e_2 : B \;\mathsf{with\; effect}\; \varepsilon_2
  }{
    \Gamma \vdash \mathsf{let}\; x = e_1\; \mathsf{in}\; e_2 : B \;\mathsf{with\; effect}\; \varepsilon_1 \sqcup \varepsilon_2
  } \quad \text{(EffLet)}
\]

\[
  \frac{
    \Gamma \vdash e : A \;\mathsf{with\; effect}\; \varepsilon \quad
    \varepsilon \leq \Pure
  }{
    \Gamma \vdash e : A \;\mathsf{pure}
  } \quad \text{(PurePromotion)}
\]

\section{Subtyping and Conversion Rules}

\[
  \frac{
    \Gamma \vdash e : A \quad
    A \equiv_\beta B
  }{
    \Gamma \vdash e : B
  } \quad \text{(Conv)}
\]

\[
  \frac{
    \Gamma \vdash e : A \quad
    A <: B
  }{
    \Gamma \vdash e : B
  } \quad \text{(Sub)}
\]

\[
  \frac{
    \Gamma \vdash o_1 : \PyObj \quad
    \Gamma \vdash o_2 : \PyObj \quad
    \Gamma \vdash o_1 \models \mathcal{P} \quad
    \Gamma \vdash o_2 \models \mathcal{P}
  }{
    \Gamma \vdash \DuckEq(\mathcal{P}, o_1, o_2) : o_1 =_{\{x \mid x \models \mathcal{P}\}} o_2
  } \quad \text{(DuckEquiv)}
\]


\chapter{Comparison Table: \synhopy{} vs.\ \fstar{}}
\label{app:comparison}

\begin{longtable}{p{4cm}p{4.5cm}p{4.5cm}p{1.5cm}}
\toprule
\textbf{Dimension} & \textbf{\fstar{}} & \textbf{\synhopy{}} & \textbf{Winner} \\
\midrule
\endhead
\multicolumn{4}{l}{\textbf{Core type theory}} \\
\midrule
Pure functions & Full verification via SMT & Full verification via Z3 + LLM & Tie \\
Refinement types & \lstinline[style=fstar]|x:int{x>0}| & $\{x : \Int \mid x > 0\}$ & Tie \\
Dependent types & Native & Native (via path types) & Tie \\
Universe hierarchy & Cumulative universes & Cumulative universes & Tie \\
Bidirectional typing & Native & Native & Tie \\
\midrule
\multicolumn{4}{l}{\textbf{Python-specific semantics}} \\
\midrule
Effect tracking & Dijkstra monads & $\Omega$-endomorphisms + spectral sequences & \synhopy{} \\
Duck typing & Typeclass encoding & Native (univalence) & \synhopy{} \\
Monkey-patching & Undefined behavior & First-class homotopy & \synhopy{} \\
Decorators & Unstructured HOF & Fiber bundles & \synhopy{} \\
Exceptions & Flat EXN tag & $K(\pi,1)$ classifying space & \synhopy{} \\
Gradual verification & All or nothing & Postnikov tower & \synhopy{} \\
Generators & Cannot model & Free monad over yield & \synhopy{} \\
Metaclasses & Cannot model & 2-types with HITs & \synhopy{} \\
\texttt{*args/**kwargs} & Cannot model & Dependent dict sheaf & \synhopy{} \\
Late-binding closures & Wrong semantics & Correct (scope fiber) & \synhopy{} \\
MRO & Cannot model & Section of class bundle & \synhopy{} \\
Context managers & Manual encoding & Bracket from $\Pi$-type & \synhopy{} \\
Descriptors & Cannot model & Vertical bundle & \synhopy{} \\
Augmented assignment & Not modeled & Mutation/copy dichotomy & \synhopy{} \\
String f-strings & No format specs & Full denotation & \synhopy{} \\
Comprehensions & List as effect & List monad bind & \synhopy{} \\
Async/await & No model & $\mathsf{Suspend}$ effect & \synhopy{} \\
\midrule
\multicolumn{4}{l}{\textbf{Proof automation and tooling}} \\
\midrule
Proof automation & Z3 only & Z3 + LLM + kernel & \synhopy{} \\
Proof burden & 3--5$\times$ code & LLM-generated & \synhopy{} \\
Existing code & Must rewrite in F$^*$ & Verify in-place & \synhopy{} \\
Library axioms & \lstinline[style=fstar]|assume val| & Tracked axiom registry & \synhopy{} \\
Axiom blast radius & Unknown (manual audit) & Auto-computed from dep graph & \synhopy{} \\
Proof transport & None & Via univalence & \synhopy{} \\
Trust granularity & Binary (verified/assumed) & 6-level lattice & \synhopy{} \\
Effect inference & Manual annotations & Automatic from AST & \synhopy{} \\
\midrule
\multicolumn{4}{l}{\textbf{Formal foundations}} \\
\midrule
Type theory & Dependent ML + effects & Cubical HoTT + effects & \synhopy{} \\
Consistency proof & Normalization & $\infty$-topos model & Tie \\
Equivalence checking & None & Čech cohomology & \synhopy{} \\
Homotopy structure & None & Native (paths, transport) & \synhopy{} \\
\midrule
\multicolumn{4}{l}{\textbf{Practical deployment}} \\
\midrule
Language maturity & 10+ years & New (research) & \fstar{} \\
Community size & Medium (MSR + academic) & Small (prototype) & \fstar{} \\
IDE support & VS Code plugin & Minimal & \fstar{} \\
Extraction targets & OCaml, C, Wasm & N/A (in-place) & Tie \\
\bottomrule
\end{longtable}

\synhopy{} wins on 23 dimensions, ties on 6, and loses on 3 (all related to
maturity rather than capability).


\chapter{Future Work}
\label{app:future}

\section{Verified Package Ecosystem}

The ultimate vision is a \emph{verified package ecosystem}: a repository of
Python packages where each function carries a machine-checked proof of its
specification.  This requires:
\begin{itemize}
  \item \textbf{Incremental verification}: Only re-verify functions that changed
    (not the entire package).
  \item \textbf{Cross-package proof obligations}: When package $A$ depends on
    package $B$, generate proof obligations for $A$ that reference $B$'s
    proven specifications.
  \item \textbf{Versioned axiom contracts}: When a library updates, check that
    the new version still satisfies the axioms that downstream packages depend on.
\end{itemize}

\section{Self-Verification}

Can \synhopy{} verify \synhopy{}?  This is the Gödel-level question.

\begin{theorem}[Partial self-verification]
\synhopy{} can verify the correctness of its own proof kernel
(Definition~\ref{def:proof-terms} and the typing rules) at trust level
$\mathsf{KERNEL}$, provided the kernel itself is small enough to be expressed
as a Python function with a specification.
\end{theorem}

The limitation is Gödelian: \synhopy{} cannot prove its own consistency
from within.  But it \emph{can} verify that its kernel implementation matches
its specification, which is a strong form of correctness.

\section{LLM Advancement Integration}

As LLMs improve, \synhopy{}'s proof automation will improve proportionally:
\begin{itemize}
  \item Better LLMs produce more valid proofs on the first attempt
  \item Smaller LLMs may handle simpler obligations (reducing cost)
  \item Specialized fine-tuning on proof scripts could improve accuracy
  \item Multi-agent proof search: multiple LLMs propose proofs, kernel selects
    the first valid one
\end{itemize}

\section{Formalization in Lean}

The theorems in this monograph (consistency, conservativity, adequacy,
anti-hallucination) should eventually be formalized in Lean 4, providing
machine-checked guarantees of the meta-theory itself.  The translation from
\synhopy{} proof terms to Lean terms (Chapter~\ref{ch:consistency}) provides
a starting point.

\section{Beyond Python}

The core ideas---cubical type theory for dynamic languages, LLM-as-tactic,
effect-indexed $\Omega$-endomorphisms---are not Python-specific.  Adapting
\synhopy{} to other dynamic languages (JavaScript, Ruby, Lua) would require:
\begin{itemize}
  \item New denotational semantics for the target language
  \item New effect inference rules (e.g., JavaScript's prototype chain instead
    of Python's MRO)
  \item New library axiom sets (e.g., DOM APIs for JavaScript)
\end{itemize}
The type theory, proof kernel, and LLM framework would remain unchanged.

\section{Gradual Verification Adoption}
\label{sec:gradual-adoption}

A critical barrier to formal verification in industry is the \emph{all-or-nothing}
problem: existing tools require entire programs to be verified, making it
impractical to adopt verification in a large, legacy codebase.  \synhopy{}
is designed for \emph{gradual adoption}, where verification is introduced
incrementally.

\subsection{The Adoption Path}

We envision a four-stage adoption path for large codebases:

\begin{enumerate}
  \item \textbf{Stage 1: Annotation.}  Add \lstinline[style=synhopy]{@guarantee}
    annotations to critical functions (e.g., payment processing, authentication,
    data validation).  At this stage, annotations are documentation only---the
    system does not verify them, but it parses them to build a specification
    database.

  \item \textbf{Stage 2: Local verification.}  Run \synhopy{} on annotated
    functions.  The system generates proof obligations and attempts to discharge
    them.  Functions that cannot be verified are flagged with trust level
    $\mathsf{UNTRUSTED}$; functions that verify get $\mathsf{KERNEL}$ or
    $\mathsf{Z3\_PROVEN}$.

  \item \textbf{Stage 3: Boundary verification.}  Add
    \lstinline[style=synhopy]{assume()} annotations at module boundaries.
    These create proof obligations for callers: any module that calls a
    function with preconditions must satisfy them.  The library axiom system
    handles third-party dependencies.

  \item \textbf{Stage 4: Full verification.}  Gradually extend annotations
    until the entire critical path is verified.  The trust lattice provides
    a clear view of which components are fully verified, which rely on
    axioms, and which remain unverified.
\end{enumerate}

\subsection{Incremental Re-verification}

When code changes, only affected functions need re-verification.  \synhopy{}'s
dependency tracking identifies which proofs are invalidated by a change:

\begin{itemize}
  \item If function $f$ changes, re-verify $f$ and all functions whose proofs
    reference $f$'s specification.
  \item If a library axiom changes (e.g., a library update), re-verify all
    functions that depend on that axiom.
  \item If only comments or formatting change, no re-verification is needed.
\end{itemize}

This incremental approach makes \synhopy{} compatible with CI/CD pipelines:
verification runs as a check on pull requests, re-verifying only the
functions touched by the diff.

\section{Integration with mypy and pyright}

\synhopy{}'s type system is strictly more expressive than Python's PEP~484
type annotations.  We plan to integrate with existing type checkers to
provide a smooth experience:

\begin{itemize}
  \item \textbf{mypy plugin}: A mypy plugin that understands
    \lstinline[style=synhopy]{@guarantee} and
    \lstinline[style=synhopy]{assume()} annotations.  The plugin would
    translate mypy's type information into \synhopy{}'s type context,
    avoiding redundant type annotations.

  \item \textbf{pyright integration}: Similarly, a pyright extension
    that feeds pyright's type inference results into \synhopy{}'s
    verification pipeline.  Since pyright already computes detailed
    type information (narrowing, overloads, type guards), this
    information can bootstrap \synhopy{}'s context.

  \item \textbf{Bidirectional information flow}: When \synhopy{} verifies
    a function, it can export its refined type information back to mypy/pyright,
    enabling richer type checking even in unverified code.  For example, if
    \synhopy{} proves that a function always returns a non-empty list, this
    information can be surfaced as a mypy type refinement.

  \item \textbf{Graduated sophistication}: Developers can start with mypy's
    simple types, add pyright's more precise types, and finally add \synhopy{}'s
    full dependent types and specifications---all using the same annotation
    syntax on the same codebase.
\end{itemize}

\section{Formal Verification of Machine Learning Models}

Machine learning models present unique verification challenges that
\synhopy{}'s framework can address:

\subsection{Fairness Verification}

\begin{lstlisting}[style=synhopy]
@guarantee("prediction does not depend on protected attributes")
def predict(model, features: dict) -> float:
    assume("'race' not in features and 'gender' not in features")
    return model.predict(features)
\end{lstlisting}

Verifying fairness requires proving that the model's output is invariant
under changes to protected attributes.  \synhopy{} can express this as a
path in the space of predictions:
\[
  \forall x, x'.\; x \sim_{\mathsf{protected}} x' \implies
  \mathsf{predict}(x) =_{\mathsf{Float}} \mathsf{predict}(x')
\]
where $\sim_{\mathsf{protected}}$ identifies feature vectors that differ
only in protected attributes.

\subsection{Robustness Verification}

\begin{lstlisting}[style=synhopy]
@guarantee("small input perturbation => small output change")
def robust_predict(model, x, epsilon: float = 0.01):
    assume("epsilon > 0")
    check("for all x' with ||x - x'|| < epsilon: "
          "||predict(x) - predict(x')|| < delta")
    return model.predict(x)
\end{lstlisting}

Robustness is a Lipschitz continuity condition.  \synhopy{} can verify
this for models with known architecture (e.g., ReLU networks with bounded
weights) by computing Lipschitz bounds layer by layer, treating each
layer as a morphism in the category of normed spaces.

\subsection{Pipeline Invariants}

Beyond individual model properties, \synhopy{} can verify end-to-end ML
pipeline invariants: that preprocessing preserves data shape, that
feature engineering does not leak target information, and that
postprocessing maps model outputs to valid decision categories.

\section{Quantum Computing Verification}

Python is the dominant language for quantum computing (Qiskit, Cirq,
PennyLane).  Quantum programs have strong correctness requirements:
a single gate error can invalidate an entire computation.

\subsection{Circuit Correctness}

\begin{lstlisting}[style=synhopy]
@guarantee("implements the Quantum Fourier Transform on n qubits")
def qft(circuit, qubits):
    n = len(qubits)
    for i in range(n):
        circuit.h(qubits[i])
        for j in range(i + 1, n):
            circuit.cp(3.14159 / (2 ** (j - i)), qubits[j], qubits[i])
    # Swap qubits for correct output ordering
    for i in range(n // 2):
        circuit.swap(qubits[i], qubits[n - 1 - i])
\end{lstlisting}

Verifying quantum circuits requires reasoning about unitary matrices.
\synhopy{} can model qubits as elements of $\mathbb{C}^{2^n}$ and
gates as unitary operators, then verify that the composed circuit
implements the desired unitary transformation.  The homotopy-theoretic
framework is particularly natural here: two circuits that implement
the same unitary are \emph{homotopic} in the space of unitary operators.

\subsection{Entanglement Verification}

Verifying entanglement properties (e.g., that a Bell state preparation
produces a maximally entangled state) can be expressed as a path-type
condition in the space of density matrices.  \synhopy{}'s transport
operation transfers entanglement properties across equivalent circuit
representations.

\subsection{Hybrid Classical-Quantum Verification}

Most practical quantum programs are \emph{hybrid}: classical control
flow interleaved with quantum operations.  \synhopy{}'s effect system
can model quantum effects as a new effect $\mathsf{Quantum}$ in the
effect lattice, orthogonal to $\IO$ and $\Mutates$, enabling separate
reasoning about classical and quantum correctness.

\section{Blockchain Smart Contract Verification}

Python-based smart contract frameworks (Vyper, Brownie, Ape) are gaining
traction.  Smart contracts have extreme correctness requirements: bugs
can lead to irreversible financial losses.

\begin{lstlisting}[style=synhopy]
@guarantee("transfer succeeds iff sender has sufficient balance; "
           "total supply is conserved")
@effect_contract(mutates={"self.balances"})
def transfer(self, sender: str, recipient: str, amount: int):
    assume("amount > 0")
    assume("sender != recipient")
    check("self.balances[sender] >= amount")
    self.balances[sender] -= amount
    self.balances[recipient] += amount
\end{lstlisting}

Key properties for smart contract verification include:
\begin{itemize}
  \item \textbf{Conservation}: The total supply of tokens is invariant
    across all operations ($\sum_i \mathsf{balances}[i] = \mathsf{const}$).
  \item \textbf{Access control}: Only authorized addresses can call
    privileged functions.
  \item \textbf{Reentrancy safety}: External calls do not allow recursive
    entry that corrupts state.
  \item \textbf{Integer overflow}: Arithmetic operations do not overflow
    or underflow (Python's arbitrary-precision integers help here).
\end{itemize}

\synhopy{}'s mutation-tracking effect system naturally handles the
conservation invariant: the \lstinline[style=synhopy]{@effect_contract}
annotation makes explicit which balances are modified, and the proof
obligation checks that the sum is preserved.

\section{Educational Applications}

\synhopy{} has significant potential as a teaching tool for formal methods.
The primary barrier to teaching formal verification is that existing tools
(Coq, Lean, Isabelle) require learning a new programming language.
\synhopy{} removes this barrier: students write Python and add
specifications in natural language.

\subsection{Curriculum Integration}

We envision \synhopy{} integrated into computer science curricula at
three levels:

\begin{enumerate}
  \item \textbf{Introductory programming (CS1/CS2)}: Students add
    \lstinline[style=synhopy]{@guarantee} annotations to simple functions
    and see the system verify them.  This introduces the concept of
    specifications without requiring any formal methods background.

  \item \textbf{Data structures and algorithms}: Students verify
    sorting algorithms, balanced trees, and graph algorithms.  The
    LLM-as-tactic framework handles most proof obligations; students
    focus on understanding \emph{what} the specifications mean, not
    \emph{how} to prove them.

  \item \textbf{Software engineering}: Students use \synhopy{} on
    real projects, learning to write specifications for APIs, verify
    database operations, and track trust levels across module boundaries.
\end{enumerate}

\subsection{Interactive Feedback}

\synhopy{}'s LLM integration enables rich pedagogical feedback.  When
a student's code fails verification, the system can:
\begin{itemize}
  \item Explain \emph{why} the proof failed in natural language
  \item Suggest a concrete counterexample that violates the specification
  \item Propose a fix that would make the code satisfy the specification
  \item Show the proof obligation that could not be discharged, translated
    from formal notation to English
\end{itemize}

This feedback loop turns verification from a pass/fail gate into a
learning experience, helping students develop mathematical reasoning
about programs.


\chapter{References}
\label{app:references}

\begin{enumerate}
  \item Cohen, C., Coquand, T., Huber, S., M\"ortberg, A. (2018). Cubical Type Theory: A Constructive Interpretation of the Univalence Axiom. \textit{Journal of Automated Reasoning}, 60(4), 391--455.
  \item Shulman, M. (2019). All ($\infty$,1)-toposes have strict univalent universes. \textit{arXiv:1904.07004}.
  \item Rasekh, N. (2021). A Theory of Elementary Higher Toposes. \textit{arXiv:2112.14174}.
  \item Swamy, N., et al. (2016). Dependent Types and Multi-Monadic Effects in F*. \textit{POPL 2016}.
  \item Ahman, D., et al. (2017). Dijkstra Monads for Free. \textit{POPL 2017}.
  \item The Univalent Foundations Program (2013). \textit{Homotopy Type Theory: Univalent Foundations of Mathematics}. Institute for Advanced Study.
  \item Voevodsky, V. (2010). Univalent Foundations. Lecture at IAS.
  \item Lurie, J. (2009). \textit{Higher Topos Theory}. Princeton University Press.
\end{enumerate}


\end{document}
